%% ============================================================
%% PARTE II -- SEQUENCIA DIDATICA
%% Sistema Kumon Expandido -- Niveis 1 a 4
%% Resolucao CNE/CEB no 2/2012 -- Principios da formacao integral
%% BNCC EF01LP01 a EF01LP09 | CAEd D1-D5 | SPAECE EF01LP01-EF01LP09
%% ============================================================

%% ============================================================
%% UNIDADE 1 -- VOGAIS (Nivel 1 -- Pre-silabico)
%% ============================================================

%% ============================================================
%% UNIDADE 0 -- PRONTIDAO (Nivel K0)
%% Rotina neuroinclusiva, motricidade fina, atencao e oralidade
%% ============================================================

\chapter{Unidade 0 --- Prontidão para a Alfabetização}

% Rotina neuroinclusiva e badges (R201)
\rotinasimples
\nivelKbadge{K0 --- Prontidao}
\rotabadge{1A --- Recomposicao Intensiva}

Esta unidade prepara o corpo e a atenção da criança para a alfabetização.
Não há letras novas: aqui se trabalham motricidade fina, direção do traçado,
escuta atenta e oralidade. É o ponto de partida da Rota 1A e pode ser usada
por qualquer aluno como aquecimento.

\section{Atividade 0.1 --- Movimentos Preparatórios}

\instrucao{Passe o dedo sobre cada linha, da esquerda para a direita.}

\begin{center}
\begin{tikzpicture}
  \draw[very thick, color=primary, dashed] (0,0) -- (8,0);
  \draw[very thick, color=secondary, dashed] (0,-1) -- (8,-1);
  \draw[very thick, color=accent, dashed] (0,-2) -- (8,-2);
  \node[right, font=\small] at (0.2,0.3) {$\rightarrow$};
  \node[right, font=\small] at (0.2,-0.7) {$\rightarrow$};
  \node[right, font=\small] at (0.2,-1.7) {$\rightarrow$};
\end{tikzpicture}
\end{center}

\caixapista{Trace primeiro no ar, com o dedo grande. Depois passe o dedo no papel.}

\caixapausa{Se cansar, feche os olhos e respire fundo três vezes.}

\section{Atividade 0.2 --- Curvas e Círculos}

\instrucao{Passe o dedo sobre cada curva.}

\begin{center}
\begin{tikzpicture}
  \draw[very thick, color=primary] (0,0) to[out=0,in=180] (3,0);
  \draw[very thick, color=secondary] (4.5,0) to[out=0,in=180] (7.5,0);
  \draw[very thick, color=accent] (9,0) circle (0.8);
\end{tikzpicture}
\end{center}

\section{Atividade 0.3 --- Escuta Atenta}

\instrucao{Ouça o som que o professor faz. Aponte a figura que combina.}

\begin{center}
\begin{tabular}{ccc}
\ilustra{Som: cachorro}{au-au} & \ilustra{Som: gato}{miau} & \ilustra{Som: sino}{dim-dom}
\end{tabular}
\end{center}

\section{Atividade 0.4 --- Fala e Ritmo}

\instrucao{Fale as palavras devagar, batendo palma em cada parte.}

\begin{center}
BO -- LA \quad CA -- SA \quad SA -- PO \quad BO -- NE -- CA
\end{center}

\caixapista{Cada palma é uma parte da palavra. Conte as palmas.}

\section{Atividade 0.5 --- Orientação Espacial}

\instrucao{Siga a ordem: para cima, para baixo, para o lado.}

\begin{center}
\begin{tikzpicture}
  \node[draw=primary, minimum width=2.5cm, minimum height=1.2cm, align=center] at (0,0) {EM CIMA};
  \node[draw=secondary, minimum width=2.5cm, minimum height=1.2cm, align=center] at (0,-2) {EM BAIXO};
  \node[draw=accent, minimum width=2.5cm, minimum height=1.2cm, align=center] at (4,-1) {AO LADO};
\end{tikzpicture}
\end{center}

\checklistdominio{Completar as 5 atividades da Unidade 0 com confiança}

\chapter{Unidade 1 --- As Vogais}

%% ------------------------------------------------------------
%% LICAO 1: LETRA A
%% ------------------------------------------------------------
\section{Licao 1: A Letra A}

% Rotina neuroinclusiva e badges (R201)
\rotinasimples
\nivelKbadge{K1 --- Letras e Grafias}
\rotabadge{1A --- Recomposicao Intensiva}


\letra{A}
% Quatro formas gráficas da letra (R201)
\quatroformasletra{A}{a}

\boq[aberta]{A}{%
  Boca bem ABERTA, maxilar relaxado.\\
  Ar sai forte e livre pela boca.\\
  Lingua repousada no assoalho bucal.\\
  Som: ``AAAAAAA'' (aberto, vibrante).\\[0.5em]
  Palavras exemplo: agua, aviao, abacaxi, anel, arco.
}

\pistaarticulatoria{Som: \'{a}\'{a}\'{a} (aberto, longo)}{Boca bem aberta, maxilar relaxado}{M\~{a}o na garganta para sentir a vibra\c{c}\~{a}o da voz}

\metadata{A-01}{Nivel 1 -- Pre-silabico}{EF01LP01}{D1}{EF01LP01}{Resolucao CNE/CEB 7/2010, art. 12}

\plancheta{Folha de Pratica 1.1 --- Letra A}
\begin{exercicio}[Folha de Pratica 1.1 --- Letra A]
\noindent\textbf{Nome:} \underline{\hspace{3.2cm}} \quad \textbf{Data:} \underline{\hspace{1.6cm}} \quad \textbf{Nota:} \underline{\hspace{0.9cm}}/10

\subsection*{PARTE 1: RECONHECIMENTO}

\begin{enumerate}
  \item Circule todas as letras ``A'' que encontrar na sequencia:
  \begin{center}
  \fbox{A} \quad B \quad \fbox{A} \quad C \quad \fbox{A} \quad D \quad \fbox{A} \quad E \quad \fbox{A} \quad F
  \end{center}

  \item Sublinhe todas as palavras que comecam com a letra ``A'':
  \begin{center}
  \uline{agua} \quad bola \quad \uline{avi\~{a}o} \quad sapo \quad \uline{abacaxi} \quad lua \quad \uline{anel}
  \end{center}

  \item Coloque um $\checkmark$ ao lado de cada linha que contiver a letra ``A'':
  \begin{center}
  \begin{tabular}{ll}
  $\checkmark$ \texttt{A R T O} & $\checkmark$ \texttt{C A S A} \\
  $\checkmark$ \texttt{B O L A} & $\square$ \texttt{L U R O} \\
  $\checkmark$ \texttt{S A P O} & $\square$ \texttt{P I N T O}
  \end{tabular}
  \end{center}

  \item Identifique a letra ``A'' no nome abaixo e circule-a:
  \begin{center}
  \texttt{M A R C E L O} $\rightarrow$ a letra A est\'a na posicao \underline{\hspace{1cm}}
  \end{center}

  \item Diga em voz alta tres palavras que comecam com ``A'' e anote aqui:
  \begin{center}
  1. \underline{\hspace{4cm}} \quad 2. \underline{\hspace{4cm}} \quad 3. \underline{\hspace{4cm}}
  \end{center}
\end{enumerate}

\subsection*{PARTE 2: PRODUCAO GUIADA}

\begin{enumerate}[resume]
  \item Complete com a letra ``A'' nas palavras abaixo:
  \begin{center}
  \underline{A}viao \quad \underline{A}gua \quad \underline{A}bacaxi \quad \underline{A}nimal \quad \underline{A}zul
  \end{center}

  \item Escreva a letra ``A'' maiuscula cinco vezes na linha pontilhada:
  \begin{center}
  \underline{\hspace{1.5cm}} \quad \underline{\hspace{1.5cm}} \quad \underline{\hspace{1.5cm}} \quad \underline{\hspace{1.5cm}} \quad \underline{\hspace{1.5cm}}
  \end{center}

  \item Escreva a letra ``A'' minuscula cinco vezes:
  \begin{center}
  \underline{\hspace{1.5cm}} \quad \underline{\hspace{1.5cm}} \quad \underline{\hspace{1.5cm}} \quad \underline{\hspace{1.5cm}} \quad \underline{\hspace{1.5cm}}
  \end{center}

  \item Copie as palavras abaixo, prestando atencao na letra ``A'':
  \begin{center}
  agua \quad aviao \quad abacaxi \quad anel \quad arco
  \end{center}

  \item Escreva o nome completo da palavra que completa a frase:
  \begin{center}
  O menino bebe \underline{\hspace{3cm}} (a agua).
  \end{center}
\end{enumerate}

\subsection*{PARTE 3: INTEGRACAO}

\begin{enumerate}[resume]
  \item Leia em voz alta para o professor:
  \begin{center}
  $\rightarrow$ A --- AGUA --- AVIAO --- ABACAXI
  \end{center}

  \item Escreva o nome de uma fruta que comeca com ``A'':
  \begin{center}
  \underline{\hspace{6cm}}
  \end{center}

  \item Desenhe algo que comeca com ``A'' e escreva o nome abaixo:
  \begin{center}
  \begin{tikzpicture}
    \draw[thick, dashed] (0,0) rectangle (5,3);
    \node[below] at (2.5,-0.3) {Nome: \underline{\hspace{4cm}}};
  \end{tikzpicture}
  \end{center}

  \item Circule a opcao correta:
  \begin{enumerate}[label=(\alph*)]
    \item A letra A tem som de: ( ) aberto ( ) fechado ( ) nasal
    \item A letra A aparece em: ( ) bolacha ( ) abacaxi ( ) limao
    \item A boca fica: ( ) fechada ( ) aberta ( ) entreaberta
  \end{enumerate}

  \item Leitura assistida: o professor aponta e o aluno le em voz alta:
  \begin{center}
  a \quad agua \quad aviao \quad abacaxi \quad anel \quad arco \quad azul
  \end{center}
\end{enumerate}
\end{exercicio}

\begin{dica}
\textbf{Dica para o Professor:} Ao apresentar a letra A, comece pela boquinha. Mostre a posicao da boca (bem aberta) e peca que os alunos repitam o som em unissono. Depois, trace a letra no quadro com setas indicando a direcao do tracado. Os alunos devem tracar no ar com o dedo antes de escrever no papel. Reforce que A e uma vogal, ou seja, tem som proprio e nao depende de outra letra.
\end{dica}

%% ------------------------------------------------------------


% ============================================================
% CONSOLIDACAO DA LETRA A (R201.1) -- Missao 1/16
% ============================================================

\plancheta{Miss\~a{}o 1 --- Letra A}

\begin{missao}
\textbf{\faGamepad~Miss\~a{}o 1 de 16:} Consolide a letra A com as três atividades. Depois descubra os quatro sons do A, conheça sua família silábica e complete a missão motora!
\end{missao}

\begin{consolidacao}
\textbf{Atividade 1 --- Ca\c{c}ada ao A}

\instrucao{Circule todas as letras A que encontrar.}

\begin{center}
\begin{tabular}{cccccc}
\fbox{A} & E & \fbox{A} & I & \fbox{A} & O \\
U & \fbox{A} & E & \fbox{A} & I & \fbox{A}
\end{tabular}
\end{center}
Agora sublinhe as palavras com A:
\begin{center}
\uline{asa} \quad \uline{ala} \quad eco \quad \uline{aba} \quad ilha \quad \uline{ata} \quad ovo
\end{center}

\caixapausa{Se precisar, pare aqui. Respire fundo e continue quando quiser.}

\textbf{Atividade 2 --- Escrita do A --- 4 formas}

\instrucao{Escreva o A cinco vezes em cada forma. Sem pressa.}

\begin{center}
\textbf{Imprensa mai\'uscula:} \underline{\hspace{1cm}} \underline{\hspace{1cm}} \underline{\hspace{1cm}} \underline{\hspace{1cm}} \underline{\hspace{1cm}}\\[0.5em]
\textbf{Imprensa min\'uscula:} \underline{\hspace{1cm}} \underline{\hspace{1cm}} \underline{\hspace{1cm}} \underline{\hspace{1cm}} \underline{\hspace{1cm}}\\[0.5em]
\textbf{Cursiva mai\'uscula:} \underline{\hspace{1cm}} \underline{\hspace{1cm}} \underline{\hspace{1cm}} \underline{\hspace{1cm}} \underline{\hspace{1cm}}\\[0.5em]
\textbf{Cursiva min\'uscula:} \underline{\hspace{1cm}} \underline{\hspace{1cm}} \underline{\hspace{1cm}} \underline{\hspace{1cm}} \underline{\hspace{1cm}}
\end{center}

\textbf{Atividade 3 --- Jogo do Som}

\instrucao{Complete cada palavra com a letra A.}

\begin{center}
\underline{~~~}n\~a \quad \underline{~~~}la \quad \underline{~~~}sa \quad \underline{~~~}ba \quad \underline{~~~}ta \quad \underline{~~~}c\'a
\end{center}
Diga em voz alta: \textbf{an\~a, ala, asa, aba, ata, ac\'a}

\caixapista{Se esquecer o som do A, olhe a pista da boca da li\c{c}\~a{}o 1: boca bem aberta, ar forte.}
\end{consolidacao}

\plancheta{Sons da Letra A --- IPA}

\begin{sonsdaletra}
Os sons da letra A no portugu\^es{} brasileiro (IPA --- Alfabeto Fon\'etico Internacional):

\somipa[aberta]{A aberto}{a}{som aberto, forte (tônico)}{p\'a, casa, sapato}
\somipa[aberta]{A fechado}{6}{som fechado, fraco (\'a{}tono)}{cama, pata, bala}
\somipa[nasal]{A nasal}{a\~}{som anasalado (com til)}{an\~a, ma\c{c}\~a, rom\~a}
\somipa[nasal]{A nasal fraco}{6\~}{nasal em s\'ilaba fraca}{canta, tampa, banda}
\end{sonsdaletra}

\plancheta{Fam\'ilia Sil\'abica da Letra A}

\begin{familiasilabica}
A fam\'ilia sil\'abica da letra A:

\begin{center}
ana \quad an\~a \quad ala \quad asa \quad aba \quad ata \quad aia \quad ac\'a \quad ara \quad ava
\end{center}
\instrucao{Leia a fam\'ilia em voz alta, devagar. Depois cubra e tente lembrar.}
\end{familiasilabica}

\plancheta{Atividade de Som e S\'ilabas da Letra A}

\begin{somletra}
\textbf{1. Ca\c{c}a ao som.} Ou\c{c}a a palavra. Se tiver o som de A, marque o quadradinho:
\begin{center}
$\square$~ABELHA \quad $\square$~AVI\~AO \quad $\square$~CASA \quad $\square$~BOLA \quad $\square$~DEDO \quad $\square$~P\'E
\end{center}

\textbf{2. Tem ou n\~ao{} tem?} Rodeie SIM ou N\~ao:
\begin{center}
ABELHA --- (\textbf{SIM}) (\textbf{N\~AO}) \quad BOLA --- (\textbf{SIM}) (\textbf{N\~AO}) \\ AVI\~AO --- (\textbf{SIM}) (\textbf{N\~AO}) \quad DEDO --- (\textbf{SIM}) (\textbf{N\~AO})
\end{center}
\caixapausa{A boca abre para o som de A.}
\end{somletra}

\begin{somsilaba}
\textbf{3. Som das s\'ilabas.} Ou\c{c}a, bata palmas e conte as s\'ilabas:
\begin{center}
A.BE.LHA $\rightarrow$ \underline{\hspace{1cm}} s\'ilabas \quad SA.PA.TO $\rightarrow$ \underline{\hspace{1cm}} s\'ilabas
\end{center}

\textbf{4. Qual s\'ilaba tem o som?} Circule a s\'ilaba com o som de A:
\begin{center}
\textbf{A}.BE.LHA \quad SA.\textbf{PA}.TO
\end{center}
\caixapista{A vogal A aparece na s\'ilaba forte (\'a) e na fraca (a).}
\end{somsilaba}


\plancheta{Miss\~a{}o Motora --- Letra A}

\begin{motricidade}
\textbf{Movimento:} trace o caminho da abelha at\'e{} a flor, sem tirar o dedo do papel. Primeiro no ar, depois no papel.
\begin{center}
\begin{tikzpicture}
  \draw[very thick, dashed, color=primary] (0,0) to[out=20,in=160] (2,0.4) to[out=-20,in=200] (4,0) to[out=20,in=160] (6,0.4) to[out=-20,in=200] (8,0);
  \node[draw=accent, circle, minimum size=0.7cm] at (0,0) {\faBug};
  \node[draw=success, circle, minimum size=0.7cm] at (8.4,0) {\faPagelines};
\end{tikzpicture}
\end{center}
\caixapista{Use o dedo indicador. V\'a{} devagar. A forma de ondas prepara o tra\c{c}ado do A cursivo.}
\end{motricidade}

\plancheta{Ponte --- Da Letra A para a Pr\'oxima}

\begin{transicao}
Compare o A com o E:
\begin{center}
\textbf{p\'a} / \textbf{p\'e} \qquad \textbf{l\'a} / \textbf{l\^e} \qquad \textbf{m\'a} / \textbf{m\^e} \qquad \textbf{casa} / \textbf{cedo}
\end{center}
\instrucao{Ou\c{c}a a palavra que o professor disser. Se tiver som de A, aponte para cima. Se tiver som de E, aponte para baixo.}
\end{transicao}

\plancheta{Recompensa --- Miss\~a{}o 1}

\begin{recompensa}
\estrelas{3}\quad \ofensiva{1}\\[0.3em]\barraprogresso{1}{16}\\[0.5em]

\checklistdominio{reconhecer o A nas 4 formas, escrever e ler palavras com A}
\end{recompensa}


%% LICAO 2: LETRA E
%% ------------------------------------------------------------
\section{Licao 2: A Letra E}

% Rotina neuroinclusiva e badges (R201)
\rotinasimples
\nivelKbadge{K1 --- Letras e Grafias}
\rotabadge{1A --- Recomposicao Intensiva}


\letra{E}
% Quatro formas gráficas da letra (R201)
\quatroformasletra{E}{e}

\boq[sorriso]{E}{%
  Boca ENTREABERTA, maxilar levemente relaxado.\\
  Ar sai suave entre os dentes.\\
  Lingua levemente elevada no assoalho.\\
  Som: ``EEEEEEE'' (medio, suave).\\[0.5em]
  Palavras exemplo: elefante, escola, escada, espada, estrela.
}

\pistaarticulatoria{Som: \'{e}\'{e}\'{e} (m\'{e}dio, suave)}{Boca entreaberta, l'{a}bios levemente esticados}{M\~{a}o na frente da boca para sentir o ar sair suave}

\metadata{A-02}{Nivel 1 -- Pre-silabico}{EF01LP01}{D1}{EF01LP01}{Resolucao CNE/CEB 7/2010, art. 12}

\plancheta{Folha de Pratica 1.2 --- Letra E}
\begin{exercicio}[Folha de Pratica 1.2 --- Letra E]
\noindent\textbf{Nome:} \underline{\hspace{3.2cm}} \quad \textbf{Data:} \underline{\hspace{1.6cm}} \quad \textbf{Nota:} \underline{\hspace{0.9cm}}/10

\subsection*{PARTE 1: RECONHECIMENTO}

\begin{enumerate}
  \item Circule todas as letras ``E'':
  \begin{center}
  \fbox{E} \quad F \quad \fbox{E} \quad G \quad \fbox{E} \quad H \quad \fbox{E} \quad I \quad \fbox{E} \quad J
  \end{center}

  \item Sublinhe as palavras que comecam com ``E'':
  \begin{center}
  \uline{elefante} \quad gato \quad \uline{escada} \quad sapo \quad \uline{escola} \quad rato \quad \uline{espada}
  \end{center}

  \item Coloque $\checkmark$ nas linhas que contiverem a letra ``E'':
  \begin{center}
  \begin{tabular}{ll}
  $\checkmark$ \texttt{E S C O L A} & $\checkmark$ \texttt{C A D E I R A} \\
  $\checkmark$ \texttt{M E S A} & $\square$ \texttt{C A R T A} \\
  $\checkmark$ \texttt{R E D E} & $\square$ \texttt{L A T A}
  \end{tabular}
  \end{center}

  \item Identifique a posicao da letra ``E'' no nome ``ELEFANTE'':
  \begin{center}
  E --- L --- \fbox{E} --- F --- A --- N --- T --- \fbox{E}
  \end{center}

  \item Diga em voz alta tres palavras com ``E'' e anote:
  \begin{center}
  1. \underline{\hspace{4cm}} \quad 2. \underline{\hspace{4cm}} \quad 3. \underline{\hspace{4cm}}
  \end{center}
\end{enumerate}

\subsection*{PARTE 2: PRODUCAO GUIADA}

\begin{enumerate}[resume]
  \item Complete com ``E'':
  \begin{center}
  \underline{E}lefante \quad \underline{E}scola \quad \underline{E}spada \quad \underline{E}studo \quad \underline{E}strela
  \end{center}

  \item Escreva a letra ``E'' maiuscula cinco vezes:
  \begin{center}
  \underline{\hspace{1.5cm}} \quad \underline{\hspace{1.5cm}} \quad \underline{\hspace{1.5cm}} \quad \underline{\hspace{1.5cm}} \quad \underline{\hspace{1.5cm}}
  \end{center}

  \item Escreva a letra ``E'' minuscula cinco vezes:
  \begin{center}
  \underline{\hspace{1.5cm}} \quad \underline{\hspace{1.5cm}} \quad \underline{\hspace{1.5cm}} \quad \underline{\hspace{1.5cm}} \quad \underline{\hspace{1.5cm}}
  \end{center}

  \item Copie as palavras abaixo:
  \begin{center}
  elefante \quad escola \quad escada \quad espada \quad estrela
  \end{center}

  \item Escreva o nome de um animal que comeca com ``E'':
  \begin{center}
  \underline{\hspace{6cm}}
  \end{center}
\end{enumerate}

\subsection*{PARTE 3: INTEGRACAO}

\begin{enumerate}[resume]
  \item Leia em voz alta:
  \begin{center}
  $\rightarrow$ E --- ELEFANTE --- ESCOLA --- ESCADA
  \end{center}

  \item Nome de um animal que comeca com ``E'':
  \begin{center}
  \underline{\hspace{6cm}}
  \end{center}

  \item Desenhe algo com ``E'' e escreva o nome:
  \begin{center}
  \begin{tikzpicture}
    \draw[thick, dashed] (0,0) rectangle (5,3);
    \node[below] at (2.5,-0.3) {Nome: \underline{\hspace{4cm}}};
  \end{tikzpicture}
  \end{center}

  \item Circule a opcao correta:
  \begin{enumerate}[label=(\alph*)]
    \item A letra E tem som: ( ) aberto ( ) medio ( ) fechado
    \item A letra E aparece em: ( ) agua ( ) elefante ( ) uva
    \item A boca fica: ( ) aberta ( ) entreaberta ( ) arredondada
  \end{enumerate}

  \item Leia as palavras e assinale a que comeca com ``E'':
  \begin{center}
  ( ) mesa \quad ( ) elefante \quad ( ) tigre \quad ( ) sapo
  \end{center}
\end{enumerate}
\end{exercicio}

\begin{dica}
\textbf{Dica para o Professor:} A letra E tem som medio, entre o A e o I. Mostre a diferenca entre a boca aberta (A) e a boca entreaberta (E). Use espelho para que os alunos observem a propria posicao bucal. Peca que repitam ``A --- E --- A --- E'' para sentir a mudanca.
\end{dica}

%% ------------------------------------------------------------


% ============================================================
% CONSOLIDACAO DA LETRA E (R201.1) -- Missao 2/16
% ============================================================

\plancheta{Miss\~a{}o 2 --- Letra E}

\begin{missao}
\textbf{\faGamepad~Miss\~a{}o 2 de 16:} Consolide a letra E. Cuidado: o E tem quatro sons diferentes! Complete tudo para ganhar suas estrelas.
\end{missao}

\begin{consolidacao}
\textbf{Atividade 1 --- Ca\c{c}ada ao E}

\instrucao{Circule todas as letras E e sublinhe as palavras com E.}

\begin{center}
\begin{tabular}{cccccc}
\fbox{E} & I & \fbox{E} & A & \fbox{E} & O \\
\fbox{E} & U & I & \fbox{E} & A & \fbox{E}
\end{tabular}
\end{center}
\begin{center}
\uline{ela} \quad \uline{eco} \quad ilha \quad \uline{era} \quad uva \quad \uline{eba}
\end{center}

\caixapausa{Se precisar, pare aqui. Respire fundo e continue quando quiser.}

\textbf{Atividade 2 --- Escrita do E --- 4 formas}

\instrucao{Escreva o E cinco vezes em cada forma.}

\begin{center}
\textbf{Imprensa mai\'uscula:} \underline{\hspace{1cm}} \underline{\hspace{1cm}} \underline{\hspace{1cm}} \underline{\hspace{1cm}} \underline{\hspace{1cm}}\\[0.5em]
\textbf{Imprensa min\'uscula:} \underline{\hspace{1cm}} \underline{\hspace{1cm}} \underline{\hspace{1cm}} \underline{\hspace{1cm}} \underline{\hspace{1cm}}\\[0.5em]
\textbf{Cursiva mai\'uscula:} \underline{\hspace{1cm}} \underline{\hspace{1cm}} \underline{\hspace{1cm}} \underline{\hspace{1cm}} \underline{\hspace{1cm}}\\[0.5em]
\textbf{Cursiva min\'uscula:} \underline{\hspace{1cm}} \underline{\hspace{1cm}} \underline{\hspace{1cm}} \underline{\hspace{1cm}} \underline{\hspace{1cm}}
\end{center}

\textbf{Atividade 3 --- Jogo dos Dois Sons}

\instrucao{Ou\c{c}a a palavra e marque: E aberto (\'e) ou E fechado (\^e)?}

\begin{center}
\begin{tabular}{ll}
p\'e \quad $\rightarrow$ \underline{\hspace{2cm}} & m\^es \quad $\rightarrow$ \underline{\hspace{2cm}} \\
caf\'e \quad $\rightarrow$ \underline{\hspace{2cm}} & voc\^e \quad $\rightarrow$ \underline{\hspace{2cm}}
\end{tabular}
\end{center}
Complete com E: \underline{~~~}la \quad \underline{~~~}co \quad \underline{~~~}ra \quad \underline{~~~}ma

\caixapista{\'E = boca entreaberta, som claro. \^E = boca quase fechada, som escuro.}
\end{consolidacao}

\plancheta{Sons da Letra E --- IPA}

\begin{sonsdaletra}
Os sons da letra E no portugu\^es{} brasileiro (IPA --- Alfabeto Fon\'etico Internacional):

\somipa[sorriso]{E aberto}{E}{som aberto (\'e)}{p\'e, f\'e, caf\'e}
\somipa[sorriso]{E fechado}{e}{som fechado (\^e)}{m\^es, voc\^e, cedo}
\somipa[sorriso]{E fraco}{i}{som de I no fim da palavra (\'a{}tono)}{noite, tarde, doce}
\somipa[nasal]{E nasal}{e\~}{nasal (varia\c{c}\~a{}o regional)}{tempo, sempre}
\end{sonsdaletra}

\plancheta{Fam\'ilia Sil\'abica da Letra E}

\begin{familiasilabica}
A fam\'ilia sil\'abica da letra E:

\begin{center}
eca \quad eco \quad era \quad ela \quad eba \quad ema \quad ene
\end{center}
\instrucao{Leia a fam\'ilia em voz alta, devagar. Depois cubra e tente lembrar.}
\end{familiasilabica}

\plancheta{Atividade de Som e S\'ilabas da Letra E}

\begin{somletra}
\textbf{1. Ca\c{c}a ao som.} Ou\c{c}a a palavra. Se tiver o som de E, marque o quadradinho:
\begin{center}
$\square$~ESCOLA \quad $\square$~ELEFANTE \quad $\square$~CEDO \quad $\square$~BOLA \quad $\square$~LUVA \quad $\square$~P\'E
\end{center}

\textbf{2. Tem ou n\~ao{} tem?} Rodeie SIM ou N\~ao:
\begin{center}
ESCOLA --- (\textbf{SIM}) (\textbf{N\~AO}) \quad BOLA --- (\textbf{SIM}) (\textbf{N\~AO}) \\ ELEFANTE --- (\textbf{SIM}) (\textbf{N\~AO}) \quad LUVA --- (\textbf{SIM}) (\textbf{N\~AO})
\end{center}
\caixapausa{A boca faz sorriso para o som de E.}
\end{somletra}

\begin{somsilaba}
\textbf{3. Som das s\'ilabas.} Ou\c{c}a, bata palmas e conte as s\'ilabas:
\begin{center}
E.SCO.LA $\rightarrow$ \underline{\hspace{1cm}} s\'ilabas \quad CE.DO $\rightarrow$ \underline{\hspace{1cm}} s\'ilabas
\end{center}

\textbf{4. Qual s\'ilaba tem o som?} Circule a s\'ilaba com o som de E:
\begin{center}
\textbf{E}.SCO.LA \quad \textbf{CE}.DO
\end{center}
\caixapista{A vogal E aparece na s\'ilaba forte (\^e) e na fraca (e).}
\end{somsilaba}


\plancheta{Miss\~a{}o Motora --- Letra E}

\begin{motricidade}
\textbf{Movimento:} ligue os pontos e forme o E. Depois trace de novo, sem os pontos.
\begin{center}
\begin{tikzpicture}
  \foreach \x in {0,0.8,1.6,2.4,3.2,4.0,4.8} {
    \fill[primary] (\x,1.2) circle (0.08);
    \fill[primary] (\x,0) circle (0.08);
    \fill[primary] (\x,-1.2) circle (0.08);
  }
  \fill[primary] (0,1.2) circle (0.08);
  \fill[primary] (0,0) circle (0.08);
  \fill[primary] (0,-1.2) circle (0.08);
\end{tikzpicture}
\end{center}
\caixapista{Tr\^es linhas horizontais + uma linha vertical na esquerda.}
\end{motricidade}

\plancheta{Ponte --- Da Letra E para a Pr\'oxima}

\begin{transicao}
Compare o E com o I:
\begin{center}
\textbf{p\'e} / \textbf{pilha} \qquad \textbf{m\^es} / \textbf{missa} \qquad \textbf{cela} / \textbf{sino}
\end{center}
\instrucao{Fale \'e e iii devagar, sentindo a boca mudar. \'E = boca mais aberta. I = sorriso esticado.}
\end{transicao}

\plancheta{Recompensa --- Miss\~a{}o 2}

\begin{recompensa}
\estrelas{3}\quad \ofensiva{2}\\[0.3em]\barraprogresso{2}{16}\\[0.5em]

\checklistdominio{reconhecer e diferenciar os sons do E (aberto, fechado, fraco)}
\end{recompensa}


%% LICAO 3: LETRA I
%% ------------------------------------------------------------
\section{Licao 3: A Letra I}

% Rotina neuroinclusiva e badges (R201)
\rotinasimples
\nivelKbadge{K1 --- Letras e Grafias}
\rotabadge{1A --- Recomposicao Intensiva}


\letra{I}
% Quatro formas gráficas da letra (R201)
\quatroformasletra{I}{i}

\boq[sorriso]{I}{%
  Boca ENTREABERTA, cantos da boca para TRAS.\\
  Ar sai estreito entre os dentes.\\
  Lingua elevada no palato duro.\\
  Som: ``IIIIIII'' (agudo, estreito).\\[0.5em]
  Palavras exemplo: igreja, iguana, icaro, isca, imel.
}

\pistaarticulatoria{Som: iii (agudo, fino)}{L'{a}bios esticados em sorriso}{M\~{a}o no canto da boca para sentir o sorriso}

\metadata{A-03}{Nivel 1 -- Pre-silabico}{EF01LP01}{D1}{EF01LP01}{Resolucao CNE/CEB 7/2010, art. 12}

\plancheta{Folha de Pratica 1.3 --- Letra I}
\begin{exercicio}[Folha de Pratica 1.3 --- Letra I]
\noindent\textbf{Nome:} \underline{\hspace{3.2cm}} \quad \textbf{Data:} \underline{\hspace{1.6cm}} \quad \textbf{Nota:} \underline{\hspace{0.9cm}}/10

\subsection*{PARTE 1: RECONHECIMENTO}

\begin{enumerate}
  \item Circule todas as letras ``I'':
  \begin{center}
  \fbox{I} \quad H \quad \fbox{I} \quad G \quad \fbox{I} \quad F \quad \fbox{I} \quad E \quad \fbox{I} \quad D
  \end{center}

  \item Sublinhe as palavras que comecam com ``I'':
  \begin{center}
  \uline{igreja} \quad bolo \quad \uline{iguana} \quad sopa \quad \uline{icaro} \quad dedo \quad \uline{isca}
  \end{center}

  \item Coloque $\checkmark$ nas linhas com a letra ``I'':
  \begin{center}
  \begin{tabular}{ll}
  $\checkmark$ \texttt{I G R E J A} & $\square$ \texttt{C A S A} \\
  $\checkmark$ \texttt{A L I M E N T O} & $\checkmark$ \texttt{M A R I N H O} \\
  $\square$ \texttt{L A T A} & $\checkmark$ \texttt{P A T I N H O}
  \end{tabular}
  \end{center}

  \item Identifique a posicao da letra ``I'' em ``IGREJA'':
  \begin{center}
  \fbox{I} --- G --- R --- E --- J --- A $\rightarrow$ posicao \underline{1}
  \end{center}

  \item Diga em voz alta tres palavras com ``I'' e anote:
  \begin{center}
  1. \underline{\hspace{4cm}} \quad 2. \underline{\hspace{4cm}} \quad 3. \underline{\hspace{4cm}}
  \end{center}
\end{enumerate}

\subsection*{PARTE 2: PRODUCAO GUIADA}

\begin{enumerate}[resume]
  \item Complete com ``I'':
  \begin{center}
  \underline{I}greja \quad \underline{I}guana \quad \underline{I}caro \quad \underline{I}sca \quad \underline{I}mel
  \end{center}

  \item Escreva a letra ``I'' maiuscula cinco vezes:
  \begin{center}
  \underline{\hspace{1.5cm}} \quad \underline{\hspace{1.5cm}} \quad \underline{\hspace{1.5cm}} \quad \underline{\hspace{1.5cm}} \quad \underline{\hspace{1.5cm}}
  \end{center}

  \item Escreva a letra ``I'' minuscula cinco vezes:
  \begin{center}
  \underline{\hspace{1.5cm}} \quad \underline{\hspace{1.5cm}} \quad \underline{\hspace{1.5cm}} \quad \underline{\hspace{1.5cm}} \quad \underline{\hspace{1.5cm}}
  \end{center}

  \item Copie as palavras:
  \begin{center}
  igreja \quad iguana \quad icaro \quad isca \quad imel
  \end{center}

  \item Escreva o nome de um lugar que comeca com ``I'':
  \begin{center}
  \underline{\hspace{6cm}}
  \end{center}
\end{enumerate}

\subsection*{PARTE 3: INTEGRACAO}

\begin{enumerate}[resume]
  \item Leia em voz alta:
  \begin{center}
  $\rightarrow$ I --- IGREJA --- IGUANA --- ICARO
  \end{center}

  \item Nome de um lugar que comeca com ``I'':
  \begin{center}
  \underline{\hspace{6cm}}
  \end{center}

  \item Desenhe algo com ``I'' e escreva o nome:
  \begin{center}
  \begin{tikzpicture}
    \draw[thick, dashed] (0,0) rectangle (5,3);
    \node[below] at (2.5,-0.3) {Nome: \underline{\hspace{4cm}}};
  \end{tikzpicture}
  \end{center}

  \item Circule a opcao correta:
  \begin{enumerate}[label=(\alph*)]
    \item A letra I tem som: ( ) grave ( ) agudo ( ) nasal
    \item A letra I aparece em: ( ) agua ( ) elefante ( ) igreja
    \item A boca fica: ( ) aberta ( ) entreaberta ( ) arredondada
  \end{enumerate}

  \item Leia e assinale a palavra com ``I'':
  \begin{center}
  ( ) mesa \quad ( ) igreja \quad ( ) tinta \quad ( ) lata
  \end{center}
\end{enumerate}
\end{exercicio}

\begin{dica}
\textbf{Dica para o Professor:} A letra I e a vogal mais aguda. Peca aos alunos que sorrirem levemente (cantos da boca para tras) e repitam ``IIIIII''. Compare com ``AAAAA'' e ``EEEEE'' para que percebam as diferencas. Use o espelho novamente para auto-observacao.
\end{dica}

%% ------------------------------------------------------------


% ============================================================
% CONSOLIDACAO DA LETRA I (R201.1) -- Missao 3/16
% ============================================================

\plancheta{Miss\~a{}o 3 --- Letra I}

\begin{missao}
\textbf{\faGamepad~Miss\~a{}o 3 de 16:} Consolide a letra I. Descubra quando o I vira som de J (como em pai).
\end{missao}

\begin{consolidacao}
\textbf{Atividade 1 --- Ca\c{c}ada ao I}

\instrucao{Circule o I e sublinhe as palavras com I.}

\begin{center}
\begin{tabular}{cccccc}
\fbox{I} & O & \fbox{I} & E & \fbox{I} & A \\
U & \fbox{I} & O & \fbox{I} & E & \fbox{I}
\end{tabular}
\end{center}
\begin{center}
\uline{ilha} \quad \uline{im\~a} \quad ovo \quad \uline{iva} \quad eco \quad \uline{iana}
\end{center}

\caixapausa{Se precisar, pare aqui. Respire fundo e continue quando quiser.}

\textbf{Atividade 2 --- Escrita do I --- 4 formas}

\instrucao{Escreva o I cinco vezes em cada forma.}

\begin{center}
\textbf{Imprensa mai\'uscula:} \underline{\hspace{1cm}} \underline{\hspace{1cm}} \underline{\hspace{1cm}} \underline{\hspace{1cm}} \underline{\hspace{1cm}}\\[0.5em]
\textbf{Imprensa min\'uscula:} \underline{\hspace{1cm}} \underline{\hspace{1cm}} \underline{\hspace{1cm}} \underline{\hspace{1cm}} \underline{\hspace{1cm}}\\[0.5em]
\textbf{Cursiva mai\'uscula:} \underline{\hspace{1cm}} \underline{\hspace{1cm}} \underline{\hspace{1cm}} \underline{\hspace{1cm}} \underline{\hspace{1cm}}\\[0.5em]
\textbf{Cursiva min\'uscula:} \underline{\hspace{1cm}} \underline{\hspace{1cm}} \underline{\hspace{1cm}} \underline{\hspace{1cm}} \underline{\hspace{1cm}}
\end{center}

\textbf{Atividade 3 --- Jogo do Som Final}

\instrucao{Complete com I e leia. Sinta o I virando J no fim.}

\begin{center}
pa\underline{~~~} \quad va\underline{~~~} \quad ca\underline{~~~} \quad s\underline{~~~}m \quad f\underline{~~~}m \quad m\underline{~~~}m
\end{center}
Diga em voz alta: \textbf{pai, vai, cai, sim, fim, mim}

\caixapista{pai = pa + i r\'apido. O i no fim vira um J leve.}
\end{consolidacao}

\plancheta{Sons da Letra I --- IPA}

\begin{sonsdaletra}
Os sons da letra I no portugu\^es{} brasileiro (IPA --- Alfabeto Fon\'etico Internacional):

\somipa[sorriso]{I forte}{i}{som agudo e claro}{ilha, pipi, bico}
\somipa[nasal]{I nasal}{i\~}{anasalado}{sim, fim, mim}
\somipa[sorriso]{I semivogal}{j}{som r\'apido de J no fim da s\'ilaba}{pai, vai, cai}
\end{sonsdaletra}

\plancheta{Fam\'ilia Sil\'abica da Letra I}

\begin{familiasilabica}
A fam\'ilia sil\'abica da letra I:

\begin{center}
iai\'a \quad iana \quad im\~a \quad ilha \quad iva
\end{center}
\instrucao{Leia a fam\'ilia em voz alta, devagar. Depois cubra e tente lembrar.}
\end{familiasilabica}

\plancheta{Atividade de Som e S\'ilabas da Letra I}

\begin{somletra}
\textbf{1. Ca\c{c}a ao som.} Ou\c{c}a a palavra. Se tiver o som de I, marque o quadradinho:
\begin{center}
$\square$~IGREJA \quad $\square$~\'INDIO \quad $\square$~AMIGO \quad $\square$~BOLA \quad $\square$~P\~AO \quad $\square$~M\~AE
\end{center}

\textbf{2. Tem ou n\~ao{} tem?} Rodeie SIM ou N\~ao:
\begin{center}
IGREJA --- (\textbf{SIM}) (\textbf{N\~AO}) \quad BOLA --- (\textbf{SIM}) (\textbf{N\~AO}) \\ \'INDIO --- (\textbf{SIM}) (\textbf{N\~AO}) \quad P\~AO --- (\textbf{SIM}) (\textbf{N\~AO})
\end{center}
\caixapausa{A boca faz sorriso para o som de I.}
\end{somletra}

\begin{somsilaba}
\textbf{3. Som das s\'ilabas.} Ou\c{c}a, bata palmas e conte as s\'ilabas:
\begin{center}
I.GRE.JA $\rightarrow$ \underline{\hspace{1cm}} s\'ilabas \quad A.MI.GO $\rightarrow$ \underline{\hspace{1cm}} s\'ilabas
\end{center}

\textbf{4. Qual s\'ilaba tem o som?} Circule a s\'ilaba com o som de I:
\begin{center}
\textbf{I}.GRE.JA \quad A.\textbf{MI}.GO
\end{center}
\caixapista{A vogal I aparece na s\'ilaba forte (\'i) e na fraca (i).}
\end{somsilaba}


\plancheta{Miss\~a{}o Motora --- Letra I}

\begin{motricidade}
\textbf{Movimento:} trace as linhas verticais de cima para baixo. Elas preparam o I.
\begin{center}
\begin{tikzpicture}
  \foreach \x in {0,1.2,2.4,3.6,4.8,6,7.2} {
    \draw[very thick, dashed, color=primary] (\x,1) -- (\x,-1);
    \draw[->, color=accent] (\x+0.15,0.8) -- (\x+0.15,0.2);
  }
\end{tikzpicture}
\end{center}
\caixapista{Sempre de cima para baixo, sem tirar o l\'apis.}
\end{motricidade}

\plancheta{Ponte --- Da Letra I para a Pr\'oxima}

\begin{transicao}
Compare o I com o O:
\begin{center}
\textbf{fim} / \textbf{fome} \qquad \textbf{ilha} / \textbf{olho} \qquad \textbf{pia} / \textbf{po\'a}
\end{center}
\instrucao{Fale iii e \'o\'o. O I estica os l\'abios. O O arredonda.}
\end{transicao}

\plancheta{Recompensa --- Miss\~a{}o 3}

\begin{recompensa}
\estrelas{3}\quad \ofensiva{3}\\[0.3em]\barraprogresso{3}{16}\\[0.5em]

\checklistdominio{reconhecer o I forte, nasal e semivogal; ler pai, vai, cai}
\end{recompensa}


%% LICAO 4: LETRA O
%% ------------------------------------------------------------
\section{Licao 4: A Letra O}

% Rotina neuroinclusiva e badges (R201)
\rotinasimples
\nivelKbadge{K1 --- Letras e Grafias}
\rotabadge{1A --- Recomposicao Intensiva}


\letra{O}
% Quatro formas gráficas da letra (R201)
\quatroformasletra{O}{o}

\boq[bico]{O}{%
  Boca ARREDONDADA, labios projetados.\\
  Ar sai em forma de circulo.\\
  Lingua no assoalho bucal.\\
  Som: ``OOOOOOO'' (redondo, aberto).\\[0.5em]
  Palavras exemplo: ovo, olho, onca, orelha, orgao.
}

\pistaarticulatoria{Som: \'{o}\'{o}\'{o} (grave, redondo)}{L'{a}bios arredondados}{M\~{a}o na frente da boca para sentir o ar em c'{i}rculo}

\metadata{A-04}{Nivel 1 -- Pre-silabico}{EF01LP01}{D1}{EF01LP01}{Resolucao CNE/CEB 7/2010, art. 12}

\plancheta{Folha de Pratica 1.4 --- Letra O}
\begin{exercicio}[Folha de Pratica 1.4 --- Letra O]
\noindent\textbf{Nome:} \underline{\hspace{3.2cm}} \quad \textbf{Data:} \underline{\hspace{1.6cm}} \quad \textbf{Nota:} \underline{\hspace{0.9cm}}/10

\subsection*{PARTE 1: RECONHECIMENTO}

\begin{enumerate}
  \item Circule todas as letras ``O'':
  \begin{center}
  \fbox{O} \quad P \quad \fbox{O} \quad Q \quad \fbox{O} \quad R \quad \fbox{O} \quad S \quad \fbox{O} \quad T
  \end{center}

  \item Sublinhe as palavras que comecam com ``O'':
  \begin{center}
  \uline{ovo} \quad gato \quad \uline{olho} \quad sapo \quad \uline{onca} \quad rato \quad \uline{orelha}
  \end{center}

  \item Coloque $\checkmark$ nas linhas com ``O'':
  \begin{center}
  \begin{tabular}{ll}
  $\checkmark$ \texttt{O L H O} & $\checkmark$ \texttt{B O L A} \\
  $\square$ \texttt{C A S A} & $\checkmark$ \texttt{O N C A} \\
  $\checkmark$ \texttt{M E S A} & $\square$ \texttt{L I R A}
  \end{tabular}
  \end{center}

  \item Identifique a posicao da letra ``O'' em ``OLHO'':
  \begin{center}
  \fbox{O} --- L --- H --- \fbox{O} $\rightarrow$ posicoes \underline{1} e \underline{4}
  \end{center}

  \item Diga em voz alta tres palavras com ``O'' e anote:
  \begin{center}
  1. \underline{\hspace{4cm}} \quad 2. \underline{\hspace{4cm}} \quad 3. \underline{\hspace{4cm}}
  \end{center}
\end{enumerate}

\subsection*{PARTE 2: PRODUCAO GUIADA}

\begin{enumerate}[resume]
  \item Complete com ``O'':
  \begin{center}
  \underline{O}vo \quad \underline{O}lho \quad \underline{O}nca \quad \underline{O}relha \quad \underline{O}rgao
  \end{center}

  \item Escreva a letra ``O'' maiuscula cinco vezes:
  \begin{center}
  \underline{\hspace{1.5cm}} \quad \underline{\hspace{1.5cm}} \quad \underline{\hspace{1.5cm}} \quad \underline{\hspace{1.5cm}} \quad \underline{\hspace{1.5cm}}
  \end{center}

  \item Escreva a letra ``O'' minuscula cinco vezes:
  \begin{center}
  \underline{\hspace{1.5cm}} \quad \underline{\hspace{1.5cm}} \quad \underline{\hspace{1.5cm}} \quad \underline{\hspace{1.5cm}} \quad \underline{\hspace{1.5cm}}
  \end{center}

  \item Copie as palavras:
  \begin{center}
  ovo \quad olho \quad onca \quad orelha \quad orgao
  \end{center}

  \item Escreva o nome de um objeto que comeca com ``O'':
  \begin{center}
  \underline{\hspace{6cm}}
  \end{center}
\end{enumerate}

\subsection*{PARTE 3: INTEGRACAO}

\begin{enumerate}[resume]
  \item Leia em voz alta:
  \begin{center}
  $\rightarrow$ O --- OVO --- OLHO --- ONCA
  \end{center}

  \item Nome de um objeto que comeca com ``O'':
  \begin{center}
  \underline{\hspace{6cm}}
  \end{center}

  \item Desenhe algo com ``O'' e escreva o nome:
  \begin{center}
  \begin{tikzpicture}
    \draw[thick, dashed] (0,0) rectangle (5,3);
    \node[below] at (2.5,-0.3) {Nome: \underline{\hspace{4cm}}};
  \end{tikzpicture}
  \end{center}

  \item Circule a opcao correta:
  \begin{enumerate}[label=(\alph*)]
    \item A letra O tem som: ( ) agudo ( ) medio ( ) grave e redondo
    \item A letra O aparece em: ( ) igreja ( ) ovo ( ) sapo
    \item A boca fica: ( ) aberta ( ) entreaberta ( ) arredondada
  \end{enumerate}

  \item Leia e assinale a palavra com ``O'':
  \begin{center}
  ( ) mesa \quad ( ) onca \quad ( ) tinta \quad ( ) igreja
  \end{center}
\end{enumerate}
\end{exercicio}

\begin{dica}
\textbf{Dica para o Professor:} A letra O e a vogal mais grave e redonda. Peca que os alunos formem os labios como se fossem beijar e repitam ``OOOOO''. Compare com ``IIIII'' para mostrar a oposicao (agudo vs. grave). Use exercicios de rotacao: ``I --- E --- A --- O --- U'' para percorrer todas as vogais.
\end{dica}

%% ------------------------------------------------------------


% ============================================================
% CONSOLIDACAO DA LETRA O (R201.1) -- Missao 4/16
% ============================================================

\plancheta{Miss\~a{}o 4 --- Letra O}

\begin{missao}
\textbf{\faGamepad~Miss\~a{}o 4 de 16:} Consolide a letra O. O O muda de som: av\'o e av\^o n\~a{}o s\~a{}o a mesma coisa!
\end{missao}

\begin{consolidacao}
\textbf{Atividade 1 --- Ca\c{c}ada ao O}

\instrucao{Circule o O e sublinhe as palavras com O.}

\begin{center}
\begin{tabular}{cccccc}
\fbox{O} & A & \fbox{O} & I & \fbox{O} & E \\
U & \fbox{O} & A & \fbox{O} & I & \fbox{O}
\end{tabular}
\end{center}
\begin{center}
\uline{ovo} \quad \uline{osso} \quad asa \quad \uline{olho} \quad im\~a \quad \uline{oca}
\end{center}

\caixapausa{Se precisar, pare aqui. Respire fundo e continue quando quiser.}

\textbf{Atividade 2 --- Escrita do O --- 4 formas}

\instrucao{Escreva o O cinco vezes em cada forma.}

\begin{center}
\textbf{Imprensa mai\'uscula:} \underline{\hspace{1cm}} \underline{\hspace{1cm}} \underline{\hspace{1cm}} \underline{\hspace{1cm}} \underline{\hspace{1cm}}\\[0.5em]
\textbf{Imprensa min\'uscula:} \underline{\hspace{1cm}} \underline{\hspace{1cm}} \underline{\hspace{1cm}} \underline{\hspace{1cm}} \underline{\hspace{1cm}}\\[0.5em]
\textbf{Cursiva mai\'uscula:} \underline{\hspace{1cm}} \underline{\hspace{1cm}} \underline{\hspace{1cm}} \underline{\hspace{1cm}} \underline{\hspace{1cm}}\\[0.5em]
\textbf{Cursiva min\'uscula:} \underline{\hspace{1cm}} \underline{\hspace{1cm}} \underline{\hspace{1cm}} \underline{\hspace{1cm}} \underline{\hspace{1cm}}
\end{center}

\textbf{Atividade 3 --- Jogo dos Dois O}

\instrucao{Ou\c{c}a e marque: O aberto (\'o) ou O fechado (\^o)?}

\begin{center}
\begin{tabular}{ll}
av\'o \quad $\rightarrow$ \underline{\hspace{2cm}} & av\^o \quad $\rightarrow$ \underline{\hspace{2cm}} \\
p\'o \quad $\rightarrow$ \underline{\hspace{2cm}} & p\^or \quad $\rightarrow$ \underline{\hspace{2cm}} \\
n\'o \quad $\rightarrow$ \underline{\hspace{2cm}} & bolo \quad $\rightarrow$ \underline{\hspace{2cm}}
\end{tabular}
\end{center}

\caixapista{av\'o = boca mais aberta. av\^o = boca mais fechada, som escuro.}
\end{consolidacao}

\plancheta{Sons da Letra O --- IPA}

\begin{sonsdaletra}
Os sons da letra O no portugu\^es{} brasileiro (IPA --- Alfabeto Fon\'etico Internacional):

\somipa[bico]{O aberto}{O}{som aberto (\'o)}{p\'o, n\'o, s\'o, av\'o}
\somipa[bico]{O fechado}{o}{som fechado (\^o)}{av\^o, p\^or, bolo}
\somipa[bico]{O fraco}{u}{som de U no fim da palavra (\'a{}tono)}{gato, bolo, pato}
\somipa[nasal]{O nasal}{o\~}{anasalado}{bom, som, tom}
\somipa[bico]{O semivogal}{w}{som r\'apido de U no fim}{pau, c\'eu, r\'eu}
\end{sonsdaletra}

\plancheta{Fam\'ilia Sil\'abica da Letra O}

\begin{familiasilabica}
A fam\'ilia sil\'abica da letra O:

\begin{center}
ovo \quad osso \quad olho \quad oca \quad onda \quad on\c{c}a \quad ouro
\end{center}
\instrucao{Leia a fam\'ilia em voz alta, devagar. Depois cubra e tente lembrar.}
\end{familiasilabica}

\plancheta{Atividade de Som e S\'ilabas da Letra O}

\begin{somletra}
\textbf{1. Ca\c{c}a ao som.} Ou\c{c}a a palavra. Se tiver o som de O, marque o quadradinho:
\begin{center}
$\square$~OVO \quad $\square$~OSSO \quad $\square$~BOLA \quad $\square$~P\'E \quad $\square$~DEDO \quad $\square$~LUA
\end{center}

\textbf{2. Tem ou n\~ao{} tem?} Rodeie SIM ou N\~ao:
\begin{center}
OVO --- (\textbf{SIM}) (\textbf{N\~AO}) \quad P\'E --- (\textbf{SIM}) (\textbf{N\~AO}) \\ OSSO --- (\textbf{SIM}) (\textbf{N\~AO}) \quad DEDO --- (\textbf{SIM}) (\textbf{N\~AO})
\end{center}
\caixapausa{A boca faz bico para o som de O.}
\end{somletra}

\begin{somsilaba}
\textbf{3. Som das s\'ilabas.} Ou\c{c}a, bata palmas e conte as s\'ilabas:
\begin{center}
O.VO $\rightarrow$ \underline{\hspace{1cm}} s\'ilabas \quad BO.LA $\rightarrow$ \underline{\hspace{1cm}} s\'ilabas
\end{center}

\textbf{4. Qual s\'ilaba tem o som?} Circule a s\'ilaba com o som de O:
\begin{center}
\textbf{O}.VO \quad \textbf{BO}.LA
\end{center}
\caixapista{A vogal O aparece na s\'ilaba forte (\'o) e na fraca (o).}
\end{somsilaba}


\plancheta{Miss\~a{}o Motora --- Letra O}

\begin{motricidade}
\textbf{Movimento:} trace os c\'irculos de dentro para fora. Comece no ponto.
\begin{center}
\begin{tikzpicture}
  \draw[very thick, dashed, color=primary] (0,0) circle (0.5);
  \draw[very thick, dashed, color=secondary] (0,0) circle (1.0);
  \draw[very thick, dashed, color=accent] (0,0) circle (1.5);
  \fill[primary] (0.5,0) circle (0.1);
\end{tikzpicture}
\end{center}
\caixapista{O c\'irculo come\c{c}a no alto, gira para a esquerda e volta ao ponto.}
\end{motricidade}

\plancheta{Ponte --- Da Letra O para a Pr\'oxima}

\begin{transicao}
Compare o O com o U:
\begin{center}
\textbf{ovo} / \textbf{uva} \qquad \textbf{osso} / \textbf{urso} \qquad \textbf{onda} / \textbf{uma}
\end{center}
\instrucao{Fale \'o\'o e uuu. O O abre um pouco mais. O U fecha em bico.}
\end{transicao}

\plancheta{Recompensa --- Miss\~a{}o 4}

\begin{recompensa}
\estrelas{3}\quad \ofensiva{4}\\[0.3em]\barraprogresso{4}{16}\\[0.5em]

\checklistdominio{diferenciar os cinco sons do O e ler av\'o/av\^o corretamente}
\end{recompensa}


%% LICAO 5: LETRA U
%% ------------------------------------------------------------
\section{Licao 5: A Letra U}

% Rotina neuroinclusiva e badges (R201)
\rotinasimples
\nivelKbadge{K1 --- Letras e Grafias}
\rotabadge{1A --- Recomposicao Intensiva}


\letra{U}
% Quatro formas gráficas da letra (R201)
\quatroformasletra{U}{u}

\boq[bico]{U}{%
  Boca BEM ARREDONDADA e pequena.\\
  Ar sai focado em jato estreito.\\
  Lingua no assoalho, labios projetados.\\
  Som: ``UUUUUUU'' (grave, fechado).\\[0.5em]
  Palavras exemplo: uva, urso, urubu, umbigo.
}

\pistaarticulatoria{Som: uuu (fechado, grave)}{L'{a}bios em bico}{M\~{a}o na frente da boca para sentir o ar focado}

\metadata{A-05}{Nivel 1 -- Pre-silabico}{EF01LP01}{D1}{EF01LP01}{Resolucao CNE/CEB 7/2010, art. 12}

\plancheta{Folha de Pratica 1.5 --- Letra U}
\begin{exercicio}[Folha de Pratica 1.5 --- Letra U]
\noindent\textbf{Nome:} \underline{\hspace{3.2cm}} \quad \textbf{Data:} \underline{\hspace{1.6cm}} \quad \textbf{Nota:} \underline{\hspace{0.9cm}}/10

\subsection*{PARTE 1: RECONHECIMENTO}

\begin{enumerate}
  \item Circule todas as letras ``U'':
  \begin{center}
  \fbox{U} \quad V \quad \fbox{U} \quad W \quad \fbox{U} \quad X \quad \fbox{U} \quad Y \quad \fbox{U} \quad Z
  \end{center}

  \item Sublinhe as palavras que comecam com ``U'':
  \begin{center}
  \uline{uva} \quad gato \quad \uline{urso} \quad sapo \quad \uline{urubu} \quad rato \quad \uline{umbigo}
  \end{center}

  \item Coloque $\checkmark$ nas linhas com ``U'':
  \begin{center}
  \begin{tabular}{ll}
  $\checkmark$ \texttt{U V A} & $\checkmark$ \texttt{U R S O} \\
  $\square$ \texttt{C A S A} & $\checkmark$ \texttt{U R U B U} \\
  $\checkmark$ \texttt{B O L A} & $\square$ \texttt{L I R A}
  \end{tabular}
  \end{center}

  \item Identifique a posicao da letra ``U'' em ``URUBU'':
  \begin{center}
  \fbox{U} --- R --- \fbox{U} --- B --- \fbox{U} $\rightarrow$ posicoes \underline{1}, \underline{3} e \underline{5}
  \end{center}

  \item Diga em voz alta tres palavras com ``U'' e anote:
  \begin{center}
  1. \underline{\hspace{4cm}} \quad 2. \underline{\hspace{4cm}} \quad 3. \underline{\hspace{4cm}}
  \end{center}
\end{enumerate}

\subsection*{PARTE 2: PRODUCAO GUIADA}

\begin{enumerate}[resume]
  \item Complete com ``U'':
  \begin{center}
  \underline{U}va \quad \underline{U}rso \quad \underline{U}rubu \quad \underline{U}mbigo \quad \underline{U}va-do-mato
  \end{center}

  \item Escreva a letra ``U'' maiuscula cinco vezes:
  \begin{center}
  \underline{\hspace{1.5cm}} \quad \underline{\hspace{1.5cm}} \quad \underline{\hspace{1.5cm}} \quad \underline{\hspace{1.5cm}} \quad \underline{\hspace{1.5cm}}
  \end{center}

  \item Escreva a letra ``U'' minuscula cinco vezes:
  \begin{center}
  \underline{\hspace{1.5cm}} \quad \underline{\hspace{1.5cm}} \quad \underline{\hspace{1.5cm}} \quad \underline{\hspace{1.5cm}} \quad \underline{\hspace{1.5cm}}
  \end{center}

  \item Copie as palavras:
  \begin{center}
  uva \quad urso \quad urubu \quad umbigo
  \end{center}

  \item Escreva o nome de uma fruta que comeca com ``U'':
  \begin{center}
  \underline{\hspace{6cm}}
  \end{center}
\end{enumerate}

\subsection*{PARTE 3: INTEGRACAO}

\begin{enumerate}[resume]
  \item Leia em voz alta:
  \begin{center}
  $\rightarrow$ U --- UVA --- URSO --- URUBU
  \end{center}

  \item Nome de uma fruta que comeca com ``U'':
  \begin{center}
  \underline{\hspace{6cm}}
  \end{center}

  \item Desenhe algo com ``U'' e escreva o nome:
  \begin{center}
  \begin{tikzpicture}
    \draw[thick, dashed] (0,0) rectangle (5,3);
    \node[below] at (2.5,-0.3) {Nome: \underline{\hspace{4cm}}};
  \end{tikzpicture}
  \end{center}

  \item Circule a opcao correta:
  \begin{enumerate}[label=(\alph*)]
    \item A letra U tem som: ( ) agudo ( ) medio ( ) o mais grave
    \item A letra U aparece em: ( ) igreja ( ) elefante ( ) uva
    \item A boca fica: ( ) aberta ( ) entreaberta ( ) bem arredondada e pequena
  \end{enumerate}

  \item Leia e assinale a palavra com ``U'':
  \begin{center}
  ( ) mesa \quad ( ) uva \quad ( ) tinta \quad ( ) igreja
  \end{center}
\end{enumerate}
\end{exercicio}

\begin{dica}
\textbf{Dica para o Professor:} A letra U e a vogal mais grave e com a boca mais fechada. Peca que os alunos facam o bico e repitam ``UUUUU''. Compare com ``AAAAA'' (boca aberta) e ``IIIII'' (boca estreita) para mostrar todo o espectro das vogais.
\end{dica}

%% ------------------------------------------------------------


% ============================================================
% CONSOLIDACAO DA LETRA U (R201.1) -- Missao 5/16
% ============================================================

\plancheta{Miss\~a{}o 5 --- Letra U}

\begin{missao}
\textbf{\faGamepad~Miss\~a{}o 5 de 16:} Consolide a letra U, a \'ultima vogal. Ao completar esta miss\~a{}o, voc\^e termina o n\'ivel K1!
\end{missao}

\begin{consolidacao}
\textbf{Atividade 1 --- Ca\c{c}ada ao U}

\instrucao{Circule o U e sublinhe as palavras com U.}

\begin{center}
\begin{tabular}{cccccc}
\fbox{U} & O & \fbox{U} & A & \fbox{U} & E \\
I & \fbox{U} & O & \fbox{U} & A & \fbox{U}
\end{tabular}
\end{center}
\begin{center}
\uline{uva} \quad \uline{unha} \quad ovo \quad \uline{urso} \quad asa \quad \uline{urubu}
\end{center}

\caixapausa{Se precisar, pare aqui. Respire fundo e continue quando quiser.}

\textbf{Atividade 2 --- Escrita do U --- 4 formas}

\instrucao{Escreva o U cinco vezes em cada forma.}

\begin{center}
\textbf{Imprensa mai\'uscula:} \underline{\hspace{1cm}} \underline{\hspace{1cm}} \underline{\hspace{1cm}} \underline{\hspace{1cm}} \underline{\hspace{1cm}}\\[0.5em]
\textbf{Imprensa min\'uscula:} \underline{\hspace{1cm}} \underline{\hspace{1cm}} \underline{\hspace{1cm}} \underline{\hspace{1cm}} \underline{\hspace{1cm}}\\[0.5em]
\textbf{Cursiva mai\'uscula:} \underline{\hspace{1cm}} \underline{\hspace{1cm}} \underline{\hspace{1cm}} \underline{\hspace{1cm}} \underline{\hspace{1cm}}\\[0.5em]
\textbf{Cursiva min\'uscula:} \underline{\hspace{1cm}} \underline{\hspace{1cm}} \underline{\hspace{1cm}} \underline{\hspace{1cm}} \underline{\hspace{1cm}}
\end{center}

\textbf{Atividade 3 --- Revis\~a{}o das Cinco Vogais}

\instrucao{Complete com a vogal que falta em cada palavra.}

\begin{center}
\begin{tabular}{llll}
\underline{~~~}la & \underline{~~~}la & \underline{~~~}la & \underline{~~~}la \\
(A) & (E) & (I) & (O) \\[0.5em]
\multicolumn{4}{c}{\underline{~~~}la \quad (U)}
\end{tabular}
\end{center}
Diga em voz alta: \textbf{ala, ela, ila, ola, ula}

\caixapista{A, E, I, O, U --- as cinco vogais. Cada uma tem sua boca.}
\end{consolidacao}

\plancheta{Sons da Letra U --- IPA}

\begin{sonsdaletra}
Os sons da letra U no portugu\^es{} brasileiro (IPA --- Alfabeto Fon\'etico Internacional):

\somipa[bico]{U forte}{u}{som grave e fechado}{uva, lua, unha, cume}
\somipa[nasal]{U nasal}{u\~}{anasalado}{um, algum, jejum}
\somipa[bico]{U semivogal}{w}{som r\'apido no fim da s\'ilaba}{pau, mau, c\'eu}
\end{sonsdaletra}

\plancheta{Fam\'ilia Sil\'abica da Letra U}

\begin{familiasilabica}
A fam\'ilia sil\'abica da letra U:

\begin{center}
uva \quad unha \quad urso \quad urubu \quad um \quad uba
\end{center}
\instrucao{Leia a fam\'ilia em voz alta, devagar. Depois cubra e tente lembrar.}
\end{familiasilabica}

\plancheta{Atividade de Som e S\'ilabas da Letra U}

\begin{somletra}
\textbf{1. Ca\c{c}a ao som.} Ou\c{c}a a palavra. Se tiver o som de U, marque o quadradinho:
\begin{center}
$\square$~UVA \quad $\square$~URSO \quad $\square$~MUNDO \quad $\square$~P\'E \quad $\square$~MA\c{c}\~A \quad $\square$~DEDO
\end{center}

\textbf{2. Tem ou n\~ao{} tem?} Rodeie SIM ou N\~ao:
\begin{center}
UVA --- (\textbf{SIM}) (\textbf{N\~AO}) \quad P\'E --- (\textbf{SIM}) (\textbf{N\~AO}) \\ URSO --- (\textbf{SIM}) (\textbf{N\~AO}) \quad MA\c{c}\~A --- (\textbf{SIM}) (\textbf{N\~AO})
\end{center}
\caixapausa{A boca faz bico pequeno para o som de U.}
\end{somletra}

\begin{somsilaba}
\textbf{3. Som das s\'ilabas.} Ou\c{c}a, bata palmas e conte as s\'ilabas:
\begin{center}
U.VA $\rightarrow$ \underline{\hspace{1cm}} s\'ilabas \quad MUN.DO $\rightarrow$ \underline{\hspace{1cm}} s\'ilabas
\end{center}

\textbf{4. Qual s\'ilaba tem o som?} Circule a s\'ilaba com o som de U:
\begin{center}
\textbf{U}.VA \quad \textbf{MUN}.DO
\end{center}
\caixapista{A vogal U aparece na s\'ilaba forte (\'u) e na fraca (u).}
\end{somsilaba}


\plancheta{Miss\~a{}o Motora --- Letra U}

\begin{motricidade}
\textbf{Movimento:} trace as ondas em U, de cima para baixo e de baixo para cima.
\begin{center}
\begin{tikzpicture}
  \draw[very thick, dashed, color=primary] (0,0) to[out=-90,in=180] (0.75,-0.8) to[out=0,in=-90] (1.5,0);
  \draw[very thick, dashed, color=secondary] (2.5,0) to[out=-90,in=180] (3.25,-0.8) to[out=0,in=-90] (4,0);
  \draw[very thick, dashed, color=accent] (5,0) to[out=-90,in=180] (5.75,-0.8) to[out=0,in=-90] (6.5,0);
\end{tikzpicture}
\end{center}
\caixapista{A onda do U \'e{} a mesma curva do U cursivo!}
\end{motricidade}

\plancheta{Ponte --- Da Letra U para a Pr\'oxima}

\begin{transicao}
As vogais terminam aqui! Compare o U com o M, a primeira consoante:
\begin{center}
\textbf{uva} / \textbf{mala} \qquad \textbf{uma} / \textbf{m\~a{}o} \qquad \textbf{urso} / \textbf{mesa}
\end{center}
\instrucao{No U, a boca fica em bico e o ar passa livre. No M, a boca fecha e o som sai pelo nariz.}
\end{transicao}

\plancheta{Recompensa --- Miss\~a{}o 5}

\begin{recompensa}
\estrelas{3}\quad \ofensiva{5}\\[0.3em]\barraprogresso{5}{16}\\[0.5em]
\subiunivel{K2 --- Consci\^encia Fonol\'ogica}\\[0.3em]
\checklistdominio{dominar as cinco vogais nas 4 formas e seus sons principais}
\end{recompensa}


%% LICAO 6: CONSOLIDACAO DAS VOGAIS
%% ------------------------------------------------------------
\section{Licao 6: Consolidacao das Vogais}

% Rotina neuroinclusiva e badges (R201)
\rotinasimples
\nivelKbadge{K1 --- Letras e Grafias}
\rotabadge{1A --- Recomposicao Intensiva}


\metadata{A-06}{Nivel 1 -- Pre-silabico}{EF01LP01}{D1}{EF01LP01}{Resolucao CNE/CEB 7/2010, art. 12}

\plancheta{Folha de Pratica 1.6 --- CONSOLIDACAO DAS VOGAIS}
\begin{exercicio}[Folha de Pratica 1.6 --- CONSOLIDACAO DAS VOGAIS]
\noindent\textbf{Nome:} \underline{\hspace{3.2cm}} \quad \textbf{Data:} \underline{\hspace{1.6cm}} \quad \textbf{Nota:} \underline{\hspace{0.9cm}}/10

\subsection*{PARTE 1: RECONHECIMENTO}

\begin{enumerate}
  \item Escreva a vogal que falta em cada palavra:
  \begin{center}
  m\underline{E}sa \quad s\underline{I}po \quad c\underline{O}sa \quad l\underline{U}vro \quad \underline{A}gua
  \end{center}

  \item Separe as letras por vogais (use cores diferentes):
  \begin{center}
  \textcolor{red}{m} \textcolor{green!70!black}{e} \textcolor{blue}{s} \textcolor{red}{a} $\rightarrow$ A, E\\[0.5em]
  \textcolor{red}{c} \textcolor{red}{a} \textcolor{blue}{s} \textcolor{red}{a} $\rightarrow$ A\\[0.5em]
  \textcolor{blue}{s} \textcolor{red}{a} \textcolor{yellow!80!black}{p} \textcolor{yellow!80!black}{o} $\rightarrow$ A, O
  \end{center}

  \item Circule a vogal que aparece mais vezes:
  \begin{center}
  A  ~  E  ~  A  ~  I  ~  A  ~  O  ~  A  ~  U  ~  A  ~  E $\rightarrow$ \fbox{A}
  \end{center}

  \item Escreva todas as vogais em ordem:
  \begin{center}
  \underline{\hspace{1cm}} \quad \underline{\hspace{1cm}} \quad \underline{\hspace{1cm}} \quad \underline{\hspace{1cm}} \quad \underline{\hspace{1cm}}
  \end{center}

  \item Identifique quantas vogais tem cada palavra:
  \begin{center}
  ``mesa'' = \underline{\hspace{0.5cm}} vogais \quad ``sapo'' = \underline{\hspace{0.5cm}} vogais \quad ``agua'' = \underline{\hspace{0.5cm}} vogais
  \end{center}
\end{enumerate}

\subsection*{PARTE 2: PRODUCAO GUIADA}

\begin{enumerate}[resume]
  \item Escreva 3 palavras com cada vogal:
  \begin{center}
  \begin{tabular}{ll}
  A: \underline{\hspace{1.8cm}} \underline{\hspace{1.8cm}} \underline{\hspace{1.8cm}} &
  E: \underline{\hspace{1.8cm}} \underline{\hspace{1.8cm}} \underline{\hspace{1.8cm}} \\[1em]
  I: \underline{\hspace{1.8cm}} \underline{\hspace{1.8cm}} \underline{\hspace{1.8cm}} &
  O: \underline{\hspace{1.8cm}} \underline{\hspace{1.8cm}} \underline{\hspace{1.8cm}} \\[1em]
  U: \underline{\hspace{1.8cm}} \underline{\hspace{1.8cm}} \underline{\hspace{1.8cm}} &
  \end{tabular}
  \end{center}

  \item Escreva o nome completo de todas as vogais:
  \begin{center}
  \underline{\hspace{6cm}}
  \end{center}

  \item Desenhe 5 coisas (uma para cada vogal) e escreva o nome:
  \begin{center}
  \begin{tikzpicture}
    \draw[thick, dashed] (0,0) rectangle (2.5,2);
    \draw[thick, dashed] (3,0) rectangle (5.5,2);
    \draw[thick, dashed] (6,0) rectangle (8.5,2);
    \draw[thick, dashed] (9,0) rectangle (11.5,2);
    \draw[thick, dashed] (12,0) rectangle (14.5,2);
    \node[below] at (1.25,-0.3) {A: \underline{\hspace{2cm}}};
    \node[below] at (4.25,-0.3) {E: \underline{\hspace{2cm}}};
    \node[below] at (7.25,-0.3) {I: \underline{\hspace{2cm}}};
    \node[below] at (10.25,-0.3) {O: \underline{\hspace{2cm}}};
    \node[below] at (13.25,-0.3) {U: \underline{\hspace{2cm}}};
  \end{tikzpicture}
  \end{center}

  \item Copie e complete a frase:
  \begin{center}
  ``A \underline{\hspace{2cm}} voa e o \underline{\hspace{2cm}} pula no \underline{\hspace{2cm}}.''
  \end{center}
\end{enumerate}

\subsection*{PARTE 3: INTEGRACAO}

\begin{enumerate}[resume]
  \item Leia em voz alta para o professor:
  \begin{center}
  $\rightarrow$ A --- E --- I --- O --- U\\[0.5em]
  $\rightarrow$ agua --- escada --- igreja --- ovo --- uva
  \end{center}

  \item Nome de um animal para cada vogal:
  \begin{center}
  A: \underline{\hspace{3cm}} \quad E: \underline{\hspace{3cm}} \quad I: \underline{\hspace{3cm}} \quad O: \underline{\hspace{3cm}} \quad U: \underline{\hspace{3cm}}
  \end{center}

  \item Circule a opcao correta:
  \begin{enumerate}[label=(\alph*)]
    \item Quantas vogais existem? ( ) 3 ( ) 5 ( ) 7
    \item Qual e a vogal mais aberta? ( ) I ( ) A ( ) O
    \item Qual e a vogal mais fechada? ( ) A ( ) I ( ) U
  \end{enumerate}

  \item Avaliacao diagnostica (marque o que o aluno ja faz):
  \begin{center}
  $\square$ Le todas as vogais isoladas\\
  $\square$ Le vogais em palavras simples\\
  $\square$ Escreve vogais corretamente\\
  $\square$ Identifica vogais em palavras\\
  $\square$ Diferencia som das vogais
  \end{center}

  \item Leitura em voz alta: o professor aponta e o aluno le:
  \begin{center}
  a \quad e \quad i \quad o \quad u \quad agua \quad elefante \quad igreja \quad ovo \quad uva
  \end{center}
\end{enumerate}
\end{exercicio}

\begin{dica}
\textbf{Dica para o Professor:} Esta licao consolida todas as vogais. Antes de comecar, revise rapidamente cada boquinha. Use o espelho para que os alunos observem a progressao da boca (aberta $\rightarrow$ entreaberta $\rightarrow$ estreita $\rightarrow$ arredondada $\rightarrow$ bem arredondada). Avalie se o aluno est\'a no n\'ivel pr\'e-sil\'abico: se identifica e escreve todas as vogais, est\'a pronto para avan\c{c}ar para as consoantes.
\end{dica}

%% ============================================================
%% UNIDADE 2 -- CONSONANTES DO BLOCO 1 (Nivel 2 -- Silabico)
%% ============================================================

\input{checkpoints/rastreio-u1}
\chapter{Unidade 2 --- Consonantes do Bloco 1}

%% ------------------------------------------------------------
%% LICAO 7: LETRA M
%% ------------------------------------------------------------
\section{Licao 7: A Letra M}

% Rotina neuroinclusiva e badges (R201)
\rotinasimples
\nivelKbadge{K1 --- Letras e Grafias}
\rotabadge{1A --- Recomposicao Intensiva}


\letra{M}
% Quatro formas gráficas da letra (R201)
\quatroformasletra{M}{m}

\boq[fechada]{M}{%
  Boca FECHADA, labios juntos.\\
  Ar sai pelo NARIZ (son nasal).\\
  Vibracao sentida no nariz.\\
  Som: ``MMMMMMM'' (nasal, continuo).\\[0.5em]
  Palavras exemplo: macaco, mesa, milho, mola, lua.
}

\pistaarticulatoria{Som: mmm (nasal, cont\'{i}nuo)}{Boca fechada, som sai pelo nariz}{Dedo no nariz para sentir a vibra\c{c}\~{a}o nasal}

\metadata{B-01}{Nivel 2 -- Silabico}{EF01LP02}{D1}{EF01LP02}{Resolucao CNE/CEB 7/2010, art. 12}

\plancheta{Folha de Pratica 2.1 --- Letra M}
\begin{exercicio}[Folha de Pratica 2.1 --- Letra M]
\noindent\textbf{Nome:} \underline{\hspace{3.2cm}} \quad \textbf{Data:} \underline{\hspace{1.6cm}} \quad \textbf{Nota:} \underline{\hspace{0.9cm}}/10

\subsection*{PARTE 1: RECONHECIMENTO}

\begin{enumerate}
  \item Formar silabas com M e vogais:
  \begin{center}
  \sil{M}{A} \quad \sil{M}{E} \quad \sil{M}{I} \quad \sil{M}{O} \quad \sil{M}{U}
  \end{center}

  \item Sublinhe palavras que comecam com ``M'':
  \begin{center}
  \uline{macaco} \quad gato \quad \uline{mesa} \quad sapo \quad \uline{milho} \quad rato \quad \uline{mola}
  \end{center}

  \item Coloque $\checkmark$ nas linhas com ``M'':
  \begin{center}
  \begin{tabular}{ll}
  $\checkmark$ \texttt{M A C A C O} & $\checkmark$ \texttt{M E S A} \\
  $\square$ \texttt{C A S A} & $\checkmark$ \texttt{M I L H O} \\
  $\checkmark$ \texttt{M O L A} & $\square$ \texttt{L A T A}
  \end{tabular}
  \end{center}

  \item Identifique a posicao do ``M'' em ``MACACO'':
  \begin{center}
  \fbox{M} --- A --- C --- A --- C --- O $\rightarrow$ posicao \underline{1}
  \end{center}

  \item Diga em voz alta tres palavras com ``M'' e anote:
  \begin{center}
  1. \underline{\hspace{4cm}} \quad 2. \underline{\hspace{4cm}} \quad 3. \underline{\hspace{4cm}}
  \end{center}
\end{enumerate}

\subsection*{PARTE 2: PRODUCAO GUIADA}

\begin{enumerate}[resume]
  \item Complete com ``M'':
  \begin{center}
  \underline{M}aca \quad \underline{M}esa \quad \underline{M}ilho \quad \underline{M}acaco \quad \underline{M}oco
  \end{center}

  \item Escreva a letra ``M'' maiuscula cinco vezes:
  \begin{center}
  \underline{\hspace{1.5cm}} \quad \underline{\hspace{1.5cm}} \quad \underline{\hspace{1.5cm}} \quad \underline{\hspace{1.5cm}} \quad \underline{\hspace{1.5cm}}
  \end{center}

  \item Escreva a letra ``M'' minuscula cinco vezes:
  \begin{center}
  \underline{\hspace{1.5cm}} \quad \underline{\hspace{1.5cm}} \quad \underline{\hspace{1.5cm}} \quad \underline{\hspace{1.5cm}} \quad \underline{\hspace{1.5cm}}
  \end{center}

  \item Escreva 3 palavras com ``M'':
  \begin{center}
  1. \underline{\hspace{4cm}} \quad 2. \underline{\hspace{4cm}} \quad 3. \underline{\hspace{4cm}}
  \end{center}

  \item Copie as silabas e palavras:
  \begin{center}
  MA --- ME --- MI --- MO --- MU\\
  macaco --- mesa --- milho --- mola --- lua
  \end{center}
\end{enumerate}

\subsection*{PARTE 3: INTEGRACAO}

\begin{enumerate}[resume]
  \item Leia as silabas em voz alta:
  \begin{center}
  $\rightarrow$ MA  ~  ME  ~  MI  ~  MO  ~  MU
  \end{center}

  \item Nome de um animal com ``M'':
  \begin{center}
  \underline{\hspace{6cm}}
  \end{center}

  \item Desenhe algo com ``M'' e escreva o nome:
  \begin{center}
  \begin{tikzpicture}
    \draw[thick, dashed] (0,0) rectangle (5,3);
    \node[below] at (2.5,-0.3) {Nome: \underline{\hspace{4cm}}};
  \end{tikzpicture}
  \end{center}

  \item Circule a opcao correta:
  \begin{enumerate}[label=(\alph*)]
    \item O M tem som: ( ) oral ( ) nasal ( ) vibrante
    \item A boca fica: ( ) aberta ( ) entreaberta ( ) fechada
    \item O M aparece em: ( ) tinta ( ) mesa ( ) uva
  \end{enumerate}

  \item Leia e complete com a silaba correta:
  \begin{center}
  \underline{MA}ca \quad \sil{M}{I}lho \quad \underline{MU}la \quad \sil{M}{O}la \quad \underline{ME}sa
  \end{center}
\end{enumerate}
\end{exercicio}

\begin{dica}
\textbf{Dica para o Professor:} O M e a primeira consonante e tem som nasal. Peca que os alunos fechem a boca e vibrem o nariz. Coloque o dedo na garganta e nariz para sentir a vibracao. Mostre que M + vogal forma silaba (MA, ME, MI, MO, MU). Este e o momento de introduzir o conceito de silaba.
\end{dica}

%% ------------------------------------------------------------


% ============================================================
% CONSOLIDACAO DA LETRA M (R201.1) -- Missao 6/16
% ============================================================

\plancheta{Miss\~a{}o 6 --- Letra M}

\begin{missao}
\textbf{\faGamepad~Miss\~a{}o 6 de 16:} Consolide a letra M, a primeira consoante nasal. Sinta o nariz vibrar!
\end{missao}

\begin{consolidacao}
\textbf{Atividade 1 --- Ca\c{c}ada ao M}

\instrucao{Circule o M e sublinhe as palavras com M.}

\begin{center}
\begin{tabular}{cccccc}
\fbox{M} & S & \fbox{M} & T & \fbox{M} & R \\
L & \fbox{M} & S & \fbox{M} & T & \fbox{M}
\end{tabular}
\end{center}
\begin{center}
\uline{mala} \quad \uline{mel} \quad sopa \quad \uline{mula} \quad tatu \quad \uline{mimo}
\end{center}

\caixapausa{Se precisar, pare aqui. Respire fundo e continue quando quiser.}

\textbf{Atividade 2 --- Escrita do M --- 4 formas}

\instrucao{Escreva o M cinco vezes em cada forma.}

\begin{center}
\textbf{Imprensa mai\'uscula:} \underline{\hspace{1cm}} \underline{\hspace{1cm}} \underline{\hspace{1cm}} \underline{\hspace{1cm}} \underline{\hspace{1cm}}\\[0.5em]
\textbf{Imprensa min\'uscula:} \underline{\hspace{1cm}} \underline{\hspace{1cm}} \underline{\hspace{1cm}} \underline{\hspace{1cm}} \underline{\hspace{1cm}}\\[0.5em]
\textbf{Cursiva mai\'uscula:} \underline{\hspace{1cm}} \underline{\hspace{1cm}} \underline{\hspace{1cm}} \underline{\hspace{1cm}} \underline{\hspace{1cm}}\\[0.5em]
\textbf{Cursiva min\'uscula:} \underline{\hspace{1cm}} \underline{\hspace{1cm}} \underline{\hspace{1cm}} \underline{\hspace{1cm}} \underline{\hspace{1cm}}
\end{center}

\textbf{Atividade 3 --- Jogo do Nariz}

\instrucao{Complete com M e sinta o nariz vibrar em cada palavra.}

\begin{center}
\underline{~~~}ala \quad \underline{~~~}el \quad \underline{~~~}ula \quad \underline{~~~}imo \quad \underline{~~~}ama
\end{center}
Diga em voz alta: \textbf{mala, mel, mula, mimo, mama} --- com a m\~a{}o no nariz!

\caixapista{M = boca fechada + nariz vibrando. Coloque o dedo no nariz e fale mmmm.}
\end{consolidacao}

\plancheta{Sons da Letra M --- IPA}

\begin{sonsdaletra}
Os sons da letra M no portugu\^es{} brasileiro (IPA --- Alfabeto Fon\'etico Internacional):

\somipa[fechada]{M no in\'icio}{m}{l\'abios fechados, som nasal}{mala, m\~a{}o, mim}
\somipa[nasal]{M no fim (nasalizador)}{\~}{fecha a s\'ilaba e anasala a vogal}{bom [bo\~], som [so\~], campo [k6\~pu]}
\end{sonsdaletra}

\plancheta{Fam\'ilia Sil\'abica da Letra M}

\begin{familiasilabica}
A fam\'ilia sil\'abica da letra M:

\begin{center}
mama \quad mam\~a{}e \quad m\~a{}o \quad mala \quad mel \quad miolo \quad mula \quad mudo \quad mimo
\end{center}
\instrucao{Leia a fam\'ilia em voz alta, devagar. Depois cubra e tente lembrar.}
\end{familiasilabica}

\plancheta{Atividade de Som e S\'ilabas da Letra M}

\begin{somletra}
\textbf{1. Ca\c{c}a ao som.} Ou\c{c}a a palavra. Se tiver o som de M, marque o quadradinho:
\begin{center}
$\square$~MAM\~AO \quad $\square$~MENINA \quad $\square$~CAMA \quad $\square$~PAPAI \quad $\square$~SAPATO \quad $\square$~DEDO
\end{center}

\textbf{2. Tem ou n\~ao{} tem?} Rodeie SIM ou N\~ao:
\begin{center}
MAM\~AO --- (\textbf{SIM}) (\textbf{N\~AO}) \quad PAPAI --- (\textbf{SIM}) (\textbf{N\~AO}) \\ MENINA --- (\textbf{SIM}) (\textbf{N\~AO}) \quad SAPATO --- (\textbf{SIM}) (\textbf{N\~AO})
\end{center}
\caixapausa{A boca fecha e o som sai pelo nariz.}
\end{somletra}

\begin{somsilaba}
\textbf{3. Som das s\'ilabas.} Ou\c{c}a, bata palmas e conte as s\'ilabas:
\begin{center}
MA.M\~AO $\rightarrow$ \underline{\hspace{1cm}} s\'ilabas \quad ME.NI.NA $\rightarrow$ \underline{\hspace{1cm}} s\'ilabas
\end{center}

\textbf{4. Qual s\'ilaba tem o som?} Circule a s\'ilaba com o som de M:
\begin{center}
\textbf{MA}.M\~AO \quad \textbf{ME}.NI.NA
\end{center}
\caixapista{A consoante M se junta com a vogal: MA, ME, MI, MO, MU.}
\end{somsilaba}


\plancheta{Miss\~a{}o Motora --- Letra M}

\begin{motricidade}
\textbf{Movimento:} trace as montanhas em zigue-zague, como o M.
\begin{center}
\begin{tikzpicture}
  \draw[very thick, dashed, color=primary] (0,0) -- (0.6,1.2) -- (1.2,0) -- (1.8,1.2) -- (2.4,0) -- (3.0,1.2) -- (3.6,0);
  \draw[very thick, dashed, color=secondary] (4.6,0) -- (5.2,1.2) -- (5.8,0) -- (6.4,1.2) -- (7.0,0) -- (7.6,1.2) -- (8.2,0);
\end{tikzpicture}
\end{center}
\caixapista{Subir, descer, subir, descer --- o M \'e{} uma montanha dupla.}
\end{motricidade}

\plancheta{Ponte --- Da Letra M para a Pr\'oxima}

\begin{transicao}
Compare o M com o S:
\begin{center}
\textbf{m\~a{}o} / \textbf{sapo} \qquad \textbf{mel} / \textbf{sol} \qquad \textbf{mala} / \textbf{sala}
\end{center}
\instrucao{No M a boca fecha e o nariz vibra. No S os dentes quase se juntam e o ar assobia.}
\end{transicao}

\plancheta{Recompensa --- Miss\~a{}o 6}

\begin{recompensa}
\estrelas{3}\quad \ofensiva{6}\\[0.3em]\barraprogresso{6}{16}\\[0.5em]

\checklistdominio{sentir a vibra\c{c}\~a{}o nasal do M e ler palavras com M no in\'icio}
\end{recompensa}


%% LICAO 8: LETRA S
%% ------------------------------------------------------------
\section{Licao 8: A Letra S}

% Rotina neuroinclusiva e badges (R201)
\rotinasimples
\nivelKbadge{K1 --- Letras e Grafias}
\rotabadge{1A --- Recomposicao Intensiva}


\letra{S}
% Quatro formas gráficas da letra (R201)
\quatroformasletra{S}{s}

\boq[dental]{S}{%
  Boca ENTREABERTA, dentes levemente juntos.\\
  Ar sai pelas laterais da boca (sibilante).\\
  Lingua proxima ao palato.\\
  Som: ``SSSSSSS'' (sibilante, suave).\\[0.5em]
  Palavras exemplo: sapo, sol, sapato, sala, saco.
}

\pistaarticulatoria{Som: sss (cont\'{i}nuo, fino)}{Dentes quase juntos, l'{i}ngua atr'{a}s dos dentes}{M\~{a}o na frente da boca para sentir o ar frio e cont'{i}nuo}

\metadata{B-02}{Nivel 2 -- Silabico}{EF01LP02}{D1}{EF01LP02}{Resolucao CNE/CEB 7/2010, art. 12}

\plancheta{Folha de Pratica 2.2 --- Letra S}
\begin{exercicio}[Folha de Pratica 2.2 --- Letra S]
\noindent\textbf{Nome:} \underline{\hspace{3.2cm}} \quad \textbf{Data:} \underline{\hspace{1.6cm}} \quad \textbf{Nota:} \underline{\hspace{0.9cm}}/10

\subsection*{PARTE 1: RECONHECIMENTO}

\begin{enumerate}
  \item Formar silabas com S e vogais:
  \begin{center}
  \sil{S}{A} \quad \sil{S}{E} \quad \sil{S}{I} \quad \sil{S}{O} \quad \sil{S}{U}
  \end{center}

  \item Sublinhe palavras que comecam com ``S'':
  \begin{center}
  \uline{sapo} \quad gato \quad \uline{sol} \quad rato \quad \uline{sapato} \quad mola \quad \uline{sala}
  \end{center}

  \item Coloque $\checkmark$ nas linhas com ``S'':
  \begin{center}
  \begin{tabular}{ll}
  $\checkmark$ \texttt{S A P O} & $\checkmark$ \texttt{S O L} \\
  $\square$ \texttt{M E S A} & $\checkmark$ \texttt{S A L A} \\
  $\checkmark$ \texttt{S A P A T O} & $\square$ \texttt{M I L H O}
  \end{tabular}
  \end{center}

  \item Identifique a posicao do ``S'' em ``SAPO'':
  \begin{center}
  \fbox{S} --- A --- P --- O $\rightarrow$ posicao \underline{1}
  \end{center}

  \item Diga em voz alta tres palavras com ``S'' e anote:
  \begin{center}
  1. \underline{\hspace{4cm}} \quad 2. \underline{\hspace{4cm}} \quad 3. \underline{\hspace{4cm}}
  \end{center}
\end{enumerate}

\subsection*{PARTE 2: PRODUCAO GUIADA}

\begin{enumerate}[resume]
  \item Complete com ``S'':
  \begin{center}
  \underline{S}apo \quad \underline{S}ol \quad \underline{S}apato \quad \underline{S}ala \quad \underline{S}aco
  \end{center}

  \item Escreva a letra ``S'' maiuscula cinco vezes:
  \begin{center}
  \underline{\hspace{1.5cm}} \quad \underline{\hspace{1.5cm}} \quad \underline{\hspace{1.5cm}} \quad \underline{\hspace{1.5cm}} \quad \underline{\hspace{1.5cm}}
  \end{center}

  \item Escreva a letra ``S'' minuscula cinco vezes:
  \begin{center}
  \underline{\hspace{1.5cm}} \quad \underline{\hspace{1.5cm}} \quad \underline{\hspace{1.5cm}} \quad \underline{\hspace{1.5cm}} \quad \underline{\hspace{1.5cm}}
  \end{center}

  \item Escreva 3 palavras com ``S'':
  \begin{center}
  1. \underline{\hspace{4cm}} \quad 2. \underline{\hspace{4cm}} \quad 3. \underline{\hspace{4cm}}
  \end{center}

  \item Copie as silabas e palavras:
  \begin{center}
  SA --- SE --- SI --- SO --- SU\\
  sapo --- sol --- sapato --- sala --- saco
  \end{center}
\end{enumerate}

\subsection*{PARTE 3: INTEGRACAO}

\begin{enumerate}[resume]
  \item Leia as silabas em voz alta:
  \begin{center}
  $\rightarrow$ SA  ~  SE  ~  SI  ~  SO  ~  SU
  \end{center}

  \item Nome de um animal com ``S'':
  \begin{center}
  \underline{\hspace{6cm}}
  \end{center}

  \item Desenhe algo com ``S'' e escreva o nome:
  \begin{center}
  \begin{tikzpicture}
    \draw[thick, dashed] (0,0) rectangle (5,3);
    \node[below] at (2.5,-0.3) {Nome: \underline{\hspace{4cm}}};
  \end{tikzpicture}
  \end{center}

  \item Circule a opcao correta:
  \begin{enumerate}[label=(\alph*)]
    \item O S tem som: ( ) nasal ( ) sibilante ( ) vibrante
    \item O ar sai: ( ) pelo nariz ( ) pelas laterais ( ) pelo fundo
    \item O S aparece em: ( ) mesa ( ) sapo ( ) ovo
  \end{enumerate}

  \item Leia e complete com a silaba correta:
  \begin{center}
  \sil{S}{A}po \quad \underline{SE}te \quad \sil{S}{I}la \quad \underline{SO}la \quad \underline{SU}co
  \end{center}
\end{enumerate}
\end{exercicio}

\begin{dica}
\textbf{Dica para o Professor:} O S e uma sibilante. Peca que os alunos coloquem a lingua proxima do palato e soprem suave, sentindo o ar pelas laterais. Compare com o M (nasal) para que percebam a diferenca. Use o espelho para observar que os dentes ficam levemente juntos.
\end{dica}

%% ------------------------------------------------------------


% ============================================================
% CONSOLIDACAO DA LETRA S (R201.1) -- Missao 7/16
% ============================================================

\plancheta{Miss\~a{}o 7 --- Letra S}

\begin{missao}
\textbf{\faGamepad~Miss\~a{}o 7 de 16:} Consolide a letra S. Aten\c{c}\~a{}o: o S muda de som no meio das palavras!
\end{missao}

\begin{consolidacao}
\textbf{Atividade 1 --- Ca\c{c}ada ao S}

\instrucao{Circule o S e sublinhe as palavras com S.}

\begin{center}
\begin{tabular}{cccccc}
\fbox{S} & T & \fbox{S} & M & \fbox{S} & R \\
L & \fbox{S} & T & \fbox{S} & M & \fbox{S}
\end{tabular}
\end{center}
\begin{center}
\uline{sapo} \quad \uline{selo} \quad tatu \quad \uline{sino} \quad mel \quad \uline{suco}
\end{center}

\caixapausa{Se precisar, pare aqui. Respire fundo e continue quando quiser.}

\textbf{Atividade 2 --- Escrita do S --- 4 formas}

\instrucao{Escreva o S cinco vezes em cada forma.}

\begin{center}
\textbf{Imprensa mai\'uscula:} \underline{\hspace{1cm}} \underline{\hspace{1cm}} \underline{\hspace{1cm}} \underline{\hspace{1cm}} \underline{\hspace{1cm}}\\[0.5em]
\textbf{Imprensa min\'uscula:} \underline{\hspace{1cm}} \underline{\hspace{1cm}} \underline{\hspace{1cm}} \underline{\hspace{1cm}} \underline{\hspace{1cm}}\\[0.5em]
\textbf{Cursiva mai\'uscula:} \underline{\hspace{1cm}} \underline{\hspace{1cm}} \underline{\hspace{1cm}} \underline{\hspace{1cm}} \underline{\hspace{1cm}}\\[0.5em]
\textbf{Cursiva min\'uscula:} \underline{\hspace{1cm}} \underline{\hspace{1cm}} \underline{\hspace{1cm}} \underline{\hspace{1cm}} \underline{\hspace{1cm}}
\end{center}

\textbf{Atividade 3 --- Jogo dos Dois Sons}

\instrucao{Ou\c{c}a e marque: som de S ou som de Z?}

\begin{center}
\begin{tabular}{ll}
sapo \quad $\rightarrow$ \underline{\hspace{2cm}} & casa \quad $\rightarrow$ \underline{\hspace{2cm}} \\
sol \quad $\rightarrow$ \underline{\hspace{2cm}} & mesa \quad $\rightarrow$ \underline{\hspace{2cm}} \\
suco \quad $\rightarrow$ \underline{\hspace{2cm}} & asa \quad $\rightarrow$ \underline{\hspace{2cm}}
\end{tabular}
\end{center}
Regra: entre duas vogais, o S vira Z! \textbf{ca-sa, me-sa, a-sa}

\caixapista{Coloque a m\~a{}o na garganta: no Z ela vibra, no S n\~a{}o.}
\end{consolidacao}

\plancheta{Sons da Letra S --- IPA}

\begin{sonsdaletra}
Os sons da letra S no portugu\^es{} brasileiro (IPA --- Alfabeto Fon\'etico Internacional):

\somipa[dental]{S no in\'icio}{s}{assobio fino e forte}{sapo, sol, sala}
\somipa[dental]{S entre vogais}{z}{vira som de Z}{casa, mesa, asa}
\somipa[palatal]{S no fim de s\'ilaba}{S}{som de X (varia por regi\~a{}o)}{mesmo, festa}
\somipa[dental]{S no fim de palavra}{s}{assobio final}{\^onibus, l\'apis, m\^es}
\end{sonsdaletra}

\plancheta{Fam\'ilia Sil\'abica da Letra S}

\begin{familiasilabica}
A fam\'ilia sil\'abica da letra S:

\begin{center}
sapo \quad saco \quad sal \quad selo \quad sino \quad sopa \quad suco \quad saia
\end{center}
\instrucao{Leia a fam\'ilia em voz alta, devagar. Depois cubra e tente lembrar.}
\end{familiasilabica}

\plancheta{Atividade de Som e S\'ilabas da Letra S}

\begin{somletra}
\textbf{1. Ca\c{c}a ao som.} Ou\c{c}a a palavra. Se tiver o som de S, marque o quadradinho:
\begin{center}
$\square$~SAPATO \quad $\square$~SOL \quad $\square$~CASA \quad $\square$~BOCA \quad $\square$~DADO \quad $\square$~M\~AO
\end{center}

\textbf{2. Tem ou n\~ao{} tem?} Rodeie SIM ou N\~ao:
\begin{center}
SAPATO --- (\textbf{SIM}) (\textbf{N\~AO}) \quad BOCA --- (\textbf{SIM}) (\textbf{N\~AO}) \\ SOL --- (\textbf{SIM}) (\textbf{N\~AO}) \quad DADO --- (\textbf{SIM}) (\textbf{N\~AO})
\end{center}
\caixapausa{Os dentes ficam juntos e o ar passa assobiando.}
\end{somletra}

\begin{somsilaba}
\textbf{3. Som das s\'ilabas.} Ou\c{c}a, bata palmas e conte as s\'ilabas:
\begin{center}
SA.PA.TO $\rightarrow$ \underline{\hspace{1cm}} s\'ilabas \quad CA.SA $\rightarrow$ \underline{\hspace{1cm}} s\'ilabas
\end{center}

\textbf{4. Qual s\'ilaba tem o som?} Circule a s\'ilaba com o som de S:
\begin{center}
\textbf{SA}.PA.TO \quad CA.\textbf{SA}
\end{center}
\caixapista{A consoante S se junta com a vogal: SA, SE, SI, SO, SU.}
\end{somsilaba}


\plancheta{Miss\~a{}o Motora --- Letra S}

\begin{motricidade}
\textbf{Movimento:} trace a cobra sinuosa, como o S.
\begin{center}
\begin{tikzpicture}
  \draw[very thick, dashed, color=primary] (0,0) to[out=0,in=180] (1,-0.6) to[out=0,in=180] (2,0.3) to[out=0,in=180] (3,-0.3) to[out=0,in=180] (4,0);
  \draw[very thick, dashed, color=secondary] (5,0) to[out=0,in=180] (6,-0.6) to[out=0,in=180] (7,0.3) to[out=0,in=180] (8,-0.3) to[out=0,in=180] (9,0);
\end{tikzpicture}
\end{center}
\caixapista{A cobra do S \'e{} uma curva cont\'inua, sem parar.}
\end{motricidade}

\plancheta{Ponte --- Da Letra S para a Pr\'oxima}

\begin{transicao}
Compare o S com o T:
\begin{center}
\textbf{sapo} / \textbf{tatu} \qquad \textbf{sol} / \textbf{teto} \qquad \textbf{sino} / \textbf{tio}
\end{center}
\instrucao{No S o ar sai cont\'inuo (ssss). No T o som \'e{} seco e r\'apido (t!).}
\end{transicao}

\plancheta{Recompensa --- Miss\~a{}o 7}

\begin{recompensa}
\estrelas{3}\quad \ofensiva{7}\\[0.3em]\barraprogresso{7}{16}\\[0.5em]

\checklistdominio{diferenciar S assobio de S-z (casa, mesa) pela vibra\c{c}\~a{}o}
\end{recompensa}


%% LICAO 9: LETRA T
%% ------------------------------------------------------------
\section{Licao 9: A Letra T}

% Rotina neuroinclusiva e badges (R201)
\rotinasimples
\nivelKbadge{K1 --- Letras e Grafias}
\rotabadge{1A --- Recomposicao Intensiva}


\letra{T}
% Quatro formas gráficas da letra (R201)
\quatroformasletra{T}{t}

\boq[dental]{T}{%
  Boca ENTREABERTA, dentes superiores visiveis.\\
  Lingua toca ATRAS dos dentes superiores.\\
  Som seco e rapido (oclusiva dental).\\
  Som: ``T-T-T-T-T'' (curto, seco).\\[0.5em]
  Palavras exemplo: tinta, teto, tigre, tira, tampa.
}

\pistaarticulatoria{Som: t (seco, r\'{a}pido)}{L'{i}ngua toca atr'{a}s dos dentes superiores}{M\~{a}o na garganta: n\~{a}o vibra (som surdo)}

\metadata{B-03}{Nivel 2 -- Silabico}{EF01LP02}{D1}{EF01LP02}{Resolucao CNE/CEB 7/2010, art. 12}

\plancheta{Folha de Pratica 2.3 --- Letra T}
\begin{exercicio}[Folha de Pratica 2.3 --- Letra T]
\noindent\textbf{Nome:} \underline{\hspace{3.2cm}} \quad \textbf{Data:} \underline{\hspace{1.6cm}} \quad \textbf{Nota:} \underline{\hspace{0.9cm}}/10

\subsection*{PARTE 1: RECONHECIMENTO}

\begin{enumerate}
  \item Formar silabas com T e vogais:
  \begin{center}
  \sil{T}{A} \quad \sil{T}{E} \quad \sil{T}{I} \quad \sil{T}{O} \quad \sil{T}{U}
  \end{center}

  \item Sublinhe palavras que comecam com ``T'':
  \begin{center}
  \uline{tinta} \quad gato \quad \uline{teto} \quad sapo \quad \uline{tigre} \quad rato \quad \uline{tira}
  \end{center}

  \item Coloque $\checkmark$ nas linhas com ``T'':
  \begin{center}
  \begin{tabular}{ll}
  $\checkmark$ \texttt{T I N T A} & $\checkmark$ \texttt{T E T O} \\
  $\square$ \texttt{C A S A} & $\checkmark$ \texttt{T I G R E} \\
  $\checkmark$ \texttt{T I R A} & $\square$ \texttt{M E S A}
  \end{tabular}
  \end{center}

  \item Identifique a posicao do ``T'' em ``TINTA'':
  \begin{center}
  \fbox{T} --- I --- N --- \fbox{T} --- A $\rightarrow$ posicoes \underline{1} e \underline{4}
  \end{center}

  \item Diga em voz alta tres palavras com ``T'' e anote:
  \begin{center}
  1. \underline{\hspace{4cm}} \quad 2. \underline{\hspace{4cm}} \quad 3. \underline{\hspace{4cm}}
  \end{center}
\end{enumerate}

\subsection*{PARTE 2: PRODUCAO GUIADA}

\begin{enumerate}[resume]
  \item Complete com ``T'':
  \begin{center}
  \underline{T}inta \quad \underline{T}eto \quad \underline{T}igre \quad \underline{T}ira \quad \underline{T}ampa
  \end{center}

  \item Escreva a letra ``T'' maiuscula cinco vezes:
  \begin{center}
  \underline{\hspace{1.5cm}} \quad \underline{\hspace{1.5cm}} \quad \underline{\hspace{1.5cm}} \quad \underline{\hspace{1.5cm}} \quad \underline{\hspace{1.5cm}}
  \end{center}

  \item Escreva a letra ``T'' minuscula cinco vezes:
  \begin{center}
  \underline{\hspace{1.5cm}} \quad \underline{\hspace{1.5cm}} \quad \underline{\hspace{1.5cm}} \quad \underline{\hspace{1.5cm}} \quad \underline{\hspace{1.5cm}}
  \end{center}

  \item Escreva 3 palavras com ``T'':
  \begin{center}
  1. \underline{\hspace{4cm}} \quad 2. \underline{\hspace{4cm}} \quad 3. \underline{\hspace{4cm}}
  \end{center}

  \item Copie as silabas e palavras:
  \begin{center}
  TA --- TE --- TI --- TO --- TU\\
  tinta --- teto --- tigre --- tira --- tampa
  \end{center}
\end{enumerate}

\subsection*{PARTE 3: INTEGRACAO}

\begin{enumerate}[resume]
  \item Leia as silabas em voz alta:
  \begin{center}
  $\rightarrow$ TA  ~  TE  ~  TI  ~  TO  ~  TU
  \end{center}

  \item Nome de um animal com ``T'':
  \begin{center}
  \underline{\hspace{6cm}}
  \end{center}

  \item Desenhe algo com ``T'' e escreva o nome:
  \begin{center}
  \begin{tikzpicture}
    \draw[thick, dashed] (0,0) rectangle (5,3);
    \node[below] at (2.5,-0.3) {Nome: \underline{\hspace{4cm}}};
  \end{tikzpicture}
  \end{center}

  \item Circule a opcao correta:
  \begin{enumerate}[label=(\alph*)]
    \item O T tem som: ( ) nasal ( ) sibilante ( ) oclusivo seco
    \item A lingua toca: ( ) o nariz ( ) o palato ( ) atras dos dentes
    \item O T aparece em: ( ) mesa ( ) sapo ( ) tigre
  \end{enumerate}

  \item Leia e complete com a silaba correta:
  \begin{center}
  \sil{T}{A}pa \quad \underline{TE}la \quad \sil{T}{I}ra \quad \underline{TO}pa \quad \underline{TU}do
  \end{center}
\end{enumerate}
\end{exercicio}

\begin{dica}
\textbf{Dica para o Professor:} O T e uma oclusiva dental. Peca que os alunos coloquem a lingua atras dos dentes superiores e soltem o ar de forma seca e rapida. Compare com o S (sibilante) e o M (nasal) para que percebam as tres formas diferentes de produzir sons com a boca.
\end{dica}

%% ------------------------------------------------------------


% ============================================================
% CONSOLIDACAO DA LETRA T (R201.1) -- Missao 8/16
% ============================================================

\plancheta{Miss\~a{}o 8 --- Letra T}

\begin{missao}
\textbf{\faGamepad~Miss\~a{}o 8 de 16:} Consolide a letra T. Descubra o som que o T faz antes do I em algumas regi\~o{}es!
\end{missao}

\begin{consolidacao}
\textbf{Atividade 1 --- Ca\c{c}ada ao T}

\instrucao{Circule o T e sublinhe as palavras com T.}

\begin{center}
\begin{tabular}{cccccc}
\fbox{T} & R & \fbox{T} & S & \fbox{T} & M \\
L & \fbox{T} & R & \fbox{T} & S & \fbox{T}
\end{tabular}
\end{center}
\begin{center}
\uline{tatu} \quad \uline{teto} \quad rato \quad \uline{tio} \quad selo \quad \uline{tudo}
\end{center}

\caixapausa{Se precisar, pare aqui. Respire fundo e continue quando quiser.}

\textbf{Atividade 2 --- Escrita do T --- 4 formas}

\instrucao{Escreva o T cinco vezes em cada forma.}

\begin{center}
\textbf{Imprensa mai\'uscula:} \underline{\hspace{1cm}} \underline{\hspace{1cm}} \underline{\hspace{1cm}} \underline{\hspace{1cm}} \underline{\hspace{1cm}}\\[0.5em]
\textbf{Imprensa min\'uscula:} \underline{\hspace{1cm}} \underline{\hspace{1cm}} \underline{\hspace{1cm}} \underline{\hspace{1cm}} \underline{\hspace{1cm}}\\[0.5em]
\textbf{Cursiva mai\'uscula:} \underline{\hspace{1cm}} \underline{\hspace{1cm}} \underline{\hspace{1cm}} \underline{\hspace{1cm}} \underline{\hspace{1cm}}\\[0.5em]
\textbf{Cursiva min\'uscula:} \underline{\hspace{1cm}} \underline{\hspace{1cm}} \underline{\hspace{1cm}} \underline{\hspace{1cm}} \underline{\hspace{1cm}}
\end{center}

\textbf{Atividade 3 --- Jogo do T Seco}

\instrucao{Complete com T e fale cada palavra bem seco.}

\begin{center}
\underline{~~~}atu \quad \underline{~~~}eto \quad \underline{~~~}io \quad \underline{~~~}udo \quad \underline{~~~}ela
\end{center}
Diga em voz alta: \textbf{tatu, teto, tio, tudo, tela}

\caixapista{T = l\'ingua toca atr\'as dos dentes e solta de uma vez. Som seco, sem ar cont\'inuo.}
\end{consolidacao}

\plancheta{Sons da Letra T --- IPA}

\begin{sonsdaletra}
Os sons da letra T no portugu\^es{} brasileiro (IPA --- Alfabeto Fon\'etico Internacional):

\somipa[dental]{T seco}{t}{l\'ingua atr\'as dos dentes, som seco}{tatu, teto, tudo, tela}
\somipa[palatal]{T antes de I (regional)}{tS}{vira tch em algumas regi\~o{}es}{tia, leite, noite}
\end{sonsdaletra}

\plancheta{Fam\'ilia Sil\'abica da Letra T}

\begin{familiasilabica}
A fam\'ilia sil\'abica da letra T:

\begin{center}
tatu \quad teto \quad tio \quad tia \quad tudo \quad tela \quad tapa
\end{center}
\instrucao{Leia a fam\'ilia em voz alta, devagar. Depois cubra e tente lembrar.}
\end{familiasilabica}

\plancheta{Atividade de Som e S\'ilabas da Letra T}

\begin{somletra}
\textbf{1. Ca\c{c}a ao som.} Ou\c{c}a a palavra. Se tiver o som de T, marque o quadradinho:
\begin{center}
$\square$~TATU \quad $\square$~TOMATE \quad $\square$~MATA \quad $\square$~SOL \quad $\square$~P\~AO \quad $\square$~BOCA
\end{center}

\textbf{2. Tem ou n\~ao{} tem?} Rodeie SIM ou N\~ao:
\begin{center}
TATU --- (\textbf{SIM}) (\textbf{N\~AO}) \quad SOL --- (\textbf{SIM}) (\textbf{N\~AO}) \\ TOMATE --- (\textbf{SIM}) (\textbf{N\~AO}) \quad P\~AO --- (\textbf{SIM}) (\textbf{N\~AO})
\end{center}
\caixapausa{A l\'ingua toca os dentes de cima.}
\end{somletra}

\begin{somsilaba}
\textbf{3. Som das s\'ilabas.} Ou\c{c}a, bata palmas e conte as s\'ilabas:
\begin{center}
TA.TU $\rightarrow$ \underline{\hspace{1cm}} s\'ilabas \quad TO.MA.TE $\rightarrow$ \underline{\hspace{1cm}} s\'ilabas
\end{center}

\textbf{4. Qual s\'ilaba tem o som?} Circule a s\'ilaba com o som de T:
\begin{center}
\textbf{TA}.TU \quad \textbf{TO}.MA.TE
\end{center}
\caixapista{A consoante T se junta com a vogal: TA, TE, TI, TO, TU.}
\end{somsilaba}


\plancheta{Miss\~a{}o Motora --- Letra T}

\begin{motricidade}
\textbf{Movimento:} trace a cruz: uma linha na horizontal, uma na vertical.
\begin{center}
\begin{tikzpicture}
  \draw[very thick, dashed, color=primary] (0,0) -- (4,0);
  \draw[very thick, dashed, color=secondary] (2,0.8) -- (2,-0.8);
  \draw[->, color=accent] (0.3,0.15) -- (1.2,0.15);
  \draw[->, color=accent] (2.15,0.5) -- (2.15,-0.1);
\end{tikzpicture}
\end{center}
\caixapista{O T \'e{} o telhado: primeiro a t\'abua, depois o pilar.}
\end{motricidade}

\plancheta{Ponte --- Da Letra T para a Pr\'oxima}

\begin{transicao}
Compare o T com o R:
\begin{center}
\textbf{tatu} / \textbf{rato} \qquad \textbf{tela} / \textbf{rua} \qquad \textbf{teto} / \textbf{rede}
\end{center}
\instrucao{O T \'e{} seco, a l\'ingua para. O R forte vibra na garganta (rrrr).}
\end{transicao}

\plancheta{Recompensa --- Miss\~a{}o 8}

\begin{recompensa}
\estrelas{3}\quad \ofensiva{8}\\[0.3em]\barraprogresso{8}{16}\\[0.5em]

\checklistdominio{reconhecer o T seco e ler palavras com T}
\end{recompensa}


%% LICAO 10: LETRA R
%% ------------------------------------------------------------
\section{Licao 10: A Letra R}

% Rotina neuroinclusiva e badges (R201)
\rotinasimples
\nivelKbadge{K1 --- Letras e Grafias}
\rotabadge{1A --- Recomposicao Intensiva}


\letra{R}
% Quatro formas gráficas da letra (R201)
\quatroformasletra{R}{r}

\boq[velar]{R}{%
  Boca ENTREABERTA, labios levemente abertos.\\
  Lingua vibra no CEU DA BOCA (vibrante multipla).\\
  Vibracao forte e perceptivel.\\
  Som: ``RRRRRRR'' (vibrante, forte).\\[0.5em]
  Palavras exemplo: rato, roda, rede, rica, rua.
}

\pistaarticulatoria{Som: rrr (vibrante)}{L'{i}ngua treme perto do c'{e}u da boca}{M\~{a}o na garganta para sentir a vibra\c{c}\~{a}o forte}

\metadata{B-04}{Nivel 2 -- Silabico}{EF01LP02}{D1}{EF01LP02}{Resolucao CNE/CEB 7/2010, art. 12}

\plancheta{Folha de Pratica 2.4 --- Letra R}
\begin{exercicio}[Folha de Pratica 2.4 --- Letra R]
\noindent\textbf{Nome:} \underline{\hspace{3.2cm}} \quad \textbf{Data:} \underline{\hspace{1.6cm}} \quad \textbf{Nota:} \underline{\hspace{0.9cm}}/10

\subsection*{PARTE 1: RECONHECIMENTO}

\begin{enumerate}
  \item Formar silabas com R e vogais:
  \begin{center}
  \sil{R}{A} \quad \sil{R}{E} \quad \sil{R}{I} \quad \sil{R}{O} \quad \sil{R}{U}
  \end{center}

  \item Sublinhe palavras que comecam com ``R'':
  \begin{center}
  \uline{rato} \quad gato \quad \uline{roda} \quad sapo \quad \uline{rede} \quad mola \quad \uline{rica}
  \end{center}

  \item Coloque $\checkmark$ nas linhas com ``R'':
  \begin{center}
  \begin{tabular}{ll}
  $\checkmark$ \texttt{R A T O} & $\checkmark$ \texttt{R O D A} \\
  $\square$ \texttt{C A S A} & $\checkmark$ \texttt{R E D E} \\
  $\checkmark$ \texttt{R I C A} & $\square$ \texttt{M E S A}
  \end{tabular}
  \end{center}

  \item Identifique a posicao do ``R'' em ``RATO'':
  \begin{center}
  \fbox{R} --- A --- T --- O $\rightarrow$ posicao \underline{1}
  \end{center}

  \item Diga em voz alta tres palavras com ``R'' e anote:
  \begin{center}
  1. \underline{\hspace{4cm}} \quad 2. \underline{\hspace{4cm}} \quad 3. \underline{\hspace{4cm}}
  \end{center}
\end{enumerate}

\subsection*{PARTE 2: PRODUCAO GUIADA}

\begin{enumerate}[resume]
  \item Complete com ``R'':
  \begin{center}
  \underline{R}ato \quad \underline{R}oda \quad \underline{R}ede \quad \underline{R}ica \quad \underline{R}ua
  \end{center}

  \item Escreva a letra ``R'' maiuscula cinco vezes:
  \begin{center}
  \underline{\hspace{1.5cm}} \quad \underline{\hspace{1.5cm}} \quad \underline{\hspace{1.5cm}} \quad \underline{\hspace{1.5cm}} \quad \underline{\hspace{1.5cm}}
  \end{center}

  \item Escreva a letra ``r'' minuscula cinco vezes:
  \begin{center}
  \underline{\hspace{1.5cm}} \quad \underline{\hspace{1.5cm}} \quad \underline{\hspace{1.5cm}} \quad \underline{\hspace{1.5cm}} \quad \underline{\hspace{1.5cm}}
  \end{center}

  \item Escreva 3 palavras com ``R'':
  \begin{center}
  1. \underline{\hspace{4cm}} \quad 2. \underline{\hspace{4cm}} \quad 3. \underline{\hspace{4cm}}
  \end{center}

  \item Copie as silabas e palavras:
  \begin{center}
  RA --- RE --- RI --- RO --- RU\\
  rato --- roda --- rede --- rica --- rua
  \end{center}
\end{enumerate}

\subsection*{PARTE 3: INTEGRACAO}

\begin{enumerate}[resume]
  \item Leia as silabas em voz alta:
  \begin{center}
  $\rightarrow$ RA  ~  RE  ~  RI  ~  RO  ~  RU
  \end{center}

  \item Nome de um animal com ``R'':
  \begin{center}
  \underline{\hspace{6cm}}
  \end{center}

  \item Desenhe algo com ``R'' e escreva o nome:
  \begin{center}
  \begin{tikzpicture}
    \draw[thick, dashed] (0,0) rectangle (5,3);
    \node[below] at (2.5,-0.3) {Nome: \underline{\hspace{4cm}}};
  \end{tikzpicture}
  \end{center}

  \item Circule a opcao correta:
  \begin{enumerate}[label=(\alph*)]
    \item O R tem som: ( ) nasal ( ) sibilante ( ) vibrante
    \item A lingua toca: ( ) o nariz ( ) os dentes ( ) o ceu da boca
    \item O R aparece em: ( ) mesa ( ) sapo ( ) rato
  \end{enumerate}

  \item Leia e complete com a silaba correta:
  \begin{center}
  \sil{R}{A}to \quad \underline{RE}de \quad \sil{R}{I}ca \quad \underline{RO}da \quad \underline{RU}a
  \end{center}
\end{enumerate}
\end{exercicio}

\begin{dica}
\textbf{Dica para o Professor:} O R e uma vibrante. Peca que os alunos coloquem a lingua no ceu da boca e vibrem com forca. Alguns alunos podem ter dificuldade com a vibracao; nesse caso, peca que digam ``T-R-R-R'' lentamente, aumentando a velocidade. Compare com T (oclusivo) e S (sibilante).
\end{dica}

%% ------------------------------------------------------------


% ============================================================
% CONSOLIDACAO DA LETRA R (R201.1) -- Missao 9/16
% ============================================================

\plancheta{Miss\~a{}o 9 --- Letra R}

\begin{missao}
\textbf{\faGamepad~Miss\~a{}o 9 de 16:} Consolide a letra R. O R tem tr\^es caras: forte no in\'icio, fraco no meio, forte no fim!
\end{missao}

\begin{consolidacao}
\textbf{Atividade 1 --- Ca\c{c}ada ao R}

\instrucao{Circule o R e sublinhe as palavras com R.}

\begin{center}
\begin{tabular}{cccccc}
\fbox{R} & L & \fbox{R} & T & \fbox{R} & S \\
M & \fbox{R} & L & \fbox{R} & T & \fbox{R}
\end{tabular}
\end{center}
\begin{center}
\uline{rato} \quad \uline{rua} \quad lata \quad \uline{rio} \quad teto \quad \uline{roda}
\end{center}

\caixapausa{Se precisar, pare aqui. Respire fundo e continue quando quiser.}

\textbf{Atividade 2 --- Escrita do R --- 4 formas}

\instrucao{Escreva o R cinco vezes em cada forma.}

\begin{center}
\textbf{Imprensa mai\'uscula:} \underline{\hspace{1cm}} \underline{\hspace{1cm}} \underline{\hspace{1cm}} \underline{\hspace{1cm}} \underline{\hspace{1cm}}\\[0.5em]
\textbf{Imprensa min\'uscula:} \underline{\hspace{1cm}} \underline{\hspace{1cm}} \underline{\hspace{1cm}} \underline{\hspace{1cm}} \underline{\hspace{1cm}}\\[0.5em]
\textbf{Cursiva mai\'uscula:} \underline{\hspace{1cm}} \underline{\hspace{1cm}} \underline{\hspace{1cm}} \underline{\hspace{1cm}} \underline{\hspace{1cm}}\\[0.5em]
\textbf{Cursiva min\'uscula:} \underline{\hspace{1cm}} \underline{\hspace{1cm}} \underline{\hspace{1cm}} \underline{\hspace{1cm}} \underline{\hspace{1cm}}
\end{center}

\textbf{Atividade 3 --- Jogo do R Forte e R Fraco}

\instrucao{Ou\c{c}a e marque: R forte (garganta) ou R fraco (batida)?}

\begin{center}
\begin{tabular}{ll}
rato \quad $\rightarrow$ \underline{\hspace{2cm}} & cara \quad $\rightarrow$ \underline{\hspace{2cm}} \\
rua \quad $\rightarrow$ \underline{\hspace{2cm}} & pera \quad $\rightarrow$ \underline{\hspace{2cm}} \\
mar \quad $\rightarrow$ \underline{\hspace{2cm}} & aro \quad $\rightarrow$ \underline{\hspace{2cm}}
\end{tabular}
\end{center}
Regra: no come\c{c}o e no fim, forte. No meio, fraco.

\caixapista{R fraco = uma batida s\'o. R forte = ronco cont\'inuo na garganta.}
\end{consolidacao}

\plancheta{Sons da Letra R --- IPA}

\begin{sonsdaletra}
Os sons da letra R no portugu\^es{} brasileiro (IPA --- Alfabeto Fon\'etico Internacional):

\somipa[velar]{R forte (in\'icio)}{h}{fricativo na garganta (h, x ou rr conforme regi\~a{}o)}{rato, rua, rede}
\somipa[dental]{R fraco (meio)}{r}{batida leve da l\'ingua}{cara, pera, aro}
\somipa[velar]{R forte (fim)}{h}{fricativo de novo}{mar, amor, flor}
\end{sonsdaletra}

\plancheta{Fam\'ilia Sil\'abica da Letra R}

\begin{familiasilabica}
A fam\'ilia sil\'abica da letra R:

\begin{center}
rato \quad rua \quad rio \quad roda \quad rede \quad raio \quad rima
\end{center}
\instrucao{Leia a fam\'ilia em voz alta, devagar. Depois cubra e tente lembrar.}
\end{familiasilabica}

\plancheta{Atividade de Som e S\'ilabas da Letra R}

\begin{somletra}
\textbf{1. Ca\c{c}a ao som.} Ou\c{c}a a palavra. Se tiver o som de R, marque o quadradinho:
\begin{center}
$\square$~RATO \quad $\square$~ROUPA \quad $\square$~CARRO \quad $\square$~SOL \quad $\square$~BOCA \quad $\square$~P\'E
\end{center}

\textbf{2. Tem ou n\~ao{} tem?} Rodeie SIM ou N\~ao:
\begin{center}
RATO --- (\textbf{SIM}) (\textbf{N\~AO}) \quad SOL --- (\textbf{SIM}) (\textbf{N\~AO}) \\ ROUPA --- (\textbf{SIM}) (\textbf{N\~AO}) \quad BOCA --- (\textbf{SIM}) (\textbf{N\~AO})
\end{center}
\caixapausa{A l\'ingua vibra no c\'eu da boca.}
\end{somletra}

\begin{somsilaba}
\textbf{3. Som das s\'ilabas.} Ou\c{c}a, bata palmas e conte as s\'ilabas:
\begin{center}
RA.TO $\rightarrow$ \underline{\hspace{1cm}} s\'ilabas \quad CAR.RO $\rightarrow$ \underline{\hspace{1cm}} s\'ilabas
\end{center}

\textbf{4. Qual s\'ilaba tem o som?} Circule a s\'ilaba com o som de R:
\begin{center}
\textbf{RA}.TO \quad \textbf{CAR}.RO
\end{center}
\caixapista{O R forte aparece no come\c{c}o e no meio: RA, ROU, CAR.}
\end{somsilaba}


\plancheta{Miss\~a{}o Motora --- Letra R}

\begin{motricidade}
\textbf{Movimento:} trace o la\c{c}o do R: suba, curve e des\c{c}a com o rabicho.
\begin{center}
\begin{tikzpicture}
  \draw[very thick, dashed, color=primary] (0,0) -- (0,1.2);
  \draw[very thick, dashed, color=primary] (0,1.2) to[out=0,in=90] (0.7,0.6);
  \draw[very thick, dashed, color=primary] (0.7,0.6) to[out=-90,in=0] (0.3,0.2);
  \draw[->, color=accent] (0.2,0.6) -- (0.2,0.1);
  \draw[->, color=accent] (0.1,1.3) -- (0.3,1.0);
\end{tikzpicture}
\end{center}
\caixapista{O R mai\'usculo tem pilar, barriga e perna. Tr\^es partes!}
\end{motricidade}

\plancheta{Ponte --- Da Letra R para a Pr\'oxima}

\begin{transicao}
Compare o R com o L:
\begin{center}
\textbf{rato} / \textbf{lata} \qquad \textbf{cara} / \textbf{cala} \qquad \textbf{rua} / \textbf{lua}
\end{center}
\instrucao{No R a l\'ingua vibra. No L a ponta da l\'ingua fica parada no c\'eu da boca.}
\end{transicao}

\plancheta{Recompensa --- Miss\~a{}o 9}

\begin{recompensa}
\estrelas{3}\quad \ofensiva{9}\\[0.3em]\barraprogresso{9}{16}\\[0.5em]

\checklistdominio{diferenciar R forte de R fraco pela posi\c{c}\~a{}o na palavra}
\end{recompensa}


%% LICAO 11: LETRA L
%% ------------------------------------------------------------
\section{Licao 11: A Letra L}

% Rotina neuroinclusiva e badges (R201)
\rotinasimples
\nivelKbadge{K1 --- Letras e Grafias}
\rotabadge{1A --- Recomposicao Intensiva}


\letra{L}
% Quatro formas gráficas da letra (R201)
\quatroformasletra{L}{l}

\boq[dental]{L}{%
  Boca ENTREABERTA, dentes levemente visiveis.\\
  Lingua toca o CEU DA BOCA na lateral.\\
  Ar sai por um dos lados da boca.\\
  Som: ``LLLLLLL'' (lateral, suave).\\[0.5em]
  Palavras exemplo: lua, laito, limao, lata, la.
}

\pistaarticulatoria{Som: lll (lateral, cont\'{i}nuo)}{Ponta da l'{i}ngua no c'{e}u da boca, ar sai pelos lados}{M\~{a}o na lateral do pesco\c{c}o para sentir a vibra\c{c}\~{a}o}

\metadata{B-05}{Nivel 2 -- Silabico}{EF01LP02}{D1}{EF01LP02}{Resolucao CNE/CEB 7/2010, art. 12}

\plancheta{Folha de Pratica 2.5 --- Letra L}
\begin{exercicio}[Folha de Pratica 2.5 --- Letra L]
\noindent\textbf{Nome:} \underline{\hspace{3.2cm}} \quad \textbf{Data:} \underline{\hspace{1.6cm}} \quad \textbf{Nota:} \underline{\hspace{0.9cm}}/10

\subsection*{PARTE 1: RECONHECIMENTO}

\begin{enumerate}
  \item Formar silabas com L e vogais:
  \begin{center}
  \sil{L}{A} \quad \sil{L}{E} \quad \sil{L}{I} \quad \sil{L}{O} \quad \sil{L}{U}
  \end{center}

  \item Sublinhe palavras que comecam com ``L'':
  \begin{center}
  \uline{lua} \quad gato \quad \uline{laito} \quad sapo \quad \uline{limao} \quad rato \quad \uline{lata}
  \end{center}

  \item Coloque $\checkmark$ nas linhas com ``L'':
  \begin{center}
  \begin{tabular}{ll}
  $\checkmark$ \texttt{L U A} & $\checkmark$ \texttt{L I M A O} \\
  $\square$ \texttt{C A S A} & $\checkmark$ \texttt{L A T A} \\
  $\checkmark$ \texttt{L A} & $\square$ \texttt{M E S A}
  \end{tabular}
  \end{center}

  \item Identifique a posicao do ``L'' em ``LUA'':
  \begin{center}
  \fbox{L} --- U --- A $\rightarrow$ posicao \underline{1}
  \end{center}

  \item Diga em voz alta tres palavras com ``L'' e anote:
  \begin{center}
  1. \underline{\hspace{4cm}} \quad 2. \underline{\hspace{4cm}} \quad 3. \underline{\hspace{4cm}}
  \end{center}
\end{enumerate}

\subsection*{PARTE 2: PRODUCAO GUIADA}

\begin{enumerate}[resume]
  \item Complete com ``L'':
  \begin{center}
  \underline{L}ua \quad \underline{L}aito \quad \underline{L}imao \quad \underline{L}ata \quad \underline{L}aranja
  \end{center}

  \item Escreva a letra ``L'' maiuscula cinco vezes:
  \begin{center}
  \underline{\hspace{1.5cm}} \quad \underline{\hspace{1.5cm}} \quad \underline{\hspace{1.5cm}} \quad \underline{\hspace{1.5cm}} \quad \underline{\hspace{1.5cm}}
  \end{center}

  \item Escreva a letra ``l'' minuscula cinco vezes:
  \begin{center}
  \underline{\hspace{1.5cm}} \quad \underline{\hspace{1.5cm}} \quad \underline{\hspace{1.5cm}} \quad \underline{\hspace{1.5cm}} \quad \underline{\hspace{1.5cm}}
  \end{center}

  \item Escreva 3 palavras com ``L'':
  \begin{center}
  1. \underline{\hspace{4cm}} \quad 2. \underline{\hspace{4cm}} \quad 3. \underline{\hspace{4cm}}
  \end{center}

  \item Copie as silabas e palavras:
  \begin{center}
  LA --- LE --- LI --- LO --- LU\\
  lua --- laito --- limao --- lata --- laranja
  \end{center}
\end{enumerate}

\subsection*{PARTE 3: INTEGRACAO}

\begin{enumerate}[resume]
  \item Leia as silabas em voz alta:
  \begin{center}
  $\rightarrow$ LA  ~  LE  ~  LI  ~  LO  ~  LU
  \end{center}

  \item Nome de uma fruta com ``L'':
  \begin{center}
  \underline{\hspace{6cm}}
  \end{center}

  \item Desenhe algo com ``L'' e escreva o nome:
  \begin{center}
  \begin{tikzpicture}
    \draw[thick, dashed] (0,0) rectangle (5,3);
    \node[below] at (2.5,-0.3) {Nome: \underline{\hspace{4cm}}};
  \end{tikzpicture}
  \end{center}

  \item Circule a opcao correta:
  \begin{enumerate}[label=(\alph*)]
    \item O L tem som: ( ) nasal ( ) lateral ( ) vibrante
    \item A lingua toca: ( ) o nariz ( ) a lateral do palato ( ) os dentes
    \item O L aparece em: ( ) mesa ( ) sapo ( ) lua
  \end{enumerate}

  \item Leia e complete com a silaba correta:
  \begin{center}
  \sil{L}{A}ta \quad \underline{LE}ite \quad \sil{L}{I}mao \quad \underline{LO}bo \quad \underline{LU}a
  \end{center}
\end{enumerate}
\end{exercicio}

\begin{dica}
\textbf{Dica para o Professor:} O L e uma lateral. Peca que os alunos coloquem a lingua no ceu da boca pela lateral e soprem o ar por um dos lados. Compare com o R (vibrante no ceu da boca) e o T (oclusivo nos dentes). A diferenca entre L e R e crucial: no L, o ar sai pela lateral; no R, a lingua vibra no centro.
\end{dica}

%% ------------------------------------------------------------


% ============================================================
% CONSOLIDACAO DA LETRA L (R201.1) -- Missao 10/16
% ============================================================

\plancheta{Miss\~a{}o 10 --- Letra L}

\begin{missao}
\textbf{\faGamepad~Miss\~a{}o 10 de 16:} Consolide a letra L, a \'ultima do Bloco 1. O L muda no fim da palavra: sol vira sou!
\end{missao}

\begin{consolidacao}
\textbf{Atividade 1 --- Ca\c{c}ada ao L}

\instrucao{Circule o L e sublinhe as palavras com L.}

\begin{center}
\begin{tabular}{cccccc}
\fbox{L} & R & \fbox{L} & S & \fbox{L} & T \\
M & \fbox{L} & R & \fbox{L} & S & \fbox{L}
\end{tabular}
\end{center}
\begin{center}
\uline{lata} \quad \uline{lua} \quad rato \quad \uline{lobo} \quad sapo \quad \uline{lupa}
\end{center}

\caixapausa{Se precisar, pare aqui. Respire fundo e continue quando quiser.}

\textbf{Atividade 2 --- Escrita do L --- 4 formas}

\instrucao{Escreva o L cinco vezes em cada forma.}

\begin{center}
\textbf{Imprensa mai\'uscula:} \underline{\hspace{1cm}} \underline{\hspace{1cm}} \underline{\hspace{1cm}} \underline{\hspace{1cm}} \underline{\hspace{1cm}}\\[0.5em]
\textbf{Imprensa min\'uscula:} \underline{\hspace{1cm}} \underline{\hspace{1cm}} \underline{\hspace{1cm}} \underline{\hspace{1cm}} \underline{\hspace{1cm}}\\[0.5em]
\textbf{Cursiva mai\'uscula:} \underline{\hspace{1cm}} \underline{\hspace{1cm}} \underline{\hspace{1cm}} \underline{\hspace{1cm}} \underline{\hspace{1cm}}\\[0.5em]
\textbf{Cursiva min\'uscula:} \underline{\hspace{1cm}} \underline{\hspace{1cm}} \underline{\hspace{1cm}} \underline{\hspace{1cm}} \underline{\hspace{1cm}}
\end{center}

\textbf{Atividade 3 --- Jogo do L no Fim}

\instrucao{Complete com L e ou\c{c}a o som virar U no fim.}

\begin{center}
so\underline{~~~} \quad me\underline{~~~} \quad faro\underline{~~~} \quad ca\underline{~~~}ma \quad a\underline{~~~}mo\c{c}o
\end{center}
Diga em voz alta: \textbf{sol, mel, farol, calma, almo\c{c}o}

\caixapista{sol soa como sou. O L no fim da s\'ilaba vira um U r\'apido.}
\end{consolidacao}

\plancheta{Sons da Letra L --- IPA}

\begin{sonsdaletra}
Os sons da letra L no portugu\^es{} brasileiro (IPA --- Alfabeto Fon\'etico Internacional):

\somipa[dental]{L no in\'icio}{l}{ponta da l\'ingua no c\'eu da boca}{lata, lua, le\~a{}o}
\somipa[bico]{L no fim de s\'ilaba}{w}{vira som de U}{sol [sOw], mel [mEw], farol}
\end{sonsdaletra}

\plancheta{Fam\'ilia Sil\'abica da Letra L}

\begin{familiasilabica}
A fam\'ilia sil\'abica da letra L:

\begin{center}
lata \quad le\~a{}o \quad lim\~a{}o \quad lua \quad lobo \quad lupa \quad leite
\end{center}
\instrucao{Leia a fam\'ilia em voz alta, devagar. Depois cubra e tente lembrar.}
\end{familiasilabica}

\plancheta{Atividade de Som e S\'ilabas da Letra L}

\begin{somletra}
\textbf{1. Ca\c{c}a ao som.} Ou\c{c}a a palavra. Se tiver o som de L, marque o quadradinho:
\begin{center}
$\square$~LUA \quad $\square$~LEITE \quad $\square$~BOLA \quad $\square$~P\'E \quad $\square$~RATO \quad $\square$~M\~AO
\end{center}

\textbf{2. Tem ou n\~ao{} tem?} Rodeie SIM ou N\~ao:
\begin{center}
LUA --- (\textbf{SIM}) (\textbf{N\~AO}) \quad P\'E --- (\textbf{SIM}) (\textbf{N\~AO}) \\ LEITE --- (\textbf{SIM}) (\textbf{N\~AO}) \quad RATO --- (\textbf{SIM}) (\textbf{N\~AO})
\end{center}
\caixapausa{A l\'ingua toca o c\'eu da boca.}
\end{somletra}

\begin{somsilaba}
\textbf{3. Som das s\'ilabas.} Ou\c{c}a, bata palmas e conte as s\'ilabas:
\begin{center}
LU.A $\rightarrow$ \underline{\hspace{1cm}} s\'ilabas \quad BO.LA $\rightarrow$ \underline{\hspace{1cm}} s\'ilabas
\end{center}

\textbf{4. Qual s\'ilaba tem o som?} Circule a s\'ilaba com o som de L:
\begin{center}
\textbf{LU}.A \quad BO.\textbf{LA}
\end{center}
\caixapista{A consoante L se junta com a vogal: LA, LE, LI, LO, LU.}
\end{somsilaba}


\plancheta{Miss\~a{}o Motora --- Letra L}

\begin{motricidade}
\textbf{Movimento:} trace o \^angulo reto do L: des\c{c}a e vire \`a{} direita.
\begin{center}
\begin{tikzpicture}
  \draw[very thick, dashed, color=primary] (0,0) -- (0,-1.4);
  \draw[very thick, dashed, color=primary] (0,-1.4) -- (2.5,-1.4);
  \draw[->, color=accent] (0.15,-0.4) -- (0.15,-1.0);
  \draw[->, color=accent] (0.6,-1.25) -- (1.6,-1.25);
\end{tikzpicture}
\end{center}
\caixapista{O L \'e{} uma esquina: desce reto e vira na horizontal.}
\end{motricidade}

\plancheta{Ponte --- Da Letra L para a Pr\'oxima}

\begin{transicao}
Compare o L com o B, a pr\'oxima consoante:
\begin{center}
\textbf{lata} / \textbf{bata} \qquad \textbf{lobo} / \textbf{bolo} \qquad \textbf{lua} / \textbf{bula}
\end{center}
\instrucao{No L a l\'ingua sobe. No B os l\'abios fecham e explodem.}
\end{transicao}

\plancheta{Recompensa --- Miss\~a{}o 10}

\begin{recompensa}
\estrelas{3}\quad \ofensiva{10}\\[0.3em]\barraprogresso{10}{16}\\[0.5em]
\subiunivel{K3 --- Decodifica\c{c}\~a{}o}\\[0.3em]
\checklistdominio{reconhecer o L no in\'icio e o L-U no fim de s\'ilaba (sol)}
\end{recompensa}


%% LICAO 12: CONSOLIDACAO DO BLOCO 1
%% ------------------------------------------------------------
\section{Licao 12: Consolidacao do Bloco 1}

% Rotina neuroinclusiva e badges (R201)
\rotinasimples
\nivelKbadge{K1 --- Letras e Grafias}
\rotabadge{1A --- Recomposicao Intensiva}


\metadata{B-06}{Nivel 2 -- Silabico}{EF01LP02}{D1}{EF01LP02}{Resolucao CNE/CEB 7/2010, art. 12}

\plancheta{Folha de Pratica 2.6 --- CONSOLIDACAO BLOCO 1}
\begin{exercicio}[Folha de Pratica 2.6 --- CONSOLIDACAO BLOCO 1]
\noindent\textbf{Nome:} \underline{\hspace{3.2cm}} \quad \textbf{Data:} \underline{\hspace{1.6cm}} \quad \textbf{Nota:} \underline{\hspace{0.9cm}}/10

\subsection*{PARTE 1: RECONHECIMENTO}

\begin{enumerate}
  \item Formar silabas com as 5 consonantes e 5 vogais:
  \begin{center}
  \resizebox{\textwidth}{!}{%
\begin{tabular}{ccccc}
  \sil{M}{A} & \sil{S}{E} & \sil{T}{I} & \sil{R}{O} & \sil{L}{U} \\
  \sil{M}{O} & \sil{S}{A} & \sil{T}{U} & \sil{R}{E} & \sil{L}{I} \\
  \sil{M}{E} & \sil{S}{I} & \sil{T}{O} & \sil{R}{A} & \sil{L}{E} \\
  \end{tabular}}
  \end{center}

  \item Separe palavras em silabas:
  \begin{center}
  MESA $\rightarrow$ \underline{ME} \underline{SA} \quad SALA $\rightarrow$ \underline{SA} \underline{LA} \quad LUTO $\rightarrow$ \underline{LU} \underline{TO}\\[0.5em]
  RESTO $\rightarrow$ \underline{RE} \underline{STO} \quad MOLA $\rightarrow$ \underline{MO} \underline{LA} \quad TIRSO $\rightarrow$ \underline{TI} \underline{RSO}
  \end{center}

  \item Identifique a primeira letra de cada palavra:
  \begin{center}
  \begin{tabular}{ccccc}
  \textbf{M}esa & \textbf{S}apo & \textbf{T}igre & \textbf{R}ato & \textbf{L}ua \\
  \end{tabular}
  \end{center}

  \item Coloque $\checkmark$ nas palavras com a letra indicada:
  \begin{center}
  \begin{tabular}{ll}
  M: $\checkmark$ mesa \quad $\checkmark$ mola \quad $\square$ tinta & S: $\checkmark$ sapo \quad $\checkmark$ sala \quad $\square$ rato \\
  T: $\checkmark$ tinta \quad $\checkmark$ teto \quad $\square$ mesa & R: $\checkmark$ rato \quad $\checkmark$ roda \quad $\square$ lata \\
  L: $\checkmark$ lua \quad $\checkmark$ lata \quad $\square$ sapo &
  \end{tabular}
  \end{center}

  \item Escreva 3 palavras para cada letra estudada:
  \begin{center}
  M: \underline{\hspace{1.8cm}} \underline{\hspace{1.8cm}} \underline{\hspace{1.8cm}} \quad S: \underline{\hspace{1.8cm}} \underline{\hspace{1.8cm}} \underline{\hspace{1.8cm}}\\[0.5em]
  T: \underline{\hspace{1.8cm}} \underline{\hspace{1.8cm}} \underline{\hspace{1.8cm}} \quad R: \underline{\hspace{1.8cm}} \underline{\hspace{1.8cm}} \underline{\hspace{1.8cm}}\\[0.5em]
  L: \underline{\hspace{1.8cm}} \underline{\hspace{1.8cm}} \underline{\hspace{1.8cm}}
  \end{center}
\end{enumerate}

\subsection*{PARTE 2: PRODUCAO GUIADA}

\begin{enumerate}[resume]
  \item Formar palavras com as silabas:
  \begin{center}
  MA + SA = \underline{\hspace{2cm}} \quad SO + LA = \underline{\hspace{2cm}} \quad TU + RO = \underline{\hspace{2cm}}\\[0.5em]
  RI + MO = \underline{\hspace{2cm}} \quad LA + VA = \underline{\hspace{2cm}} \quad TO + RA = \underline{\hspace{2cm}}
  \end{center}

  \item Escreva as silabas de cada palavra:
  \begin{center}
  MACACO = \underline{\hspace{1cm}} \underline{\hspace{1cm}} \underline{\hspace{1cm}}\\
  SAPATO = \underline{\hspace{1cm}} \underline{\hspace{1cm}} \underline{\hspace{1cm}}\\
  MILHO = \underline{\hspace{1cm}} \underline{\hspace{1cm}}
  \end{center}

  \item Copie e complete:
  \begin{center}
  \underline{MA}sa \quad \sil{S}{A}po \quad \underline{TI}nta \quad \sil{R}{A}to \quad \underline{LU}a
  \end{center}

  \item Escreva o nome de uma pessoa para cada letra:
  \begin{center}
  M: \underline{\hspace{3cm}} \quad S: \underline{\hspace{3cm}} \quad T: \underline{\hspace{3cm}} \quad R: \underline{\hspace{3cm}} \quad L: \underline{\hspace{3cm}}
  \end{center}

  \item Forme 5 palavras novas usando as silabas estudadas:
  \begin{center}
  1. \underline{\hspace{3cm}} \quad 2. \underline{\hspace{3cm}} \quad 3. \underline{\hspace{3cm}} \quad 4. \underline{\hspace{3cm}} \quad 5. \underline{\hspace{3cm}}
  \end{center}
\end{enumerate}

\subsection*{PARTE 3: INTEGRACAO}

\begin{enumerate}[resume]
  \item Leia palavras completas em voz alta:
  \begin{center}
  $\rightarrow$ mesa  ~  sala  ~  luto  ~  mola  ~  tira  ~  rato  ~  sapo  ~  lua
  \end{center}

  \item Leia e circule a palavra correta:
  \begin{enumerate}[label=(\alph*)]
    \item ( ) mesa \quad ( ) mesa \quad ( ) tinta
    \item ( ) sapo \quad ( ) rato \quad ( ) lata
    \item ( ) tigre \quad ( ) mola \quad ( ) limao
  \end{enumerate}

  \item Avaliacao diagnostica:
  \begin{center}
  $\square$ Forma todas as silabas diretas\\
  $\square$ Separa palavras em silabas\\
  $\square$ Le palavras simples com fluencia\\
  $\square$ Identifica fonema-grafema corretamente\\
  $\square$ Diferencia as 5 consonantes estudadas
  \end{center}

  \item Leitura em voz alta: o professor aponta e o aluno le:
  \begin{center}
  mesa \quad sala \quad luto \quad mola \quad tira \quad rato \quad sapo \quad lua \quad tinta \quad teto \quad limao \quad lata
  \end{center}

  \item Escrita ditada: o professor dita e o aluno escreve:
  \begin{center}
  1. \underline{\hspace{4cm}} \quad 2. \underline{\hspace{4cm}} \quad 3. \underline{\hspace{4cm}} \quad 4. \underline{\hspace{4cm}} \quad 5. \underline{\hspace{4cm}}
  \end{center}
\end{enumerate}
\end{exercicio}

\begin{dica}
\textbf{Dica para o Professor:} Esta e a consolidacao do Bloco 1. Avalie se o aluno est\'a no n\'ivel sil\'abico: se forma s\'ilabas, separa palavras e l\^e palavras simples. Se houver dificuldades, retome as licoes individuais antes de avancar para as silabas diretas. Use o jogo ``Bingo das Silabas'' para tornar a avaliacao mais ludica.
\end{dica}

%% ============================================================
%% UNIDADE 3 -- SILABAS DIRETAS (Nivel 3 -- Silabico-Alfabetico)
%% ============================================================

\input{checkpoints/rastreio-u2}
\chapter{Unidade 3 --- Silabas Diretas}

%% ------------------------------------------------------------
%% LICAO 13: SILABAS COM A
%% ------------------------------------------------------------
\section{Licao 13: Silabas com A}

% Rotina neuroinclusiva e badges (R201)
\rotinasimples
\nivelKbadge{K2 --- Consciencia Fonologica}
\rotabadge{1B --- Segmentacao e Leitura Funcional}


\metadata{C-01}{Nivel 3 -- Silabico-Alfabetico}{EF01LP03}{D1}{EF01LP03}{Resolucao CNE/CEB 7/2010, art. 12}

\plancheta{Folha de Pratica 3.1 --- Silabas com A}
\begin{exercicio}[Folha de Pratica 3.1 --- Silabas com A]
\noindent\textbf{Nome:} \underline{\hspace{3.2cm}} \quad \textbf{Data:} \underline{\hspace{1.6cm}} \quad \textbf{Nota:} \underline{\hspace{0.9cm}}/10

\subsection*{PARTE 1: RECONHECIMENTO}

\begin{enumerate}
  \item Forme todas as silabas com A usando as consonantes:
  \begin{center}
  \resizebox{\textwidth}{!}{%
\begin{tabular}{ccccc}
  \sil{M}{A} & \sil{S}{A} & \sil{T}{A} & \sil{R}{A} & \sil{L}{A} \\
  \sil{B}{A} & \sil{C}{A} & \sil{D}{A} & \sil{F}{A} & \sil{G}{A} \\
  \sil{P}{A} & \sil{N}{A} & \sil{V}{A} & \sil{Z}{A} & \sil{J}{A} \\
  \end{tabular}}
  \end{center}

  \item Sublinhe as palavras que contem a silaba ``MA'':
  \begin{center}
  \uline{mesa} \quad gato \quad \uline{macaco} \quad sapo \quad \uline{mala} \quad rato \quad \uline{mapa}
  \end{center}

  \item Coloque $\checkmark$ nas palavras que comecam com ``SA'':
  \begin{center}
  \begin{tabular}{ll}
  $\checkmark$ \texttt{SAPO} & $\square$ \texttt{MESA} \\
  $\checkmark$ \texttt{SALA} & $\square$ \texttt{TINTA} \\
  $\checkmark$ \texttt{SACO} & $\square$ \texttt{RATO}
  \end{tabular}
  \end{center}

  \item Identifique a primeira silaba de cada palavra:
  \begin{center}
  MESA = \underline{\hspace{1cm}} \quad SAPO = \underline{\hspace{1cm}} \quad TAPA = \underline{\hspace{1cm}} \quad RASA = \underline{\hspace{1cm}}
  \end{center}

  \item Separe em silabas:
  \begin{center}
  MACACO = \underline{\hspace{1cm}} \underline{\hspace{1cm}} \underline{\hspace{1cm}} \quad SAPATO = \underline{\hspace{1cm}} \underline{\hspace{1cm}} \underline{\hspace{1cm}}
  \end{center}
\end{enumerate}

\subsection*{PARTE 2: PRODUCAO GUIADA}

\begin{enumerate}[resume]
  \item Forme palavras usando as silabas:
  \begin{center}
  MA + SA = \underline{\hspace{2cm}} \quad CA + SA = \underline{\hspace{2cm}} \quad TA + PA = \underline{\hspace{2cm}}\\[0.5em]
  RA + SA = \underline{\hspace{2cm}} \quad LA + VA = \underline{\hspace{2cm}} \quad SA + PO = \underline{\hspace{2cm}}
  \end{center}

  \item Escreva 3 palavras com a silaba ``MA'':
  \begin{center}
  1. \underline{\hspace{3cm}} \quad 2. \underline{\hspace{3cm}} \quad 3. \underline{\hspace{3cm}}
  \end{center}

  \item Escreva 3 palavras com a silaba ``SA'':
  \begin{center}
  1. \underline{\hspace{3cm}} \quad 2. \underline{\hspace{3cm}} \quad 3. \underline{\hspace{3cm}}
  \end{center}

  \item Copie as palavras e destaque a silaba ``A'':
  \begin{center}
  \textbf{M}\underline{\textbf{A}}sa \quad \textbf{S}\underline{\textbf{A}}po \quad \textbf{T}\underline{\textbf{A}}pa \quad \textbf{R}\underline{\textbf{A}}sa \quad \textbf{L}\underline{\textbf{A}}va
  \end{center}

  \item Escreva palavras que comecam com ``TA'':
  \begin{center}
  1. \underline{\hspace{3cm}} \quad 2. \underline{\hspace{3cm}} \quad 3. \underline{\hspace{3cm}}
  \end{center}
\end{enumerate}

\subsection*{PARTE 3: INTEGRACAO}

\begin{enumerate}[resume]
  \item Leia as palavras formadas em voz alta:
  \begin{center}
  $\rightarrow$ mesa  ~  casa  ~  tapa  ~  rasa  ~  lavo  ~  sapo
  \end{center}

  \item Escreva o nome de uma fruta que tem a silaba ``A'':
  \begin{center}
  \underline{\hspace{6cm}}
  \end{center}

  \item Desenhe algo e escreva o nome, destacando a silaba ``A'':
  \begin{center}
  \begin{tikzpicture}
    \draw[thick, dashed] (0,0) rectangle (5,3);
    \node[below] at (2.5,-0.3) {Nome: \underline{\hspace{4cm}}};
  \end{tikzpicture}
  \end{center}

  \item Circule a opcao correta:
  \begin{enumerate}[label=(\alph*)]
    \item A silaba ``MA'' aparece em: ( ) tinta ( ) mesa ( ) sapo
    \item A silaba ``SA'' aparece em: ( ) rato ( ) sapo ( ) lua
    \item A silaba ``TA'' aparece em: ( ) mesa ( ) lata ( ) sapo
  \end{enumerate}

  \item Leitura de palavras com A:
  \begin{center}
  $\rightarrow$ mesa  ~  casa  ~  tapa  ~  rasa  ~  sapo  ~  lata  ~  mala  ~  saco
  \end{center}
\end{enumerate}
\end{exercicio}

\begin{dica}
\textbf{Dica para o Professor:} Esta e a primeira licao de silabas diretas. Mostre que ao juntar uma consonante com A, formamos uma silaba que tem som proprio. Use o cubo de silabas ou cartoes para montar palavras. Comece sempre com a silaba ``MA'' e va adicionando outras consonantes.
\end{dica}

%% ------------------------------------------------------------
%% LICAO 14: SILABAS COM E
%% ------------------------------------------------------------
\section{Licao 14: Silabas com E}

% Rotina neuroinclusiva e badges (R201)
\rotinasimples
\nivelKbadge{K2 --- Consciencia Fonologica}
\rotabadge{1B --- Segmentacao e Leitura Funcional}


\metadata{C-02}{Nivel 3 -- Silabico-Alfabetico}{EF01LP03}{D1}{EF01LP03}{Resolucao CNE/CEB 7/2010, art. 12}

\plancheta{Folha de Pratica 3.2 --- Silabas com E}
\begin{exercicio}[Folha de Pratica 3.2 --- Silabas com E]
\noindent\textbf{Nome:} \underline{\hspace{3.2cm}} \quad \textbf{Data:} \underline{\hspace{1.6cm}} \quad \textbf{Nota:} \underline{\hspace{0.9cm}}/10

\subsection*{PARTE 1: RECONHECIMENTO}

\begin{enumerate}
  \item Forme todas as silabas com E:
  \begin{center}
  \resizebox{\textwidth}{!}{%
\begin{tabular}{ccccc}
  \sil{M}{E} & \sil{S}{E} & \sil{T}{E} & \sil{R}{E} & \sil{L}{E} \\
  \sil{B}{E} & \sil{C}{E} & \sil{D}{E} & \sil{F}{E} & \sil{G}{E} \\
  \sil{P}{E} & \sil{N}{E} & \sil{V}{E} & \sil{Z}{E} & \sil{J}{E} \\
  \end{tabular}}
  \end{center}

  \item Sublinhe as palavras que contem a silaba ``ME'':
  \begin{center}
  \uline{mesa} \quad gato \quad \uline{melo} \quad sapo \quad \uline{mete} \quad rato \quad \uline{meio}
  \end{center}

  \item Coloque $\checkmark$ nas palavras que comecam com ``SE'':
  \begin{center}
  \begin{tabular}{ll}
  $\checkmark$ \texttt{SETE} & $\square$ \texttt{MESA} \\
  $\checkmark$ \texttt{SECO} & $\square$ \texttt{TINTA} \\
  $\checkmark$ \texttt{SELO} & $\square$ \texttt{RATO}
  \end{tabular}
  \end{center}

  \item Identifique a primeira silaba de cada palavra:
  \begin{center}
  MESA = \underline{\hspace{1cm}} \quad SETE = \underline{\hspace{1cm}} \quad REDE = \underline{\hspace{1cm}} \quad TELA = \underline{\hspace{1cm}}
  \end{center}

  \item Separe em silabas:
  \begin{center}
  ELEFANTE = \underline{\hspace{1cm}} \underline{\hspace{1cm}} \underline{\hspace{1cm}} \underline{\hspace{1cm}}\\
  ESCOLA = \underline{\hspace{1cm}} \underline{\hspace{1cm}} \underline{\hspace{1cm}}
  \end{center}
\end{enumerate}

\subsection*{PARTE 2: PRODUCAO GUIADA}

\begin{enumerate}[resume]
  \item Forme palavras usando as silabas:
  \begin{center}
  ME + SA = \underline{\hspace{2cm}} \quad SE + TE = \underline{\hspace{2cm}} \quad RE + DE = \underline{\hspace{2cm}}\\[0.5em]
  LE + TE = \underline{\hspace{2cm}} \quad TE + LA = \underline{\hspace{2cm}} \quad DE + DO = \underline{\hspace{2cm}}
  \end{center}

  \item Escreva 3 palavras com a silaba ``ME'':
  \begin{center}
  1. \underline{\hspace{3cm}} \quad 2. \underline{\hspace{3cm}} \quad 3. \underline{\hspace{3cm}}
  \end{center}

  \item Escreva 3 palavras com a silaba ``SE'':
  \begin{center}
  1. \underline{\hspace{3cm}} \quad 2. \underline{\hspace{3cm}} \quad 3. \underline{\hspace{3cm}}
  \end{center}

  \item Copie as palavras e destaque a silaba ``E'':
  \begin{center}
  \textbf{M}\underline{\textbf{E}}sa \quad \textbf{S}\underline{\textbf{E}}te \quad \textbf{T}\underline{\textbf{E}}la \quad \textbf{R}\underline{\textbf{E}}de \quad \textbf{D}\underline{\textbf{E}}do
  \end{center}

  \item Escreva palavras que comecam com ``TE'':
  \begin{center}
  1. \underline{\hspace{3cm}} \quad 2. \underline{\hspace{3cm}} \quad 3. \underline{\hspace{3cm}}
  \end{center}
\end{enumerate}

\subsection*{PARTE 3: INTEGRACAO}

\begin{enumerate}[resume]
  \item Leia as palavras formadas em voz alta:
  \begin{center}
  $\rightarrow$ mesa  ~  sete  ~  rede  ~  leite  ~  tela  ~  dedo
  \end{center}

  \item Escreva o nome de um animal que tem a silaba ``E'':
  \begin{center}
  \underline{\hspace{6cm}}
  \end{center}

  \item Desenhe algo e escreva o nome, destacando a silaba ``E'':
  \begin{center}
  \begin{tikzpicture}
    \draw[thick, dashed] (0,0) rectangle (5,3);
    \node[below] at (2.5,-0.3) {Nome: \underline{\hspace{4cm}}};
  \end{tikzpicture}
  \end{center}

  \item Circule a opcao correta:
  \begin{enumerate}[label=(\alph*)]
    \item A silaba ``ME'' aparece em: ( ) tinta ( ) mesa ( ) sapo
    \item A silaba ``SE'' aparece em: ( ) rato ( ) sete ( ) lua
    \item A silaba ``TE'' aparece em: ( ) mesa ( ) tela ( ) sapo
  \end{enumerate}

  \item Leitura de palavras com E:
  \begin{center}
  $\rightarrow$ mesa  ~  sete  ~  rede  ~  leite  ~  tela  ~  dedo  ~  mela  ~  belo
  \end{center}
\end{enumerate}
\end{exercicio}

\begin{dica}
\textbf{Dica para o Professor:} Compare as palavras com ``A'' e ``E'': ``MASA'' vs ``MESA'', ``SAPA'' vs ``SEPA''. Mostre que a mudanca da vogal muda completamente o significado. Use rimas para fixar: ``ELE, BELE, DELE, MELE''. Os alunos devem perceber que ``E'' tem som mais fechado que ``A''.
\end{dica}

%% ------------------------------------------------------------
%% LICAO 15: SILABAS COM I
%% ------------------------------------------------------------
\section{Licao 15: Silabas com I}

% Rotina neuroinclusiva e badges (R201)
\rotinasimples
\nivelKbadge{K2 --- Consciencia Fonologica}
\rotabadge{1B --- Segmentacao e Leitura Funcional}


\metadata{C-03}{Nivel 3 -- Silabico-Alfabetico}{EF01LP03}{D1}{EF01LP03}{Resolucao CNE/CEB 7/2010, art. 12}

\plancheta{Folha de Pratica 3.3 --- Silabas com I}
\begin{exercicio}[Folha de Pratica 3.3 --- Silabas com I]
\noindent\textbf{Nome:} \underline{\hspace{3.2cm}} \quad \textbf{Data:} \underline{\hspace{1.6cm}} \quad \textbf{Nota:} \underline{\hspace{0.9cm}}/10

\subsection*{PARTE 1: RECONHECIMENTO}

\begin{enumerate}
  \item Forme todas as silabas com I:
  \begin{center}
  \resizebox{\textwidth}{!}{%
\begin{tabular}{ccccc}
  \sil{M}{I} & \sil{S}{I} & \sil{T}{I} & \sil{R}{I} & \sil{L}{I} \\
  \sil{B}{I} & \sil{C}{I} & \sil{D}{I} & \sil{F}{I} & \sil{G}{I} \\
  \sil{P}{I} & \sil{N}{I} & \sil{V}{I} & \sil{Z}{I} & \sil{J}{I} \\
  \end{tabular}}
  \end{center}

  \item Sublinhe as palavras que contem a silaba ``MI'':
  \begin{center}
  \uline{milo} \quad gato \quad \uline{mica} \quad sapo \quad \uline{mira} \quad rato \quad \uline{mio}
  \end{center}

  \item Coloque $\checkmark$ nas palavras que comecam com ``TI'':
  \begin{center}
  \begin{tabular}{ll}
  $\checkmark$ \texttt{TIPA} & $\square$ \texttt{MESA} \\
  $\checkmark$ \texttt{TIRA} & $\square$ \texttt{SAPATO} \\
  $\checkmark$ \texttt{TINTA} & $\square$ \texttt{RATO}
  \end{tabular}
  \end{center}

  \item Identifique a primeira silaba de cada palavra:
  \begin{center}
  MILO = \underline{\hspace{1cm}} \quad SILA = \underline{\hspace{1cm}} \quad TIPA = \underline{\hspace{1cm}} \quad RICA = \underline{\hspace{1cm}}
  \end{center}

  \item Separe em silabas:
  \begin{center}
  MELANCIA = \underline{\hspace{1cm}} \underline{\hspace{1cm}} \underline{\hspace{1cm}} \underline{\hspace{1cm}}\\
  SILABARIO = \underline{\hspace{1cm}} \underline{\hspace{1cm}} \underline{\hspace{1cm}} \underline{\hspace{1cm}}
  \end{center}
\end{enumerate}

\subsection*{PARTE 2: PRODUCAO GUIADA}

\begin{enumerate}[resume]
  \item Forme palavras usando as silabas:
  \begin{center}
  MI + LO = \underline{\hspace{2cm}} \quad SI + LA = \underline{\hspace{2cm}} \quad TI + PA = \underline{\hspace{2cm}}\\[0.5em]
  RI + CA = \underline{\hspace{2cm}} \quad LI + VA = \underline{\hspace{2cm}} \quad DI + CA = \underline{\hspace{2cm}}
  \end{center}

  \item Escreva 3 palavras com a silaba ``MI'':
  \begin{center}
  1. \underline{\hspace{3cm}} \quad 2. \underline{\hspace{3cm}} \quad 3. \underline{\hspace{3cm}}
  \end{center}

  \item Escreva 3 palavras com a silaba ``TI'':
  \begin{center}
  1. \underline{\hspace{3cm}} \quad 2. \underline{\hspace{3cm}} \quad 3. \underline{\hspace{3cm}}
  \end{center}

  \item Copie as palavras e destaque a silaba ``I'':
  \begin{center}
  \textbf{M}\underline{\textbf{I}}lo \quad \textbf{S}\underline{\textbf{I}}la \quad \textbf{T}\underline{\textbf{I}}pa \quad \textbf{R}\underline{\textbf{I}}ca \quad \textbf{D}\underline{\textbf{I}}ca
  \end{center}

  \item Escreva palavras que comecam com ``RI'':
  \begin{center}
  1. \underline{\hspace{3cm}} \quad 2. \underline{\hspace{3cm}} \quad 3. \underline{\hspace{3cm}}
  \end{center}
\end{enumerate}

\subsection*{PARTE 3: INTEGRACAO}

\begin{enumerate}[resume]
  \item Leia as palavras formadas em voz alta:
  \begin{center}
  $\rightarrow$ milo  ~  sila  ~  tipa  ~  rica  ~  liva  ~  dica
  \end{center}

  \item Escreva o nome de uma pessoa que tem a silaba ``I'' no nome:
  \begin{center}
  \underline{\hspace{6cm}}
  \end{center}

  \item Desenhe algo e escreva o nome, destacando a silaba ``I'':
  \begin{center}
  \begin{tikzpicture}
    \draw[thick, dashed] (0,0) rectangle (5,3);
    \node[below] at (2.5,-0.3) {Nome: \underline{\hspace{4cm}}};
  \end{tikzpicture}
  \end{center}

  \item Circule a opcao correta:
  \begin{enumerate}[label=(\alph*)]
    \item A silaba ``MI'' aparece em: ( ) tinta ( ) milo ( ) sapo
    \item A silaba ``TI'' aparece em: ( ) mesa ( ) tinta ( ) sapo
    \item A silaba ``RI'' aparece em: ( ) rato ( ) rica ( ) lua
  \end{enumerate}

  \item Leitura de palavras com I:
  \begin{center}
  $\rightarrow$ milo  ~  sila  ~  tipa  ~  rica  ~  liva  ~  dica  ~  mica  ~  pita
  \end{center}
\end{enumerate}
\end{exercicio}

\begin{dica}
\textbf{Dica para o Professor:} A silaba ``I'' e a mais aguda. Compare ``MA'' (A aberto) com ``MI'' (I estreito). Use pares minimos: ``MALA'' vs ``MILA'', ``PATA'' vs ``PITA''. Os alunos devem perceber que a vogal ``I'' deixa a silaba mais aguda e estreita.
\end{dica}

%% ------------------------------------------------------------
%% LICAO 16: SILABAS COM O
%% ------------------------------------------------------------
\section{Licao 16: Silabas com O}

% Rotina neuroinclusiva e badges (R201)
\rotinasimples
\nivelKbadge{K2 --- Consciencia Fonologica}
\rotabadge{1B --- Segmentacao e Leitura Funcional}


\metadata{C-04}{Nivel 3 -- Silabico-Alfabetico}{EF01LP03}{D1}{EF01LP03}{Resolucao CNE/CEB 7/2010, art. 12}

\plancheta{Folha de Pratica 3.4 --- Silabas com O}
\begin{exercicio}[Folha de Pratica 3.4 --- Silabas com O]
\noindent\textbf{Nome:} \underline{\hspace{3.2cm}} \quad \textbf{Data:} \underline{\hspace{1.6cm}} \quad \textbf{Nota:} \underline{\hspace{0.9cm}}/10

\subsection*{PARTE 1: RECONHECIMENTO}

\begin{enumerate}
  \item Forme todas as silabas com O:
  \begin{center}
  \resizebox{\textwidth}{!}{%
\begin{tabular}{ccccc}
  \sil{M}{O} & \sil{S}{O} & \sil{T}{O} & \sil{R}{O} & \sil{L}{O} \\
  \sil{B}{O} & \sil{C}{O} & \sil{D}{O} & \sil{F}{O} & \sil{G}{O} \\
  \sil{P}{O} & \sil{N}{O} & \sil{V}{O} & \sil{Z}{O} & \sil{J}{O} \\
  \end{tabular}}
  \end{center}

  \item Sublinhe as palavras que contem a silaba ``MO'':
  \begin{center}
  \uline{mola} \quad gato \quad \uline{moto} \quad sapo \quad \uline{modo} \quad rato \quad \uline{mopo}
  \end{center}

  \item Coloque $\checkmark$ nas palavras que comecam com ``SO'':
  \begin{center}
  \begin{tabular}{ll}
  $\checkmark$ \texttt{SOLA} & $\square$ \texttt{MESA} \\
  $\checkmark$ \texttt{SOPO} & $\square$ \texttt{TINTA} \\
  $\checkmark$ \texttt{SOM} & $\square$ \texttt{RATO}
  \end{tabular}
  \end{center}

  \item Identifique a primeira silaba de cada palavra:
  \begin{center}
  MOLA = \underline{\hspace{1cm}} \quad SOLA = \underline{\hspace{1cm}} \quad TOPA = \underline{\hspace{1cm}} \quad ROSA = \underline{\hspace{1cm}}
  \end{center}

  \item Separe em silabas:
  \begin{center}
  BONDE = \underline{\hspace{1cm}} \underline{\hspace{1cm}} \quad MOTOR = \underline{\hspace{1cm}} \underline{\hspace{1cm}}\\
  COPOS = \underline{\hspace{1cm}} \underline{\hspace{1cm}} \quad TOCO = \underline{\hspace{1cm}} \underline{\hspace{1cm}}
  \end{center}
\end{enumerate}

\subsection*{PARTE 2: PRODUCAO GUIADA}

\begin{enumerate}[resume]
  \item Forme palavras usando as silabas:
  \begin{center}
  MO + LA = \underline{\hspace{2cm}} \quad SO + LA = \underline{\hspace{2cm}} \quad TO + PA = \underline{\hspace{2cm}}\\[0.5em]
  RO + SA = \underline{\hspace{2cm}} \quad LO + CA = \underline{\hspace{2cm}} \quad BO + LA = \underline{\hspace{2cm}}
  \end{center}

  \item Escreva 3 palavras com a silaba ``MO'':
  \begin{center}
  1. \underline{\hspace{3cm}} \quad 2. \underline{\hspace{3cm}} \quad 3. \underline{\hspace{3cm}}
  \end{center}

  \item Escreva 3 palavras com a silaba ``RO'':
  \begin{center}
  1. \underline{\hspace{3cm}} \quad 2. \underline{\hspace{3cm}} \quad 3. \underline{\hspace{3cm}}
  \end{center}

  \item Copie as palavras e destaque a silaba ``O'':
  \begin{center}
  \textbf{M}\underline{\textbf{O}}la \quad \textbf{S}\underline{\textbf{O}}la \quad \textbf{T}\underline{\textbf{O}}pa \quad \textbf{R}\underline{\textbf{O}}sa \quad \textbf{B}\underline{\textbf{O}}la
  \end{center}

  \item Escreva palavras que comecam com ``LO'':
  \begin{center}
  1. \underline{\hspace{3cm}} \quad 2. \underline{\hspace{3cm}} \quad 3. \underline{\hspace{3cm}}
  \end{center}
\end{enumerate}

\subsection*{PARTE 3: INTEGRACAO}

\begin{enumerate}[resume]
  \item Leia as palavras formadas em voz alta:
  \begin{center}
  $\rightarrow$ mola  ~  sola  ~  topa  ~  rosa  ~  loca  ~  bola
  \end{center}

  \item Escreva o nome de uma cor que tem a silaba ``O'':
  \begin{center}
  \underline{\hspace{6cm}}
  \end{center}

  \item Desenhe algo e escreva o nome, destacando a silaba ``O'':
  \begin{center}
  \begin{tikzpicture}
    \draw[thick, dashed] (0,0) rectangle (5,3);
    \node[below] at (2.5,-0.3) {Nome: \underline{\hspace{4cm}}};
  \end{tikzpicture}
  \end{center}

  \item Circule a opcao correta:
  \begin{enumerate}[label=(\alph*)]
    \item A silaba ``MO'' aparece em: ( ) tinta ( ) mola ( ) sapo
    \item A silaba ``SO'' aparece em: ( ) rato ( ) sola ( ) lua
    \item A silaba ``RO'' aparece em: ( ) rosa ( ) mesa ( ) lata
  \end{enumerate}

  \item Leitura de palavras com O:
  \begin{center}
  $\rightarrow$ mola  ~  sola  ~  topa  ~  rosa  ~  loca  ~  bola  ~  moto  ~  copo
  \end{center}
\end{enumerate}
\end{exercicio}

\begin{dica}
\textbf{Dica para o Professor:} A silaba ``O'' e a mais redonda. Compare ``MA'' (A aberto) com ``MO'' (O redondo). Use pares minimos: ``MALA'' vs ``MOLA'', ``PATA'' vs ``POTA''. Os alunos devem perceber que ``O'' redonda a silaba. Use o espelho para mostrar a forma dos labios.
\end{dica}

%% ------------------------------------------------------------
%% LICAO 17: SILABAS COM U
%% ------------------------------------------------------------
\section{Licao 17: Silabas com U}

% Rotina neuroinclusiva e badges (R201)
\rotinasimples
\nivelKbadge{K2 --- Consciencia Fonologica}
\rotabadge{1B --- Segmentacao e Leitura Funcional}


\metadata{C-05}{Nivel 3 -- Silabico-Alfabetico}{EF01LP03}{D1}{EF01LP03}{Resolucao CNE/CEB 7/2010, art. 12}

\plancheta{Folha de Pratica 3.5 --- Silabas com U}
\begin{exercicio}[Folha de Pratica 3.5 --- Silabas com U]
\noindent\textbf{Nome:} \underline{\hspace{3.2cm}} \quad \textbf{Data:} \underline{\hspace{1.6cm}} \quad \textbf{Nota:} \underline{\hspace{0.9cm}}/10

\subsection*{PARTE 1: RECONHECIMENTO}

\begin{enumerate}
  \item Forme todas as silabas com U:
  \begin{center}
  \resizebox{\textwidth}{!}{%
\begin{tabular}{ccccc}
  \sil{M}{U} & \sil{S}{U} & \sil{T}{U} & \sil{R}{U} & \sil{L}{U} \\
  \sil{B}{U} & \sil{C}{U} & \sil{D}{U} & \sil{F}{U} & \sil{G}{U} \\
  \sil{P}{U} & \sil{N}{U} & \sil{V}{U} & \sil{Z}{U} & \sil{J}{U} \\
  \end{tabular}}
  \end{center}

  \item Sublinhe as palavras que contem a silaba ``MU'':
  \begin{center}
  \uline{mula} \quad gato \quad \uline{muro} \quad sapo \quad \uline{mudo} \quad rato \quad \uline{musgo}
  \end{center}

  \item Coloque $\checkmark$ nas palavras que comecam com ``SU'':
  \begin{center}
  \begin{tabular}{ll}
  $\checkmark$ \texttt{SUCO} & $\square$ \texttt{MESA} \\
  $\checkmark$ \texttt{SUMO} & $\square$ \texttt{TINTA} \\
  $\checkmark$ \texttt{SUSTO} & $\square$ \texttt{RATO}
  \end{tabular}
  \end{center}

  \item Identifique a primeira silaba de cada palavra:
  \begin{center}
  MULA = \underline{\hspace{1cm}} \quad SUCO = \underline{\hspace{1cm}} \quad TUDO = \underline{\hspace{1cm}} \quad RUDA = \underline{\hspace{1cm}}
  \end{center}

  \item Separe em silabas:
  \begin{center}
  UNIVERSO = \underline{\hspace{1cm}} \underline{\hspace{1cm}} \underline{\hspace{1cm}} \underline{\hspace{1cm}}\\
  MUSICA = \underline{\hspace{1cm}} \underline{\hspace{1cm}} \underline{\hspace{1cm}}
  \end{center}
\end{enumerate}

\subsection*{PARTE 2: PRODUCAO GUIADA}

\begin{enumerate}[resume]
  \item Forme palavras usando as silabas:
  \begin{center}
  MU + LA = \underline{\hspace{2cm}} \quad SU + CO = \underline{\hspace{2cm}} \quad TU + DO = \underline{\hspace{2cm}}\\[0.5em]
  RU + DA = \underline{\hspace{2cm}} \quad LU + TA = \underline{\hspace{2cm}} \quad BU + CA = \underline{\hspace{2cm}}
  \end{center}

  \item Escreva 3 palavras com a silaba ``MU'':
  \begin{center}
  1. \underline{\hspace{3cm}} \quad 2. \underline{\hspace{3cm}} \quad 3. \underline{\hspace{3cm}}
  \end{center}

  \item Escreva 3 palavras com a silaba ``TU'':
  \begin{center}
  1. \underline{\hspace{3cm}} \quad 2. \underline{\hspace{3cm}} \quad 3. \underline{\hspace{3cm}}
  \end{center}

  \item Copie as palavras e destaque a silaba ``U'':
  \begin{center}
  \textbf{M}\underline{\textbf{U}}la \quad \textbf{S}\underline{\textbf{U}}co \quad \textbf{T}\underline{\textbf{U}}do \quad \textbf{R}\underline{\textbf{U}}da \quad \textbf{B}\underline{\textbf{U}}ca
  \end{center}

  \item Escreva palavras que comecam com ``LU'':
  \begin{center}
  1. \underline{\hspace{3cm}} \quad 2. \underline{\hspace{3cm}} \quad 3. \underline{\hspace{3cm}}
  \end{center}
\end{enumerate}

\subsection*{PARTE 3: INTEGRACAO}

\begin{enumerate}[resume]
  \item Leia as palavras formadas em voz alta:
  \begin{center}
  $\rightarrow$ mula  ~  suco  ~  tudo  ~  ruda  ~  luta  ~  buca
  \end{center}

  \item Escreva o nome de um animal que tem a silaba ``U'':
  \begin{center}
  \underline{\hspace{6cm}}
  \end{center}

  \item Desenhe algo e escreva o nome, destacando a silaba ``U'':
  \begin{center}
  \begin{tikzpicture}
    \draw[thick, dashed] (0,0) rectangle (5,3);
    \node[below] at (2.5,-0.3) {Nome: \underline{\hspace{4cm}}};
  \end{tikzpicture}
  \end{center}

  \item Circule a opcao correta:
  \begin{enumerate}[label=(\alph*)]
    \item A silaba ``MU'' aparece em: ( ) tinta ( ) mula ( ) sapo
    \item A silaba ``SU'' aparece em: ( ) rato ( ) suco ( ) lua
    \item A silaba ``TU'' aparece em: ( ) mesa ( ) tudo ( ) lata
  \end{enumerate}

  \item Leitura de palavras com U:
  \begin{center}
  $\rightarrow$ mula  ~  suco  ~  tudo  ~  ruda  ~  luta  ~  buca  ~  muro  ~  punho
  \end{center}
\end{enumerate}
\end{exercicio}

\begin{dica}
\textbf{Dica para o Professor:} A silaba ``U'' e a mais grave e com a boca mais fechada. Compare ``MA'' (A aberto) com ``MU'' (U fechado). Use pares minimos: ``MALA'' vs ``MULA'', ``PATA'' vs ``CURA''. Os alunos devem perceber que ``U'' fecha e torna grave a silaba.
\end{dica}

%% ------------------------------------------------------------
%% LICAO 18: JOGO DE SILABAS
%% ------------------------------------------------------------
\section{Licao 18: Jogo de Silabas --- Atividade Ludica}

% Rotina neuroinclusiva e badges (R201)
\rotinasimples
\nivelKbadge{K2 --- Consciencia Fonologica}
\rotabadge{1B --- Segmentacao e Leitura Funcional}


\metadata{C-06}{Nivel 3 -- Silabico-Alfabetico}{EF01LP03}{D1}{EF01LP03}{Resolucao CNE/CEB 7/2010, art. 12}

\plancheta{Caca-Palavras das Silabas}
\begin{exercicio}[Caca-Palavras das Silabas]
Encontre palavras formadas por silabas diretas na tabela abaixo:

\begin{center}
\begin{tabular}{|c|c|c|c|c|c|c|c|c|c|c|}
\hline
M & A & S & T & R & L & B & C & D & F & G \\
\hline
E & O & A & E & I & U & O & A & E & I & U \\
\hline
S & P & L & T & R & M & S & C & D & F & G \\
\hline
A & O & A & O & O & U & A & A & E & I & U \\
\hline
\end{tabular}
\end{center}

\textbf{Palavras encontradas:}
\begin{enumerate}
  \item \underline{\hspace{4cm}}
  \item \underline{\hspace{4cm}}
  \item \underline{\hspace{4cm}}
  \item \underline{\hspace{4cm}}
  \item \underline{\hspace{4cm}}
  \item \underline{\hspace{4cm}}
  \item \underline{\hspace{4cm}}
  \item \underline{\hspace{4cm}}
  \item \underline{\hspace{4cm}}
  \item \underline{\hspace{4cm}}
\end{enumerate}

\textbf{Dica do jogo:} Procure silabas na horizontal (da esquerda para a direita) e monte palavras com 2 silabas. Exemplo: M + A = MA, S + A = SA, depois MA + SA = MASA.

\vspace{1em}

\textbf{Desafio adicional:} Forme palavras de 3 silabas usando as silabas encontradas.
\begin{enumerate}[resume]
  \item \underline{\hspace{4cm}}
  \item \underline{\hspace{4cm}}
  \item \underline{\hspace{4cm}}
\end{enumerate}
\end{exercicio}

\begin{dica}
\textbf{Dica para o Professor:} Este jogo consolida as silabas diretas de forma ludica. Permita que os alunos trabalhem em duplas e usem canetas coloridas para marcar as silabas encontradas. Ao final, cada dupla apresenta as palavras encontradas e le em voz alta. Premie as duplas que encontrarem mais palavras corretas.
\end{dica}

%% ============================================================
%% UNIDADE 4 -- PRIMEIRAS FRASES (Nivel 4 -- Silabico-Alfabetico Avancado)
%% ============================================================

\input{checkpoints/rastreio-u3}
\chapter{Unidade 4 --- Primeiras Frases}

%% ------------------------------------------------------------
%% LICAO 19: FORMANDO FRASES SIMPLES
%% ------------------------------------------------------------
\section{Licao 19: Formando Frases Simples}

% Rotina neuroinclusiva e badges (R201)
\rotinasimples
\nivelKbadge{K3 --- Decodificacao}
\rotabadge{1C --- Consolidacao e Expansao}


\metadata{D-01}{Nivel 4 -- Silabico-Alfabetico Avancado}{EF01LP04}{D2}{EF01LP04}{Resolucao CNE/CEB 7/2010, art. 12}

\textbf{Estrutura de uma frase simples:}
\begin{center}
\textbf{ARTIGO} + \textbf{SUBSTANTIVO} + \textbf{VERBO}
\end{center}

\textbf{Exemplos:}
\begin{enumerate}[leftmargin=2cm]
  \item \textcolor{primary}{O} + \textcolor{secondary}{menino} + \textcolor{accent}{corre.}
  \item \textcolor{primary}{A} + \textcolor{secondary}{menina} + \textcolor{accent}{pula.}
  \item \textcolor{primary}{O} + \textcolor{secondary}{gato} + \textcolor{accent}{mia.}
  \item \textcolor{primary}{A} + \textcolor{secondary}{cachorra} + \textcolor{accent}{late.}
\end{enumerate}

\plancheta{Folha de Pratica 4.1 --- Frases Simples}
\begin{exercicio}[Folha de Pratica 4.1 --- Frases Simples]
\noindent\textbf{Nome:} \underline{\hspace{3.2cm}} \quad \textbf{Data:} \underline{\hspace{1.6cm}} \quad \textbf{Nota:} \underline{\hspace{0.9cm}}/10

\subsection*{PARTE 1: RECONHECIMENTO}

\begin{enumerate}
  \item Leia as frases e identifique o sujeito e o verbo:
  \begin{enumerate}[label=\alph*)]
    \item O menino corre. $\rightarrow$ Sujeito: \underline{\hspace{3cm}} Verbo: \underline{\hspace{2cm}}
    \item A menina pula. $\rightarrow$ Sujeito: \underline{\hspace{3cm}} Verbo: \underline{\hspace{2cm}}
    \item O gato mia. $\rightarrow$ Sujeito: \underline{\hspace{3cm}} Verbo: \underline{\hspace{2cm}}
  \end{enumerate}

  \item Sublinhe o verbo em cada frase:
  \begin{enumerate}[label=\alph*)]
    \item O menino corre.
    \item A cachorra late.
    \item O peixe nada.
  \end{enumerate}

  \item Identifique se a frase usa ``O'' ou ``A'':
  \begin{enumerate}[label=\alph*)]
    \item \underline{\hspace{0.5cm}} menino corre. $\rightarrow$ ( ) O ( ) A
    \item \underline{\hspace{0.5cm}} menina pula. $\rightarrow$ ( ) O ( ) A
    \item \underline{\hspace{0.5cm}} gato mia. $\rightarrow$ ( ) O ( ) A
  \end{enumerate}

  \item Separe as palavras de cada frase em silabas:
  \begin{enumerate}[label=\alph*)]
    \item O menino corre. $\rightarrow$ \underline{\hspace{1cm}} | \underline{\hspace{1cm}} \underline{\hspace{1cm}} \underline{\hspace{1cm}} | \underline{\hspace{1cm}} \underline{\hspace{1cm}}
    \item A menina pula. $\rightarrow$ \underline{\hspace{1cm}} | \underline{\hspace{1cm}} \underline{\hspace{1cm}} \underline{\hspace{1cm}} | \underline{\hspace{1cm}} \underline{\hspace{1cm}}
  \end{enumerate}

  \item Leia em voz alta e conte quantas palavras tem cada frase:
  \begin{enumerate}[label=\alph*)]
    \item ``O menino corre.'' = \underline{\hspace{0.5cm}} palavras
    \item ``A cachorra late.'' = \underline{\hspace{0.5cm}} palavras
    \item ``O peixe nada.'' = \underline{\hspace{0.5cm}} palavras
  \end{enumerate}
\end{enumerate}

\subsection*{PARTE 2: PRODUCAO GUIADA}

\begin{enumerate}[resume]
  \item Complete a frase escolhendo a palavra correta:
  \begin{enumerate}[label=\alph*)]
    \item O menino \underline{\hspace{3cm}} (corre / pula / mia)
    \item A menina \underline{\hspace{3cm}} (late / corre / pula)
    \item O gato \underline{\hspace{3cm}} (pula / corre / mia)
    \item A cachorra \underline{\hspace{3cm}} (mia / corre / late)
  \end{enumerate}

  \item Forme frases juntando as palavras:
  \begin{enumerate}[label=\alph*)]
    \item O + sapo + pula = \underline{\hspace{6cm}}
    \item A + rata + corre = \underline{\hspace{6cm}}
    \item O + peixe + nada = \underline{\hspace{6cm}}
  \end{enumerate}

  \item Escreva 3 frases usando ``O'' ou ``A'' + um animal + um verbo:
  \begin{enumerate}[label=\alph*)]
    \item \underline{\hspace{8cm}}
    \item \underline{\hspace{8cm}}
    \item \underline{\hspace{8cm}}
  \end{enumerate}

  \item Copie as frases abaixo:
  \begin{enumerate}[label=\alph*)]
    \item O menino corre.
    \item A menina pula.
    \item O gato mia.
    \item A cachorra late.
  \end{enumerate}

  \item Escreva a frase completa com a pontuacao:
  \begin{enumerate}[label=\alph*)]
    \item o menino corre \underline{\hspace{1cm}}
    \item a menina pula \underline{\hspace{1cm}}
    \item o gato mia \underline{\hspace{1cm}}
  \end{enumerate}
\end{enumerate}

\subsection*{PARTE 3: INTEGRACAO}

\begin{enumerate}[resume]
  \item Leia as frases em voz alta para o professor:
  \begin{center}
  $\rightarrow$ O menino corre.\\
  $\rightarrow$ A menina pula.\\
  $\rightarrow$ O gato mia.\\
  $\rightarrow$ A cachorra late.
  \end{center}

  \item Desenhe uma cena para cada frase e escreva a legenda:
  \begin{center}
  \begin{tikzpicture}
    \draw[thick, dashed] (0,0) rectangle (4,3);
    \draw[thick, dashed] (4.5,0) rectangle (8.5,3);
    \draw[thick, dashed] (9,0) rectangle (13,3);
    \node[below] at (2,-0.3) {Legenda: \underline{\hspace{3cm}}};
    \node[below] at (6.5,-0.3) {Legenda: \underline{\hspace{3cm}}};
    \node[below] at (11,-0.3) {Legenda: \underline{\hspace{3cm}}};
  \end{tikzpicture}
  \end{center}

  \item Circule a opcao correta:
  \begin{enumerate}[label=(\alph*)]
    \item Na frase ``O menino corre'', o verbo e: ( ) menino ( ) corre ( ) o
    \item ``A menina pula'' comeca com: ( ) O ( ) A ( ) E
    \item Uma frase precisa de pelo menos: ( ) 1 ( ) 2 ( ) 3 palavras
  \end{enumerate}

  \item Avaliacao diagnostica:
  \begin{center}
  $\square$ Le frases simples com fluencia\\
  $\square$ Identifica sujeito e verbo\\
  $\square$ Usa ``O'' e ``A'' corretamente\\
  $\square$ Escreve frases com pontuacao\\
  $\square$ Compreende o significado da frase
  \end{center}

  \item Leitura em voz alta: o professor le uma frase e o aluno desenha:
  \begin{enumerate}[label=\alph*)]
    \item ``O gato pula.''
    \item ``A menina corre.''
    \item ``O sapo mia.''
  \end{enumerate}
\end{enumerate}
\end{exercicio}

\begin{dica}
\textbf{Dica para o Professor:} Esta e a primeira licao de frases. Comece com frases muito simples (3 palavras). Use a estrutura ``O/A + substantivo + verbo''. Peca que os alunos montem frases com cartoes antes de escrever. Use historias curtas para motivar a leitura de frases.
\end{dica}

%% ------------------------------------------------------------
%% LICAO 20: FRASES COM PREDICADO
%% ------------------------------------------------------------
\section{Licao 20: Frases com Predicado}

% Rotina neuroinclusiva e badges (R201)
\rotinasimples
\nivelKbadge{K3 --- Decodificacao}
\rotabadge{1C --- Consolidacao e Expansao}


\metadata{D-02}{Nivel 4 -- Silabico-Alfabetico Avancado}{EF01LP04}{D2}{EF01LP04}{Resolucao CNE/CEB 7/2010, art. 12}

\textbf{Estrutura com predicado:}
\begin{center}
\textbf{SUJEITO} + \textbf{PREDICADO}
\end{center}

\textbf{Exemplos:}
\begin{enumerate}[leftmargin=2cm]
  \item O menino + \textcolor{accent}{est\'a brincando.}
  \item A menina + \textcolor{accent}{esta estudando.}
  \item O gato + \textcolor{accent}{esta dormindo.}
  \item A cachorra + \textcolor{accent}{esta comendo.}
\end{enumerate}

\plancheta{Folha de Pratica 4.2 --- Frases com Predicado}
\begin{exercicio}[Folha de Pratica 4.2 --- Frases com Predicado]
\noindent\textbf{Nome:} \underline{\hspace{3.2cm}} \quad \textbf{Data:} \underline{\hspace{1.6cm}} \quad \textbf{Nota:} \underline{\hspace{0.9cm}}/10

\subsection*{PARTE 1: RECONHECIMENTO}

\begin{enumerate}
  \item Leia as frases e identifique o predicado:
  \begin{enumerate}[label=\alph*)]
    \item O menino est\'a brincando. $\rightarrow$ Predicado: \underline{\hspace{4cm}}
    \item A menina est\'a estudando. $\rightarrow$ Predicado: \underline{\hspace{4cm}}
    \item O gato est\'a dormindo. $\rightarrow$ Predicado: \underline{\hspace{4cm}}
  \end{enumerate}

  \item Sublinhe o sujeito em cada frase:
  \begin{enumerate}[label=\alph*)]
    \item O menino est\'a brincando.
    \item A cachorra est\'a comendo.
    \item O peixe est\'a nadando.
  \end{enumerate}

  \item Identifique se o predicado usa ``esta'' ou ``esta'':
  \begin{enumerate}[label=\alph*)]
    \item O menino \underline{\hspace{1cm}} brincando. $\rightarrow$ ( ) est\'a ( ) estavam
    \item A menina \underline{\hspace{1cm}} estudando. $\rightarrow$ ( ) est\'a ( ) estavam
    \item O gato \underline{\hspace{1cm}} dormindo. $\rightarrow$ ( ) est\'a ( ) estavam
  \end{enumerate}

  \item Separe sujeito e predicado:
  \begin{enumerate}[label=\alph*)]
    \item ``O menino est\'a brincando.'' $\rightarrow$ Sujeito: \underline{\hspace{3cm}} Predicado: \underline{\hspace{3cm}}
    \item ``A menina esta estudando.'' $\rightarrow$ Sujeito: \underline{\hspace{3cm}} Predicado: \underline{\hspace{3cm}}
  \end{enumerate}

  \item Leia em voz alta e conte quantas palavras tem cada frase:
  \begin{enumerate}[label=\alph*)]
    \item ``O menino esta brincando.'' = \underline{\hspace{0.5cm}} palavras
    \item ``A menina esta estudando.'' = \underline{\hspace{0.5cm}} palavras
    \item ``O gato esta dormindo.'' = \underline{\hspace{0.5cm}} palavras
  \end{enumerate}
\end{enumerate}

\subsection*{PARTE 2: PRODUCAO GUIADA}

\begin{enumerate}[resume]
  \item Unir sujeito e predicado:
  \begin{enumerate}[label=\alph*)]
    \item O menino $\rightarrow$ \underline{\hspace{4cm}} brincando
    \item A menina $\rightarrow$ \underline{\hspace{4cm}} estudando
    \item O gato $\rightarrow$ \underline{\hspace{4cm}} dormindo
    \item A cachorra $\rightarrow$ \underline{\hspace{4cm}} comendo
  \end{enumerate}

  \item Complete com o predicado adequado:
  \begin{enumerate}[label=\alph*)]
    \item O sapo \underline{\hspace{3cm}} (esta pulando / esta comendo / esta voando)
    \item A formiga \underline{\hspace{3cm}} (esta andando / esta dormindo / esta voando)
    \item O peixe \underline{\hspace{3cm}} (esta correndo / esta nadando / esta pULando)
  \end{enumerate}

  \item Escreva 3 frases com predicado usando ``esta'':
  \begin{enumerate}[label=\alph*)]
    \item \underline{\hspace{8cm}}
    \item \underline{\hspace{8cm}}
    \item \underline{\hspace{8cm}}
  \end{enumerate}

  \item Copie as frases abaixo:
  \begin{enumerate}[label=\alph*)]
    \item O menino esta brincando.
    \item A menina esta estudando.
    \item O gato esta dormindo.
    \item A cachorra esta comendo.
  \end{enumerate}

  \item Escreva a frase completa com pontuacao:
  \begin{enumerate}[label=\alph*)]
    \item o menino esta brincando \underline{\hspace{1cm}}
    \item a menina esta estudando \underline{\hspace{1cm}}
    \item o gato esta dormindo \underline{\hspace{1cm}}
  \end{enumerate}
\end{enumerate}

\subsection*{PARTE 3: INTEGRACAO}

\begin{enumerate}[resume]
  \item Leia as frases em voz alta para o professor:
  \begin{center}
  $\rightarrow$ O menino esta brincando.\\
  $\rightarrow$ A menina esta estudando.\\
  $\rightarrow$ O gato esta dormindo.\\
  $\rightarrow$ A cachorra esta comendo.
  \end{center}

  \item Desenhe uma cena para cada frase e escreva a legenda:
  \begin{center}
  \begin{tikzpicture}
    \draw[thick, dashed] (0,0) rectangle (4,3);
    \draw[thick, dashed] (4.5,0) rectangle (8.5,3);
    \draw[thick, dashed] (9,0) rectangle (13,3);
    \node[below] at (2,-0.3) {Legenda: \underline{\hspace{3cm}}};
    \node[below] at (6.5,-0.3) {Legenda: \underline{\hspace{3cm}}};
    \node[below] at (11,-0.3) {Legenda: \underline{\hspace{3cm}}};
  \end{tikzpicture}
  \end{center}

  \item Circule a opcao correta:
  \begin{enumerate}[label=(\alph*)]
    \item Na frase ``O menino esta brincando'', o sujeito e: ( ) menino ( ) brincando ( ) esta
    \item O predicado indica: ( ) quem ( ) o que acontece ( ) onde
    \item ``Esta'' indica que a acao: ( ) ja terminou ( ) esta acontecendo ( ) vai acontecer
  \end{enumerate}

  \item Avaliacao diagnostica:
  \begin{center}
  $\square$ Le frases com predicado com fluencia\\
  $\square$ Identifica sujeito e predicado\\
  $\square$ Diferencia frase simples de frase com predicado\\
  $\square$ Escreve frases com predicado e pontuacao\\
  $\square$ Compreende que o predicado descreve a acao
  \end{center}

  \item Leitura em voz alta: o professor le e o aluno desenha:
  \begin{enumerate}[label=\alph*)]
    \item ``O sapo esta pulando.''
    \item ``A menina esta estudando.''
    \item ``O peixe esta nadando.''
  \end{enumerate}
\end{enumerate}
\end{exercicio}

\begin{dica}
\textbf{Dica para o Professor:} Frases com predicado sao mais longas e complexas. O predicado indica o que o sujeito esta fazendo. Use o verbo ``estar'' para indicar acoes em andamento. Peca que os alunos observem a diferenca entre ``O menino corre.'' (frase simples) e ``O menino esta brincando.'' (frase com predicado). Use historias curtas com frases com predicado para motivar a leitura.
\end{dica}

%% ============================================================
%% UNIDADE 5 -- CONSONANTES DO BLOCO 2 (Nivel 3 -- Silabico-Alfabetico)
%% ============================================================

\chapter{Unidade 5 --- Consonantes do Bloco 2}

%% ------------------------------------------------------------
%% LICAO 21: LETRA B
%% ------------------------------------------------------------
\section{Licao 21: A Letra B}

% Rotina neuroinclusiva e badges (R201)
\rotinasimples
\nivelKbadge{K2 --- Consciencia Fonologica}
\rotabadge{1B --- Segmentacao e Leitura Funcional}


\letra{B}
% Quatro formas gráficas da letra (R201)
\quatroformasletra{B}{b}

\boq[fechada]{B}{%
  Boca ENTREABERTA, labios levemente juntos.\\
  Ar sai em explosao suave pelos labios (oclusiva bilabial).\\
  Labios se tocam e se abrem rapidamente.\\
  Som: ``B-B-B-B-B'' (curto, explosivo).\\[0.5em]
  Palavras exemplo: bola, boi, boca, braco, beijo.
}

\pistaarticulatoria{Som: b (explosivo, sonoro)}{L'{a}bios fechados que se abrem de repente}{M\~{a}o na garganta: vibra (som sonoro)}

\metadata{E-01}{Nivel 3 -- Silabico-Alfabetico}{EF01LP02}{D1}{EF01LP02}{Resolucao CNE/CEB 7/2010, art. 12}

\plancheta{Folha de Pratica 5.1 --- Letra B}
\begin{exercicio}[Folha de Pratica 5.1 --- Letra B]
\noindent\textbf{Nome:} \underline{\hspace{3.2cm}} \quad \textbf{Data:} \underline{\hspace{1.6cm}} \quad \textbf{Nota:} \underline{\hspace{0.9cm}}/10

\subsection*{PARTE 1: RECONHECIMENTO}

\begin{enumerate}
  \item Formar silabas com B e vogais:
  \begin{center}
  \sil{B}{A} \quad \sil{B}{E} \quad \sil{B}{I} \quad \sil{B}{O} \quad \sil{B}{U}
  \end{center}

  \item Sublinhe palavras que comecam com ``B'':
  \begin{center}
  \uline{bola} \quad gato \quad \uline{boi} \quad sapo \quad \uline{boca} \quad rato \quad \uline{braco}
  \end{center}

  \item Coloque $\checkmark$ nas linhas com ``B'':
  \begin{center}
  \begin{tabular}{ll}
  $\checkmark$ \texttt{B O L A} & $\checkmark$ \texttt{B O C A} \\
  $\square$ \texttt{C A S A} & $\checkmark$ \texttt{B R A C O} \\
  $\checkmark$ \texttt{B E I J O} & $\square$ \texttt{M E S A}
  \end{tabular}
  \end{center}

  \item Identifique a posicao do ``B'' em ``BOLA'':
  \begin{center}
  \fbox{B} --- O --- L --- A $\rightarrow$ posicao \underline{1}
  \end{center}

  \item Diga em voz alta tres palavras com ``B'' e anote:
  \begin{center}
  1. \underline{\hspace{4cm}} \quad 2. \underline{\hspace{4cm}} \quad 3. \underline{\hspace{4cm}}
  \end{center}
\end{enumerate}

\subsection*{PARTE 2: PRODUCAO GUIADA}

\begin{enumerate}[resume]
  \item Complete com ``B'':
  \begin{center}
  \underline{B}ola \quad \underline{B}oi \quad \underline{B}oca \quad \underline{B}raco \quad \underline{B}eijo
  \end{center}

  \item Escreva a letra ``B'' maiuscula cinco vezes:
  \begin{center}
  \underline{\hspace{1.5cm}} \quad \underline{\hspace{1.5cm}} \quad \underline{\hspace{1.5cm}} \quad \underline{\hspace{1.5cm}} \quad \underline{\hspace{1.5cm}}
  \end{center}

  \item Escreva a letra ``b'' minuscula cinco vezes:
  \begin{center}
  \underline{\hspace{1.5cm}} \quad \underline{\hspace{1.5cm}} \quad \underline{\hspace{1.5cm}} \quad \underline{\hspace{1.5cm}} \quad \underline{\hspace{1.5cm}}
  \end{center}

  \item Copie as silabas e palavras:
  \begin{center}
  BA --- BE --- BI --- BO --- BU\\
  bola --- boi --- boca --- braco --- beijo
  \end{center}

  \item Escreva 3 palavras com ``B'':
  \begin{center}
  1. \underline{\hspace{4cm}} \quad 2. \underline{\hspace{4cm}} \quad 3. \underline{\hspace{4cm}}
  \end{center}
\end{enumerate}

\subsection*{PARTE 3: INTEGRACAO}

\begin{enumerate}[resume]
  \item Leia as silabas em voz alta:
  \begin{center}
  $\rightarrow$ BA  ~  BE  ~  BI  ~  BO  ~  BU
  \end{center}

  \item Nome de um animal com ``B'':
  \begin{center}
  \underline{\hspace{6cm}}
  \end{center}

  \item Desenhe algo com ``B'' e escreva o nome:
  \begin{center}
  \begin{tikzpicture}
    \draw[thick, dashed] (0,0) rectangle (5,3);
    \node[below] at (2.5,-0.3) {Nome: \underline{\hspace{4cm}}};
  \end{tikzpicture}
  \end{center}

  \item Circule a opcao correta:
  \begin{enumerate}[label=(\alph*)]
    \item O B tem som: ( ) nasal ( ) sibilante ( ) oclusivo
    \item Os labios ficam: ( ) fechados ( ) abertos ( ) entreabertos
    \item O B aparece em: ( ) mesa ( ) bola ( ) sapo
  \end{enumerate}

  \item Leia e complete com a silaba correta:
  \begin{center}
  \sil{B}{A}la \quad \underline{BE}ijo \quad \sil{B}{I}co \quad \underline{BO}la \quad \underline{BU}ca
  \end{center}
\end{enumerate}
\end{exercicio}

\begin{dica}
\textbf{Dica para o Professor:} O B e uma oclusiva bilabial. Peca que os alunos encostem os labios e soltem o ar com forca. Compare com o P (que e surdo, sem vibracao na garganta) para que percebam a diferenca. Use o espelho para observar que os labios se tocam.
\end{dica}

%% ------------------------------------------------------------


% ============================================================
% CONSOLIDACAO DA LETRA B (R201.1) -- Missao 11/16
% ============================================================

\plancheta{Miss\~a{}o 11 --- Letra B}

\begin{missao}
\textbf{\faGamepad~Miss\~a{}o 11 de 16:} Consolide a letra B com a fam\'ilia do B: bab\'a, baba, beb\^e, bibi, bob\^o, bobo!
\end{missao}

\begin{consolidacao}
\textbf{Atividade 1 --- Ca\c{c}ada ao B}

\instrucao{Circule o B e sublinhe as palavras com B.}

\begin{center}
\begin{tabular}{cccccc}
\fbox{B} & C & \fbox{B} & D & \fbox{B} & F \\
G & \fbox{B} & C & \fbox{B} & D & \fbox{B}
\end{tabular}
\end{center}
\begin{center}
\uline{bola} \quad \uline{baba} \quad casa \quad \uline{beb\^e} \quad dado \quad \uline{bibi}
\end{center}

\caixapausa{Se precisar, pare aqui. Respire fundo e continue quando quiser.}

\textbf{Atividade 2 --- Escrita do B --- 4 formas}

\instrucao{Escreva o B cinco vezes em cada forma.}

\begin{center}
\textbf{Imprensa mai\'uscula:} \underline{\hspace{1cm}} \underline{\hspace{1cm}} \underline{\hspace{1cm}} \underline{\hspace{1cm}} \underline{\hspace{1cm}}\\[0.5em]
\textbf{Imprensa min\'uscula:} \underline{\hspace{1cm}} \underline{\hspace{1cm}} \underline{\hspace{1cm}} \underline{\hspace{1cm}} \underline{\hspace{1cm}}\\[0.5em]
\textbf{Cursiva mai\'uscula:} \underline{\hspace{1cm}} \underline{\hspace{1cm}} \underline{\hspace{1cm}} \underline{\hspace{1cm}} \underline{\hspace{1cm}}\\[0.5em]
\textbf{Cursiva min\'uscula:} \underline{\hspace{1cm}} \underline{\hspace{1cm}} \underline{\hspace{1cm}} \underline{\hspace{1cm}} \underline{\hspace{1cm}}
\end{center}

\textbf{Atividade 3 --- A Fam\'ilia do B}

\instrucao{Leia a fam\'ilia do B em voz alta. Depois copie as quatro palavras que o professor ditar.}

\begin{center}
\textbf{bab\'a} \quad \textbf{baba} \quad \textbf{beb\^e} \quad \textbf{bibi} \quad \textbf{bob\^o} \quad \textbf{bobo}\\[0.5em]
\textbf{bolo} \quad \textbf{bola} \quad \textbf{bala} \quad \textbf{bula} \quad \textbf{bico} \quad \textbf{boca}
\end{center}
\instrucao{Ditado: escreva aqui.}
\begin{center}
1. \underline{\hspace{3cm}} \quad 2. \underline{\hspace{3cm}} \quad 3. \underline{\hspace{3cm}} \quad 4. \underline{\hspace{3cm}}
\end{center}

\caixapista{A fam\'ilia do B troca s\'o a vogal: BA, BE, BI, BO, BU.}
\end{consolidacao}

\plancheta{Sons da Letra B --- IPA}

\begin{sonsdaletra}
Os sons da letra B no portugu\^es{} brasileiro (IPA --- Alfabeto Fon\'etico Internacional):

\somipa[fechada]{B \'unico}{b}{l\'abios fechados que explodem com voz}{bola, baba, beb\^e}
\end{sonsdaletra}

\plancheta{Fam\'ilia Sil\'abica da Letra B}

\begin{familiasilabica}
A fam\'ilia sil\'abica da letra B:

\begin{center}
bab\'a \quad baba \quad beb\^e \quad bibi \quad bob\^o \quad bobo \quad bolo \quad bola \quad bala \quad bula \quad bico \quad boca
\end{center}
\instrucao{Leia a fam\'ilia em voz alta, devagar. Depois cubra e tente lembrar.}
\end{familiasilabica}

\plancheta{Atividade de Som e S\'ilabas da Letra B}

\begin{somletra}
\textbf{1. Ca\c{c}a ao som.} Ou\c{c}a a palavra. Se tiver o som de B, marque o quadradinho:
\begin{center}
$\square$~BOLA \quad $\square$~BEB\^E \quad $\square$~CABO \quad $\square$~M\~AO \quad $\square$~P\'E \quad $\square$~SOL
\end{center}

\textbf{2. Tem ou n\~ao{} tem?} Rodeie SIM ou N\~ao:
\begin{center}
BOLA --- (\textbf{SIM}) (\textbf{N\~AO}) \quad M\~AO --- (\textbf{SIM}) (\textbf{N\~AO}) \\ BEB\^E --- (\textbf{SIM}) (\textbf{N\~AO}) \quad P\'E --- (\textbf{SIM}) (\textbf{N\~AO})
\end{center}
\caixapausa{A boca fecha, com som e voz.}
\end{somletra}

\begin{somsilaba}
\textbf{3. Som das s\'ilabas.} Ou\c{c}a, bata palmas e conte as s\'ilabas:
\begin{center}
BO.LA $\rightarrow$ \underline{\hspace{1cm}} s\'ilabas \quad BE.B\^E $\rightarrow$ \underline{\hspace{1cm}} s\'ilabas
\end{center}

\textbf{4. Qual s\'ilaba tem o som?} Circule a s\'ilaba com o som de B:
\begin{center}
\textbf{BO}.LA \quad \textbf{BE}.B\^E
\end{center}
\caixapista{A consoante B se junta com a vogal: BA, BE, BI, BO, BU.}
\end{somsilaba}


\plancheta{Miss\~a{}o Motora --- Letra B}

\begin{motricidade}
\textbf{Movimento:} trace as duas barrigas do B: uma em cima, uma embaixo.
\begin{center}
\begin{tikzpicture}
  \draw[very thick, dashed, color=primary] (0,0) -- (0,1.6);
  \draw[very thick, dashed, color=primary] (0,1.3) to[out=0,in=180] (0.6,1.0) to[out=0,in=180] (0,0.7);
  \draw[very thick, dashed, color=secondary] (0,0.55) to[out=0,in=180] (0.7,0.3) to[out=0,in=180] (0,0);
  \draw[->, color=accent] (0.15,0.9) -- (0.3,0.9);
  \draw[->, color=accent] (0.15,0.25) -- (0.3,0.25);
\end{tikzpicture}
\end{center}
\caixapista{O B \'e{} um pilar com duas barrigas. Barriga de cima menor, barriga de baixo maior.}
\end{motricidade}

\plancheta{Ponte --- Da Letra B para a Pr\'oxima}

\begin{transicao}
Compare o B com o C:
\begin{center}
\textbf{bala} / \textbf{cala} \qquad \textbf{boca} / \textbf{coca} \qquad \textbf{bola} / \textbf{cola}
\end{center}
\instrucao{No B os l\'abios explodem com voz. No C o som sai da garganta, sem l\'abios.}
\end{transicao}

\plancheta{Recompensa --- Miss\~a{}o 11}

\begin{recompensa}
\estrelas{3}\quad \ofensiva{11}\\[0.3em]\barraprogresso{11}{16}\\[0.5em]

\checklistdominio{ler e escrever toda a fam\'ilia do B (bab\'a, baba, beb\^e, bibi, bob\^o, bobo)}
\end{recompensa}


%% LICAO 22: LETRA C
%% ------------------------------------------------------------
\section{Licao 22: A Letra C (dois sons: K e S) e a Cedilha}

% Rotina neuroinclusiva e badges (R201)
\rotinasimples
\nivelKbadge{K2 --- Consciencia Fonologica}
\rotabadge{1B --- Segmentacao e Leitura Funcional}


\letra{C}
% Quatro formas gráficas da letra (R201)
\quatroformasletra{C}{c}

\boq[velar]{C}{%
  Boca ENTREABERTA, lingua toca o palato mole.\\
  Ar sai em explosao suave pelo fundo da boca.\\
  Som: ``C-C-C-C-C'' (oclusiva velar, antes de A, O, U).\\[0.5em]
  Palavras exemplo: casa, cavalo, cesto, copo, cumo.
}

\pistaarticulatoria{Som: k (seco, na garganta)}{Boca aberta, som sai do fundo da boca}{M\~{a}o na garganta: n\~{a}o vibra (som surdo)}

\metadata{E-02}{Nivel 3 -- Silabico-Alfabetico}{EF01LP02}{D1}{EF01LP02}{Resolucao CNE/CEB 7/2010, art. 12}

\plancheta{Folha de Pratica 5.2 --- Letra C (CA, CO, CU)}
\begin{exercicio}[Folha de Pratica 5.2 --- Letra C (CA, CO, CU)]
\noindent\textbf{Nome:} \underline{\hspace{3.2cm}} \quad \textbf{Data:} \underline{\hspace{1.6cm}} \quad \textbf{Nota:} \underline{\hspace{0.9cm}}/10

\subsection*{PARTE 1: RECONHECIMENTO}

\begin{enumerate}
  \item Formar silabas com C (antes de A, O, U):
  \begin{center}
  \sil{C}{A} \quad \sil{C}{O} \quad \sil{C}{U}
  \end{center}

  \item Sublinhe palavras que comecam com ``C'':
  \begin{center}
  \uline{casa} \quad gato \quad \uline{cavalo} \quad sapo \quad \uline{copo} \quad rato \quad \uline{cumo}
  \end{center}

  \item Coloque $\checkmark$ nas linhas com ``C'':
  \begin{center}
  \begin{tabular}{ll}
  $\checkmark$ \texttt{C A S A} & $\checkmark$ \texttt{C O P O} \\
  $\square$ \texttt{M E S A} & $\checkmark$ \texttt{C A V A L O} \\
  $\checkmark$ \texttt{C U M O} & $\square$ \texttt{L A T A}
  \end{tabular}
  \end{center}

  \item Identifique a posicao do ``C'' em ``CASA'':
  \begin{center}
  \fbox{C} --- A --- S --- A $\rightarrow$ posicao \underline{1}
  \end{center}

  \item Diga em voz alta tres palavras com ``C'' e anote:
  \begin{center}
  1. \underline{\hspace{4cm}} \quad 2. \underline{\hspace{4cm}} \quad 3. \underline{\hspace{4cm}}
  \end{center}
\end{enumerate}

\subsection*{PARTE 2: PRODUCAO GUIADA}

\begin{enumerate}[resume]
  \item Complete com ``C'':
  \begin{center}
  \underline{C}asa \quad \underline{C}avalo \quad \underline{C}opo \quad \underline{C}umo \quad \underline{C}esto
  \end{center}

  \item Escreva a letra ``C'' maiuscula cinco vezes:
  \begin{center}
  \underline{\hspace{1.5cm}} \quad \underline{\hspace{1.5cm}} \quad \underline{\hspace{1.5cm}} \quad \underline{\hspace{1.5cm}} \quad \underline{\hspace{1.5cm}}
  \end{center}

  \item Escreva a letra ``c'' minuscula cinco vezes:
  \begin{center}
  \underline{\hspace{1.5cm}} \quad \underline{\hspace{1.5cm}} \quad \underline{\hspace{1.5cm}} \quad \underline{\hspace{1.5cm}} \quad \underline{\hspace{1.5cm}}
  \end{center}

  \item Copie as silabas e palavras:
  \begin{center}
  CA --- CO --- CU\\
  casa --- cavalo --- copo --- cesto --- cumo
  \end{center}

  \item Escreva 3 palavras com ``C'':
  \begin{center}
  1. \underline{\hspace{4cm}} \quad 2. \underline{\hspace{4cm}} \quad 3. \underline{\hspace{4cm}}
  \end{center}
\end{enumerate}

\subsection*{PARTE 3: INTEGRACAO}

\begin{enumerate}[resume]
  \item Leia as silabas em voz alta:
  \begin{center}
  $\rightarrow$ CA \quad CO \quad CU
  \end{center}

  \item Nome de um objeto com ``C'':
  \begin{center}
  \underline{\hspace{6cm}}
  \end{center}

  \item Circule a opcao correta:
  \begin{enumerate}[label=(\alph*)]
    \item O C (antes de A) tem som: ( ) sibilante ( ) oclusivo velar ( ) nasal
    \item A lingua toca: ( ) o nariz ( ) os dentes ( ) o palato mole
    \item O C aparece em: ( ) mesa ( ) casa ( ) sapo
  \end{enumerate}

  \item Leia e complete com a silaba correta:
  \begin{center}
  \sil{C}{A}sa \quad \underline{CO}po \quad \sil{C}{U}la \quad \underline{CA}valo \quad \underline{CO}mida
  \end{center}

  \item Observacao: ``CE'' e ``CI'' tem som de ``S'' (aprendido na proxima licao).
  \begin{center}
  \textbf{C} + A, O, U = som de ``K'' (casa, copo, cumo)\\
  \textbf{C} + E, I = som de ``S'' (cesto, cima)
  \end{center}
\end{enumerate}
\end{exercicio}

\begin{dica}
\textbf{Dica para o Professor:} O C tem dois sons: ``K'' (antes de A, O, U) e ``S'' (antes de E, I). Comece pelo som ``K'' e reforce a regra. Use pares minimos: ``CA'' vs ``CE'', ``CO'' vs ``CI''. Os alunos devem memorizar que C + E/I = S.
\end{dica}

%% ------------------------------------------------------------


% ============================================================
% CONSOLIDACAO DA LETRA C (R201.1) -- Missao 12/16
% ============================================================

\plancheta{Miss\~a{}o 12 --- Letra C}

\begin{missao}
\textbf{\faGamepad~Miss\~a{}o 12 de 16:} Consolide a letra C, a letra de dois sons: ca, co, cu (K) e ce, ci (S)!
\end{missao}

\begin{consolidacao}
\textbf{Atividade 1 --- Ca\c{c}ada ao C}

\instrucao{Circule o C e sublinhe as palavras com C.}

\begin{center}
\begin{tabular}{cccccc}
\fbox{C} & B & \fbox{C} & D & \fbox{C} & F \\
G & \fbox{C} & B & \fbox{C} & D & \fbox{C}
\end{tabular}
\end{center}
\begin{center}
\uline{casa} \quad \uline{copo} \quad bola \quad \uline{cedo} \quad dado \quad \uline{cena}
\end{center}

\caixapausa{Se precisar, pare aqui. Respire fundo e continue quando quiser.}

\textbf{Atividade 2 --- Escrita do C --- 4 formas}

\instrucao{Escreva o C cinco vezes em cada forma.}

\begin{center}
\textbf{Imprensa mai\'uscula:} \underline{\hspace{1cm}} \underline{\hspace{1cm}} \underline{\hspace{1cm}} \underline{\hspace{1cm}} \underline{\hspace{1cm}}\\[0.5em]
\textbf{Imprensa min\'uscula:} \underline{\hspace{1cm}} \underline{\hspace{1cm}} \underline{\hspace{1cm}} \underline{\hspace{1cm}} \underline{\hspace{1cm}}\\[0.5em]
\textbf{Cursiva mai\'uscula:} \underline{\hspace{1cm}} \underline{\hspace{1cm}} \underline{\hspace{1cm}} \underline{\hspace{1cm}} \underline{\hspace{1cm}}\\[0.5em]
\textbf{Cursiva min\'uscula:} \underline{\hspace{1cm}} \underline{\hspace{1cm}} \underline{\hspace{1cm}} \underline{\hspace{1cm}} \underline{\hspace{1cm}}
\end{center}

\textbf{Atividade 3 --- Jogo do C de Dois Sons}

\instrucao{Ou\c{c}a e marque: som de K ou som de S?}

\begin{center}
\begin{tabular}{ll}
casa \quad $\rightarrow$ \underline{\hspace{2cm}} & cena \quad $\rightarrow$ \underline{\hspace{2cm}} \\
copo \quad $\rightarrow$ \underline{\hspace{2cm}} & cinto \quad $\rightarrow$ \underline{\hspace{2cm}} \\
cubo \quad $\rightarrow$ \underline{\hspace{2cm}} & doce \quad $\rightarrow$ \underline{\hspace{2cm}}
\end{tabular}
\end{center}
Regra: C + A, O, U = K. C + E, I = S.

\caixapista{ca, co, cu = K. ce, ci = S. Olhe a vogal que vem depois!}
\end{consolidacao}

\plancheta{Sons da Letra C --- IPA}

\begin{sonsdaletra}
Os sons da letra C no portugu\^es{} brasileiro (IPA --- Alfabeto Fon\'etico Internacional):

\somipa[velar]{C antes de A, O, U}{k}{som de K}{casa, copo, cubo}
\somipa[dental]{C antes de E, I}{s}{som de S}{cena, cinto, doce, cip\'o}
\somipa[dental]{Cedilha (C com rabinho)}{s}{som de S antes de A, O e U}{ma\c{c}\~a, a\c{c}\'ucar, pa\c{c}oca}
\end{sonsdaletra}

\plancheta{Fam\'ilia Sil\'abica da Letra C}

\begin{familiasilabica}
A fam\'ilia sil\'abica da letra C:

\begin{center}
ca\c{c}ada \quad ca\c{c}ador \quad ma\c{c}\~a \quad a\c{c}\'ucar \quad pa\c{c}oca \quad casa \quad copo \quad cubo \quad caco \quad coco \quad cuca \quad cedo
\end{center}
\instrucao{Leia a fam\'ilia em voz alta, devagar. Depois cubra e tente lembrar.}
\end{familiasilabica}

\plancheta{Atividade de Som e S\'ilabas da Letra C}

\begin{somletra}
\textbf{1. Ca\c{c}a ao som.} Ou\c{c}a a palavra. Se tiver o som de C, marque o quadradinho:
\begin{center}
$\square$~CASA \quad $\square$~CADEIRA \quad $\square$~SACO \quad $\square$~M\~AO \quad $\square$~BOI \quad $\square$~P\'E
\end{center}

\textbf{2. Tem ou n\~ao{} tem?} Rodeie SIM ou N\~ao:
\begin{center}
CASA --- (\textbf{SIM}) (\textbf{N\~AO}) \quad M\~AO --- (\textbf{SIM}) (\textbf{N\~AO}) \\ CADEIRA --- (\textbf{SIM}) (\textbf{N\~AO}) \quad BOI --- (\textbf{SIM}) (\textbf{N\~AO})
\end{center}
\caixapausa{O som sai do fundo da boca.}
\end{somletra}

\begin{somsilaba}
\textbf{3. Som das s\'ilabas.} Ou\c{c}a, bata palmas e conte as s\'ilabas:
\begin{center}
CA.SA $\rightarrow$ \underline{\hspace{1cm}} s\'ilabas \quad CA.DEI.RA $\rightarrow$ \underline{\hspace{1cm}} s\'ilabas
\end{center}

\textbf{4. Qual s\'ilaba tem o som?} Circule a s\'ilaba com o som de C:
\begin{center}
\textbf{CA}.SA \quad \textbf{CA}.DEI.RA
\end{center}
\caixapista{O C com A, O e U tem som de K (\textbf{CA}, \textbf{CO}, \textbf{CU}); com E e I tem som de S (\textbf{CE}, \textbf{CI}); com cedilha (\c{C}) tamb\'em tem som de S antes de A, O e U (\textbf{\c{C}A}, \textbf{\c{C}O}, \textbf{\c{C}U}).}
\end{somsilaba}


\plancheta{Cedilha --- o C com Rabinho}

\begin{somletra}
\textbf{Palavra geradora: \textsc{MA\c{C}\~A}.} A ma\c{c}\~a \'e doce e redonda, e o ``\c{C}A'' dela tem o som de S. Compare com CASA: o C de CASA tem som de K; o \c{C} de MA\c{C}\~A tem som de S.

\textbf{1. Onde tem \c{C}?} Ou\c{c}a a palavra e marque se tem \c{C}:
\begin{center}
$\square$~MA\c{C}\~A \quad $\square$~CASA \quad $\square$~A\c{C}\'UCAR \quad $\square$~COPO \quad $\square$~PA\c{C}OCA \quad $\square$~BOLA
\end{center}

\textbf{2. Tem ou n\~ao{} tem o som de S?} Rodeie SIM ou N\~AO:
\begin{center}
MA\c{C}\~A --- (\textbf{SIM}) (\textbf{N\~AO}) \quad CASA --- (\textbf{SIM}) (\textbf{N\~AO}) \\ A\c{C}\'UCAR --- (\textbf{SIM}) (\textbf{N\~AO}) \quad CA\c{C}A --- (\textbf{SIM}) (\textbf{N\~AO})
\end{center}
\caixapausa{O som de \c{C} \'e o mesmo de S: o ar passa assobiando.}
\end{somletra}

\begin{somsilaba}
\textbf{3. Recombina\c{c}\~ao (descoberta).} Forme novas palavras com as s\'ilabas \textbf{\c{C}A}, \textbf{\c{C}O}, \textbf{\c{C}U}:
\begin{center}
CA + \c{C}A $\rightarrow$ \underline{\hspace{2cm}} \quad MA + \c{C}\~A $\rightarrow$ \underline{\hspace{2cm}} \\
A + \c{C}\'U + CAR $\rightarrow$ \underline{\hspace{2cm}} \quad PA + \c{C}O + CA $\rightarrow$ \underline{\hspace{2cm}}
\end{center}
\textbf{4. Leia e separe.} Separe as s\'ilabas:
\begin{center}
MA\c{C}\~A $\rightarrow$ \underline{\hspace{2cm}} \quad A\c{C}\'UCAR $\rightarrow$ \underline{\hspace{2cm}} \quad PA\c{C}OCA $\rightarrow$ \underline{\hspace{2cm}}
\end{center}
\caixapista{Regra da cedilha: \c{C} s\'o aparece antes de A, O e U. Antes de E e I, quem faz o som de S \'e o pr\'oprio C: CE-DO, CIN-TO.}
\end{somsilaba}

\plancheta{Miss\~a{}o Motora --- Letra C}

\begin{motricidade}
\textbf{Movimento:} trace a meia-lua do C, de cima para baixo.
\begin{center}
\begin{tikzpicture}
  \draw[very thick, dashed, color=primary] (1.2,0.6) to[out=180,in=90] (0.2,0) to[out=-90,in=180] (1.2,-0.6);
  \draw[very thick, dashed, color=secondary] (3.2,0.6) to[out=180,in=90] (2.2,0) to[out=-90,in=180] (3.2,-0.6);
  \draw[->, color=accent] (1.0,0.45) -- (0.7,0.15);
  \draw[->, color=accent] (0.5,-0.25) -- (0.8,-0.45);
\end{tikzpicture}
\end{center}
\caixapista{O C \'e{} uma meia-lua aberta para a direita.}
\end{motricidade}

\plancheta{Ponte --- Da Letra C para a Pr\'oxima}

\begin{transicao}
Compare o C com o D:
\begin{center}
\textbf{casa} / \textbf{data} \qquad \textbf{coco} / \textbf{dado} \qquad \textbf{cama} / \textbf{dama}
\end{center}
\instrucao{No C o som sai da garganta. No D a l\'ingua toca os dentes e o som vibra.}
\end{transicao}

\plancheta{Recompensa --- Miss\~a{}o 12}

\begin{recompensa}
\estrelas{3}\quad \ofensiva{12}\\[0.3em]\barraprogresso{12}{16}\\[0.5em]

\checklistdominio{aplicar a regra do C: K antes de A/O/U e S antes de E/I}
\end{recompensa}


%% LICAO 23: LETRA D
%% ------------------------------------------------------------
\section{Licao 23: A Letra D}

% Rotina neuroinclusiva e badges (R201)
\rotinasimples
\nivelKbadge{K2 --- Consciencia Fonologica}
\rotabadge{1B --- Segmentacao e Leitura Funcional}


\letra{D}
% Quatro formas gráficas da letra (R201)
\quatroformasletra{D}{d}

\boq[dental]{D}{%
  Boca ENTREABERTA, lingua toca ATRAS dos dentes superiores.\\
  Ar sai em explosao seca (oclusiva dental sonora).\\
  Vibracao leve na garganta.\\
  Som: ``D-D-D-D-D'' (curto, seco).\\[0.5em]
  Palavras exemplo: dente, dardo, dedo, disco, duro.
}

\pistaarticulatoria{Som: d (seco, sonoro)}{L'{i}ngua toca atr'{a}s dos dentes superiores}{M\~{a}o na garganta: vibra (som sonoro)}

\metadata{E-03}{Nivel 3 -- Silabico-Alfabetico}{EF01LP02}{D1}{EF01LP02}{Resolucao CNE/CEB 7/2010, art. 12}

\plancheta{Folha de Pratica 5.3 --- Letra D}
\begin{exercicio}[Folha de Pratica 5.3 --- Letra D]
\noindent\textbf{Nome:} \underline{\hspace{3.2cm}} \quad \textbf{Data:} \underline{\hspace{1.6cm}} \quad \textbf{Nota:} \underline{\hspace{0.9cm}}/10

\subsection*{PARTE 1: RECONHECIMENTO}

\begin{enumerate}
  \item Formar silabas com D e vogais:
  \begin{center}
  \sil{D}{A} \quad \sil{D}{E} \quad \sil{D}{I} \quad \sil{D}{O} \quad \sil{D}{U}
  \end{center}

  \item Sublinhe palavras que comecam com ``D'':
  \begin{center}
  \uline{dente} \quad gato \quad \uline{dardo} \quad sapo \quad \uline{dedo} \quad rato \quad \uline{disco}
  \end{center}

  \item Coloque $\checkmark$ nas linhas com ``D'':
  \begin{center}
  \begin{tabular}{ll}
  $\checkmark$ \texttt{D E N T E} & $\checkmark$ \texttt{D E D O} \\
  $\square$ \texttt{C A S A} & $\checkmark$ \texttt{D A R D O} \\
  $\checkmark$ \texttt{D I S C O} & $\square$ \texttt{M E S A}
  \end{tabular}
  \end{center}

  \item Identifique a posicao do ``D'' em ``DENTE'':
  \begin{center}
  \fbox{D} --- E --- N --- T --- E $\rightarrow$ posicao \underline{1}
  \end{center}

  \item Diga em voz alta tres palavras com ``D'' e anote:
  \begin{center}
  1. \underline{\hspace{4cm}} \quad 2. \underline{\hspace{4cm}} \quad 3. \underline{\hspace{4cm}}
  \end{center}
\end{enumerate}

\subsection*{PARTE 2: PRODUCAO GUIADA}

\begin{enumerate}[resume]
  \item Complete com ``D'':
  \begin{center}
  \underline{D}ente \quad \underline{D}ardo \quad \underline{D}edo \quad \underline{D}isco \quad \underline{D}uro
  \end{center}

  \item Escreva a letra ``D'' maiuscula cinco vezes:
  \begin{center}
  \underline{\hspace{1.5cm}} \quad \underline{\hspace{1.5cm}} \quad \underline{\hspace{1.5cm}} \quad \underline{\hspace{1.5cm}} \quad \underline{\hspace{1.5cm}}
  \end{center}

  \item Escreva a letra ``d'' minuscula cinco vezes:
  \begin{center}
  \underline{\hspace{1.5cm}} \quad \underline{\hspace{1.5cm}} \quad \underline{\hspace{1.5cm}} \quad \underline{\hspace{1.5cm}} \quad \underline{\hspace{1.5cm}}
  \end{center}

  \item Copie as silabas e palavras:
  \begin{center}
  DA --- DE --- DI --- DO --- DU\\
  dente --- dardo --- dedo --- disco --- duro
  \end{center}

  \item Escreva 3 palavras com ``D'':
  \begin{center}
  1. \underline{\hspace{4cm}} \quad 2. \underline{\hspace{4cm}} \quad 3. \underline{\hspace{4cm}}
  \end{center}
\end{enumerate}

\subsection*{PARTE 3: INTEGRACAO}

\begin{enumerate}[resume]
  \item Leia as silabas em voz alta:
  \begin{center}
  $\rightarrow$ DA  ~  DE  ~  DI  ~  DO  ~  DU
  \end{center}

  \item Nome de um animal com ``D'':
  \begin{center}
  \underline{\hspace{6cm}}
  \end{center}

  \item Circule a opcao correta:
  \begin{enumerate}[label=(\alph*)]
    \item O D tem som: ( ) nasal ( ) sibilante ( ) oclusivo dental
    \item A lingua toca: ( ) o nariz ( ) os dentes ( ) o palato
    \item O D aparece em: ( ) mesa ( ) dente ( ) sapo
  \end{enumerate}

  \item Leia e complete com a silaba correta:
  \begin{center}
  \sil{D}{A}do \quad \underline{DE}nte \quad \sil{D}{I}co \quad \underline{DO}uro \quad \underline{DU}da
  \end{center}
\end{enumerate}
\end{exercicio}

\begin{dica}
\textbf{Dica para o Professor:} O D e uma oclusiva dental sonora. Compare com o T (surdo) para que percebam a vibracao da garganta. Use o espelho para observar que a lingua toca atras dos dentes. Os alunos devem sentir a diferenca entre D (com som) e T (sem som na garganta).
\end{dica}

%% ------------------------------------------------------------


% ============================================================
% CONSOLIDACAO DA LETRA D (R201.1) -- Missao 13/16
% ============================================================

\plancheta{Miss\~a{}o 13 --- Letra D}

\begin{missao}
\textbf{\faGamepad~Miss\~a{}o 13 de 16:} Consolide a letra D. Complete e voc\^e fecha o Bloco 2!
\end{missao}

\begin{consolidacao}
\textbf{Atividade 1 --- Ca\c{c}ada ao D}

\instrucao{Circule o D e sublinhe as palavras com D.}

\begin{center}
\begin{tabular}{cccccc}
\fbox{D} & B & \fbox{D} & C & \fbox{D} & F \\
G & \fbox{D} & B & \fbox{D} & C & \fbox{D}
\end{tabular}
\end{center}
\begin{center}
\uline{dado} \quad \uline{dedo} \quad bola \quad \uline{doca} \quad casa \quad \uline{duna}
\end{center}

\caixapausa{Se precisar, pare aqui. Respire fundo e continue quando quiser.}

\textbf{Atividade 2 --- Escrita do D --- 4 formas}

\instrucao{Escreva o D cinco vezes em cada forma.}

\begin{center}
\textbf{Imprensa mai\'uscula:} \underline{\hspace{1cm}} \underline{\hspace{1cm}} \underline{\hspace{1cm}} \underline{\hspace{1cm}} \underline{\hspace{1cm}}\\[0.5em]
\textbf{Imprensa min\'uscula:} \underline{\hspace{1cm}} \underline{\hspace{1cm}} \underline{\hspace{1cm}} \underline{\hspace{1cm}} \underline{\hspace{1cm}}\\[0.5em]
\textbf{Cursiva mai\'uscula:} \underline{\hspace{1cm}} \underline{\hspace{1cm}} \underline{\hspace{1cm}} \underline{\hspace{1cm}} \underline{\hspace{1cm}}\\[0.5em]
\textbf{Cursiva min\'uscula:} \underline{\hspace{1cm}} \underline{\hspace{1cm}} \underline{\hspace{1cm}} \underline{\hspace{1cm}} \underline{\hspace{1cm}}
\end{center}

\textbf{Atividade 3 --- Jogo da Fam\'ilia do D}

\instrucao{Complete a fam\'ilia do D com a vogal que falta.}

\begin{center}
dad\underline{~~~} \quad d\underline{~~~}do \quad d\underline{~~~}ca \quad d\underline{~~~}na \quad d\underline{~~~}ca
\end{center}
Diga em voz alta: \textbf{dad\'a, dado, doca, duna, dica}

\caixapista{A fam\'ilia do D: DA, DE, DI, DO, DU.}
\end{consolidacao}

\plancheta{Sons da Letra D --- IPA}

\begin{sonsdaletra}
Os sons da letra D no portugu\^es{} brasileiro (IPA --- Alfabeto Fon\'etico Internacional):

\somipa[dental]{D seco}{d}{l\'ingua nos dentes, som sonoro}{dado, dedo, doca, duna}
\somipa[palatal]{D antes de I (regional)}{dZ}{vira dj em algumas regi\~o{}es}{dia, dica}
\end{sonsdaletra}

\plancheta{Fam\'ilia Sil\'abica da Letra D}

\begin{familiasilabica}
A fam\'ilia sil\'abica da letra D:

\begin{center}
dad\'a \quad dado \quad dedo \quad doca \quad duna \quad dica \quad dama
\end{center}
\instrucao{Leia a fam\'ilia em voz alta, devagar. Depois cubra e tente lembrar.}
\end{familiasilabica}

\plancheta{Atividade de Som e S\'ilabas da Letra D}

\begin{somletra}
\textbf{1. Ca\c{c}a ao som.} Ou\c{c}a a palavra. Se tiver o som de D, marque o quadradinho:
\begin{center}
$\square$~DEDO \quad $\square$~DADO \quad $\square$~MODO \quad $\square$~SOL \quad $\square$~P\~AO \quad $\square$~BOCA
\end{center}

\textbf{2. Tem ou n\~ao{} tem?} Rodeie SIM ou N\~ao:
\begin{center}
DEDO --- (\textbf{SIM}) (\textbf{N\~AO}) \quad SOL --- (\textbf{SIM}) (\textbf{N\~AO}) \\ DADO --- (\textbf{SIM}) (\textbf{N\~AO}) \quad P\~AO --- (\textbf{SIM}) (\textbf{N\~AO})
\end{center}
\caixapausa{A l\'ingua toca os dentes de cima, com voz.}
\end{somletra}

\begin{somsilaba}
\textbf{3. Som das s\'ilabas.} Ou\c{c}a, bata palmas e conte as s\'ilabas:
\begin{center}
DE.DO $\rightarrow$ \underline{\hspace{1cm}} s\'ilabas \quad DA.DO $\rightarrow$ \underline{\hspace{1cm}} s\'ilabas
\end{center}

\textbf{4. Qual s\'ilaba tem o som?} Circule a s\'ilaba com o som de D:
\begin{center}
\textbf{DE}.DO \quad \textbf{DA}.DO
\end{center}
\caixapista{A consoante D se junta com a vogal: DA, DE, DI, DO, DU.}
\end{somsilaba}


\plancheta{Miss\~a{}o Motora --- Letra D}

\begin{motricidade}
\textbf{Movimento:} trace a barriga grande do D.
\begin{center}
\begin{tikzpicture}
  \draw[very thick, dashed, color=primary] (0,0) -- (0,1.6);
  \draw[very thick, dashed, color=primary] (0,1.6) to[out=0,in=90] (1.1,0.8) to[out=-90,in=0] (0,0);
  \draw[->, color=accent] (0.3,1.45) -- (0.7,1.3);
  \draw[->, color=accent] (0.7,0.35) -- (0.35,0.15);
\end{tikzpicture}
\end{center}
\caixapista{O D \'e{} um pilar com uma barriga grande. S\'o uma!}
\end{motricidade}

\plancheta{Ponte --- Da Letra D para a Pr\'oxima}

\begin{transicao}
Compare o D com o X, da Unidade 8:
\begin{center}
\textbf{dado} / \textbf{xale} \qquad \textbf{dedo} / \textbf{t\'axi} \qquad \textbf{dama} / \textbf{caixa}
\end{center}
\instrucao{O D vibra na garganta. O X tem v\'arios sons que voc\^e vai descobrir na Unidade 8.}
\end{transicao}

\plancheta{Recompensa --- Miss\~a{}o 13}

\begin{recompensa}
\estrelas{3}\quad \ofensiva{13}\\[0.3em]\barraprogresso{13}{16}\\[0.5em]

\checklistdominio{ler a fam\'ilia do D e diferenciar D de B (dado/bola)}
\end{recompensa}


%% LICAO 24: LETRAS F, G, N, P, V, Z, J
%% ------------------------------------------------------------
\section{Licao 24: Consonantes do Bloco 2 --- F, G, J, K, N, P, V, Z, W, Y}

% Rotina neuroinclusiva e badges (R201)
\rotinasimples
\nivelKbadge{K2 --- Consciencia Fonologica}
\rotabadge{1B --- Segmentacao e Leitura Funcional}


\metadata{E-04}{Nivel 3 -- Silabico-Alfabetico}{EF01LP02}{D1}{EF01LP02}{Resolucao CNE/CEB 7/2010, art. 12}

\plancheta{Folha de Pratica 5.4 --- Silabas com F, G, N, P, V, Z, J}
\begin{exercicio}[Folha de Pratica 5.4 --- Silabas com F, G, N, P, V, Z, J]
\noindent\textbf{Nome:} \underline{\hspace{3.2cm}} \quad \textbf{Data:} \underline{\hspace{1.6cm}} \quad \textbf{Nota:} \underline{\hspace{0.9cm}}/10

\subsection*{PARTE 1: RECONHECIMENTO}

\begin{enumerate}
  \item Forme todas as silabas com F, G, N, P, V, Z, J e vogais:
  \begin{center}
  \resizebox{\textwidth}{!}{%
\begin{tabular}{ccccc}
  \sil{F}{A} & \sil{F}{E} & \sil{F}{I} & \sil{F}{O} & \sil{F}{U} \\
  \sil{G}{A} & \sil{G}{E} & \sil{G}{I} & \sil{G}{O} & \sil{G}{U} \\
  \sil{N}{A} & \sil{N}{E} & \sil{N}{I} & \sil{N}{O} & \sil{N}{U} \\
  \sil{P}{A} & \sil{P}{E} & \sil{P}{I} & \sil{P}{O} & \sil{P}{U} \\
  \sil{V}{A} & \sil{V}{E} & \sil{V}{I} & \sil{V}{O} & \sil{V}{U} \\
  \sil{Z}{A} & \sil{Z}{E} & \sil{Z}{I} & \sil{Z}{O} & \sil{Z}{U} \\
  \sil{J}{A} & \sil{J}{E} & \sil{J}{I} & \sil{J}{O} & \sil{J}{U} \\
  \end{tabular}}
  \end{center}

  \item Sublinhe palavras que comecam com cada letra:
  \begin{center}
  \uline{faca} \quad \uline{gato} \quad \uline{nose} \quad \uline{pato} \quad \uline{vaca} \quad \uline{zebra} \quad \uline{jato}
  \end{center}

  \item Separe por letras:
  \begin{center}
  FACAS $\rightarrow$ \underline{\hspace{2cm}} \quad GATOS $\rightarrow$ \underline{\hspace{2cm}} \quad PATOS $\rightarrow$ \underline{\hspace{2cm}} \\
  VACAS $\rightarrow$ \underline{\hspace{2cm}} \quad ZEBRAS $\rightarrow$ \underline{\hspace{2cm}} \quad JATOS $\rightarrow$ \underline{\hspace{2cm}}
  \end{center}

  \item Identifique a primeira letra de cada palavra:
  \begin{center}
  \begin{tabular}{ccccccc}
  \textbf{F}aca & \textbf{G}ato & \textbf{N}aso & \textbf{P}ato & \textbf{V}aca & \textbf{Z}ebra & \textbf{J}ato
  \end{tabular}
  \end{center}

  \item Diga em voz alta tres palavras com ``P'' e anote:
  \begin{center}
  1. \underline{\hspace{4cm}} \quad 2. \underline{\hspace{4cm}} \quad 3. \underline{\hspace{4cm}}
  \end{center}
\end{enumerate}

\subsection*{PARTE 2: PRODUCAO GUIADA}

\begin{enumerate}[resume]
  \item Complete com a letra correta:
  \begin{center}
  \underline{F}aca \quad \underline{G}ato \quad \underline{N}aso \quad \underline{P}ato \quad \underline{V}aca \quad \underline{Z}ebra \quad \underline{J}ato
  \end{center}

  \item Forme palavras com as silabas:
  \begin{center}
  FA + CA = \underline{\hspace{2cm}} \quad GA + TO = \underline{\hspace{2cm}} \quad PA + TO = \underline{\hspace{2cm}} \\
  VA + CA = \underline{\hspace{2cm}} \quad ZE + BRA = \underline{\hspace{2cm}} \quad JA + TO = \underline{\hspace{2cm}}
  \end{center}

  \item Escreva 2 palavras para cada letra:
  \begin{center}
  F: \underline{\hspace{3cm}} \underline{\hspace{3cm}} \quad G: \underline{\hspace{3cm}} \underline{\hspace{3cm}} \quad P: \underline{\hspace{3cm}} \underline{\hspace{3cm}} \\
  V: \underline{\hspace{3cm}} \underline{\hspace{3cm}} \quad Z: \underline{\hspace{3cm}} \underline{\hspace{3cm}} \quad J: \underline{\hspace{3cm}} \underline{\hspace{3cm}}
  \end{center}

  \item Copie as palavras:
  \begin{center}
  faca --- gato --- pato --- vaca --- zebra --- jato
  \end{center}

  \item Separe em silabas:
  \begin{center}
  FACAS $\rightarrow$ \underline{\hspace{1cm}} \underline{\hspace{1cm}} \quad GATOS $\rightarrow$ \underline{\hspace{1cm}} \underline{\hspace{1cm}} \quad PATOS $\rightarrow$ \underline{\hspace{1cm}} \underline{\hspace{1cm}}
  \end{center}
\end{enumerate}

\subsection*{PARTE 3: INTEGRACAO}

\begin{enumerate}[resume]
  \item Leia as palavras em voz alta:
  \begin{center}
  $\rightarrow$ faca  ~  gato  ~  pato  ~  vaca  ~  zebra  ~  jato  ~  nabo  ~  pano
  \end{center}

  \item Circule a opcao correta:
  \begin{enumerate}[label=(\alph*)]
    \item O F tem som: ( ) nasal ( ) fricativo ( ) vibrante
    \item O P tem som: ( ) sonoro ( ) surdo ( ) nasal
    \item O J tem som: ( ) oclusivo ( ) fricativo ( ) lateral
  \end{enumerate}

  \item Leia e complete:
  \begin{center}
  \underline{F}aca \quad \underline{G}ato \quad \underline{P}ato \quad \underline{V}aca \quad \underline{Z}ebra
  \end{center}

  \item Avaliacao diagnostica:
  \begin{center}
  $\square$ Forma silabas com todas as consonantes\\
  $\square$ Separa palavras em silabas\\
  $\square$ Le palavras com fluencia\\
  $\square$ Identifica fonema-grafema corretamente\\
  $\square$ Diferencia as 12 consonantes estudadas
  \end{center}

  \item Leitura em voz alta: o professor aponta e o aluno le:
  \begin{center}
  faca \quad gato \quad pato \quad vaca \quad zebra \quad jato \quad mesa \quad sapo \quad tinta \quad rato \quad lua \quad bola
  \end{center}
\end{enumerate}
\end{exercicio}

\begin{dica}
\textbf{Dica para o Professor:} Esta licao consolida todas as consonantes do Bloco 2. Revisa rapidamente a boquinha de cada letra. Use o jogo ``Bingo das Silabas'' com todas as consonantes para tornar a avaliacao ludica. Avalie se o aluno esta no nivel silabico-alfabetico: se forma silabas com todas as consonantes e le palavras simples.
\end{dica}

%% ============================================================

%% ============================================================
%% ALFABETO COMPLETO -- CONSOLIDACAO DAS LETRAS F, G, J, K, N, P, V, Z, W, Y (R201.2)
%% ============================================================

\section{Alfabeto Completo --- Consolida\c{c}\~a{}o das Letras F, G, J, K, N, P, V, Z, W, Y}

% Rotina neuroinclusiva e badges (R201)
\rotinasimples
\nivelKbadge{K1 --- Letras e Grafias}
\rotabadge{1A --- Recomposicao Intensiva}

Estas miss\~o{}es completam o alfabeto inteiro de A a Z. Cada letra tem
suas atividades de consolida\c{c}\~a{}o, seus sons com IPA, sua fam\'ilia
sil\'abica e sua miss\~a{}o motora --- o mesmo caminho das outras letras.



% ============================================================
% CONSOLIDACAO DA LETRA F (R201.2) -- Missao Alfabeto 1/10
% ============================================================

\subsection*{Consolida\c{c}\~a{}o da Letra F --- Alfabeto Completo}

\plancheta{Miss\~a{}o Alfabeto 1/10 --- Letra F}

\begin{missao}
\textbf{\faGamepad~Miss\~a{}o Alfabeto 1 de 10:} Consolide a letra F, o sopro forte. O F \'e{} o irm\~a{}o surdo do V!
\end{missao}


\boq[labiodental]{F}{%
  Dentes de cima encostam no labio de baixo.\\Ar sai forte e continuo, SEM voz.\\Som: fff (assobio suave).
}

\pistaarticulatoria{fff (fricativo, sem voz)}{Dentes de cima no labio de baixo}{Mao na frente da boca: sente o ar frio}

\begin{consolidacao}
\textbf{Atividade 1 --- Ca\c{c}ada ao F}

\instrucao{Circule o F e sublinhe as palavras com F.}

\begin{center}
\begin{tabular}{cccccc}
\fbox{F} & G & \fbox{F} & J & \fbox{F} & N \\
P & \fbox{F} & G & \fbox{F} & J & \fbox{F}
\end{tabular}
\end{center}
\begin{center}
\uline{fada} \quad \uline{foca} \quad gato \quad \uline{figo} \quad jaca \quad \uline{fumo}
\end{center}

\caixapausa{Se precisar, pare aqui. Respire fundo e continue quando quiser.}

\textbf{Atividade 2 --- Escrita do F --- 4 formas}

\instrucao{Escreva o F cinco vezes em cada forma.}

\begin{center}
\textbf{Imprensa mai\'uscula:} \underline{\hspace{1cm}} \underline{\hspace{1cm}} \underline{\hspace{1cm}} \underline{\hspace{1cm}} \underline{\hspace{1cm}}\\[0.5em]
\textbf{Imprensa min\'uscula:} \underline{\hspace{1cm}} \underline{\hspace{1cm}} \underline{\hspace{1cm}} \underline{\hspace{1cm}} \underline{\hspace{1cm}}\\[0.5em]
\textbf{Cursiva mai\'uscula:} \underline{\hspace{1cm}} \underline{\hspace{1cm}} \underline{\hspace{1cm}} \underline{\hspace{1cm}} \underline{\hspace{1cm}}\\[0.5em]
\textbf{Cursiva min\'uscula:} \underline{\hspace{1cm}} \underline{\hspace{1cm}} \underline{\hspace{1cm}} \underline{\hspace{1cm}} \underline{\hspace{1cm}}
\end{center}

\textbf{Atividade 3 --- Jogo do Sopro}

\instrucao{Complete com F e sinta o ar sair forte pela boca.}

\begin{center}
\underline{~~~}ada \quad \underline{~~~}oca \quad \underline{~~~}igo \quad \underline{~~~}umo \quad \underline{~~~}aca
\end{center}
Diga em voz alta: \textbf{fada, foca, figo, fumo, faca}

\caixapista{F = dentes no l\'abio de baixo + ar forte. M\~a{}o na frente da boca: sente o vento.}
\end{consolidacao}

\plancheta{Sons da Letra F --- IPA}

\begin{sonsdaletra}
Os sons da letra F no portugu\^es{} brasileiro (IPA --- Alfabeto Fon\'etico Internacional):

\somipa[labiodental]{F \'unico}{f}{dentes no l\'abio, ar forte, sem voz}{faca, fada, foca, fuma\c{c}a}
\end{sonsdaletra}

\plancheta{Fam\'ilia Sil\'abica da Letra F}

\begin{familiasilabica}
A fam\'ilia sil\'abica completa da letra F:

\begin{center}
FA \quad FE \quad FI \quad FO \quad FU \\[0.5em] fada \quad f\'e \quad figo \quad foca \quad fumo
\end{center}
\instrucao{Leia a fam\'ilia em voz alta, devagar. Depois cubra e tente lembrar.}
\end{familiasilabica}

\plancheta{Atividade de Som e S\'ilabas da Letra F}

\begin{somletra}
\textbf{1. Ca\c{c}a ao som.} Ou\c{c}a a palavra. Se tiver o som de F, marque o quadradinho:
\begin{center}
$\square$~FACA \quad $\square$~FOGO \quad $\square$~CAF\'E \quad $\square$~BOLA \quad $\square$~M\~AO \quad $\square$~P\'E
\end{center}

\textbf{2. Tem ou n\~ao{} tem?} Rodeie SIM ou N\~ao:
\begin{center}
FACA --- (\textbf{SIM}) (\textbf{N\~AO}) \quad BOLA --- (\textbf{SIM}) (\textbf{N\~AO}) \\ FOGO --- (\textbf{SIM}) (\textbf{N\~AO}) \quad M\~AO --- (\textbf{SIM}) (\textbf{N\~AO})
\end{center}
\caixapausa{O l\'abio de baixo toca os dentes de cima.}
\end{somletra}

\begin{somsilaba}
\textbf{3. Som das s\'ilabas.} Ou\c{c}a, bata palmas e conte as s\'ilabas:
\begin{center}
FA.CA $\rightarrow$ \underline{\hspace{1cm}} s\'ilabas \quad FO.GO $\rightarrow$ \underline{\hspace{1cm}} s\'ilabas
\end{center}

\textbf{4. Qual s\'ilaba tem o som?} Circule a s\'ilaba com o som de F:
\begin{center}
\textbf{FA}.CA \quad \textbf{FO}.GO
\end{center}
\caixapista{A consoante F se junta com a vogal: FA, FE, FI, FO, FU.}
\end{somsilaba}


\plancheta{Miss\~a{}o Motora --- Letra F}

\begin{motricidade}
\textbf{Movimento:} trace o poste e os dois bra\c{c}os do F.
\begin{center}
\begin{tikzpicture}
  \draw[very thick, dashed, color=primary] (0,0) -- (0,1.6);
  \draw[very thick, dashed, color=secondary] (0,1.6) -- (1.4,1.6);
  \draw[very thick, dashed, color=accent] (0,0.9) -- (1.1,0.9);
  \draw[->, color=accent] (0.15,0.6) -- (0.15,0.2);
\end{tikzpicture}
\end{center}
\caixapista{O F \'e{} um poste com dois bra\c{c}os apontando para a direita.}
\end{motricidade}

\plancheta{Ponte --- Da Letra F para a Pr\'oxima}

\begin{transicao}
Compare o F com o G:
\begin{center}
\textbf{fada} / \textbf{gato} \qquad \textbf{foca} / \textbf{gola} \qquad \textbf{figo} / \textbf{gelo}
\end{center}
\instrucao{No F o ar sai forte entre os dentes e o l\'abio. No G o som nasce na garganta.}
\end{transicao}

\plancheta{Recompensa --- Miss\~a{}o Alfabeto 1/10}

\begin{recompensa}
\estrelas{3}\quad \ofensiva{1}\\[0.3em]\barraprogresso{1}{10}\\[0.5em]

\checklistdominio{ler e escrever palavras com F sentindo o sopro forte}
\end{recompensa}



% ============================================================
% CONSOLIDACAO DA LETRA G (R201.2) -- Missao Alfabeto 2/10
% ============================================================

\subsection*{Consolida\c{c}\~a{}o da Letra G --- Alfabeto Completo}

\plancheta{Miss\~a{}o Alfabeto 2/10 --- Letra G}

\begin{missao}
\textbf{\faGamepad~Miss\~a{}o Alfabeto 2 de 10:} Consolide a letra G, a letra de dois sons: ga, go, gu (g forte) e ge, gi (som de j)!
\end{missao}


\boq[velar]{G}{%
  Boca aberta, som sai do fundo.\\Som forte COM voz.\\gato, gola, gude.
}

\pistaarticulatoria{ggg (forte, com voz)}{Fundo da boca, lingua encostada}{Mao na garganta: vibra}

\begin{consolidacao}
\textbf{Atividade 1 --- Ca\c{c}ada ao G}

\instrucao{Circule o G e sublinhe as palavras com G.}

\begin{center}
\begin{tabular}{cccccc}
\fbox{G} & F & \fbox{G} & J & \fbox{G} & N \\
P & \fbox{G} & F & \fbox{G} & J & \fbox{G}
\end{tabular}
\end{center}
\begin{center}
\uline{gato} \quad \uline{gelo} \quad fada \quad \uline{gola} \quad jaca \quad \uline{gude}
\end{center}

\caixapausa{Se precisar, pare aqui. Respire fundo e continue quando quiser.}

\textbf{Atividade 2 --- Escrita do G --- 4 formas}

\instrucao{Escreva o G cinco vezes em cada forma.}

\begin{center}
\textbf{Imprensa mai\'uscula:} \underline{\hspace{1cm}} \underline{\hspace{1cm}} \underline{\hspace{1cm}} \underline{\hspace{1cm}} \underline{\hspace{1cm}}\\[0.5em]
\textbf{Imprensa min\'uscula:} \underline{\hspace{1cm}} \underline{\hspace{1cm}} \underline{\hspace{1cm}} \underline{\hspace{1cm}} \underline{\hspace{1cm}}\\[0.5em]
\textbf{Cursiva mai\'uscula:} \underline{\hspace{1cm}} \underline{\hspace{1cm}} \underline{\hspace{1cm}} \underline{\hspace{1cm}} \underline{\hspace{1cm}}\\[0.5em]
\textbf{Cursiva min\'uscula:} \underline{\hspace{1cm}} \underline{\hspace{1cm}} \underline{\hspace{1cm}} \underline{\hspace{1cm}} \underline{\hspace{1cm}}
\end{center}

\textbf{Atividade 3 --- Jogo do G de Dois Sons}

\instrucao{Ou\c{c}a e marque: G forte (ga, go, gu) ou G de J (ge, gi)?}

\begin{center}
\begin{tabular}{ll}
gato \quad $\rightarrow$ \underline{\hspace{2cm}} & gelo \quad $\rightarrow$ \underline{\hspace{2cm}} \\
gola \quad $\rightarrow$ \underline{\hspace{2cm}} & girafa \quad $\rightarrow$ \underline{\hspace{2cm}} \\
gude \quad $\rightarrow$ \underline{\hspace{2cm}} & gema \quad $\rightarrow$ \underline{\hspace{2cm}}
\end{tabular}
\end{center}
Regra: GA, GO, GU = G forte. GE, GI = som de J.

\caixapista{ge = je, gi = ji. Olhe a vogal que vem depois do G!}
\end{consolidacao}

\plancheta{Sons da Letra G --- IPA}

\begin{sonsdaletra}
Os sons da letra G no portugu\^es{} brasileiro (IPA --- Alfabeto Fon\'etico Internacional):

\somipa[velar]{G antes de A, O, U}{g}{som forte na garganta}{gato, gola, gude}
\somipa[aberta]{G antes de E, I}{Z}{vira som de J}{gelo, girafa, gema, p\'agina}
\end{sonsdaletra}

\plancheta{Fam\'ilia Sil\'abica da Letra G}

\begin{familiasilabica}
A fam\'ilia sil\'abica completa da letra G:

\begin{center}
GA \quad GE \quad GI \quad GO \quad GU \\[0.5em] gato \quad gelo \quad girafa \quad gola \quad gude
\end{center}
\instrucao{Leia a fam\'ilia em voz alta, devagar. Depois cubra e tente lembrar.}
\end{familiasilabica}

\plancheta{Atividade de Som e S\'ilabas da Letra G}

\begin{somletra}
\textbf{1. Ca\c{c}a ao som.} Ou\c{c}a a palavra. Se tiver o som de G, marque o quadradinho:
\begin{center}
$\square$~GATO \quad $\square$~GOMA \quad $\square$~FIGO \quad $\square$~BOLA \quad $\square$~DADO \quad $\square$~P\'E
\end{center}

\textbf{2. Tem ou n\~ao{} tem?} Rodeie SIM ou N\~ao:
\begin{center}
GATO --- (\textbf{SIM}) (\textbf{N\~AO}) \quad BOLA --- (\textbf{SIM}) (\textbf{N\~AO}) \\ GOMA --- (\textbf{SIM}) (\textbf{N\~AO}) \quad DADO --- (\textbf{SIM}) (\textbf{N\~AO})
\end{center}
\caixapausa{O som sai do fundo da boca, com voz.}
\end{somletra}

\begin{somsilaba}
\textbf{3. Som das s\'ilabas.} Ou\c{c}a, bata palmas e conte as s\'ilabas:
\begin{center}
GA.TO $\rightarrow$ \underline{\hspace{1cm}} s\'ilabas \quad FI.GO $\rightarrow$ \underline{\hspace{1cm}} s\'ilabas
\end{center}

\textbf{4. Qual s\'ilaba tem o som?} Circule a s\'ilaba com o som de G:
\begin{center}
\textbf{GA}.TO \quad FI.\textbf{GO}
\end{center}
\caixapista{O G com A, O e U tem som de G (\textbf{GA}, \textbf{GO}, \textbf{GU}).}
\end{somsilaba}


\plancheta{Miss\~a{}o Motora --- Letra G}

\begin{motricidade}
\textbf{Movimento:} trace a meia-lua com a perna interna do G.
\begin{center}
\begin{tikzpicture}
  \draw[very thick, dashed, color=primary] (1.2,0.6) to[out=180,in=90] (0.2,0) to[out=-90,in=180] (1.2,-0.6);
  \draw[very thick, dashed, color=secondary] (0.8,0) -- (1.4,0);
  \draw[->, color=accent] (0.5,-0.35) -- (0.8,-0.35);
\end{tikzpicture}
\end{center}
\caixapista{O G \'e{} um C com uma perna para dentro.}
\end{motricidade}

\plancheta{Ponte --- Da Letra G para a Pr\'oxima}

\begin{transicao}
Compare o G com o J:
\begin{center}
\textbf{gela} / \textbf{jaca} \qquad \textbf{gira} / \textbf{jipe} \qquad \textbf{gelo} / \textbf{jogo}
\end{center}
\instrucao{O G antes de E e I tem o MESMO som do J! Compare: ge = je, gi = ji.}
\end{transicao}

\plancheta{Recompensa --- Miss\~a{}o Alfabeto 2/10}

\begin{recompensa}
\estrelas{3}\quad \ofensiva{2}\\[0.3em]\barraprogresso{2}{10}\\[0.5em]

\checklistdominio{aplicar a regra do G: forte antes de A/O/U, som de J antes de E/I}
\end{recompensa}



% ============================================================
% CONSOLIDACAO DA LETRA J (R201.2) -- Missao Alfabeto 3/10
% ============================================================

\subsection*{Consolida\c{c}\~a{}o da Letra J --- Alfabeto Completo}

\plancheta{Miss\~a{}o Alfabeto 3/10 --- Letra J}

\begin{missao}
\textbf{\faGamepad~Miss\~a{}o Alfabeto 3 de 10:} Consolide a letra J. O J tem um som s\'o, mas aparece no G tamb\'em (ge, gi)!
\end{missao}


\boq[palatal]{J}{%
  Lingua encosta no ceu da boca.\\Som molhado COM voz.\\jaca, jogo, jipe.
}

\pistaarticulatoria{jjj (molhado, com voz)}{Lingua no ceu da boca}{Mao na garganta: vibra}

\begin{consolidacao}
\textbf{Atividade 1 --- Ca\c{c}ada ao J}

\instrucao{Circule o J e sublinhe as palavras com J.}

\begin{center}
\begin{tabular}{cccccc}
\fbox{J} & G & \fbox{J} & F & \fbox{J} & N \\
P & \fbox{J} & G & \fbox{J} & F & \fbox{J}
\end{tabular}
\end{center}
\begin{center}
\uline{jaca} \quad \uline{jogo} \quad gato \quad \uline{jipe} \quad fada \quad \uline{juba}
\end{center}

\caixapausa{Se precisar, pare aqui. Respire fundo e continue quando quiser.}

\textbf{Atividade 2 --- Escrita do J --- 4 formas}

\instrucao{Escreva o J cinco vezes em cada forma.}

\begin{center}
\textbf{Imprensa mai\'uscula:} \underline{\hspace{1cm}} \underline{\hspace{1cm}} \underline{\hspace{1cm}} \underline{\hspace{1cm}} \underline{\hspace{1cm}}\\[0.5em]
\textbf{Imprensa min\'uscula:} \underline{\hspace{1cm}} \underline{\hspace{1cm}} \underline{\hspace{1cm}} \underline{\hspace{1cm}} \underline{\hspace{1cm}}\\[0.5em]
\textbf{Cursiva mai\'uscula:} \underline{\hspace{1cm}} \underline{\hspace{1cm}} \underline{\hspace{1cm}} \underline{\hspace{1cm}} \underline{\hspace{1cm}}\\[0.5em]
\textbf{Cursiva min\'uscula:} \underline{\hspace{1cm}} \underline{\hspace{1cm}} \underline{\hspace{1cm}} \underline{\hspace{1cm}} \underline{\hspace{1cm}}
\end{center}

\textbf{Atividade 3 --- Jogo J x G}

\instrucao{Ou\c{c}a e marque: escreve-se com J ou com G?}

\begin{center}
\begin{tabular}{ll}
jaca \quad $\rightarrow$ \underline{\hspace{2cm}} & gelo \quad $\rightarrow$ \underline{\hspace{2cm}} \\
jogo \quad $\rightarrow$ \underline{\hspace{2cm}} & girafa \quad $\rightarrow$ \underline{\hspace{2cm}} \\
jipe \quad $\rightarrow$ \underline{\hspace{2cm}} & gema \quad $\rightarrow$ \underline{\hspace{2cm}}
\end{tabular}
\end{center}
Dica: JA, JO, JU s\'o com J. JE, JI com J; GE, GI com G.

\caixapista{ja, jo, ju = J. ge, gi = G. S\'o decora\c{c}\~a{}o de algumas palavras!}
\end{consolidacao}

\plancheta{Sons da Letra J --- IPA}

\begin{sonsdaletra}
Os sons da letra J no portugu\^es{} brasileiro (IPA --- Alfabeto Fon\'etico Internacional):

\somipa[aberta]{J \'unico}{Z}{som molhado, com voz}{jaca, jogo, jipe, juba, jejum}
\end{sonsdaletra}

\plancheta{Fam\'ilia Sil\'abica da Letra J}

\begin{familiasilabica}
A fam\'ilia sil\'abica completa da letra J:

\begin{center}
JA \quad JE \quad JI \quad JO \quad JU \\[0.5em] jaca \quad jeca \quad jipe \quad jogo \quad juba
\end{center}
\instrucao{Leia a fam\'ilia em voz alta, devagar. Depois cubra e tente lembrar.}
\end{familiasilabica}

\plancheta{Atividade de Som e S\'ilabas da Letra J}

\begin{somletra}
\textbf{1. Ca\c{c}a ao som.} Ou\c{c}a a palavra. Se tiver o som de J, marque o quadradinho:
\begin{center}
$\square$~JACA \quad $\square$~JANELA \quad $\square$~CAJU \quad $\square$~BOLA \quad $\square$~P\'E \quad $\square$~DADO
\end{center}

\textbf{2. Tem ou n\~ao{} tem?} Rodeie SIM ou N\~ao:
\begin{center}
JACA --- (\textbf{SIM}) (\textbf{N\~AO}) \quad BOLA --- (\textbf{SIM}) (\textbf{N\~AO}) \\ JANELA --- (\textbf{SIM}) (\textbf{N\~AO}) \quad P\'E --- (\textbf{SIM}) (\textbf{N\~AO})
\end{center}
\caixapausa{A l\'ingua toca o c\'eu da boca, com voz.}
\end{somletra}

\begin{somsilaba}
\textbf{3. Som das s\'ilabas.} Ou\c{c}a, bata palmas e conte as s\'ilabas:
\begin{center}
JA.CA $\rightarrow$ \underline{\hspace{1cm}} s\'ilabas \quad JA.NE.LA $\rightarrow$ \underline{\hspace{1cm}} s\'ilabas
\end{center}

\textbf{4. Qual s\'ilaba tem o som?} Circule a s\'ilaba com o som de J:
\begin{center}
\textbf{JA}.CA \quad \textbf{JA}.NE.LA
\end{center}
\caixapista{A consoante J se junta com a vogal: JA, JE, JI, JO, JU.}
\end{somsilaba}


\plancheta{Miss\~a{}o Motora --- Letra J}

\begin{motricidade}
\textbf{Movimento:} trace o gancho do J.
\begin{center}
\begin{tikzpicture}
  \draw[very thick, dashed, color=primary] (0.8,1.2) -- (0.8,0);
  \draw[very thick, dashed, color=secondary] (0.8,0) to[out=-90,in=180] (1.4,-0.5) to[out=0,in=-90] (0.8,-0.9);
  \draw[->, color=accent] (0.95,0.9) -- (0.95,0.4);
\end{tikzpicture}
\end{center}
\caixapista{O J \'e{} um gancho de pescar: desce reto e curva na ponta.}
\end{motricidade}

\plancheta{Ponte --- Da Letra J para a Pr\'oxima}

\begin{transicao}
Compare o J com o K:
\begin{center}
\textbf{jaca} / \textbf{kiwi} \qquad \textbf{jogo} / \textbf{kart} \qquad \textbf{jipe} / \textbf{ketchup}
\end{center}
\instrucao{O J vibra com voz. O K \'e{} seco, como o C de casa.}
\end{transicao}

\plancheta{Recompensa --- Miss\~a{}o Alfabeto 3/10}

\begin{recompensa}
\estrelas{3}\quad \ofensiva{3}\\[0.3em]\barraprogresso{3}{10}\\[0.5em]

\checklistdominio{escrever com J as palavras de som \textbf{ja, jo, ju}}
\end{recompensa}



% ============================================================
% CONSOLIDACAO DA LETRA K (R201.2) -- Missao Alfabeto 4/10
% ============================================================

\subsection*{Consolida\c{c}\~a{}o da Letra K --- Alfabeto Completo}

\plancheta{Miss\~a{}o Alfabeto 4/10 --- Letra K}

\begin{missao}
\textbf{\faGamepad~Miss\~a{}o Alfabeto 4 de 10:} Consolide a letra K, a letra do kiwi! O K vive em palavras estrangeiras.
\end{missao}


\boq[velar]{K}{%
  Igual ao C de casa: som seco.\\Som SEM voz.\\kiwi, kart, ketchup.
}

\pistaarticulatoria{kkk (seco, sem voz)}{Fundo da boca}{Mao na garganta: nao vibra}

\begin{consolidacao}
\textbf{Atividade 1 --- Ca\c{c}ada ao K}

\instrucao{Circule o K e sublinhe as palavras com K.}

\begin{center}
\begin{tabular}{cccccc}
\fbox{K} & J & \fbox{K} & N & \fbox{K} & P \\
V & \fbox{K} & J & \fbox{K} & N & \fbox{K}
\end{tabular}
\end{center}
\begin{center}
\uline{kiwi} \quad \uline{kart} \quad jaca \quad \uline{ketchup} \quad nada \quad \uline{karaok\^e}
\end{center}

\caixapausa{Se precisar, pare aqui. Respire fundo e continue quando quiser.}

\textbf{Atividade 2 --- Escrita do K --- 4 formas}

\instrucao{Escreva o K cinco vezes em cada forma.}

\begin{center}
\textbf{Imprensa mai\'uscula:} \underline{\hspace{1cm}} \underline{\hspace{1cm}} \underline{\hspace{1cm}} \underline{\hspace{1cm}} \underline{\hspace{1cm}}\\[0.5em]
\textbf{Imprensa min\'uscula:} \underline{\hspace{1cm}} \underline{\hspace{1cm}} \underline{\hspace{1cm}} \underline{\hspace{1cm}} \underline{\hspace{1cm}}\\[0.5em]
\textbf{Cursiva mai\'uscula:} \underline{\hspace{1cm}} \underline{\hspace{1cm}} \underline{\hspace{1cm}} \underline{\hspace{1cm}} \underline{\hspace{1cm}}\\[0.5em]
\textbf{Cursiva min\'uscula:} \underline{\hspace{1cm}} \underline{\hspace{1cm}} \underline{\hspace{1cm}} \underline{\hspace{1cm}} \underline{\hspace{1cm}}
\end{center}

\textbf{Atividade 3 --- Jogo K x C}

\instrucao{Ou\c{c}a e marque: som de K com K ou com C?}

\begin{center}
\begin{tabular}{ll}
kiwi \quad $\rightarrow$ \underline{\hspace{2cm}} & casa \quad $\rightarrow$ \underline{\hspace{2cm}} \\
kart \quad $\rightarrow$ \underline{\hspace{2cm}} & copo \quad $\rightarrow$ \underline{\hspace{2cm}}
\end{tabular}
\end{center}
Regra: o K aparece em palavras que vieram de outras l\'inguas.

\caixapista{K = C duro. S\'o muda a letra, o som \'e{} o mesmo!}
\end{consolidacao}

\plancheta{Sons da Letra K --- IPA}

\begin{sonsdaletra}
Os sons da letra K no portugu\^es{} brasileiro (IPA --- Alfabeto Fon\'etico Internacional):

\somipa[velar]{K \'unico}{k}{som seco, igual ao C de casa}{kiwi, kart, ketchup, karaok\^e}
\end{sonsdaletra}

\plancheta{Fam\'ilia Sil\'abica da Letra K}

\begin{familiasilabica}
A fam\'ilia sil\'abica completa da letra K:

\begin{center}
KA \quad KE \quad KI \quad KO \quad KU \\[0.5em] karaok\^e \quad ketchup \quad kiwi \quad --- \quad kung fu
\end{center}
\instrucao{Leia a fam\'ilia em voz alta, devagar. Depois cubra e tente lembrar.}
\end{familiasilabica}

\plancheta{Atividade de Som e S\'ilabas da Letra K}

\begin{somletra}
\textbf{1. Ca\c{c}a ao som.} Ou\c{c}a a palavra. Se tiver o som de K, marque o quadradinho:
\begin{center}
$\square$~KIWI \quad $\square$~KARAOK\^E \quad $\square$~KOMBI \quad $\square$~BOLA \quad $\square$~P\'E \quad $\square$~M\~AO
\end{center}

\textbf{2. Tem ou n\~ao{} tem?} Rodeie SIM ou N\~ao:
\begin{center}
KIWI --- (\textbf{SIM}) (\textbf{N\~AO}) \quad BOLA --- (\textbf{SIM}) (\textbf{N\~AO}) \\ KARAOK\^E --- (\textbf{SIM}) (\textbf{N\~AO}) \quad P\'E --- (\textbf{SIM}) (\textbf{N\~AO})
\end{center}
\caixapausa{O som sai do fundo da boca, sem voz.}
\end{somletra}

\begin{somsilaba}
\textbf{3. Som das s\'ilabas.} Ou\c{c}a, bata palmas e conte as s\'ilabas:
\begin{center}
KI.WI $\rightarrow$ \underline{\hspace{1cm}} s\'ilabas \quad KA.RA.O.K\^E $\rightarrow$ \underline{\hspace{1cm}} s\'ilabas
\end{center}

\textbf{4. Qual s\'ilaba tem o som?} Circule a s\'ilaba com o som de K:
\begin{center}
\textbf{KI}.WI \quad \textbf{KA}.RA.O.K\^E
\end{center}
\caixapista{O K se junta com a vogal: KA, KE, KI, KO, KU.}
\end{somsilaba}


\plancheta{Miss\~a{}o Motora --- Letra K}

\begin{motricidade}
\textbf{Movimento:} trace o poste com os dois bra\c{c}os em V do K.
\begin{center}
\begin{tikzpicture}
  \draw[very thick, dashed, color=primary] (0,0) -- (0,1.6);
  \draw[very thick, dashed, color=secondary] (1.2,1.6) -- (0,0.8) -- (1.2,0);
  \draw[->, color=accent] (0.9,1.45) -- (0.5,1.0);
  \draw[->, color=accent] (0.5,0.55) -- (0.9,0.15);
\end{tikzpicture}
\end{center}
\caixapista{O K \'e{} um poste que recebe dois bra\c{c}os em V no meio.}
\end{motricidade}

\plancheta{Ponte --- Da Letra K para a Pr\'oxima}

\begin{transicao}
Compare o K com o N:
\begin{center}
\textbf{kiwi} / \textbf{nada} \qquad \textbf{kart} / \textbf{ninho} \qquad \textbf{ketchup} / \textbf{nuvem}
\end{center}
\instrucao{O K \'e{} seco, sem voz. O N vibra e sai pelo nariz.}
\end{transicao}

\plancheta{Recompensa --- Miss\~a{}o Alfabeto 4/10}

\begin{recompensa}
\estrelas{3}\quad \ofensiva{4}\\[0.3em]\barraprogresso{4}{10}\\[0.5em]

\checklistdominio{reconhecer o K e saber que ele soa como o C de casa}
\end{recompensa}



% ============================================================
% CONSOLIDACAO DA LETRA N (R201.2) -- Missao Alfabeto 5/10
% ============================================================

\subsection*{Consolida\c{c}\~a{}o da Letra N --- Alfabeto Completo}

\plancheta{Miss\~a{}o Alfabeto 5/10 --- Letra N}

\begin{missao}
\textbf{\faGamepad~Miss\~a{}o Alfabeto 5 de 10:} Consolide a letra N. O N vibra no nariz, como o M!
\end{missao}


\boq[dental]{N}{%
  Lingua no ceu da boca, perto dos dentes.\\Som nasal COM voz.\\nada, nena, ninho, no, nuvem.
}

\pistaarticulatoria{nnn (nasal, com voz)}{Lingua no ceu da boca}{Mao no nariz: sente a vibracao}

\begin{consolidacao}
\textbf{Atividade 1 --- Ca\c{c}ada ao N}

\instrucao{Circule o N e sublinhe as palavras com N.}

\begin{center}
\begin{tabular}{cccccc}
\fbox{N} & P & \fbox{N} & V & \fbox{N} & Z \\
K & \fbox{N} & P & \fbox{N} & V & \fbox{N}
\end{tabular}
\end{center}
\begin{center}
\uline{nada} \quad \uline{ninho} \quad pato \quad \uline{nuvem} \quad vaca \quad \uline{nena}
\end{center}

\caixapausa{Se precisar, pare aqui. Respire fundo e continue quando quiser.}

\textbf{Atividade 2 --- Escrita do N --- 4 formas}

\instrucao{Escreva o N cinco vezes em cada forma.}

\begin{center}
\textbf{Imprensa mai\'uscula:} \underline{\hspace{1cm}} \underline{\hspace{1cm}} \underline{\hspace{1cm}} \underline{\hspace{1cm}} \underline{\hspace{1cm}}\\[0.5em]
\textbf{Imprensa min\'uscula:} \underline{\hspace{1cm}} \underline{\hspace{1cm}} \underline{\hspace{1cm}} \underline{\hspace{1cm}} \underline{\hspace{1cm}}\\[0.5em]
\textbf{Cursiva mai\'uscula:} \underline{\hspace{1cm}} \underline{\hspace{1cm}} \underline{\hspace{1cm}} \underline{\hspace{1cm}} \underline{\hspace{1cm}}\\[0.5em]
\textbf{Cursiva min\'uscula:} \underline{\hspace{1cm}} \underline{\hspace{1cm}} \underline{\hspace{1cm}} \underline{\hspace{1cm}} \underline{\hspace{1cm}}
\end{center}

\textbf{Atividade 3 --- Jogo N x M}

\instrucao{Complete com N ou M. Sinta o nariz vibrar!}

\begin{center}
\underline{~~~}ada \quad ca\underline{~~~}to \quad \underline{~~~}inho \quad o\underline{~~~}da \quad \underline{~~~}uvem \quad ba\underline{~~~}da
\end{center}
Diga em voz alta: \textbf{nada, canto, ninho, onda, nuvem, banda}

\caixapista{N e M s\~a{}o irm\~a{}os do nariz. No fim da s\'ilaba, anasalam a vogal.}
\end{consolidacao}

\plancheta{Sons da Letra N --- IPA}

\begin{sonsdaletra}
Os sons da letra N no portugu\^es{} brasileiro (IPA --- Alfabeto Fon\'etico Internacional):

\somipa[dental]{N no in\'icio}{n}{l\'ingua no c\'eu da boca, som nasal}{nada, nena, ninho, n\'o, nuvem}
\somipa[nasal]{N no fim (nasalizador)}{\~}{anasala a vogal antes dele}{canto [k6\~tu], onda [o\~da], banda}
\end{sonsdaletra}

\plancheta{Fam\'ilia Sil\'abica da Letra N}

\begin{familiasilabica}
A fam\'ilia sil\'abica completa da letra N:

\begin{center}
NA \quad NE \quad NI \quad NO \quad NU \\[0.5em] nada \quad nena \quad ninho \quad n\'o \quad nuvem
\end{center}
\instrucao{Leia a fam\'ilia em voz alta, devagar. Depois cubra e tente lembrar.}
\end{familiasilabica}

\plancheta{Atividade de Som e S\'ilabas da Letra N}

\begin{somletra}
\textbf{1. Ca\c{c}a ao som.} Ou\c{c}a a palavra. Se tiver o som de N, marque o quadradinho:
\begin{center}
$\square$~NAVIO \quad $\square$~NINHO \quad $\square$~CANA \quad $\square$~BOLA \quad $\square$~GATO \quad $\square$~P\'E
\end{center}

\textbf{2. Tem ou n\~ao{} tem?} Rodeie SIM ou N\~ao:
\begin{center}
NAVIO --- (\textbf{SIM}) (\textbf{N\~AO}) \quad BOLA --- (\textbf{SIM}) (\textbf{N\~AO}) \\ NINHO --- (\textbf{SIM}) (\textbf{N\~AO}) \quad GATO --- (\textbf{SIM}) (\textbf{N\~AO})
\end{center}
\caixapausa{O som sai pelo nariz.}
\end{somletra}

\begin{somsilaba}
\textbf{3. Som das s\'ilabas.} Ou\c{c}a, bata palmas e conte as s\'ilabas:
\begin{center}
NA.VI.O $\rightarrow$ \underline{\hspace{1cm}} s\'ilabas \quad NI.NHO $\rightarrow$ \underline{\hspace{1cm}} s\'ilabas
\end{center}

\textbf{4. Qual s\'ilaba tem o som?} Circule a s\'ilaba com o som de N:
\begin{center}
\textbf{NA}.VI.O \quad \textbf{NI}.NHO
\end{center}
\caixapista{A consoante N se junta com a vogal: NA, NE, NI, NO, NU.}
\end{somsilaba}


\plancheta{Miss\~a{}o Motora --- Letra N}

\begin{motricidade}
\textbf{Movimento:} trace as duas pontes do N.
\begin{center}
\begin{tikzpicture}
  \draw[very thick, dashed, color=primary] (0,0) -- (0,1.4);
  \draw[very thick, dashed, color=primary] (0,1.4) -- (1.1,0);
  \draw[very thick, dashed, color=primary] (1.1,0) -- (1.1,1.4);
  \draw[->, color=accent] (0.15,0.9) -- (0.15,0.3);
  \draw[->, color=accent] (0.3,1.2) -- (0.75,0.55);
\end{tikzpicture}
\end{center}
\caixapista{O N \'e{} um poste, uma descida em diagonal e outro poste.}
\end{motricidade}

\plancheta{Ponte --- Da Letra N para a Pr\'oxima}

\begin{transicao}
Compare o N com o P:
\begin{center}
\textbf{nada} / \textbf{pato} \qquad \textbf{ninho} / \textbf{pipa} \qquad \textbf{nuvem} / \textbf{pote}
\end{center}
\instrucao{No N o som sai pelo nariz. No P os l\'abios fecham e explodem.}
\end{transicao}

\plancheta{Recompensa --- Miss\~a{}o Alfabeto 5/10}

\begin{recompensa}
\estrelas{3}\quad \ofensiva{5}\\[0.3em]\barraprogresso{5}{10}\\[0.5em]

\checklistdominio{ler palavras com N e perceber a nasaliza\c{c}\~a{}o (canto, onda)}
\end{recompensa}



% ============================================================
% CONSOLIDACAO DA LETRA P (R201.2) -- Missao Alfabeto 6/10
% ============================================================

\subsection*{Consolida\c{c}\~a{}o da Letra P --- Alfabeto Completo}

\plancheta{Miss\~a{}o Alfabeto 6/10 --- Letra P}

\begin{missao}
\textbf{\faGamepad~Miss\~a{}o Alfabeto 6 de 10:} Consolide a letra P, a explos\~a{}o dos l\'abios! O P \'e{} o irm\~a{}o surdo do B.
\end{missao}


\boq[fechada]{P}{%
  Labios fechados que explodem de repente.\\Som seco, SEM voz.\\pato, papai, pipa, pote, pulo.
}

\pistaarticulatoria{ppp (explosivo, sem voz)}{Labios fechados que abrem}{Mao na garganta: nao vibra}

\begin{consolidacao}
\textbf{Atividade 1 --- Ca\c{c}ada ao P}

\instrucao{Circule o P e sublinhe as palavras com P.}

\begin{center}
\begin{tabular}{cccccc}
\fbox{P} & V & \fbox{P} & Z & \fbox{P} & N \\
K & \fbox{P} & V & \fbox{P} & Z & \fbox{P}
\end{tabular}
\end{center}
\begin{center}
\uline{pato} \quad \uline{pipa} \quad vaca \quad \uline{pote} \quad zebra \quad \uline{papai}
\end{center}

\caixapausa{Se precisar, pare aqui. Respire fundo e continue quando quiser.}

\textbf{Atividade 2 --- Escrita do P --- 4 formas}

\instrucao{Escreva o P cinco vezes em cada forma.}

\begin{center}
\textbf{Imprensa mai\'uscula:} \underline{\hspace{1cm}} \underline{\hspace{1cm}} \underline{\hspace{1cm}} \underline{\hspace{1cm}} \underline{\hspace{1cm}}\\[0.5em]
\textbf{Imprensa min\'uscula:} \underline{\hspace{1cm}} \underline{\hspace{1cm}} \underline{\hspace{1cm}} \underline{\hspace{1cm}} \underline{\hspace{1cm}}\\[0.5em]
\textbf{Cursiva mai\'uscula:} \underline{\hspace{1cm}} \underline{\hspace{1cm}} \underline{\hspace{1cm}} \underline{\hspace{1cm}} \underline{\hspace{1cm}}\\[0.5em]
\textbf{Cursiva min\'uscula:} \underline{\hspace{1cm}} \underline{\hspace{1cm}} \underline{\hspace{1cm}} \underline{\hspace{1cm}} \underline{\hspace{1cm}}
\end{center}

\textbf{Atividade 3 --- Jogo P x B}

\instrucao{Ou\c{c}a e marque: P (sem voz) ou B (com voz)?}

\begin{center}
\begin{tabular}{ll}
pato \quad $\rightarrow$ \underline{\hspace{2cm}} & bola \quad $\rightarrow$ \underline{\hspace{2cm}} \\
pipa \quad $\rightarrow$ \underline{\hspace{2cm}} & baba \quad $\rightarrow$ \underline{\hspace{2cm}} \\
pote \quad $\rightarrow$ \underline{\hspace{2cm}} & beb\^e \quad $\rightarrow$ \underline{\hspace{2cm}}
\end{tabular}
\end{center}
Teste: m\~a{}o na garganta. No B vibra, no P n\~a{}o!

\caixapista{P e B usam os mesmos l\'abios. A diferen\c{c}a \'e{} a voz: P surdo, B sonoro.}
\end{consolidacao}

\plancheta{Sons da Letra P --- IPA}

\begin{sonsdaletra}
Os sons da letra P no portugu\^es{} brasileiro (IPA --- Alfabeto Fon\'etico Internacional):

\somipa[fechada]{P \'unico}{p}{l\'abios fechados que explodem sem voz}{pato, papai, pipa, pote, pulo}
\end{sonsdaletra}

\plancheta{Fam\'ilia Sil\'abica da Letra P}

\begin{familiasilabica}
A fam\'ilia sil\'abica completa da letra P:

\begin{center}
PA \quad PE \quad PI \quad PO \quad PU \\[0.5em] pato \quad pele \quad pipa \quad pote \quad pula
\end{center}
\instrucao{Leia a fam\'ilia em voz alta, devagar. Depois cubra e tente lembrar.}
\end{familiasilabica}

\plancheta{Atividade de Som e S\'ilabas da Letra P}

\begin{somletra}
\textbf{1. Ca\c{c}a ao som.} Ou\c{c}a a palavra. Se tiver o som de P, marque o quadradinho:
\begin{center}
$\square$~PATO \quad $\square$~PIPA \quad $\square$~SAPO \quad $\square$~BOLA \quad $\square$~M\~AO \quad $\square$~SOL
\end{center}

\textbf{2. Tem ou n\~ao{} tem?} Rodeie SIM ou N\~ao:
\begin{center}
PATO --- (\textbf{SIM}) (\textbf{N\~AO}) \quad BOLA --- (\textbf{SIM}) (\textbf{N\~AO}) \\ PIPA --- (\textbf{SIM}) (\textbf{N\~AO}) \quad M\~AO --- (\textbf{SIM}) (\textbf{N\~AO})
\end{center}
\caixapausa{A boca fecha e o ar estoura.}
\end{somletra}

\begin{somsilaba}
\textbf{3. Som das s\'ilabas.} Ou\c{c}a, bata palmas e conte as s\'ilabas:
\begin{center}
PA.TO $\rightarrow$ \underline{\hspace{1cm}} s\'ilabas \quad PI.PA $\rightarrow$ \underline{\hspace{1cm}} s\'ilabas
\end{center}

\textbf{4. Qual s\'ilaba tem o som?} Circule a s\'ilaba com o som de P:
\begin{center}
\textbf{PA}.TO \quad \textbf{PI}.PA
\end{center}
\caixapista{A consoante P se junta com a vogal: PA, PE, PI, PO, PU.}
\end{somsilaba}


\plancheta{Miss\~a{}o Motora --- Letra P}

\begin{motricidade}
\textbf{Movimento:} trace o poste com a barriga de cima do P.
\begin{center}
\begin{tikzpicture}
  \draw[very thick, dashed, color=primary] (0,0) -- (0,1.6);
  \draw[very thick, dashed, color=secondary] (0,1.3) to[out=0,in=180] (0.6,1.0) to[out=0,in=180] (0,0.7);
  \draw[->, color=accent] (0.15,0.5) -- (0.15,0.1);
\end{tikzpicture}
\end{center}
\caixapista{O P \'e{} um poste com uma barriga s\'o, l\'a{} em cima.}
\end{motricidade}

\plancheta{Ponte --- Da Letra P para a Pr\'oxima}

\begin{transicao}
Compare o P com o V:
\begin{center}
\textbf{pato} / \textbf{vaca} \qquad \textbf{pipa} / \textbf{vida} \qquad \textbf{pote} / \textbf{vela}
\end{center}
\instrucao{No P os l\'abios explodem sem voz. No V os dentes tocam o l\'abio e a garganta vibra.}
\end{transicao}

\plancheta{Recompensa --- Miss\~a{}o Alfabeto 6/10}

\begin{recompensa}
\estrelas{3}\quad \ofensiva{6}\\[0.3em]\barraprogresso{6}{10}\\[0.5em]

\checklistdominio{diferenciar P de B pela vibra\c{c}\~a{}o da garganta}
\end{recompensa}



% ============================================================
% CONSOLIDACAO DA LETRA V (R201.2) -- Missao Alfabeto 7/10
% ============================================================

\subsection*{Consolida\c{c}\~a{}o da Letra V --- Alfabeto Completo}

\plancheta{Miss\~a{}o Alfabeto 7/10 --- Letra V}

\begin{missao}
\textbf{\faGamepad~Miss\~a{}o Alfabeto 7 de 10:} Consolide a letra V, o sopro com voz. O V \'e{} o irm\~a{}o sonoro do F!
\end{missao}


\boq[labiodental]{V}{%
  Dentes de cima no labio de baixo.\\Ar forte COM voz.\\vaca, vela, vida, vovo, vulcao.
}

\pistaarticulatoria{vvv (fricativo, com voz)}{Dentes de cima no labio de baixo}{Mao na garganta: vibra}

\begin{consolidacao}
\textbf{Atividade 1 --- Ca\c{c}ada ao V}

\instrucao{Circule o V e sublinhe as palavras com V.}

\begin{center}
\begin{tabular}{cccccc}
\fbox{V} & Z & \fbox{V} & P & \fbox{V} & N \\
K & \fbox{V} & Z & \fbox{V} & P & \fbox{V}
\end{tabular}
\end{center}
\begin{center}
\uline{vaca} \quad \uline{vela} \quad zebra \quad \uline{vida} \quad pato \quad \uline{vov\'o}
\end{center}

\caixapausa{Se precisar, pare aqui. Respire fundo e continue quando quiser.}

\textbf{Atividade 2 --- Escrita do V --- 4 formas}

\instrucao{Escreva o V cinco vezes em cada forma.}

\begin{center}
\textbf{Imprensa mai\'uscula:} \underline{\hspace{1cm}} \underline{\hspace{1cm}} \underline{\hspace{1cm}} \underline{\hspace{1cm}} \underline{\hspace{1cm}}\\[0.5em]
\textbf{Imprensa min\'uscula:} \underline{\hspace{1cm}} \underline{\hspace{1cm}} \underline{\hspace{1cm}} \underline{\hspace{1cm}} \underline{\hspace{1cm}}\\[0.5em]
\textbf{Cursiva mai\'uscula:} \underline{\hspace{1cm}} \underline{\hspace{1cm}} \underline{\hspace{1cm}} \underline{\hspace{1cm}} \underline{\hspace{1cm}}\\[0.5em]
\textbf{Cursiva min\'uscula:} \underline{\hspace{1cm}} \underline{\hspace{1cm}} \underline{\hspace{1cm}} \underline{\hspace{1cm}} \underline{\hspace{1cm}}
\end{center}

\textbf{Atividade 3 --- Jogo V x F}

\instrucao{Ou\c{c}a e marque: V (com voz) ou F (sem voz)?}

\begin{center}
\begin{tabular}{ll}
vaca \quad $\rightarrow$ \underline{\hspace{2cm}} & fada \quad $\rightarrow$ \underline{\hspace{2cm}} \\
vela \quad $\rightarrow$ \underline{\hspace{2cm}} & foca \quad $\rightarrow$ \underline{\hspace{2cm}} \\
vida \quad $\rightarrow$ \underline{\hspace{2cm}} & figo \quad $\rightarrow$ \underline{\hspace{2cm}}
\end{tabular}
\end{center}
Teste: m\~a{}o na garganta. No V vibra, no F n\~a{}o!

\caixapista{V e F usam os mesmos dentes. A diferen\c{c}a \'e{} a voz: V sonoro, F surdo.}
\end{consolidacao}

\plancheta{Sons da Letra V --- IPA}

\begin{sonsdaletra}
Os sons da letra V no portugu\^es{} brasileiro (IPA --- Alfabeto Fon\'etico Internacional):

\somipa[labiodental]{V \'unico}{v}{dentes no l\'abio, com voz}{vaca, vela, vida, vov\'o, vulc\~a{}o}
\end{sonsdaletra}

\plancheta{Fam\'ilia Sil\'abica da Letra V}

\begin{familiasilabica}
A fam\'ilia sil\'abica completa da letra V:

\begin{center}
VA \quad VE \quad VI \quad VO \quad VU \\[0.5em] vaca \quad vela \quad vida \quad vov\'o \quad vulc\~a{}o
\end{center}
\instrucao{Leia a fam\'ilia em voz alta, devagar. Depois cubra e tente lembrar.}
\end{familiasilabica}

\plancheta{Atividade de Som e S\'ilabas da Letra V}

\begin{somletra}
\textbf{1. Ca\c{c}a ao som.} Ou\c{c}a a palavra. Se tiver o som de V, marque o quadradinho:
\begin{center}
$\square$~VACA \quad $\square$~VELA \quad $\square$~NOVO \quad $\square$~BOLA \quad $\square$~P\'E \quad $\square$~M\~AO
\end{center}

\textbf{2. Tem ou n\~ao{} tem?} Rodeie SIM ou N\~ao:
\begin{center}
VACA --- (\textbf{SIM}) (\textbf{N\~AO}) \quad BOLA --- (\textbf{SIM}) (\textbf{N\~AO}) \\ VELA --- (\textbf{SIM}) (\textbf{N\~AO}) \quad P\'E --- (\textbf{SIM}) (\textbf{N\~AO})
\end{center}
\caixapausa{O l\'abio de baixo toca os dentes de cima, com voz.}
\end{somletra}

\begin{somsilaba}
\textbf{3. Som das s\'ilabas.} Ou\c{c}a, bata palmas e conte as s\'ilabas:
\begin{center}
VA.CA $\rightarrow$ \underline{\hspace{1cm}} s\'ilabas \quad VE.LA $\rightarrow$ \underline{\hspace{1cm}} s\'ilabas
\end{center}

\textbf{4. Qual s\'ilaba tem o som?} Circule a s\'ilaba com o som de V:
\begin{center}
\textbf{VA}.CA \quad \textbf{VE}.LA
\end{center}
\caixapista{A consoante V se junta com a vogal: VA, VE, VI, VO, VU.}
\end{somsilaba}


\plancheta{Miss\~a{}o Motora --- Letra V}

\begin{motricidade}
\textbf{Movimento:} trace os dois postes inclinados do V.
\begin{center}
\begin{tikzpicture}
  \draw[very thick, dashed, color=primary] (0,1.2) -- (0.5,0);
  \draw[very thick, dashed, color=primary] (0.5,0) -- (1.0,1.2);
  \draw[->, color=accent] (0.15,0.9) -- (0.35,0.5);
  \draw[->, color=accent] (0.65,0.5) -- (0.85,0.9);
\end{tikzpicture}
\end{center}
\caixapista{O V \'e{} um bico: desce de um lado e sobe do outro.}
\end{motricidade}

\plancheta{Ponte --- Da Letra V para a Pr\'oxima}

\begin{transicao}
Compare o V com o Z:
\begin{center}
\textbf{vaca} / \textbf{zebra} \qquad \textbf{vida} / \textbf{zelo} \qquad \textbf{vela} / \textbf{z\'iper}
\end{center}
\instrucao{O V vibra com os dentes no l\'abio. O Z vibra com a l\'ingua perto dos dentes.}
\end{transicao}

\plancheta{Recompensa --- Miss\~a{}o Alfabeto 7/10}

\begin{recompensa}
\estrelas{3}\quad \ofensiva{7}\\[0.3em]\barraprogresso{7}{10}\\[0.5em]

\checklistdominio{diferenciar V de F pela vibra\c{c}\~a{}o da garganta}
\end{recompensa}



% ============================================================
% CONSOLIDACAO DA LETRA Z (R201.2) -- Missao Alfabeto 8/10
% ============================================================

\subsection*{Consolida\c{c}\~a{}o da Letra Z --- Alfabeto Completo}

\plancheta{Miss\~a{}o Alfabeto 8/10 --- Letra Z}

\begin{missao}
\textbf{\faGamepad~Miss\~a{}o Alfabeto 8 de 10:} Consolide a letra Z, a letra que zune como abelha!
\end{missao}


\boq[dental]{Z}{%
  Lingua atras dos dentes de cima.\\Som de abelha, COM voz.\\zebra, zelo, ziper, zoologico.
}

\pistaarticulatoria{zzz (de abelha, com voz)}{Lingua atras dos dentes}{Mao na garganta: vibra}

\begin{consolidacao}
\textbf{Atividade 1 --- Ca\c{c}ada ao Z}

\instrucao{Circule o Z e sublinhe as palavras com Z.}

\begin{center}
\begin{tabular}{cccccc}
\fbox{Z} & W & \fbox{Z} & V & \fbox{Z} & P \\
N & \fbox{Z} & W & \fbox{Z} & V & \fbox{Z}
\end{tabular}
\end{center}
\begin{center}
\uline{zebra} \quad \uline{zelo} \quad waffle \quad \uline{z\'iper} \quad vaca \quad \uline{zumbido}
\end{center}

\caixapausa{Se precisar, pare aqui. Respire fundo e continue quando quiser.}

\textbf{Atividade 2 --- Escrita do Z --- 4 formas}

\instrucao{Escreva o Z cinco vezes em cada forma.}

\begin{center}
\textbf{Imprensa mai\'uscula:} \underline{\hspace{1cm}} \underline{\hspace{1cm}} \underline{\hspace{1cm}} \underline{\hspace{1cm}} \underline{\hspace{1cm}}\\[0.5em]
\textbf{Imprensa min\'uscula:} \underline{\hspace{1cm}} \underline{\hspace{1cm}} \underline{\hspace{1cm}} \underline{\hspace{1cm}} \underline{\hspace{1cm}}\\[0.5em]
\textbf{Cursiva mai\'uscula:} \underline{\hspace{1cm}} \underline{\hspace{1cm}} \underline{\hspace{1cm}} \underline{\hspace{1cm}} \underline{\hspace{1cm}}\\[0.5em]
\textbf{Cursiva min\'uscula:} \underline{\hspace{1cm}} \underline{\hspace{1cm}} \underline{\hspace{1cm}} \underline{\hspace{1cm}} \underline{\hspace{1cm}}
\end{center}

\textbf{Atividade 3 --- Jogo do Z que Vira S}

\instrucao{Ou\c{c}a e marque: Z de abelha ou Z que vira assobio?}

\begin{center}
\begin{tabular}{ll}
zebra \quad $\rightarrow$ \underline{\hspace{2cm}} & luz \quad $\rightarrow$ \underline{\hspace{2cm}} \\
zelo \quad $\rightarrow$ \underline{\hspace{2cm}} & voz \quad $\rightarrow$ \underline{\hspace{2cm}} \\
z\'iper \quad $\rightarrow$ \underline{\hspace{2cm}} & paz \quad $\rightarrow$ \underline{\hspace{2cm}}
\end{tabular}
\end{center}
Regra: no come\c{c}o, zumbe (z). No fim, assobia (s).

\caixapista{luz, voz, paz --- o Z do fim assobia como S.}
\end{consolidacao}

\plancheta{Sons da Letra Z --- IPA}

\begin{sonsdaletra}
Os sons da letra Z no portugu\^es{} brasileiro (IPA --- Alfabeto Fon\'etico Internacional):

\somipa[dental]{Z no in\'icio}{z}{som de abelha, com voz}{zebra, zelo, z\'iper, zool\'ogico}
\somipa[dental]{Z no fim de palavra}{s}{vira assobio}{luz, voz, paz, raiz}
\end{sonsdaletra}

\plancheta{Fam\'ilia Sil\'abica da Letra Z}

\begin{familiasilabica}
A fam\'ilia sil\'abica completa da letra Z:

\begin{center}
ZA \quad ZE \quad ZI \quad ZO \quad ZU \\[0.5em] zebra \quad zelo \quad z\'iper \quad zool\'ogico \quad zumbido
\end{center}
\instrucao{Leia a fam\'ilia em voz alta, devagar. Depois cubra e tente lembrar.}
\end{familiasilabica}

\plancheta{Atividade de Som e S\'ilabas da Letra Z}

\begin{somletra}
\textbf{1. Ca\c{c}a ao som.} Ou\c{c}a a palavra. Se tiver o som de Z, marque o quadradinho:
\begin{center}
$\square$~ZEBRA \quad $\square$~ZERO \quad $\square$~AMIZADE \quad $\square$~BOLA \quad $\square$~P\'E \quad $\square$~DADO
\end{center}

\textbf{2. Tem ou n\~ao{} tem?} Rodeie SIM ou N\~ao:
\begin{center}
ZEBRA --- (\textbf{SIM}) (\textbf{N\~AO}) \quad BOLA --- (\textbf{SIM}) (\textbf{N\~AO}) \\ ZERO --- (\textbf{SIM}) (\textbf{N\~AO}) \quad P\'E --- (\textbf{SIM}) (\textbf{N\~AO})
\end{center}
\caixapausa{Os dentes juntos, com voz ligada.}
\end{somletra}

\begin{somsilaba}
\textbf{3. Som das s\'ilabas.} Ou\c{c}a, bata palmas e conte as s\'ilabas:
\begin{center}
ZE.BRA $\rightarrow$ \underline{\hspace{1cm}} s\'ilabas \quad A.MI.ZA.DE $\rightarrow$ \underline{\hspace{1cm}} s\'ilabas
\end{center}

\textbf{4. Qual s\'ilaba tem o som?} Circule a s\'ilaba com o som de Z:
\begin{center}
\textbf{ZE}.BRA \quad A.MI.\textbf{ZA}.DE
\end{center}
\caixapista{A consoante Z se junta com a vogal: ZA, ZE, ZI, ZO, ZU.}
\end{somsilaba}


\plancheta{Miss\~a{}o Motora --- Letra Z}

\begin{motricidade}
\textbf{Movimento:} trace o zigue-zague de tr\^es linhas do Z.
\begin{center}
\begin{tikzpicture}
  \draw[very thick, dashed, color=primary] (0,1.2) -- (1.4,1.2);
  \draw[very thick, dashed, color=primary] (1.4,1.2) -- (0,0);
  \draw[very thick, dashed, color=primary] (0,0) -- (1.4,0);
  \draw[->, color=accent] (0.3,1.05) -- (0.9,1.05);
  \draw[->, color=accent] (1.1,0.8) -- (0.6,0.35);
\end{tikzpicture}
\end{center}
\caixapista{O Z \'e{} um raio: teto, descida em diagonal e ch\~a{}o.}
\end{motricidade}

\plancheta{Ponte --- Da Letra Z para a Pr\'oxima}

\begin{transicao}
Compare o Z com o W:
\begin{center}
\textbf{zebra} / \textbf{web} \qquad \textbf{z\'iper} / \textbf{waffle} \qquad \textbf{zumbido} / \textbf{kiwi}
\end{center}
\instrucao{O Z zumbe com voz. O W s\'o aparece em palavras estrangeiras e pode soar como U ou V.}
\end{transicao}

\plancheta{Recompensa --- Miss\~a{}o Alfabeto 8/10}

\begin{recompensa}
\estrelas{3}\quad \ofensiva{8}\\[0.3em]\barraprogresso{8}{10}\\[0.5em]

\checklistdominio{reconhecer o Z que zumbe no in\'icio e assobia no fim (luz, voz)}
\end{recompensa}



% ============================================================
% CONSOLIDACAO DA LETRA W (R201.2) -- Missao Alfabeto 9/10
% ============================================================

\subsection*{Consolida\c{c}\~a{}o da Letra W --- Alfabeto Completo}

\plancheta{Miss\~a{}o Alfabeto 9/10 --- Letra W}

\begin{missao}
\textbf{\faGamepad~Miss\~a{}o Alfabeto 9 de 10:} Consolide a letra W, a letra estrangeira do duplo V!
\end{missao}


\boq[bico]{W}{%
  Letra de outras linguas.\\Soa como U (web) ou V (waffle).\\Boca em bico para o som de U.
}

\pistaarticulatoria{uuu ou vvv (emprestimo)}{Boca em bico (U) ou labio-dental (V)}{Escolha a palavra: web ou waffle}

\begin{consolidacao}
\textbf{Atividade 1 --- Ca\c{c}ada ao W}

\instrucao{Circule o W e sublinhe as palavras com W.}

\begin{center}
\begin{tabular}{cccccc}
\fbox{W} & Z & \fbox{W} & Y & \fbox{W} & V \\
P & \fbox{W} & Z & \fbox{W} & Y & \fbox{W}
\end{tabular}
\end{center}
\begin{center}
\uline{web} \quad \uline{waffle} \quad zebra \quad \uline{windsurf} \quad yoga \quad \uline{kiwi}
\end{center}

\caixapausa{Se precisar, pare aqui. Respire fundo e continue quando quiser.}

\textbf{Atividade 2 --- Escrita do W --- 4 formas}

\instrucao{Escreva o W cinco vezes em cada forma.}

\begin{center}
\textbf{Imprensa mai\'uscula:} \underline{\hspace{1cm}} \underline{\hspace{1cm}} \underline{\hspace{1cm}} \underline{\hspace{1cm}} \underline{\hspace{1cm}}\\[0.5em]
\textbf{Imprensa min\'uscula:} \underline{\hspace{1cm}} \underline{\hspace{1cm}} \underline{\hspace{1cm}} \underline{\hspace{1cm}} \underline{\hspace{1cm}}\\[0.5em]
\textbf{Cursiva mai\'uscula:} \underline{\hspace{1cm}} \underline{\hspace{1cm}} \underline{\hspace{1cm}} \underline{\hspace{1cm}} \underline{\hspace{1cm}}\\[0.5em]
\textbf{Cursiva min\'uscula:} \underline{\hspace{1cm}} \underline{\hspace{1cm}} \underline{\hspace{1cm}} \underline{\hspace{1cm}} \underline{\hspace{1cm}}
\end{center}

\textbf{Atividade 3 --- Jogo do W de Dois Sons}

\instrucao{Ou\c{c}a e marque: W de U ou W de V?}

\begin{center}
\begin{tabular}{ll}
web \quad $\rightarrow$ \underline{\hspace{2cm}} & waffle \quad $\rightarrow$ \underline{\hspace{2cm}} \\
windsurf \quad $\rightarrow$ \underline{\hspace{2cm}} & Wagner \quad $\rightarrow$ \underline{\hspace{2cm}}
\end{tabular}
\end{center}
O W vive em palavras que vieram de outras l\'inguas.

\caixapista{Na d\'uvida, pergunte: \'e{} som de U ou de V?}
\end{consolidacao}

\plancheta{Sons da Letra W --- IPA}

\begin{sonsdaletra}
Os sons da letra W no portugu\^es{} brasileiro (IPA --- Alfabeto Fon\'etico Internacional):

\somipa[bico]{W de U}{w}{som de U r\'apido}{web, windsurf}
\somipa[labiodental]{W de V}{v}{som de V}{waffle, Wagner}
\end{sonsdaletra}

\plancheta{Fam\'ilia Sil\'abica da Letra W}

\begin{familiasilabica}
A fam\'ilia sil\'abica completa da letra W:

\begin{center}
WA \quad WE \quad WI \quad WO \quad WU \\[0.5em] waffle \quad web \quad windsurf \quad --- \quad ---
\end{center}
\instrucao{Leia a fam\'ilia em voz alta, devagar. Depois cubra e tente lembrar.}
\end{familiasilabica}

\plancheta{Atividade de Som e S\'ilabas da Letra W}

\begin{somletra}
\textbf{1. Ca\c{c}a ao som.} Ou\c{c}a a palavra. Se tiver o som de W, marque o quadradinho:
\begin{center}
$\square$~WAFER \quad $\square$~WINDOW \quad $\square$~WAGNER \quad $\square$~BOLA \quad $\square$~P\'E \quad $\square$~M\~AO
\end{center}

\textbf{2. Tem ou n\~ao{} tem?} Rodeie SIM ou N\~ao:
\begin{center}
WAFER --- (\textbf{SIM}) (\textbf{N\~AO}) \quad BOLA --- (\textbf{SIM}) (\textbf{N\~AO}) \\ WINDOW --- (\textbf{SIM}) (\textbf{N\~AO}) \quad P\'E --- (\textbf{SIM}) (\textbf{N\~AO})
\end{center}
\caixapausa{A boca faz bico (som de U).}
\end{somletra}

\begin{somsilaba}
\textbf{3. Som das s\'ilabas.} Ou\c{c}a, bata palmas e conte as s\'ilabas:
\begin{center}
WA.FER $\rightarrow$ \underline{\hspace{1cm}} s\'ilabas \quad WA.GNER $\rightarrow$ \underline{\hspace{1cm}} s\'ilabas
\end{center}

\textbf{4. Qual s\'ilaba tem o som?} Circule a s\'ilaba com o som de W:
\begin{center}
\textbf{WA}.FER \quad \textbf{WA}.GNER
\end{center}
\caixapista{O W aparece em palavras de outros idiomas: WA, WE, WI.}
\end{somsilaba}


\plancheta{Miss\~a{}o Motora --- Letra W}

\begin{motricidade}
\textbf{Movimento:} trace os dois bicos do W.
\begin{center}
\begin{tikzpicture}
  \draw[very thick, dashed, color=primary] (0,1.2) -- (0.35,0) -- (0.7,1.2) -- (1.05,0) -- (1.4,1.2);
  \draw[->, color=accent] (0.1,0.9) -- (0.25,0.5);
  \draw[->, color=accent] (0.5,0.9) -- (0.6,0.5);
\end{tikzpicture}
\end{center}
\caixapista{O W \'e{} um V duplo: dois bicos seguidos.}
\end{motricidade}

\plancheta{Ponte --- Da Letra W para a Pr\'oxima}

\begin{transicao}
Compare o W com o Y:
\begin{center}
\textbf{web} / \textbf{yoga} \qquad \textbf{waffle} / \textbf{yakisoba} \qquad \textbf{kiwi} / \textbf{yin-yang}
\end{center}
\instrucao{W e Y s\~a{}o letras estrangeiras. O W soa como U ou V; o Y soa como I.}
\end{transicao}

\plancheta{Recompensa --- Miss\~a{}o Alfabeto 9/10}

\begin{recompensa}
\estrelas{3}\quad \ofensiva{9}\\[0.3em]\barraprogresso{9}{10}\\[0.5em]

\checklistdominio{reconhecer o W em palavras estrangeiras com som de U ou V}
\end{recompensa}



% ============================================================
% CONSOLIDACAO DA LETRA Y (R201.2) -- Missao Alfabeto 10/10
% ============================================================

\subsection*{Consolida\c{c}\~a{}o da Letra Y --- Alfabeto Completo}

\plancheta{Miss\~a{}o Alfabeto 10/10 --- Letra Y}

\begin{missao}
\textbf{\faGamepad~Miss\~a{}o Alfabeto 10 de 10:} Consolide a letra Y, a \'ultima miss\~a{}o do alfabeto completo!
\end{missao}


\boq[sorriso]{Y}{%
  Letra de outras linguas.\\Soa como I (yoga) ou J (Yan).\\Labios esticados em sorriso.
}

\pistaarticulatoria{iii ou jjj (emprestimo)}{Labios esticados em sorriso}{Escolha a palavra: yoga ou Yan}

\begin{consolidacao}
\textbf{Atividade 1 --- Ca\c{c}ada ao Y}

\instrucao{Circule o Y e sublinhe as palavras com Y.}

\begin{center}
\begin{tabular}{cccccc}
\fbox{Y} & W & \fbox{Y} & Z & \fbox{Y} & X \\
Q & \fbox{Y} & W & \fbox{Y} & Z & \fbox{Y}
\end{tabular}
\end{center}
\begin{center}
\uline{yoga} \quad \uline{yakisoba} \quad waffle \quad \uline{yin-yang} \quad zebra \quad \uline{Yan}
\end{center}

\caixapausa{Se precisar, pare aqui. Respire fundo e continue quando quiser.}

\textbf{Atividade 2 --- Escrita do Y --- 4 formas}

\instrucao{Escreva o Y cinco vezes em cada forma.}

\begin{center}
\textbf{Imprensa mai\'uscula:} \underline{\hspace{1cm}} \underline{\hspace{1cm}} \underline{\hspace{1cm}} \underline{\hspace{1cm}} \underline{\hspace{1cm}}\\[0.5em]
\textbf{Imprensa min\'uscula:} \underline{\hspace{1cm}} \underline{\hspace{1cm}} \underline{\hspace{1cm}} \underline{\hspace{1cm}} \underline{\hspace{1cm}}\\[0.5em]
\textbf{Cursiva mai\'uscula:} \underline{\hspace{1cm}} \underline{\hspace{1cm}} \underline{\hspace{1cm}} \underline{\hspace{1cm}} \underline{\hspace{1cm}}\\[0.5em]
\textbf{Cursiva min\'uscula:} \underline{\hspace{1cm}} \underline{\hspace{1cm}} \underline{\hspace{1cm}} \underline{\hspace{1cm}} \underline{\hspace{1cm}}
\end{center}

\textbf{Atividade 3 --- Revis\~a{}o do Alfabeto Completo}

\instrucao{Escreva as letras que faltam na ordem do alfabeto.}

\begin{center}
A \quad \underline{~~~} \quad C \quad \underline{~~~} \quad E \quad \underline{~~~} \quad G \quad \underline{~~~} \quad I \quad \underline{~~~} \quad K \\[0.5em]
L \quad \underline{~~~} \quad N \quad \underline{~~~} \quad P \quad \underline{~~~} \quad R \quad \underline{~~~} \quad T \quad \underline{~~~} \quad V \\[0.5em]
W \quad \underline{~~~} \quad Y \quad \underline{~~~}
\end{center}
Diga em voz alta o alfabeto inteiro, do A ao Z!

\caixapista{Vogais: A E I O U. O resto \'e{} consoante.}
\end{consolidacao}

\plancheta{Sons da Letra Y --- IPA}

\begin{sonsdaletra}
Os sons da letra Y no portugu\^es{} brasileiro (IPA --- Alfabeto Fon\'etico Internacional):

\somipa[sorriso]{Y de I}{i}{som de I}{yoga, yakisoba}
\somipa[sorriso]{Y de J (come\c{c}o)}{j}{som de J r\'apido}{yin-yang, Yan}
\end{sonsdaletra}

\plancheta{Fam\'ilia Sil\'abica da Letra Y}

\begin{familiasilabica}
A fam\'ilia sil\'abica completa da letra Y:

\begin{center}
YA \quad YE \quad YI \quad YO \quad YU \\[0.5em] yakisoba \quad yin-yang \quad --- \quad yoga \quad ---
\end{center}
\instrucao{Leia a fam\'ilia em voz alta, devagar. Depois cubra e tente lembrar.}
\end{familiasilabica}

\plancheta{Atividade de Som e S\'ilabas da Letra Y}

\begin{somletra}
\textbf{1. Ca\c{c}a ao som.} Ou\c{c}a a palavra. Se tiver o som de Y, marque o quadradinho:
\begin{center}
$\square$~YOGA \quad $\square$~YOUTUBE \quad $\square$~YAGO \quad $\square$~BOLA \quad $\square$~P\'E \quad $\square$~DADO
\end{center}

\textbf{2. Tem ou n\~ao{} tem?} Rodeie SIM ou N\~ao:
\begin{center}
YOGA --- (\textbf{SIM}) (\textbf{N\~AO}) \quad BOLA --- (\textbf{SIM}) (\textbf{N\~AO}) \\ YOUTUBE --- (\textbf{SIM}) (\textbf{N\~AO}) \quad P\'E --- (\textbf{SIM}) (\textbf{N\~AO})
\end{center}
\caixapausa{A boca faz sorriso (som de I).}
\end{somletra}

\begin{somsilaba}
\textbf{3. Som das s\'ilabas.} Ou\c{c}a, bata palmas e conte as s\'ilabas:
\begin{center}
YO.GA $\rightarrow$ \underline{\hspace{1cm}} s\'ilabas \quad YA.GO $\rightarrow$ \underline{\hspace{1cm}} s\'ilabas
\end{center}

\textbf{4. Qual s\'ilaba tem o som?} Circule a s\'ilaba com o som de Y:
\begin{center}
\textbf{YO}.GA \quad \textbf{YA}.GO
\end{center}
\caixapista{O Y aparece em palavras de outros idiomas: YA, YO, YU.}
\end{somsilaba}


\plancheta{Miss\~a{}o Motora --- Letra Y}

\begin{motricidade}
\textbf{Movimento:} trace os dois bra\c{c}os e o poste do Y.
\begin{center}
\begin{tikzpicture}
  \draw[very thick, dashed, color=primary] (0,1.4) -- (0.5,0.7);
  \draw[very thick, dashed, color=primary] (1.0,1.4) -- (0.5,0.7);
  \draw[very thick, dashed, color=primary] (0.5,0.7) -- (0.5,0);
  \draw[->, color=accent] (0.5,1.2) -- (0.5,0.9);
\end{tikzpicture}
\end{center}
\caixapista{O Y \'e{} uma forquilha: dois bra\c{c}os que se juntam em um poste.}
\end{motricidade}

\plancheta{Ponte --- Da Letra Y para a Pr\'oxima}

\begin{transicao}
O Y fecha o alfabeto! Revise a ponte completa:
\begin{center}
\textbf{A E I O U} --- vogais \quad$\cdot$\quad \textbf{B C D F G H J K L M N P Q R S T V W X Y Z} --- consoantes
\end{center}
\instrucao{Digam juntos o alfabeto em voz alta, apontando cada letra.}
\end{transicao}

\plancheta{Recompensa --- Miss\~a{}o Alfabeto 10/10}

\begin{recompensa}
\estrelas{3}\quad \ofensiva{10}\\[0.3em]\barraprogresso{10}{10}\\[0.5em]
\subiunivel{K4 --- Frases e Compreens\~a{}o --- Alfabeto completo!}\\[0.3em]
\checklistdominio{recitar e escrever o alfabeto completo de A a Z}
\end{recompensa}


%% UNIDADE 6 -- SILABAS NASAIS (Nivel 3 -- Silabico-Alfabetico)
%% ============================================================

\input{checkpoints/rastreio-u4}
\chapter{Unidade 6 --- Silabas Nasais}

%% ------------------------------------------------------------
%% LICAO 25: SILABAS COM M NO FINAL (AM, EM, IM, OM, UM)
%% ------------------------------------------------------------
\section{Licao 25: Silabas com M no Final}

% Rotina neuroinclusiva e badges (R201)
\rotinasimples
\nivelKbadge{K3 --- Decodificacao}
\rotabadge{1C --- Consolidacao e Expansao}


\metadata{F-01}{Nivel 3 -- Silabico-Alfabetico}{EF01LP03}{D1}{EF01LP03}{Resolucao CNE/CEB 7/2010, art. 12}

\textbf{Regra:} Quando o ``M'' aparece no FINAL de uma silaba ou palavra, ele nasaliza a vogal anterior.

\plancheta{Folha de Pratica 6.1 --- Silabas Nasais com M}
\begin{exercicio}[Folha de Pratica 6.1 --- Silabas Nasais com M]
\noindent\textbf{Nome:} \underline{\hspace{3.2cm}} \quad \textbf{Data:} \underline{\hspace{1.6cm}} \quad \textbf{Nota:} \underline{\hspace{0.9cm}}/10

\subsection*{PARTE 1: RECONHECIMENTO}

\begin{enumerate}
  \item Leia as silabas nasais em voz alta:
  \begin{center}
  AM \quad EM \quad IM \quad OM \quad UM
  \end{center}

  \item Sublinhe palavras que terminam em ``M'':
  \begin{center}
  \uline{colher} \quad \uline{com} \quad \uline{tem} \quad \uline{bom} \quad \uline{lim}
  \end{center}

  \item Identifique a silaba nasal em cada palavra:
  \begin{center}
  COM $\rightarrow$ silaba nasal: \underline{\hspace{1cm}} \quad BOM $\rightarrow$ silaba nasal: \underline{\hspace{1cm}} \quad TEM $\rightarrow$ silaba nasal: \underline{\hspace{1cm}}
  \end{center}

  \item Separe em silabas e identifique a nasal:
  \begin{center}
  CAMPO $\rightarrow$ \underline{\hspace{1cm}} \underline{\hspace{1cm}} \quad COMIDA $\rightarrow$ \underline{\hspace{1cm}} \underline{\hspace{1cm}} \underline{\hspace{1cm}} \quad LIM\~AO $\rightarrow$ \underline{\hspace{1cm}} \underline{\hspace{1cm}}
  \end{center}

  \item Conte quantas silabas nasais tem cada palavra:
  \begin{center}
  CAMPO $\rightarrow$ \underline{\hspace{1cm}} silaba(s) nasal(is) \quad COMIDA $\rightarrow$ \underline{\hspace{1cm}} \quad LIM\~AO $\rightarrow$ \underline{\hspace{1cm}}
  \end{center}
\end{enumerate}

\subsection*{PARTE 2: PRODUCAO GUIADA}

\begin{enumerate}[resume]
  \item Forme palavras com as silabas nasais:
  \begin{center}
  COM = \underline{\hspace{2cm}} \quad BOM = \underline{\hspace{2cm}} \quad TEM = \underline{\hspace{2cm}} \quad LIM = \underline{\hspace{2cm}}
  \end{center}

  \item Complete com ``AM'', ``EM'', ``IM'', ``OM'' ou ``UM'':
  \begin{center}
  \underline{AM}plo \quad \underline{EM}pfim \quad \underline{IM}vel \quad \underline{OM}briga \quad \underline{UM}ilde
  \end{center}

  \item Escreva 3 palavras que terminam em ``AM'':
  \begin{center}
  1. \underline{\hspace{3cm}} \quad 2. \underline{\hspace{3cm}} \quad 3. \underline{\hspace{3cm}}
  \end{center}

  \item Escreva 3 palavras que terminam em ``EM'':
  \begin{center}
  1. \underline{\hspace{3cm}} \quad 2. \underline{\hspace{3cm}} \quad 3. \underline{\hspace{3cm}}
  \end{center}

  \item Copie as palavras e destaque a silaba nasal:
  \begin{center}
  \textbf{C}\underline{\textbf{OM}} \quad \textbf{B}\underline{\textbf{OM}} \quad \textbf{T}\underline{\textbf{EM}} \quad \textbf{L}\underline{\textbf{IM}} \quad \textbf{S}\underline{\textbf{OM}}
  \end{center}
\end{enumerate}

\subsection*{PARTE 3: INTEGRACAO}

\begin{enumerate}[resume]
  \item Leia as palavras em voz alta:
  \begin{center}
  $\rightarrow$ com  ~  bom  ~  tem  ~  lim  ~  som  ~  campo  ~  lim\~{a}o  ~  Bombom
  \end{center}

  \item Circule a opcao correta:
  \begin{enumerate}[label=(\alph*)]
    \item O M no final nasaliza: ( ) a vogal anterior ( ) a consoante ( ) nada
    \item COM termina em: ( ) silaba nasal ( ) silaba oral ( ) ditongo
    \item A silaba nasal ``EM'' aparece em: ( ) campo ( ) enfim ( ) limao
  \end{enumerate}

  \item Leia e complete:
  \begin{center}
  O c\underline{~~~}po \quad A m\underline{~~~}o \quad O li\underline{~~~}ao \quad B\underline{~~~}m \quad T\underline{~~~}m
  \end{center}

  \item Avaliacao diagnostica:
  \begin{center}
  $\square$ Identifica silabas nasais com M\\
  $\square$ Diferencia silaba oral de nasal\\
  $\square$ Le palavras com silabas nasais\\
  $\square$ Separa palavras em silabas\\
  $\square$ Escreve palavras com silabas nasais
  \end{center}
\end{enumerate}
\end{exercicio}

\begin{dica}
\textbf{Dica para o Professor:} As silabas nasais sao produzidas com a boca FECHADA (pelo M). Peca que os alunos fechem a boca e digam ``AM'', ``EM'', ``IM'', ``OM'', ``UM''. Compare com ``A'', ``E'', ``I'', ``O'', ``U'' (boca aberta) para perceberem a diferenca. Use espelho para observar que os labios se fecham.
\end{dica}

%% ------------------------------------------------------------
%% LICAO 26: SILABAS COM N NO FINAL (AN, EN, IN, ON, UN)
%% ------------------------------------------------------------
\section{Licao 26: Silabas com N no Final}

% Rotina neuroinclusiva e badges (R201)
\rotinasimples
\nivelKbadge{K3 --- Decodificacao}
\rotabadge{1C --- Consolidacao e Expansao}


\metadata{F-02}{Nivel 3 -- Silabico-Alfabetico}{EF01LP03}{D1}{EF01LP03}{Resolucao CNE/CEB 7/2010, art. 12}

\textbf{Regra:} Quando o ``N'' aparece no FINAL de uma silaba ou palavra, ele nasaliza a vogal anterior.

\plancheta{Folha de Pratica 6.2 --- Silabas Nasais com N}
\begin{exercicio}[Folha de Pratica 6.2 --- Silabas Nasais com N]
\noindent\textbf{Nome:} \underline{\hspace{3.2cm}} \quad \textbf{Data:} \underline{\hspace{1.6cm}} \quad \textbf{Nota:} \underline{\hspace{0.9cm}}/10

\subsection*{PARTE 1: RECONHECIMENTO}

\begin{enumerate}
  \item Leia as silabas nasais em voz alta:
  \begin{center}
  AN \quad EN \quad IN \quad ON \quad UN
  \end{center}

  \item Sublinhe palavras que terminam em ``N'':
  \begin{center}
  \uline{pan} \quad \uline{bem} \quad \uline{sim} \quad \uline{bom} \quad \uline{com}
  \end{center}

  \item Identifique a silaba nasal em cada palavra:
  \begin{center}
  PAN $\rightarrow$ silaba nasal: \underline{\hspace{1cm}} \quad BEM $\rightarrow$ silaba nasal: \underline{\hspace{1cm}} \quad SIM $\rightarrow$ silaba nasal: \underline{\hspace{1cm}}
  \end{center}

  \item Separe em silabas e identifique a nasal:
  \begin{center}
  PANO $\rightarrow$ \underline{\hspace{1cm}} \underline{\hspace{1cm}} \quad BEMOL $\rightarrow$ \underline{\hspace{1cm}} \underline{\hspace{1cm}} \quad SIMBOLO $\rightarrow$ \underline{\hspace{1cm}} \underline{\hspace{1cm}} \underline{\hspace{1cm}}
  \end{center}

  \item Conte quantas silabas nasais tem cada palavra:
  \begin{center}
  PANO $\rightarrow$ \underline{\hspace{1cm}} silaba(s) nasal(is) \quad BEMOL $\rightarrow$ \underline{\hspace{1cm}} \quad SIMBOLO $\rightarrow$ \underline{\hspace{1cm}}
  \end{center}
\end{enumerate}

\subsection*{PARTE 2: PRODUCAO GUIADA}

\begin{enumerate}[resume]
  \item Forme palavras com as silabas nasais:
  \begin{center}
  PAN = \underline{\hspace{2cm}} \quad BEM = \underline{\hspace{2cm}} \quad SIM = \underline{\hspace{2cm}} \quad BON = \underline{\hspace{2cm}}
  \end{center}

  \item Complete com ``AN'', ``EN'', ``IN'', ``ON'' ou ``UN'':
  \begin{center}
  \underline{AN}tes \quad \underline{EN}tao \quad \underline{IN}icio \quad \underline{ON}de \quad \underline{UN}ica
  \end{center}

  \item Escreva 3 palavras que terminam com ``N'' (como PAN, BEM, SIM):
  \begin{center}
  1. \underline{\hspace{3cm}} \quad 2. \underline{\hspace{3cm}} \quad 3. \underline{\hspace{3cm}}
  \end{center}

  \item Escreva 3 palavras que terminam com ``M'' (como BOM, COM, SEM):
  \begin{center}
  1. \underline{\hspace{3cm}} \quad 2. \underline{\hspace{3cm}} \quad 3. \underline{\hspace{3cm}}
  \end{center}

  \item Copie as palavras e destaque a silaba nasal:
  \begin{center}
  \textbf{P}\underline{\textbf{AN}} \quad \textbf{B}\underline{\textbf{EM}} \quad \textbf{S}\underline{\textbf{IM}} \quad \textbf{B}\underline{\textbf{ON}} \quad \textbf{C}\underline{\textbf{OM}}
  \end{center}
\end{enumerate}

\subsection*{PARTE 3: INTEGRACAO}

\begin{enumerate}[resume]
  \item Leia as palavras em voz alta:
  \begin{center}
  $\rightarrow$ pan  ~  bem  ~  sim  ~  bon  ~  com  ~  pano  ~  bemol  ~  simbolo
  \end{center}

  \item Circule a opcao correta:
  \begin{enumerate}[label=(\alph*)]
    \item O N no final nasaliza: ( ) a vogal anterior ( ) a consoante ( ) nada
    \item PAN termina em: ( ) silaba nasal ( ) silaba oral ( ) ditongo
    \item A silaba nasal ``AN'' aparece em: ( ) pano ( ) antes ( ) onda
  \end{enumerate}

  \item Leia e complete:
  \begin{center}
  \underline{AN}tes \quad \underline{EN}tao \quad \underline{IN}icio \quad \underline{ON}de \quad \underline{UN}ica
  \end{center}

  \item Leitura em voz alta:
  \begin{center}
  $\rightarrow$ pan  ~  bem  ~  sim  ~  bon  ~  com  ~  pano  ~  bemol  ~  simbolo  ~  antes  ~  entao  ~  inicio
  \end{center}
\end{enumerate}
\end{exercicio}

\begin{dica}
\textbf{Dica para o Professor:} As silabas nasais com N sao produzidas com a lingua tocando o palato (pelo N). Compare com as nasais com M (boca fechada). Peca que os alunos alternem entre ``AM'' (boca fechada) e ``AN'' (boca aberta, lingua no palato) para perceberem a diferenca.
\end{dica}

%% ============================================================
%% UNIDADE 7 -- DITONGOS NASAIS (Nivel 4 -- Silabico-Alfabetico Avancado)
%% ============================================================

\chapter{Unidade 7 --- Ditongos Nasais}

%% ------------------------------------------------------------
%% LICAO 27: DITONGO AO
%% ------------------------------------------------------------
\section{Licao 27: Ditongo Nasal \~AO}

% Rotina neuroinclusiva e badges (R201)
\rotinasimples
\nivelKbadge{K3 --- Decodificacao}
\rotabadge{1C --- Consolidacao e Expansao}


\metadata{G-01}{Nivel 4 -- Silabico-Alfabetico Avancado}{EF01LP03}{D1}{EF01LP03}{Resolucao CNE/CEB 7/2010, art. 12}

\textbf{Regra:} O ditongo ``\~AO'' \'e nasal: aparece no fim de palavras como m\~ao, p\~ao e corac\~ao, e como ``AM'' no fim dos verbos (falam, cantam). A boca abre e o ar sai pelo nariz.

\textbf{Palavra geradora do bloco:} \textbf{M\~AO} (m\~ao) --- a crian\c{c}a olha, toca e fala a pr\'opria m\~ao antes de ler a palavra escrita.

\plancheta{Folha de Pratica 7.1 --- Ditongo \~AO}
\begin{exercicio}[Folha de Pratica 7.1 --- Ditongo AO]
\noindent\textbf{Nome:} \underline{\hspace{3.2cm}} \quad \textbf{Data:} \underline{\hspace{1.6cm}} \quad \textbf{Nota:} \underline{\hspace{0.9cm}}/10

\subsection*{PARTE 1: RECONHECIMENTO}

\begin{enumerate}
  \item Leia as palavras com ditongo ``\~AO'' em voz alta:
  \begin{center}
  $\rightarrow$ m\~ao  ~  n\~ao  ~  p\~ao  ~  t\~ao  ~  s\~ao  ~  corac\~ao  ~  mam\~ao
  \end{center}

  \item Sublinhe palavras que terminam em ``\~AO'':
  \begin{center}
  \uline{m\~ao} \quad gato \quad \uline{n\~ao} \quad sapo \quad \uline{p\~ao} \quad rato \quad \uline{t\~ao}
  \end{center}

  \item Identifique o ditongo ``\~AO'' em cada palavra:
  \begin{center}
  M\~AO $\rightarrow$ ditongo: \underline{\hspace{1cm}} \quad N\~AO $\rightarrow$ ditongo: \underline{\hspace{1cm}} \quad CORAC\~AO $\rightarrow$ ditongo: \underline{\hspace{1cm}}
  \end{center}

  \item Separe em silabas:
  \begin{center}
  M\~AO $\rightarrow$ \underline{\hspace{1cm}} \quad N\~AO $\rightarrow$ \underline{\hspace{1cm}} \quad CORAC\~AO $\rightarrow$ \underline{\hspace{1cm}} \underline{\hspace{1cm}} \underline{\hspace{1cm}}
  \end{center}

  \item Conte as silabas:
  \begin{center}
  M\~AO $\rightarrow$ \underline{\hspace{1cm}} silaba(s) \quad CORAC\~AO $\rightarrow$ \underline{\hspace{1cm}} silabas \quad MAM\~AO $\rightarrow$ \underline{\hspace{1cm}} silabas
  \end{center}
\end{enumerate}

\subsection*{PARTE 2: PRODUCAO GUIADA}

\begin{enumerate}[resume]
  \item Complete com ``\~AO'' ou ``AM'':
  \begin{center}
  m\underline{~~~} \quad n\underline{~~~} \quad p\underline{~~~} \quad s\underline{~~~} \quad cora\underline{~~~}
  \end{center}

  \item Escreva 3 palavras que terminam em ``\~AO'':
  \begin{center}
  1. \underline{\hspace{3cm}} \quad 2. \underline{\hspace{3cm}} \quad 3. \underline{\hspace{3cm}}
  \end{center}

  \item Escreva 3 frases usando palavras com ``\~AO'':
  \begin{enumerate}[label=\alph*)]
    \item \underline{\hspace{8cm}}
    \item \underline{\hspace{8cm}}
    \item \underline{\hspace{8cm}}
  \end{enumerate}

  \item Copie as palavras:
  \begin{center}
  m\~ao \quad n\~ao \quad p\~ao \quad t\~ao \quad s\~ao \quad corac\~ao
  \end{center}

  \item Separe em silabas:
  \begin{center}
  CORAC\~AO = \underline{\hspace{1cm}} \underline{\hspace{1cm}} \underline{\hspace{1cm}} \quad MAM\~AO = \underline{\hspace{1cm}} \underline{\hspace{1cm}} \underline{\hspace{1cm}} \quad LIM\~AO = \underline{\hspace{1cm}} \underline{\hspace{1cm}}
  \end{center}
\end{enumerate}

\subsection*{PARTE 3: INTEGRACAO}

\begin{enumerate}[resume]
  \item Leia as frases em voz alta:
  \begin{center}
  $\rightarrow$ A m\~ao \'e grande.\\
  $\rightarrow$ N\~ao \'e bom.\\
  $\rightarrow$ O p\~ao est\'a quente.
  \end{center}

  \item Circule a opcao correta:
  \begin{enumerate}[label=(\alph*)]
    \item O ditongo ``\~AO'' \'e: ( ) oral ( ) nasal
    \item ``N\~AO'' tem: ( ) 1 silaba ( ) 2 silabas ( ) 3 silabas
    \item O ditongo ``\~AO'' aparece em: ( ) gato ( ) m\~ao ( ) sapo
  \end{enumerate}

  \item Leia e complete:
  \begin{center}
  A m\underline{~~~} \quad N\underline{~~~} \quad O p\underline{~~~} \quad O s\underline{~~~} \quad O cora\underline{~~~}
  \end{center}

  \item Avaliacao diagnostica:
  \begin{center}
  $\square$ Identifica ditongo ``\~AO''\\
  $\square$ Diferencia ditongo ``\~AO'' de silaba com ``O''\\
  $\square$ Le palavras com ditongo ``\~AO''\\
  $\square$ Separa palavras em silabas\\
  $\square$ Escreve frases com ditongo ``\~AO''
  \end{center}
\end{enumerate}
\end{exercicio}

\begin{dica}
\textbf{Dica para o Professor:} O ditongo ``\~AO'' \'e nasal. Compare ``M\~AO'' (nasal) com ``MAU'' (oral) para que percebam a diferen\c{c}a. Use o espelho: em ``M\~AO'', a boca abre e o ar sai pelo nariz; em ``MAU'', a boca abre e o ar sai pela boca.
\end{dica}

%% ------------------------------------------------------------
%% LICAO 28: DITONGO AE
%% ------------------------------------------------------------
\section{Licao 28: Ditongo Nasal \~AE}

% Rotina neuroinclusiva e badges (R201)
\rotinasimples
\nivelKbadge{K3 --- Decodificacao}
\rotabadge{1C --- Consolidacao e Expansao}


\metadata{G-02}{Nivel 4 -- Silabico-Alfabetico Avancado}{EF01LP03}{D1}{EF01LP03}{Resolucao CNE/CEB 7/2010, art. 12}

\textbf{Regra:} O ditongo ``\~AE'' \'e nasal: aparece no fim de palavras como m\~ae, p\~aes e c\~aes. A boca entreabre e o ar sai pelo nariz.

\textbf{Palavra geradora do bloco:} \textbf{M\~AE} (m\~ae) --- a palavra mais afetiva da alfabetiza\c{c}\~ao: quem chama ``mam\~ae''? Aqui o ``\~AE'' fecha a palavra com carinho.

\plancheta{Folha de Pratica 7.2 --- Ditongo \~AE}
\begin{exercicio}[Folha de Pratica 7.2 --- Ditongo AE]
\noindent\textbf{Nome:} \underline{\hspace{3.2cm}} \quad \textbf{Data:} \underline{\hspace{1.6cm}} \quad \textbf{Nota:} \underline{\hspace{0.9cm}}/10

\subsection*{PARTE 1: RECONHECIMENTO}

\begin{enumerate}
  \item Leia as palavras com ditongo ``\~AE'' em voz alta:
  \begin{center}
  $\rightarrow$ m\~ae  ~  p\~aes  ~  c\~aes  ~  m\~aes  ~  mam\~ae
  \end{center}

  \item Sublinhe palavras que terminam em ``\~AE'':
  \begin{center}
  \uline{m\~ae} \quad gato \quad \uline{p\~aes} \quad sapo \quad \uline{c\~aes} \quad rato \quad \uline{mam\~ae}
  \end{center}

  \item Identifique o ditongo ``\~AE'' em cada palavra:
  \begin{center}
  M\~AE $\rightarrow$ ditongo: \underline{\hspace{1cm}} \quad P\~AES $\rightarrow$ ditongo: \underline{\hspace{1cm}} \quad C\~AES $\rightarrow$ ditongo: \underline{\hspace{1cm}}
  \end{center}

  \item Separe em silabas:
  \begin{center}
  M\~AE $\rightarrow$ \underline{\hspace{1cm}} \quad P\~AES $\rightarrow$ \underline{\hspace{1cm}} \quad C\~AES $\rightarrow$ \underline{\hspace{1cm}} \underline{\hspace{1cm}}
  \end{center}

  \item Conte as silabas:
  \begin{center}
  M\~AE $\rightarrow$ \underline{\hspace{1cm}} silaba(s) \quad MAM\~AE $\rightarrow$ \underline{\hspace{1cm}} silabas \quad CAVALO $\rightarrow$ \underline{\hspace{1cm}} silabas
  \end{center}
\end{enumerate}

\subsection*{PARTE 2: PRODUCAO GUIADA}

\begin{enumerate}[resume]
  \item Complete com ``\~AE'' ou ``\~AES'':
  \begin{center}
  m\underline{~~~} \quad p\underline{~~~} \quad c\underline{~~~} \quad mam\underline{~~~}
  \end{center}

  \item Escreva 3 palavras que terminam em ``\~AE'' ou ``\~AES'':
  \begin{center}
  1. \underline{\hspace{3cm}} \quad 2. \underline{\hspace{3cm}} \quad 3. \underline{\hspace{3cm}}
  \end{center}

  \item Escreva 3 frases usando palavras com ``\~AE'':
  \begin{enumerate}[label=\alph*)]
    \item \underline{\hspace{8cm}}
    \item \underline{\hspace{8cm}}
    \item \underline{\hspace{8cm}}
  \end{enumerate}

  \item Copie as palavras:
  \begin{center}
  m\~ae \quad p\~aes \quad c\~aes \quad m\~aes \quad mam\~ae
  \end{center}

  \item Separe em silabas:
  \begin{center}
  M\~AE = \underline{\hspace{1cm}} \quad P\~AES = \underline{\hspace{1cm}} \quad MAM\~AE = \underline{\hspace{1cm}} \underline{\hspace{1cm}}
  \end{center}
\end{enumerate}

\subsection*{PARTE 3: INTEGRACAO}

\begin{enumerate}[resume]
  \item Leia as frases em voz alta:
  \begin{center}
  $\rightarrow$ A m\~ae \'e carinhosa.\\
  $\rightarrow$ N\~ao \'e f\'acil.\\
  $\rightarrow$ Os p\~aes s\~ao legais.
  \end{center}

  \item Circule a opcao correta:
  \begin{enumerate}[label=(\alph*)]
    \item O ditongo ``\~AE'' \'e: ( ) oral ( ) nasal
    \item ``M\~AE'' tem: ( ) 1 silaba ( ) 2 silabas ( ) 3 silabas
    \item O ditongo ``\~AE'' aparece em: ( ) gato ( ) m\~ae ( ) sapo
  \end{enumerate}

  \item Leia e complete:
  \begin{center}
  A m\underline{~~~} \quad Os p\underline{~~~} \quad Os c\underline{~~~} \quad A mam\underline{~~~}
  \end{center}

  \item Leitura em voz alta:
  \begin{center}
  $\rightarrow$ m\~ae  ~  p\~aes  ~  c\~aes  ~  m\~aes  ~  mam\~ae  ~  cavalo  ~  m\'armore
  \end{center}
\end{enumerate}
\end{exercicio}

\begin{dica}
\textbf{Dica para o Professor:} O ditongo ``\~AE'' \'e nasal. Compare ``M\~AE'' (nasal) com ``MAU'' (oral): s\~ao pares m\'inimos --- em ``M\~AE'' o ar sai pelo nariz e em ``MAU'' pela boca. Use o espelho para observar a diferen\c{c}a.
\end{dica}

%% ------------------------------------------------------------
%% LICAO 29: DITONGO OE
%% ------------------------------------------------------------
\section{Licao 29: Ditongo Nasal \~OE}

% Rotina neuroinclusiva e badges (R201)
\rotinasimples
\nivelKbadge{K3 --- Decodificacao}
\rotabadge{1C --- Consolidacao e Expansao}


\metadata{G-03}{Nivel 4 -- Silabico-Alfabetico Avancado}{EF01LP03}{D1}{EF01LP03}{Resolucao CNE/CEB 7/2010, art. 12}

\textbf{Regra:} O ditongo ``\~OE'' \'e nasal: aparece no fim de palavras como p\~oe, lim\~oes e avi\~oes. A boca se arredonda e o ar sai pelo nariz.

\textbf{Palavra geradora do bloco:} \textbf{P\~OE} (p\~oe) --- quem p\~oe a mesa? O ``\~OE'' fecha a a\c{c}\~ao com som anasalado.

\plancheta{Folha de Pratica 7.3 --- Ditongo \~OE}
\begin{exercicio}[Folha de Pratica 7.3 --- Ditongo OE]
\noindent\textbf{Nome:} \underline{\hspace{3.2cm}} \quad \textbf{Data:} \underline{\hspace{1.6cm}} \quad \textbf{Nota:} \underline{\hspace{0.9cm}}/10

\subsection*{PARTE 1: RECONHECIMENTO}

\begin{enumerate}
  \item Leia as palavras com ditongo ``\~OE'' em voz alta:
  \begin{center}
  $\rightarrow$ p\~oe  ~  lim\~oes  ~  avi\~oes  ~  bal\~oes  ~  cora\c{c}\~oes
  \end{center}

  \item Sublinhe palavras que terminam em ``\~OE'':
  \begin{center}
  \uline{p\~oe} \quad gato \quad \uline{lim\~oes} \quad sapo \quad \uline{avi\~oes} \quad rato \quad \uline{bal\~oes}
  \end{center}

  \item Identifique o ditongo ``\~OE'' em cada palavra:
  \begin{center}
  P\~OE $\rightarrow$ ditongo: \underline{\hspace{1cm}} \quad LIM\~OES $\rightarrow$ ditongo: \underline{\hspace{1cm}} \quad AVI\~OES $\rightarrow$ ditongo: \underline{\hspace{1cm}}
  \end{center}

  \item Separe em silabas:
  \begin{center}
  P\~OE $\rightarrow$ \underline{\hspace{1cm}} \quad LIM\~OES $\rightarrow$ \underline{\hspace{1cm}} \quad AVI\~OES $\rightarrow$ \underline{\hspace{1cm}} \underline{\hspace{1cm}} \underline{\hspace{1cm}}
  \end{center}

  \item Conte as silabas:
  \begin{center}
  P\~OE $\rightarrow$ \underline{\hspace{1cm}} silaba(s) \quad LIM\~OES $\rightarrow$ \underline{\hspace{1cm}} silabas \quad BAL\~OES $\rightarrow$ \underline{\hspace{1cm}} silabas
  \end{center}
\end{enumerate}

\subsection*{PARTE 2: PRODUCAO GUIADA}

\begin{enumerate}[resume]
  \item Complete com ``\~OE'' ou ``\~OES'':
  \begin{center}
  p\underline{~~~} \quad lim\underline{~~~} \quad avi\underline{~~~} \quad bal\underline{~~~}
  \end{center}

  \item Escreva 3 palavras que terminam em ``\~OE'' ou ``\~OES'':
  \begin{center}
  1. \underline{\hspace{3cm}} \quad 2. \underline{\hspace{3cm}} \quad 3. \underline{\hspace{3cm}}
  \end{center}

  \item Escreva 3 frases usando palavras com ``\~OE'':
  \begin{enumerate}[label=\alph*)]
    \item \underline{\hspace{8cm}}
    \item \underline{\hspace{8cm}}
    \item \underline{\hspace{8cm}}
  \end{enumerate}

  \item Copie as palavras:
  \begin{center}
  p\~oe \quad lim\~oes \quad avi\~oes \quad bal\~oes \quad cora\c{c}\~oes
  \end{center}

  \item Separe em silabas:
  \begin{center}
  LIM\~OES = \underline{\hspace{1cm}} \underline{\hspace{1cm}} \quad AVI\~OES = \underline{\hspace{1cm}} \underline{\hspace{1cm}} \underline{\hspace{1cm}} \quad BAL\~OES = \underline{\hspace{1cm}} \underline{\hspace{1cm}}
  \end{center}
\end{enumerate}

\subsection*{PARTE 3: INTEGRACAO}

\begin{enumerate}[resume]
  \item Leia as frases em voz alta:
  \begin{center}
  $\rightarrow$ A mam\~ae p\~oe a mesa.\\
  $\rightarrow$ Os avi\~oes voam.\\
  $\rightarrow$ Os lim\~oes s\~ao azedos.
  \end{center}

  \item Circule a opcao correta:
  \begin{enumerate}[label=(\alph*)]
    \item O ditongo ``\~OE'' \'e: ( ) oral ( ) nasal
    \item ``P\~OE'' tem: ( ) 1 silaba ( ) 2 silabas ( ) 3 silabas
    \item O ditongo ``\~OE'' aparece em: ( ) gato ( ) p\~oe ( ) sapo
  \end{enumerate}

  \item Leia e complete:
  \begin{center}
  A po\underline{~~~}ira \quad O co\underline{~~~}tro \quad A men\underline{~~~}ira \quad O p\underline{~~~} \quad A m\underline{~~~}
  \end{center}

  \item Leitura em voz alta:
  \begin{center}
  $\rightarrow$ moe  ~  noe  ~  poe  ~  toe  ~  soe  ~  poeira  ~  coentro  ~  mentira
  \end{center}
\end{enumerate}
\end{exercicio}

\begin{dica}
\textbf{Dica para o Professor:} O ditongo ``\~OE'' \'e nasal. Compare ``P\~OE'' (nasal) com ``POEIRA'' (oral) para que percebam a diferen\c{c}a: em ``P\~OE'' a boca se arredonda e o ar sai pelo nariz.
\end{dica}

%% ============================================================
%% UNIDADE 8 -- LETRAS X, Q, H (ULTIMAS CONSOANTES)
%% ============================================================

\input{checkpoints/rastreio-u5}
\chapter{Unidade 8 --- Letras X, Q, H}

%% ------------------------------------------------------------
%% LICAO 30: LETRA X
%% ------------------------------------------------------------
\section{Licao 30: A Letra X}

% Rotina neuroinclusiva e badges (R201)
\rotinasimples
\nivelKbadge{K2 --- Consciencia Fonologica}
\rotabadge{1B --- Segmentacao e Leitura Funcional}


\letra{X}
% Quatro formas gráficas da letra (R201)
\quatroformasletra{X}{x}

\boq[palatal]{X}{%
  Boca ENTREABERTA, dentes levemente juntos.\\
  Ar sai pelas laterais (sibilante similar ao S).\\
  Lingua proxima ao palato.\\
  Som: ``X-X-X-X-X'' (sibilante, forte).\\[0.5em]
  Palavras exemplo: xicara, xadrez, ixé, exame, xeque.
}

\pistaarticulatoria{Som: ch, ks, s ou z (quatro sons)}{Dentes quase juntos, l'{i}ngua atr'{a}s dos dentes}{M\~{a}o na frente da boca para sentir o ar forte}

\metadata{H-01}{Nivel 4 -- Silabico-Alfabetico Avancado}{EF01LP02}{D1}{EF01LP02}{Resolucao CNE/CEB 7/2010, art. 12}

\plancheta{Folha de Pratica 8.1 --- Letra X}
\begin{exercicio}[Folha de Pratica 8.1 --- Letra X]
\noindent\textbf{Nome:} \underline{\hspace{3.2cm}} \quad \textbf{Data:} \underline{\hspace{1.6cm}} \quad \textbf{Nota:} \underline{\hspace{0.9cm}}/10

\subsection*{PARTE 1: RECONHECIMENTO}

\begin{enumerate}
  \item Formar silabas com X e vogais:
  \begin{center}
  \sil{X}{A} \quad \sil{X}{E} \quad \sil{X}{I} \quad \sil{X}{O} \quad \sil{X}{U}
  \end{center}

  \item Sublinhe palavras que comecam com ``X'':
  \begin{center}
  \uline{xicara} \quad gato \quad \uline{xadrez} \quad sapo \quad \uline{xale} \quad rato
  \end{center}

  \item Identifique o ``X'' em cada palavra:
  \begin{center}
  XICARA $\rightarrow$ posicao \underline{\hspace{1cm}} \quad EXAME $\rightarrow$ posicao \underline{\hspace{1cm}} \quad IXE $\rightarrow$ posicao \underline{\hspace{1cm}}
  \end{center}

  \item Separe em silabas:
  \begin{center}
  XICARA $\rightarrow$ \underline{\hspace{1cm}} \underline{\hspace{1cm}} \underline{\hspace{1cm}} \quad XADREZ $\rightarrow$ \underline{\hspace{1cm}} \underline{\hspace{1cm}} \quad EXAME $\rightarrow$ \underline{\hspace{1cm}} \underline{\hspace{1cm}} \underline{\hspace{1cm}}
  \end{center}

  \item Diga em voz alta tres palavras com ``X'' e anote:
  \begin{center}
  1. \underline{\hspace{4cm}} \quad 2. \underline{\hspace{4cm}} \quad 3. \underline{\hspace{4cm}}
  \end{center}
\end{enumerate}

\subsection*{PARTE 2: PRODUCAO GUIADA}

\begin{enumerate}[resume]
  \item Complete com ``X'':
  \begin{center}
  \underline{X}icara \quad \underline{X}adrez \quad \underline{X}ale \quad e\underline{X}ame \quad i\underline{X}e
  \end{center}

  \item Escreva a letra ``X'' maiuscula cinco vezes:
  \begin{center}
  \underline{\hspace{1.5cm}} \quad \underline{\hspace{1.5cm}} \quad \underline{\hspace{1.5cm}} \quad \underline{\hspace{1.5cm}} \quad \underline{\hspace{1.5cm}}
  \end{center}

  \item Escreva a letra ``x'' minuscula cinco vezes:
  \begin{center}
  \underline{\hspace{1.5cm}} \quad \underline{\hspace{1.5cm}} \quad \underline{\hspace{1.5cm}} \quad \underline{\hspace{1.5cm}} \quad \underline{\hspace{1.5cm}}
  \end{center}

  \item Copie as palavras:
  \begin{center}
  xicara \quad xadrez \quad xale \quad exame \quad ixé
  \end{center}

  \item Escreva 3 palavras com ``X'':
  \begin{center}
  1. \underline{\hspace{4cm}} \quad 2. \underline{\hspace{4cm}} \quad 3. \underline{\hspace{4cm}}
  \end{center}
\end{enumerate}

\subsection*{PARTE 3: INTEGRACAO}

\begin{enumerate}[resume]
  \item Leia as silabas em voz alta:
  \begin{center}
  $\rightarrow$ XA  ~  XE  ~  XI  ~  XO  ~  XU
  \end{center}

  \item Nome de uma coisa com ``X'':
  \begin{center}
  \underline{\hspace{6cm}}
  \end{center}

  \item Circule a opcao correta:
  \begin{enumerate}[label=(\alph*)]
    \item O X tem som: ( ) sibilante ( ) nasal ( ) vibrante
    \item O X aparece em: ( ) mesa ( ) xicara ( ) sapo
    \item O X pode estar no: ( ) inicio ( ) meio ( ) fim da palavra
  \end{enumerate}

  \item Leia e complete:
  \begin{center}
  \sil{X}{A}drez \quad \underline{XI}cara \quad \sil{X}{E}que \quad e\underline{X}ercicio \quad \underline{X}ale
  \end{center}

  \item Leitura em voz alta:
  \begin{center}
  xicara \quad xadrez \quad xale \quad exame \quad ixé \quad exerccio
  \end{center}
\end{enumerate}
\end{exercicio}

\begin{dica}
\textbf{Dica para o Professor:} O X e uma sibilante forte, similar ao ``CH'' do ingles ``church''. Compare com o S (sibilante suave) para que percebam a diferenca. O X pode aparecer no inicio (xicara), meio (exame) ou fim (ifix) da palavra.
\end{dica}

%% ------------------------------------------------------------


% ============================================================
% CONSOLIDACAO DA LETRA X (R201.1) -- Missao 14/16
% ============================================================

\plancheta{Miss\~a{}o 14 --- Letra X}

\begin{missao}
\textbf{\faGamepad~Miss\~a{}o 14 de 16:} Consolide a letra X, a letra camale\~a{}o: quatro sons diferentes em portugu\^es!
\end{missao}

\begin{consolidacao}
\textbf{Atividade 1 --- Ca\c{c}ada ao X}

\instrucao{Circule o X e sublinhe as palavras com X.}

\begin{center}
\begin{tabular}{cccccc}
\fbox{X} & Q & \fbox{X} & H & \fbox{X} & Z \\
W & \fbox{X} & Q & \fbox{X} & H & \fbox{X}
\end{tabular}
\end{center}
\begin{center}
\uline{x\'icara} \quad \uline{xale} \quad queijo \quad \uline{t\'axi} \quad hora \quad \uline{exame}
\end{center}

\caixapausa{Se precisar, pare aqui. Respire fundo e continue quando quiser.}

\textbf{Atividade 2 --- Escrita do X --- 4 formas}

\instrucao{Escreva o X cinco vezes em cada forma.}

\begin{center}
\textbf{Imprensa mai\'uscula:} \underline{\hspace{1cm}} \underline{\hspace{1cm}} \underline{\hspace{1cm}} \underline{\hspace{1cm}} \underline{\hspace{1cm}}\\[0.5em]
\textbf{Imprensa min\'uscula:} \underline{\hspace{1cm}} \underline{\hspace{1cm}} \underline{\hspace{1cm}} \underline{\hspace{1cm}} \underline{\hspace{1cm}}\\[0.5em]
\textbf{Cursiva mai\'uscula:} \underline{\hspace{1cm}} \underline{\hspace{1cm}} \underline{\hspace{1cm}} \underline{\hspace{1cm}} \underline{\hspace{1cm}}\\[0.5em]
\textbf{Cursiva min\'uscula:} \underline{\hspace{1cm}} \underline{\hspace{1cm}} \underline{\hspace{1cm}} \underline{\hspace{1cm}} \underline{\hspace{1cm}}
\end{center}

\textbf{Atividade 3 --- Jogo dos Quatro Sons do X}

\instrucao{Ou\c{c}a e ligue cada palavra ao som correto.}

\begin{center}
\begin{tabular}{ll}
x\'icara & $\rightarrow$ CH \\
t\'axi & $\rightarrow$ KS \\
pr\'oximo & $\rightarrow$ S \\
exame & $\rightarrow$ Z
\end{tabular}
\end{center}
\instrucao{Fale cada palavra em voz alta e sinta o som do X.}

\caixapista{X = camale\~a{}o. O som depende da palavra. Na d\'uvida, ou\c{c}a e repita.}
\end{consolidacao}

\plancheta{Sons da Letra X --- IPA}

\begin{sonsdaletra}
Os sons da letra X no portugu\^es{} brasileiro (IPA --- Alfabeto Fon\'etico Internacional):

\somipa[palatal]{X = CH}{S}{som de ch}{x\'icara, xale, caixa, peixe}
\somipa[velar]{X = KS}{ks}{dois sons juntos}{t\'axi, fixo, anexo}
\somipa[dental]{X = S}{s}{som de s simples}{pr\'oximo, m\'aximo}
\somipa[dental]{X = Z}{z}{som de z no come\c{c}o de palavra}{exame, exemplo, exerc\'icio}
\end{sonsdaletra}

\plancheta{Fam\'ilia Sil\'abica da Letra X}

\begin{familiasilabica}
A fam\'ilia sil\'abica da letra X:

\begin{center}
x\'icara \quad xale \quad xadrez \quad caixa \quad peixe \quad t\'axi \quad exame
\end{center}
\instrucao{Leia a fam\'ilia em voz alta, devagar. Depois cubra e tente lembrar.}
\end{familiasilabica}

\plancheta{Atividade de Som e S\'ilabas da Letra X}

\begin{somletra}
\textbf{1. Ca\c{c}a ao som.} Ou\c{c}a a palavra. Se tiver o som de X, marque o quadradinho:
\begin{center}
$\square$~X\'ICARA \quad $\square$~CAIXA \quad $\square$~PEIXE \quad $\square$~BOLA \quad $\square$~SOL \quad $\square$~M\~AO
\end{center}

\textbf{2. Tem ou n\~ao{} tem?} Rodeie SIM ou N\~ao:
\begin{center}
X\'ICARA --- (\textbf{SIM}) (\textbf{N\~AO}) \quad BOLA --- (\textbf{SIM}) (\textbf{N\~AO}) \\ CAIXA --- (\textbf{SIM}) (\textbf{N\~AO}) \quad SOL --- (\textbf{SIM}) (\textbf{N\~AO})
\end{center}
\caixapausa{A l\'ingua sobe no c\'eu da boca (som de CH).}
\end{somletra}

\begin{somsilaba}
\textbf{3. Som das s\'ilabas.} Ou\c{c}a, bata palmas e conte as s\'ilabas:
\begin{center}
X\'I.CA.RA $\rightarrow$ \underline{\hspace{1cm}} s\'ilabas \quad CAI.XA $\rightarrow$ \underline{\hspace{1cm}} s\'ilabas
\end{center}

\textbf{4. Qual s\'ilaba tem o som?} Circule a s\'ilaba com o som de X:
\begin{center}
\textbf{X\'I}.CA.RA \quad CAI.\textbf{XA}
\end{center}
\caixapista{O X com som de CH: X\'I, XA, XE (x\'icara, caixa, peixe).}
\end{somsilaba}


\plancheta{Miss\~a{}o Motora --- Letra X}

\begin{motricidade}
\textbf{Movimento:} trace as duas diagonais do X.
\begin{center}
\begin{tikzpicture}
  \draw[very thick, dashed, color=primary] (0,0.8) -- (1.8,-0.8);
  \draw[very thick, dashed, color=secondary] (0,-0.8) -- (1.8,0.8);
  \draw[->, color=accent] (0.4,0.5) -- (0.9,0.15);
  \draw[->, color=accent] (0.4,-0.5) -- (0.9,-0.15);
\end{tikzpicture}
\end{center}
\caixapista{O X \'e{} uma cruz em diagonal. Duas linhas que se cruzam.}
\end{motricidade}

\plancheta{Ponte --- Da Letra X para a Pr\'oxima}

\begin{transicao}
Compare o X com o Q:
\begin{center}
\textbf{xale} / \textbf{queijo} \qquad \textbf{caixa} / \textbf{quero} \qquad \textbf{peixe} / \textbf{quinta}
\end{center}
\instrucao{O X tem quatro sons. O Q tem sempre som de K, com o U mudo ou sonoro.}
\end{transicao}

\plancheta{Recompensa --- Miss\~a{}o 14}

\begin{recompensa}
\estrelas{3}\quad \ofensiva{14}\\[0.3em]\barraprogresso{14}{16}\\[0.5em]

\checklistdominio{identificar os quatro sons do X em palavras conhecidas}
\end{recompensa}


%% LICAO 31: LETRA Q (QUE, QUI)
%% ------------------------------------------------------------
\section{Licao 31: A Letra Q (QUE, QUI)}

% Rotina neuroinclusiva e badges (R201)
\rotinasimples
\nivelKbadge{K2 --- Consciencia Fonologica}
\rotabadge{1B --- Segmentacao e Leitura Funcional}


\letra{Q}
% Quatro formas gráficas da letra (R201)
\quatroformasletra{Q}{q}

\boq[velar]{Q}{%
  Boca ENTREABERTA, lingua toca o palato mole.\\
  Ar sai em explosao surda (oclusiva velar surda).\\
  Sempre aparece com U antes de E ou I.\\
  Som: ``QUE'' = ``KE'', ``QUI'' = ``KI''.\\[0.5em]
  Palavras exemplo: queijo, quero, quinta, quem, aqui.
}

\pistaarticulatoria{Som: k (com u mudo depois)}{Boca como no K; U n\~{a}o tem som}{M\~{a}o na garganta: n\~{a}o vibra (som surdo)}

\metadata{H-02}{Nivel 4 -- Silabico-Alfabetico Avancado}{EF01LP02}{D1}{EF01LP02}{Resolucao CNE/CEB 7/2010, art. 12}

\plancheta{Folha de Pratica 8.2 --- Letra Q (QUE, QUI)}
\begin{exercicio}[Folha de Pratica 8.2 --- Letra Q (QUE, QUI)]
\noindent\textbf{Nome:} \underline{\hspace{3.2cm}} \quad \textbf{Data:} \underline{\hspace{1.6cm}} \quad \textbf{Nota:} \underline{\hspace{0.9cm}}/10

\subsection*{PARTE 1: RECONHECIMENTO}

\begin{enumerate}
  \item Formar silabas QUE e QUI:
  \begin{center}
  \sil{QU}{E} \quad \sil{QU}{I}
  \end{center}

  \item Sublinhe palavras que comecam com ``QU'':
  \begin{center}
  \uline{queijo} \quad gato \quad \uline{quero} \quad sapo \quad \uline{quinta} \quad rato \quad \uline{quem}
  \end{center}

  \item Identifique QUE ou QUI em cada palavra:
  \begin{center}
  QUEIJO $\rightarrow$ \underline{\hspace{1cm}} \quad QUERO $\rightarrow$ \underline{\hspace{1cm}} \quad QUINTA $\rightarrow$ \underline{\hspace{1cm}} \quad AQUI $\rightarrow$ \underline{\hspace{1cm}}
  \end{center}

  \item Separe em silabas:
  \begin{center}
  QUEIJO $\rightarrow$ \underline{\hspace{1cm}} \underline{\hspace{1cm}} \quad QUERO $\rightarrow$ \underline{\hspace{1cm}} \underline{\hspace{1cm}} \quad QUINTA $\rightarrow$ \underline{\hspace{1cm}} \underline{\hspace{1cm}}
  \end{center}

  \item Diga em voz alta tres palavras com ``QU'' e anote:
  \begin{center}
  1. \underline{\hspace{4cm}} \quad 2. \underline{\hspace{4cm}} \quad 3. \underline{\hspace{4cm}}
  \end{center}
\end{enumerate}

\subsection*{PARTE 2: PRODUCAO GUIADA}

\begin{enumerate}[resume]
  \item Complete com QUE ou QUI:
  \begin{center}
  \underline{QUE}ijo \quad \underline{QUER}o \quad \underline{QUIN}ta \quad a\underline{QUI} \quad \underline{QUE}m
  \end{center}

  \item Escreva a letra ``Q'' maiuscula cinco vezes:
  \begin{center}
  \underline{\hspace{1.5cm}} \quad \underline{\hspace{1.5cm}} \quad \underline{\hspace{1.5cm}} \quad \underline{\hspace{1.5cm}} \quad \underline{\hspace{1.5cm}}
  \end{center}

  \item Escreva a letra ``q'' minuscula cinco vezes:
  \begin{center}
  \underline{\hspace{1.5cm}} \quad \underline{\hspace{1.5cm}} \quad \underline{\hspace{1.5cm}} \quad \underline{\hspace{1.5cm}} \quad \underline{\hspace{1.5cm}}
  \end{center}

  \item Copie as palavras:
  \begin{center}
  queijo \quad quero \quad quinta \quad quem \quad aqui
  \end{center}

  \item Escreva 3 palavras com ``QU'':
  \begin{center}
  1. \underline{\hspace{4cm}} \quad 2. \underline{\hspace{4cm}} \quad 3. \underline{\hspace{4cm}}
  \end{center}
\end{enumerate}

\subsection*{PARTE 3: INTEGRACAO}

\begin{enumerate}[resume]
  \item Leia as silabas em voz alta:
  \begin{center}
  $\rightarrow$ QUE \quad QUI
  \end{center}

  \item Regra importante:
  \begin{center}
  \textbf{Q} + \textbf{UE} = ``KE'' (queijo, quero)\\
  \textbf{Q} + \textbf{UI} = ``KI'' (quinta, quiabo)\\
  O ``U'' do ``QU'' e MUDO (nao tem som)!
  \end{center}

  \item Circule a opcao correta:
  \begin{enumerate}[label=(\alph*)]
    \item No ``QUE'', o U: ( ) tem som ( ) e mudo ( ) nasaliza
    \item ``QU'' aparece antes de: ( ) A e O ( ) E e I ( ) sozinho
    \item ``AQUI'' tem: ( ) 2 silabas ( ) 3 silabas ( ) 1 silaba
  \end{enumerate}

  \item Leia e complete:
  \begin{center}
  \underline{QUE}ijo \quad \underline{QUI}abo \quad a\underline{QUI} \quad \underline{QUE}m \quad \underline{QUIN}ta
  \end{center}

  \item Leitura em voz alta:
  \begin{center}
  queijo \quad quero \quad quinta \quad quem \quad aqui \quad quiabo \quad quota
  \end{center}
\end{enumerate}
\end{exercicio}

\begin{dica}
\textbf{Dica para o Professor:} O Q sempre aparece com U antes de E ou I. O U do ``QU'' e mudo! Peca que os alunos digam ``QUE'' e percebam que o ``U'' nao e pronunciado. Compare ``QUE'' (K-E) com ``CUE'' (n existe em portugues) para reforar a regra.
\end{dica}

%% ------------------------------------------------------------


% ============================================================
% CONSOLIDACAO DA LETRA Q (R201.1) -- Missao 15/16
% ============================================================

\plancheta{Miss\~a{}o 15 --- Letra Q}

\begin{missao}
\textbf{\faGamepad~Miss\~a{}o 15 de 16:} Consolide a letra Q. O Q anda sempre de m\~a{}o dada com o U!
\end{missao}

\begin{consolidacao}
\textbf{Atividade 1 --- Ca\c{c}ada ao Q}

\instrucao{Circule o Q e sublinhe as palavras com QU.}

\begin{center}
\begin{tabular}{cccccc}
\fbox{Q} & H & \fbox{Q} & X & \fbox{Q} & Z \\
W & \fbox{Q} & H & \fbox{Q} & X & \fbox{Q}
\end{tabular}
\end{center}
\begin{center}
\uline{queijo} \quad \uline{quero} \quad hora \quad \uline{quinta} \quad xale \quad \uline{aqui}
\end{center}

\caixapausa{Se precisar, pare aqui. Respire fundo e continue quando quiser.}

\textbf{Atividade 2 --- Escrita do Q --- 4 formas}

\instrucao{Escreva o Q cinco vezes em cada forma.}

\begin{center}
\textbf{Imprensa mai\'uscula:} \underline{\hspace{1cm}} \underline{\hspace{1cm}} \underline{\hspace{1cm}} \underline{\hspace{1cm}} \underline{\hspace{1cm}}\\[0.5em]
\textbf{Imprensa min\'uscula:} \underline{\hspace{1cm}} \underline{\hspace{1cm}} \underline{\hspace{1cm}} \underline{\hspace{1cm}} \underline{\hspace{1cm}}\\[0.5em]
\textbf{Cursiva mai\'uscula:} \underline{\hspace{1cm}} \underline{\hspace{1cm}} \underline{\hspace{1cm}} \underline{\hspace{1cm}} \underline{\hspace{1cm}}\\[0.5em]
\textbf{Cursiva min\'uscula:} \underline{\hspace{1cm}} \underline{\hspace{1cm}} \underline{\hspace{1cm}} \underline{\hspace{1cm}} \underline{\hspace{1cm}}
\end{center}

\textbf{Atividade 3 --- Jogo do U Mudo}

\instrucao{Ou\c{c}a e marque: o U soa ou \'e{} mudo?}

\begin{center}
\begin{tabular}{ll}
queijo \quad $\rightarrow$ \underline{\hspace{2cm}} & qual \quad $\rightarrow$ \underline{\hspace{2cm}} \\
quero \quad $\rightarrow$ \underline{\hspace{2cm}} & quatro \quad $\rightarrow$ \underline{\hspace{2cm}} \\
aqui \quad $\rightarrow$ \underline{\hspace{2cm}} & quando \quad $\rightarrow$ \underline{\hspace{2cm}}
\end{tabular}
\end{center}
Regra: QUE, QUI = U mudo. QUA, QUO = U sonoro.

\caixapista{QUE = KE. QUI = KI. QUA = KUA. Sinta o U sonoro em qual!}
\end{consolidacao}

\plancheta{Sons da Letra Q --- IPA}

\begin{sonsdaletra}
Os sons da letra Q no portugu\^es{} brasileiro (IPA --- Alfabeto Fon\'etico Internacional):

\somipa[velar]{QU com U mudo}{k}{o U n\~a{}o soa: QUE = KE, QUI = KI}{queijo, quero, aqui, quibe}
\somipa[velar]{QU com U sonoro}{kw}{o U soa: QUA = KUA}{qual, quando, quatro, quase}
\end{sonsdaletra}

\plancheta{Fam\'ilia Sil\'abica da Letra Q}

\begin{familiasilabica}
A fam\'ilia sil\'abica da letra Q:

\begin{center}
queijo \quad quero \quad quinta \quad quem \quad aqui \quad quatro \quad quando
\end{center}
\instrucao{Leia a fam\'ilia em voz alta, devagar. Depois cubra e tente lembrar.}
\end{familiasilabica}

\plancheta{Atividade de Som e S\'ilabas da Letra Q}

\begin{somletra}
\textbf{1. Ca\c{c}a ao som.} Ou\c{c}a a palavra. Se tiver o som de Q, marque o quadradinho:
\begin{center}
$\square$~QUEIJO \quad $\square$~QUADRO \quad $\square$~TOQUE \quad $\square$~BOLA \quad $\square$~DADO \quad $\square$~P\'E
\end{center}

\textbf{2. Tem ou n\~ao{} tem?} Rodeie SIM ou N\~ao:
\begin{center}
QUEIJO --- (\textbf{SIM}) (\textbf{N\~AO}) \quad BOLA --- (\textbf{SIM}) (\textbf{N\~AO}) \\ QUADRO --- (\textbf{SIM}) (\textbf{N\~AO}) \quad DADO --- (\textbf{SIM}) (\textbf{N\~AO})
\end{center}
\caixapausa{O som sai do fundo da boca, junto com U.}
\end{somletra}

\begin{somsilaba}
\textbf{3. Som das s\'ilabas.} Ou\c{c}a, bata palmas e conte as s\'ilabas:
\begin{center}
QUEI.JO $\rightarrow$ \underline{\hspace{1cm}} s\'ilabas \quad QUA.DRO $\rightarrow$ \underline{\hspace{1cm}} s\'ilabas
\end{center}

\textbf{4. Qual s\'ilaba tem o som?} Circule a s\'ilaba com o som de Q:
\begin{center}
\textbf{QUEI}.JO \quad \textbf{QUA}.DRO
\end{center}
\caixapista{O Q sempre vem com U: QUE, QUI, QUA.}
\end{somsilaba}


\plancheta{Miss\~a{}o Motora --- Letra Q}

\begin{motricidade}
\textbf{Movimento:} trace o c\'irculo e o rabicho do Q.
\begin{center}
\begin{tikzpicture}
  \draw[very thick, dashed, color=primary] (0.5,0) circle (0.6);
  \draw[very thick, dashed, color=secondary] (0.9,-0.4) -- (1.4,-0.9);
  \fill[primary] (1.1,0) circle (0.1);
  \draw[->, color=accent] (0.8,0.5) -- (1.0,0.4);
\end{tikzpicture}
\end{center}
\caixapista{O Q \'e{} um O com rabicho. O rabicho desce para a direita.}
\end{motricidade}

\plancheta{Ponte --- Da Letra Q para a Pr\'oxima}

\begin{transicao}
Compare o Q com o H, a \'ultima letra:
\begin{center}
\textbf{quero} / \textbf{hoje} \qquad \textbf{queijo} / \textbf{chave} \qquad \textbf{quinta} / \textbf{chuva}
\end{center}
\instrucao{O Q tem sempre som de K. O H \'e{} mudo, mas forma os times CH, LH e NH.}
\end{transicao}

\plancheta{Recompensa --- Miss\~a{}o 15}

\begin{recompensa}
\estrelas{3}\quad \ofensiva{15}\\[0.3em]\barraprogresso{15}{16}\\[0.5em]

\checklistdominio{ler QUE e QUI com U mudo e QUA com U sonoro}
\end{recompensa}


%% LICAO 32: LETRA H
%% ------------------------------------------------------------
\section{Licao 32: A Letra H (Muda e nos Digrafos)}

% Rotina neuroinclusiva e badges (R201)
\rotinasimples
\nivelKbadge{K2 --- Consciencia Fonologica}
\rotabadge{1B --- Segmentacao e Leitura Funcional}


\letra{H}
% Quatro formas gráficas da letra (R201)
\quatroformasletra{H}{h}

\boq[muda]{H}{%
  O H N\~AO TEM SOM pr\'oprio em portugu\^es: \'e sempre mudo.\\
  HOJE se le ``OJE''; HORA se le ``ORA''.\\
  O H aparece nos times CH, LH e NH.\\
  Som: a boca fica neutra e soa a vogal a seguir.\\[0.5em]
  Palavras exemplo: hoje, heroi, hotel, hortela.
}

\pistaarticulatoria{Sem som proprio (mudo)}{Boca neutra, sem sopro; soa a vogal a seguir}{Mostre que HOJE se le OJE}

\metadata{H-03}{Nivel 4 -- Silabico-Alfabetico Avancado}{EF01LP02}{D1}{EF01LP02}{Resolucao CNE/CEB 7/2010, art. 12}

\plancheta{Folha de Pratica 8.3 --- Letra H}
\begin{exercicio}[Folha de Pratica 8.3 --- Letra H]
\noindent\textbf{Nome:} \underline{\hspace{3.2cm}} \quad \textbf{Data:} \underline{\hspace{1.6cm}} \quad \textbf{Nota:} \underline{\hspace{0.9cm}}/10

\subsection*{PARTE 1: RECONHECIMENTO}

\begin{enumerate}
  \item Ler familias com H (o H nao tem som: leia so a vogal):
  \begin{center}
  HA \quad HE \quad HI \quad HO \quad HU
  \end{center}

  \item Sublinhe palavras que comecam com ``H'':
  \begin{center}
  \uline{hoje} \quad gato \quad \uline{heroi} \quad sapo \quad \uline{hotel} \quad rato \quad \uline{hortela}
  \end{center}

  \item Identifique o ``H'' em cada palavra:
  \begin{center}
  HOJE $\rightarrow$ posicao \underline{\hspace{1cm}} \quad HEROI $\rightarrow$ posicao \underline{\hspace{1cm}} \quad HOTEL $\rightarrow$ posicao \underline{\hspace{1cm}}
  \end{center}

  \item Separe em silabas:
  \begin{center}
  HOJE $\rightarrow$ \underline{\hspace{1cm}} \underline{\hspace{1cm}} \quad HEROICO $\rightarrow$ \underline{\hspace{1cm}} \underline{\hspace{1cm}} \underline{\hspace{1cm}} \underline{\hspace{1cm}} \quad HORTELA $\rightarrow$ \underline{\hspace{1cm}} \underline{\hspace{1cm}} \underline{\hspace{1cm}}
  \end{center}

  \item Diga em voz alta tres palavras com ``H'' e anote:
  \begin{center}
  1. \underline{\hspace{4cm}} \quad 2. \underline{\hspace{4cm}} \quad 3. \underline{\hspace{4cm}}
  \end{center}
\end{enumerate}

\subsection*{PARTE 2: PRODUCAO GUIADA}

\begin{enumerate}[resume]
  \item Complete com ``H'':
  \begin{center}
  \underline{H}oje \quad \underline{H}eroi \quad \underline{H}otel \quad \underline{H}eroico \quad \underline{H}ortela
  \end{center}

  \item Escreva a letra ``H'' maiuscula cinco vezes:
  \begin{center}
  \underline{\hspace{1.5cm}} \quad \underline{\hspace{1.5cm}} \quad \underline{\hspace{1.5cm}} \quad \underline{\hspace{1.5cm}} \quad \underline{\hspace{1.5cm}}
  \end{center}

  \item Escreva a letra ``h'' minuscula cinco vezes:
  \begin{center}
  \underline{\hspace{1.5cm}} \quad \underline{\hspace{1.5cm}} \quad \underline{\hspace{1.5cm}} \quad \underline{\hspace{1.5cm}} \quad \underline{\hspace{1.5cm}}
  \end{center}

  \item Copie as palavras:
  \begin{center}
  hoje \quad heroi \quad hotel \quad heroico \quad hortela
  \end{center}

  \item Escreva 3 palavras com ``H'':
  \begin{center}
  1. \underline{\hspace{4cm}} \quad 2. \underline{\hspace{4cm}} \quad 3. \underline{\hspace{4cm}}
  \end{center}
\end{enumerate}

\subsection*{PARTE 3: INTEGRACAO}

\begin{enumerate}[resume]
  \item Leia as silabas em voz alta:
  \begin{center}
  $\rightarrow$ HA  ~  HE  ~  HI  ~  HO  ~  HU
  \end{center}

  \item Regra importante:
  \begin{center}
  \textbf{H} no inicio = MUDO (hoje se le OJE)\\
  \textbf{H} no meio = MUDO (bahia se le baia)\\
  Digrafos com H: CH (chuva), LH (milho), NH (ninho)
  \end{center}

  \item Circule a opcao correta:
  \begin{enumerate}[label=(\alph*)]
    \item O H tem som proprio? ( ) sim ( ) nao
    \item O H aparece nos times: ( ) CH, LH e NH ( ) XH, WH e ZH
    \item O H aparece em: ( ) mesa ( ) hoje ( ) sapo
  \end{enumerate}

  \item Leia e complete:
  \begin{center}
  \underline{H}oje \quad \underline{H}eroi \quad \underline{H}otel \quad \underline{H}eroico \quad \underline{H}ortela
  \end{center}

  \item Leitura em voz alta:
  \begin{center}
  hoje \quad heroi \quad hotel \quad heroico \quad hortela \quad Hospital
  \end{center}
\end{enumerate}
\end{exercicio}

\begin{dica}
\textbf{Dica para o Professor:} No portugues brasileiro o H nao tem som proprio: HOJE se pronuncia ``OJE'' e HORA se pronuncia ``ORA''. O H so aparece ``soando'' dentro dos digrafos CH (chuva), LH (milho) e NH (ninho). Peca que os alunos comparem HOJE e OJE e percebam que o H sozinho e mudo: o som vem da vogal seguinte ou do digrafo.
\end{dica}

%% ============================================================


% ============================================================
% CONSOLIDACAO DA LETRA H (R201.1) -- Missao 16/16
% ============================================================

\plancheta{Miss\~a{}o 16 --- Letra H}

\begin{missao}
\textbf{\faGamepad~Miss\~a{}o 16 de 16:} Consolide a letra H, a \'ultima do alfabeto! O H \'e{} mudo, mas \'e{} o chefe dos times CH, LH e NH.
\end{missao}

\begin{consolidacao}
\textbf{Atividade 1 --- Ca\c{c}ada ao H}

\instrucao{Circule o H e sublinhe as palavras com H (sozinho ou em time).}

\begin{center}
\begin{tabular}{cccccc}
\fbox{H} & Q & \fbox{H} & X & \fbox{H} & Z \\
W & \fbox{H} & Q & \fbox{H} & X & \fbox{H}
\end{tabular}
\end{center}
\begin{center}
\uline{hoje} \quad \uline{chave} \quad queijo \quad \uline{milho} \quad xale \quad \uline{unha}
\end{center}

\caixapausa{Se precisar, pare aqui. Respire fundo e continue quando quiser.}

\textbf{Atividade 2 --- Escrita do H --- 4 formas}

\instrucao{Escreva o H cinco vezes em cada forma.}

\begin{center}
\textbf{Imprensa mai\'uscula:} \underline{\hspace{1cm}} \underline{\hspace{1cm}} \underline{\hspace{1cm}} \underline{\hspace{1cm}} \underline{\hspace{1cm}}\\[0.5em]
\textbf{Imprensa min\'uscula:} \underline{\hspace{1cm}} \underline{\hspace{1cm}} \underline{\hspace{1cm}} \underline{\hspace{1cm}} \underline{\hspace{1cm}}\\[0.5em]
\textbf{Cursiva mai\'uscula:} \underline{\hspace{1cm}} \underline{\hspace{1cm}} \underline{\hspace{1cm}} \underline{\hspace{1cm}} \underline{\hspace{1cm}}\\[0.5em]
\textbf{Cursiva min\'uscula:} \underline{\hspace{1cm}} \underline{\hspace{1cm}} \underline{\hspace{1cm}} \underline{\hspace{1cm}} \underline{\hspace{1cm}}
\end{center}

\textbf{Atividade 3 --- Jogo dos Times do H}

\instrucao{Complete com CH, LH ou NH.}

\begin{center}
\underline{~~~}ave \quad mi\underline{~~~}o \quad u\underline{~~~}a \quad \underline{~~~}uva \quad ni\underline{~~~}o \quad so\underline{~~~}o
\end{center}
Diga em voz alta: \textbf{chave, milho, unha, chuva, ninho, sonho}

\caixapista{CH = ch. LH = lh molhado. NH = nh molhado. O H muda o som da letra!}
\end{consolidacao}

\plancheta{Sons da Letra H --- IPA}

\begin{sonsdaletra}
Os sons da letra H no portugu\^es{} brasileiro (IPA --- Alfabeto Fon\'etico Internacional):

\somipa[muda]{H mudo}{}{n\~a{}o tem som pr\'oprio}{hoje, hora, hotel, her\'oi}
\somipa[palatal]{CH}{S}{som de ch}{chave, chuva, cacho}
\somipa[palatal]{LH}{\textturny}{som molhado de lh}{milho, alho, filho}
\somipa[palatal]{NH}{\textltailn}{som molhado de nh}{unha, ninho, sonho}
\end{sonsdaletra}

\plancheta{Fam\'ilia Sil\'abica da Letra H}

\begin{familiasilabica}
A fam\'ilia sil\'abica da letra H:

\begin{center}
hoje \quad hora \quad chave \quad chuva \quad milho \quad alho \quad unha \quad ninho
\end{center}
\instrucao{Leia a fam\'ilia em voz alta, devagar. Depois cubra e tente lembrar.}
\end{familiasilabica}

\plancheta{Atividade de Som e S\'ilabas da Letra H}

\begin{somletra}
\textbf{1. O H \'e{} mudo.} Marque as palavras que t\^e{}m H, mesmo sem som:
\begin{center}
$\square$~HOJE \quad $\square$~BOLA \quad $\square$~HORA \quad $\square$~P\'E \quad $\square$~HOTEL \quad $\square$~DADO
\end{center}

\textbf{2. Tem ou n\~ao{} tem?} Rodeie SIM se a palavra tem H, N\~ao{} se n\~ao{} tem:
\begin{center}
HOJE --- (\textbf{SIM}) (\textbf{N\~AO}) \quad BOLA --- (\textbf{SIM}) (\textbf{N\~AO})\\
HORA --- (\textbf{SIM}) (\textbf{N\~AO}) \quad DADO --- (\textbf{SIM}) (\textbf{N\~AO})
\end{center}
\caixapausa{A boca abre, mas o H n\~ao{} tem som.}
\end{somletra}

\begin{somsilaba}
\textbf{3. Compare:} HOJE l\^e-se OJE; HORA l\^e-se ORA; HOTEL l\^e-se OTEL.

\textbf{4. O segredo:} o H n\~ao{} tem som pr\'oprio. Ele acompanha outras letras: CH, LH, NH.
\caixapista{NH tem som de nariz: NINHO. LH: MILHO. CH tem som de x: CHUVA.}
\end{somsilaba}


\plancheta{Miss\~a{}o Motora --- Letra H}

\begin{motricidade}
\textbf{Movimento:} trace os dois trilhos e a ponte do H.
\begin{center}
\begin{tikzpicture}
  \draw[very thick, dashed, color=primary] (0,0.8) -- (0,-0.8);
  \draw[very thick, dashed, color=primary] (2,0.8) -- (2,-0.8);
  \draw[very thick, dashed, color=secondary] (0,0) -- (2,0);
  \draw[->, color=accent] (0.15,0.5) -- (0.15,-0.1);
  \draw[->, color=accent] (0.5,0.15) -- (1.2,0.15);
\end{tikzpicture}
\end{center}
\caixapista{O H \'e{} uma ponte entre dois trilhos.}
\end{motricidade}

\plancheta{Ponte --- Da Letra H para a Pr\'oxima}

\begin{transicao}
O H fecha o alfabeto! Compare com os sons nasais que voc\^e j\'a{} aprendeu:
\begin{center}
\textbf{chave} / \textbf{m\~a{}e} \qquad \textbf{unha} / \textbf{lim\~a{}o} \qquad \textbf{ninho} / \textbf{p\~a{}o}
\end{center}
\instrucao{Fale cada par em voz alta. O som nasal vem pelo nariz, o som do H vem pela boca.}
\end{transicao}

\plancheta{Recompensa --- Miss\~a{}o 16}

\begin{recompensa}
\estrelas{3}\quad \ofensiva{16}\\[0.3em]\barraprogresso{16}{16}\\[0.5em]
\subiunivel{K4 --- Frases e Compreens\~a{}o}\\[0.3em]
\checklistdominio{usar os times CH, LH e NH e saber que o H sozinho \'e{} mudo}
\end{recompensa}


%% UNIDADE 9 -- SILEBAS COM ÃO PARA TODAS AS CONSOANTES
%% ============================================================

\chapter{Unidade 9 --- Sílabas com ÃO (BÃO, CÃO, DÃO, FÃO, GÃO, JÃO, MÃO, NÃO, PÃO, VÃO, ZÃO)}

\metadata{I-01}{Nivel 4 -- Silabico-Alfabetico Avancado}{EF01LP03}{D1}{EF01LP03}{Resolucao CNE/CEB 7/2010, art. 12}

\textbf{Regra:} A silaba ``ÃO'' (ditongo nasal) pode ser precedida por QUALQUER consoante. Basta juntar a consoante + ``ÃO''.

\plancheta{Folha de Pratica 9.1 --- Silabas com ÃO}
\begin{exercicio}[Folha de Pratica 9.1 --- Silabas com ÃO]
\noindent\textbf{Nome:} \underline{\hspace{3.2cm}} \quad \textbf{Data:} \underline{\hspace{1.6cm}} \quad \textbf{Nota:} \underline{\hspace{0.9cm}}/10

\subsection*{PARTE 1: RECONHECIMENTO}

\begin{enumerate}
  \item Forme todas as silabas com ÃO:
  \begin{center}
  \resizebox{\textwidth}{!}{%
\begin{tabular}{ccccc}
  \sil{B}{ÃO} & \sil{C}{ÃO} & \sil{D}{ÃO} & \sil{F}{ÃO} & \sil{G}{ÃO} \\
  \sil{J}{ÃO} & \sil{M}{ÃO} & \sil{N}{ÃO} & \sil{P}{ÃO} & \sil{V}{ÃO} \\
  \sil{Z}{ÃO} & \sil{S}{ÃO} & \sil{T}{ÃO} & \sil{R}{ÃO} & \sil{L}{ÃO}
  \end{tabular}}
  \end{center}

  \item Sublinhe palavras que terminam em ``ÃO'':
  \begin{center}
  \uline{pao} \quad gato \quad \uline{nao} \quad sapo \quad \uline{mao} \quad rato \quad \uline{cao}
  \end{center}

  \item Identifique a silaba ``ÃO'' em cada palavra:
  \begin{center}
  BÃO $\rightarrow$ silaba: \underline{\hspace{1cm}} \quad CÃO $\rightarrow$ silaba: \underline{\hspace{1cm}} \quad DÃO $\rightarrow$ silaba: \underline{\hspace{1cm}} \\
  FÃO $\rightarrow$ silaba: \underline{\hspace{1cm}} \quad GÃO $\rightarrow$ silaba: \underline{\hspace{1cm}} \quad JÃO $\rightarrow$ silaba: \underline{\hspace{1cm}} \\
  MÃO $\rightarrow$ silaba: \underline{\hspace{1cm}} \quad NÃO $\rightarrow$ silaba: \underline{\hspace{1cm}} \quad PÃO $\rightarrow$ silaba: \underline{\hspace{1cm}} \\
  VÃO $\rightarrow$ silaba: \underline{\hspace{1cm}} \quad ZÃO $\rightarrow$ silaba: \underline{\hspace{1cm}}
  \end{center}

  \item Separe em silabas:
  \begin{center}
  PÃO $\rightarrow$ \underline{\hspace{1cm}} \quad CÃO $\rightarrow$ \underline{\hspace{1cm}} \quad MÃO $\rightarrow$ \underline{\hspace{1cm}} \quad CORAÇÃO $\rightarrow$ \underline{\hspace{1cm}} \underline{\hspace{1cm}} \underline{\hspace{1cm}}
  \end{center}

  \item Conte as silabas:
  \begin{center}
  PÃO $\rightarrow$ \underline{\hspace{1cm}} \quad CACHORRO $\rightarrow$ \underline{\hspace{1cm}} \quad CORAÇÃO $\rightarrow$ \underline{\hspace{1cm}}
  \end{center}
\end{enumerate}

\subsection*{PARTE 2: PRODUCAO GUIADA}

\begin{enumerate}[resume]
  \item Complete com ``ÃO'':
  \begin{center}
  b\underline{~~~} \quad c\underline{~~~} \quad d\underline{~~~} \quad f\underline{~~~} \quad g\underline{~~~} \quad j\underline{~~~} \quad m\underline{~~~} \quad n\underline{~~~} \quad p\underline{~~~} \quad v\underline{~~~} \quad z\underline{~~~}
  \end{center}

  \item Escreva 2 palavras para cada silaba ``ÃO'':
  \begin{center}
  BÃO: \underline{\hspace{3cm}} \underline{\hspace{3cm}} \quad CÃO: \underline{\hspace{3cm}} \underline{\hspace{3cm}} \quad DÃO: \underline{\hspace{3cm}} \underline{\hspace{3cm}} \\
  FÃO: \underline{\hspace{3cm}} \underline{\hspace{3cm}} \quad GÃO: \underline{\hspace{3cm}} \underline{\hspace{3cm}} \quad MÃO: \underline{\hspace{3cm}} \underline{\hspace{3cm}} \\
  NÃO: \underline{\hspace{3cm}} \underline{\hspace{3cm}} \quad PÃO: \underline{\hspace{3cm}} \underline{\hspace{3cm}} \quad VÃO: \underline{\hspace{3cm}} \underline{\hspace{3cm}}
  \end{center}

  \item Copie as palavras:
  \begin{center}
  bão \quad cão \quad dão \quad fão \quad gão \quad jão \quad mão \quad não \quad pão \quad vão \quad zão
  \end{center}

  \item Separe em silabas:
  \begin{center}
  CORAÇÃO = \underline{\hspace{1cm}} \underline{\hspace{1cm}} \underline{\hspace{1cm}} \quad CACHORRO = \underline{\hspace{1cm}} \underline{\hspace{1cm}} \underline{\hspace{1cm}} \quad BARRIGÃO = \underline{\hspace{1cm}} \underline{\hspace{1cm}} \underline{\hspace{1cm}}
  \end{center}

  \item Escreva 3 frases usando palavras com ``ÃO'':
  \begin{enumerate}[label=\alph*)]
    \item \underline{\hspace{8cm}}
    \item \underline{\hspace{8cm}}
    \item \underline{\hspace{8cm}}
  \end{enumerate}
\end{enumerate}

\subsection*{PARTE 3: INTEGRACAO}

\begin{enumerate}[resume]
  \item Leia as palavras em voz alta:
  \begin{center}
  $\rightarrow$ bão  ~  cão  ~  dão  ~  fão  ~  gão  ~  jão  ~  mão  ~  não  ~  pão  ~  vão  ~  zão
  \end{center}

  \item Leia as frases em voz alta:
  \begin{center}
  $\rightarrow$ A mao e grande.\\
  $\rightarrow$ Nao e bom.\\
  $\rightarrow$ O pao esta quente.\\
  $\rightarrow$ O cao late.\\
  $\rightarrow$ Eles vao para a escola.
  \end{center}

  \item Circule a opcao correta:
  \begin{enumerate}[label=(\alph*)]
    \item ``ÃO'' e: ( ) oral ( ) nasal ( ) vibrante
    \item ``PÃO'' tem: ( ) 1 silaba ( ) 2 silabas ( ) 3 silabas
    \item ``ÃO'' pode ser precedido por: ( ) so vogais ( ) consoantes e vogais ( ) so consoantes
  \end{enumerate}

  \item Leia e complete:
  \begin{center}
  A m\underline{~~~} \quad Na\underline{~~~} \quad O p\underline{~~~} \quad O c\underline{~~~} \quad Eles v\underline{~~~}
  \end{center}

  \item Avaliacao diagnostica:
  \begin{center}
  $\square$ Forma silabas com ÃO para todas as consoantes\\
  $\square$ Identifica ditongo ÃO em palavras\\
  $\square$ Le palavras com ÃO fluentemente\\
  $\square$ Separa palavras em silabas\\
  $\square$ Escreve frases com ÃO
  \end{center}
\end{enumerate}
\end{exercicio}

\begin{dica}
\textbf{Dica para o Professor:} Esta e a licao mais importante do Volume 1! A silaba ``ÃO'' com QUALQUER consoante e a base da escrita em portugues. Use cartoes com as consoantes (B, C, D, F, G, J, M, N, P, V, Z) e um cartao ``ÃO''. Peca que os alunos montem as silabas: B+ÃO=BÃO, C+ÃO=CÃO, etc. Reforce que ``ÃO'' e nasal (boca fechada pelo M antes do O).
\end{dica}

%% ============================================================
%% UNIDADE 10 -- DITONGOS NASAIS ÃE E ÕE
%% ============================================================

\chapter{Unidade 10 --- Ditongos Nasais ÃE e ÕE (MÃE, PÃES, LIÇÕES, CANÕES)}

\metadata{I-02}{Nivel 4 -- Silabico-Alfabetico Avancado}{EF01LP03}{D1}{EF01LP03}{Resolucao CNE/CEB 7/2010, art. 12}

\textbf{Regra:} Os ditongos nasais ``ÃE'' (como em m\~ae e p\~aes) e ``ÕE'' (como em li\c{c}\~oes e can\c{c}\~oes) seguem a mesma logica do ``ÃO'' --- podem ser precedidos por QUALQUER consoante. Nao confundir com ``AM'' e ``OM'', que sao desinencias verbais (falam, comem) e nao ditongos com E.

\plancheta{Folha de Pratica 10.1 --- Ditongo ÃE}
\begin{exercicio}[Folha de Pratica 10.1 --- Ditongo ÃE]
\noindent\textbf{Nome:} \underline{\hspace{3.2cm}} \quad \textbf{Data:} \underline{\hspace{1.6cm}} \quad \textbf{Nota:} \underline{\hspace{0.9cm}}/10

\subsection*{PARTE 1: RECONHECIMENTO}

\begin{enumerate}
  \item Forme silabas com ÃE:
  \begin{center}
  \resizebox{\textwidth}{!}{%
\begin{tabular}{ccccc}
  \sil{M}{ÃE} & \sil{P}{ÃE} & \sil{V}{ÃE} & \sil{C}{ÃE} & \sil{S}{ÃE}
  \end{tabular}}
  \end{center}

  \item Sublinhe palavras que terminam em ``ÃE'':
  \begin{center}
  \uline{mae} \quad gato \quad \uline{pae} \quad sapo \quad \uline{nae} \quad rato \quad \uline{tae}
  \end{center}

  \item Identifique o ditongo ``ÃE'' em cada palavra:
  \begin{center}
  MÃE $\rightarrow$ ditongo: \underline{\hspace{1cm}} \quad PÃE $\rightarrow$ ditongo: \underline{\hspace{1cm}} \quad NÃO $\rightarrow$ ditongo: \underline{\hspace{1cm}}
  \end{center}

  \item Separe em silabas:
  \begin{center}
  MÃE $\rightarrow$ \underline{\hspace{1cm}} \quad PÃES $\rightarrow$ \underline{\hspace{1cm}} \underline{\hspace{1cm}} \quad CANÇÃO $\rightarrow$ \underline{\hspace{1cm}} \underline{\hspace{1cm}} \underline{\hspace{1cm}}
  \end{center}

  \item Conte as silabas:
  \begin{center}
  MÃE $\rightarrow$ \underline{\hspace{1cm}} \quad CANÇÕES $\rightarrow$ \underline{\hspace{1cm}} \quad CORAÇÕES $\rightarrow$ \underline{\hspace{1cm}}
  \end{center}
\end{enumerate}

\subsection*{PARTE 2: PRODUCAO GUIADA}

\begin{enumerate}[resume]
  \item Complete com ``ÃE'' ou ``ÃES'':
  \begin{center}
  m\underline{~~~} \quad p\underline{~~~} \quad v\underline{~~~} \quad c\underline{~~~} \quad n\underline{~~~}
  \end{center}

  \item Escreva 2 palavras com ditongo ``ÃE'':
  \begin{center}
  1. \underline{\hspace{3cm}} \quad 2. \underline{\hspace{3cm}} \quad 3. \underline{\hspace{3cm}}
  \end{center}

  \item Copie as palavras:
  \begin{center}
  mãe \quad pães \quad mães \quad vães \quad nães
  \end{center}

  \item Separe em silabas:
  \begin{center}
  PÃES = \underline{\hspace{1cm}} \underline{\hspace{1cm}} \quad CANÇÕES = \underline{\hspace{1cm}} \underline{\hspace{1cm}} \underline{\hspace{1cm}} \quad CORAÇÕES = \underline{\hspace{1cm}} \underline{\hspace{1cm}} \underline{\hspace{1cm}}
  \end{center}

  \item Escreva 3 frases usando palavras com ``ÃE'':
  \begin{enumerate}[label=\alph*)]
    \item \underline{\hspace{8cm}}
    \item \underline{\hspace{8cm}}
    \item \underline{\hspace{8cm}}
  \end{enumerate}
\end{enumerate}

\subsection*{PARTE 3: INTEGRACAO}

\begin{enumerate}[resume]
  \item Leia as palavras em voz alta:
  \begin{center}
  $\rightarrow$ mãe  ~  pães  ~  mães  ~  vães  ~  nães
  \end{center}

  \item Leia as frases em voz alta:
  \begin{center}
  $\rightarrow$ A mae e carinhosa.\\
  $\rightarrow$ Os paes estao no forno.\\
  $\rightarrow$ As maes trabalham.
  \end{center}

  \item Circule a opcao correta:
  \begin{enumerate}[label=(\alph*)]
    \item ``ÃE'' e: ( ) oral ( ) nasal ( ) vibrante
    \item ``MÃE'' tem: ( ) 1 silaba ( ) 2 silabas ( ) 3 silabas
    \item ``PÃES'' tem ditongo: ( ) ÃE ( ) ÃO ( ) ÕE
  \end{enumerate}

  \item Leia e complete:
  \begin{center}
  A m\underline{~~~} \quad Os p\underline{~~~} \quad As m\underline{~~~}
  \end{center}

  \item Leitura em voz alta:
  \begin{center}
  mãe \quad pães \quad mães \quad vães \quad nães \quad canções \quad corações
  \end{center}
\end{enumerate}
\end{exercicio}

\begin{dica}
\textbf{Dica para o Professor:} O ditongo ``ÃE'' segue a mesma regra do ``ÃO'' --- pode ser precedido por qualquer consoante. Compare ``MÃE'' (ÃE) com ``MÃO'' (ÃO) para que percebam a diferenca na abertura da boca.
\end{dica}

\plancheta{Folha de Pratica 10.2 --- Ditongo ÕE}
\begin{exercicio}[Folha de Pratica 10.2 --- Ditongo ÕE]
\noindent\textbf{Nome:} \underline{\hspace{3.2cm}} \quad \textbf{Data:} \underline{\hspace{1.6cm}} \quad \textbf{Nota:} \underline{\hspace{0.9cm}}/10

\subsection*{PARTE 1: RECONHECIMENTO}

\begin{enumerate}
  \item Forme silabas com ÕE:
  \begin{center}
  \resizebox{\textwidth}{!}{%
\begin{tabular}{ccccc}
  \sil{L}{IÇÕES} & \sil{C}{ANÕES} & \sil{S}{IÇÕES} & \sil{L}{IÇÕES} & \sil{F}{UNÇÕES}
  \end{tabular}}
  \end{center}

  \item Sublinhe palavras que terminam em ``ÕES'':
  \begin{center}
  \uline{lições} \quad gato \quad \uline{canções} \quad sapo \quad \uline{funções} \quad rato \quad \uline{opções}
  \end{center}

  \item Identifique o ditongo ``ÕE'' em cada palavra:
  \begin{center}
  LIÇÕES $\rightarrow$ ditongo: \underline{\hspace{1cm}} \quad CANÕES $\rightarrow$ ditongo: \underline{\hspace{1cm}} \quad FUNÇÕES $\rightarrow$ ditongo: \underline{\hspace{1cm}}
  \end{center}

  \item Separe em silabas:
  \begin{center}
  LIÇÕES $\rightarrow$ \underline{\hspace{1cm}} \underline{\hspace{1cm}} \underline{\hspace{1cm}} \quad CANÕES $\rightarrow$ \underline{\hspace{1cm}} \underline{\hspace{1cm}} \underline{\hspace{1cm}} \quad FUNÇÕES $\rightarrow$ \underline{\hspace{1cm}} \underline{\hspace{1cm}} \underline{\hspace{1cm}}
  \end{center}

  \item Conte as silabas:
  \begin{center}
  LIÇÕES $\rightarrow$ \underline{\hspace{1cm}} \quad CANÕES $\rightarrow$ \underline{\hspace{1cm}} \quad FUNÇÕES $\rightarrow$ \underline{\hspace{1cm}}
  \end{center}
\end{enumerate}

\subsection*{PARTE 2: PRODUCAO GUIADA}

\begin{enumerate}[resume]
  \item Complete com ``ÕES'':
  \begin{center}
  li\underline{~~~} \quad can\underline{~~~} \quad fun\underline{~~~} \quad o\underline{~~~} \quad le\~{o}es \underline{\hspace{0.5cm}}
  \end{center}

  \item Escreva 3 palavras que terminam em ``ÕES'':
  \begin{center}
  1. \underline{\hspace{3cm}} \quad 2. \underline{\hspace{3cm}} \quad 3. \underline{\hspace{3cm}}
  \end{center}

  \item Copie as palavras:
  \begin{center}
  lições \quad canções \quad funções \quad opções \quad lições
  \end{center}

  \item Separe em silabas:
  \begin{center}
  LIÇÕES = \underline{\hspace{1cm}} \underline{\hspace{1cm}} \underline{\hspace{1cm}} \quad CANÕES = \underline{\hspace{1cm}} \underline{\hspace{1cm}} \underline{\hspace{1cm}} \quad FUNÇÕES = \underline{\hspace{1cm}} \underline{\hspace{1cm}} \underline{\hspace{1cm}}
  \end{center}

  \item Escreva 3 frases usando palavras com ``ÕES'':
  \begin{enumerate}[label=\alph*)]
    \item \underline{\hspace{8cm}}
    \item \underline{\hspace{8cm}}
    \item \underline{\hspace{8cm}}
  \end{enumerate}
\end{enumerate}

\subsection*{PARTE 3: INTEGRACAO}

\begin{enumerate}[resume]
  \item Leia as palavras em voz alta:
  \begin{center}
  $\rightarrow$ lições  ~  canções  ~  funções  ~  opções  ~  lições
  \end{center}

  \item Leia as frases em voz alta:
  \begin{center}
  $\rightarrow$ As licoes sao importantes.\\
  $\rightarrow$ As cancoes sao bonitas.\\
  $\rightarrow$ As funcoes do corpo.
  \end{center}

  \item Circule a opcao correta:
  \begin{enumerate}[label=(\alph*)]
    \item ``ÕES'' e: ( ) oral ( ) nasal ( ) vibrante
    \item ``LIÇÕES'' tem: ( ) 2 silabas ( ) 3 silabas ( ) 4 silabas
    \item ``ÕES'' aparece: ( ) no inicio ( ) no meio ( ) no fim da palavra
  \end{enumerate}

  \item Leia e complete:
  \begin{center}
  As li\underline{~~~} \quad As can\underline{~~~} \quad As fun\underline{~~~}
  \end{center}

  \item Leitura em voz alta:
  \begin{center}
lições \quad canções \quad funções \quad opções \quad lições \quad uniões \quad sanções
\end{center}
\end{enumerate}
\end{exercicio}

\begin{dica}
\textbf{Dica para o Professor:} O ditongo ``ÕE'' (escrito ``ÕES'' no plural) e nasal. Compare ``LIÇÃO'' (singular, ÃO) com ``LIÇÕES'' (plural, ÕES) para que percebam a mudanca. Use o espelho: em ``LIÇÃO'', a boca abre; em ``LIÇÕES'', a boca se arredonda e o ar sai pelo nariz.
\end{dica}

%% ============================================================
\input{checkpoints/rastreio-u6}
%% FIM DA PARTE II -- SEQUENCIA DIDATICA
%% ============================================================
